\hypertarget{home}{%
\section{Home}\label{home}}

\hypertarget{vuxe4lkommen}{%
\subsection{Välkommen}\label{vuxe4lkommen}}

Välkommen till Campus Mölndals fantastiska kodwebbplats för JIN23. Här
hittar du allt du behöver för att komma igång med programmering i Java.

Java är ett objektorienterat programmeringsspråk som utvecklades av Sun
Microsystems (numera ägt av Oracle). Det används ofta för att skapa
programvaror och plattformar som sträcker sig från små mobila enheter
till stora datorsystem. Java har blivit populärt på grund av sin
plattformsoberoende natur och förmåga att köra på olika operativsystem.

Java har en enkel och lättförståelig syntax som gör det till ett
idealiskt språk för nybörjare. Samtidigt erbjuder det avancerade
funktioner och bibliotek för att hantera komplexa program. Generiska
typer, anonyma inre klasser, lambda-uttryck och trådhantering är några
av de kraftfulla funktionerna som Java erbjuder.

En av de stora fördelarna med Java är dess stora ekosystem av bibliotek
och ramverk. Det finns ett brett utbud av bibliotek för att hjälpa till
med allt från grafisk användargränssnittsutveckling till
databashantering och nätverksprogrammering. Ramverk som Spring och
Hibernate förenklar utvecklingsprocessen och främjar goda designmönster.

Java används inom en mängd olika områden, inklusive webbutveckling,
mobilapputveckling, spelutveckling och företagsapplikationer. Det är ett
mångsidigt språk som ger utvecklare möjlighet att skapa kraftfulla och
skalbara program.

\hypertarget{java}{%
\section{Java}\label{java}}

En introduktion till ämnet Java inom programmering.

\hypertarget{introduktion}{%
\subsection{Introduktion}\label{introduktion}}

Java är ett objektorienterat programmeringsspråk som utvecklades av Sun
Microsystems. Det används ofta för att skapa applikationer för olika
plattformar. Java är ett modern, flexibel och mångsidig
programmeringsspråk som ger utvecklare möjlighet att skapa allt från
enkla webbapplikationer till komplexa distribuerade system.

Java har en stark och tydlig syntax som gör det lätt att lära sig och
använda. Språket har ett brett utbud av funktioner som hjälper
utvecklare att skapa avancerade system. Funktioner som generiska typer,
anonyma metoder, LINQ, lambda-uttryck och parallell programmering gör
Java till ett mycket kraftfullt programmeringsspråk. Java stöder också
många olika typer av programutveckling, inklusive webbappar, mobilappar,
spel och systemprogramvara.

\hypertarget{skolan}{%
\section{Skolan}\label{skolan}}

Här är en lista på webbtjänster vi använder i skolan.

\hypertarget{huxe4r-uxe4r-skolans-webbplatser}{%
\subsection{Här är skolans
webbplatser}\label{huxe4r-uxe4r-skolans-webbplatser}}

\begin{itemize}
\item
  \href{https://campus.molndal.se}{Campus Mölndals hemsida}
\item
  \href{https://www.linkedin.com/school/campus-molndal/}{LinkedIn}
\item
  \href{https://www.instagram.com/yh_campusmolndal_it/}{Instagram}
\item
  \href{https://www.facebook.com/campusmolndal/}{Facebook}
\item
  \href{https://www.linkedin.com/in/marcusmedina/?locale=sv_SE}{Marcus
  Medina}
\item
  \href{https://www.linkedin.com/in/annika-lund-profile/}{Annika Lund}
\item
  \href{https://www.linkedin.com/search/results/all/?keywords=\%23campusm\%C3\%B6lndal&origin=GLOBAL_SEARCH_HEADER&sid=gtu}{\#CampusMölndal}
\end{itemize}

LinkedIn är ett måste när man vill skapa sig en karriär. Så tveka inte
och skapa ett konto. Det är gratis och du kan skapa ett konto med din
skolmail.

\hypertarget{skolportalen}{%
\section{Skolportalen}\label{skolportalen}}

\href{https://www.molndal.se/skolportal?authmech=Personal\%20och\%20elever}{Öppna
skolportalen}\{: .btn .btn-blue \}

Skolportalen är din ingång till skolan och dina kurser. Här kan du se
dina betyg, läsa om dina kurser och se dina scheman. Här finns även
länkar till Google Classroom och Google Drive. Plus en massa annat.

\hypertarget{google-classroom}{%
\section{Google Classroom}\label{google-classroom}}

\href{https://classroom.google.com/}{Öppna Google Classroom}\{: .btn
.btn-blue \}

Vi använder Google Classroom på Campus Mölndal.

I ditt klassrum hittar du information, länkar och planering för kursen.
Du kan också lämna in uppgifter och inlämningar där.

Varje kurs har sin egen Meet kanal där.

\hypertarget{google-calendar}{%
\section{Google Calendar}\label{google-calendar}}

\href{https://calendar.google.com/}{Öppna Google Calendar}\{: .btn
.btn-blue \}

Lektioner, tentor, inlämningar och annat skoj hamnar i klassens google
calendar. Du kan prenumerera på kalendern och få alla händelser direkt i
din telefon.

\hypertarget{google-documents}{%
\section{Google Documents}\label{google-documents}}

\href{https://docs.google.com/}{Öppna Google Documents}\{: .btn
.btn-blue \}

Vi använder Google Documents för att skriva uppsatser, föra anteckningar
och annat. Då du får ett Google konto av skolan kan du samla all din
material där. Det är också enkelt att dela med sig av dokument till
andra.

\hypertarget{discord}{%
\section{Discord}\label{discord}}

\href{https://discord.com/download}{Installera Discord}\{: .btn
.btn-blue \}

När du blir intagen till kursen får du en länk till en Discordgrupp som
är skapad för klassen.

I den postas det information som behöver komma ut snabbt, exempelvis
påminnelser om möten och annat. Där sker även det mesta av online
kommunikationen mellan studenter och utbildare.

\emph{(Email är så 90tal\ldots)}

\hypertarget{info}{%
\section{Info}\label{info}}

Här kommer jag att posta information om sidan, som t.ex.
versionshistorik, användarvillkor, hur man använder den mm. Denna sidan
är baserad på
\href{https://campusmolndaleducation.github.io/CLO22/}{CLO22} som är en
sida som jag skapat för att samla information om C\# och .NET. All
information ``översattes'' till Java.

\hypertarget{versionshistorik}{%
\section{Versionshistorik}\label{versionshistorik}}

\begin{longtable}[]{@{}
  >{\raggedright\arraybackslash}p{(\columnwidth - 4\tabcolsep) * \real{0.0800}}
  >{\raggedright\arraybackslash}p{(\columnwidth - 4\tabcolsep) * \real{0.1200}}
  >{\raggedright\arraybackslash}p{(\columnwidth - 4\tabcolsep) * \real{0.7900}}@{}}
\toprule\noalign{}
\begin{minipage}[b]{\linewidth}\raggedright
Version
\end{minipage} & \begin{minipage}[b]{\linewidth}\raggedright
Datum
\end{minipage} & \begin{minipage}[b]{\linewidth}\raggedright
Beskrivning
\end{minipage} \\
\midrule\noalign{}
\endhead
\bottomrule\noalign{}
\endlastfoot
1.0.1 & 2022-12-09 & Uppdatering av sidan. Mermaid funkar inte,
diagrammen syns inte 😞 \\
1.0.0 & 2022-12-09 & Första publicering. Första publicering, \\
\end{longtable}

\hypertarget{om-sidan}{%
\section{Om sidan}\label{om-sidan}}

Denna sidan skapades för att kunna samla alla dokument på en plats, för
att göra det enklare för studenter att hitta material från lektioner och
annat.

Skicka gärna in förslag på förbättringar eller nya dokument till mig.

Ni hittar mig på \href{https://discord.com/users/mrmarcus1337}{Discord}
eller \href{marcus.medina-ramirez@molndal.se}{skolans mail}.

/Marcus

\hypertarget{licens-fuxf6r-materialet-puxe5-denna-webbsida}{%
\section{Licens för materialet på denna
webbsida}\label{licens-fuxf6r-materialet-puxe5-denna-webbsida}}

Materialet på denna webbsida är licensierat under en
\href{https://creativecommons.org/licenses/by-nc-sa/4.0/}{Creative
Commons BY-NC-SA (Attribution-NonCommercial-ShareAlike)} licens. Det
innebär att du fritt får använda materialet för
\textbf{icke-kommersiella ändamål}, under förutsättning att du
\textbf{anger källan} och delar \textbf{ditt material under samma
licens}.

Denna licens har valts för att ge studenterna friheten att använda koden
på det sätt de önskar, samtidigt som materialet inte tillåts kopieras
eller användas i kommersiella syften av andra utbildningscenter än
Campus Mölndal.

För mer information om Creative Commons BY-NC-SA-licensen och dess
villkor kan du besöka
\href{https://creativecommons.org/licenses/by-nc-sa/}{Creative Commons
webbplats}.

\hypertarget{sammanfattning}{%
\subsection{Sammanfattning}\label{sammanfattning}}

\begin{itemize}
\tightlist
\item
  \textbf{Licens:} Creative Commons BY-NC-SA
  (Attribution-NonCommercial-ShareAlike)
\item
  \textbf{Användning:} Materialet får användas för icke-kommersiella
  ändamål.
\item
  \textbf{Krav:} Du måste ange källan och dela ditt material under samma
  licens.
\item
  \textbf{Undantag:} Materialet får inte kopieras eller användas i
  kommersiella syften av andra utbildningscenter än Campus Mölndal.
\end{itemize}

Vänligen observera att licensvalet är noggrant genomtänkt för att uppnå
de önskade användningsvillkoren för materialet. Vid eventuella tvivel
eller frågor rekommenderas det att du konsulterar en juridisk expert för
att säkerställa att licensen är korrekt och uppfyller dina behov och
jurisdiktionens krav.

\hypertarget{enkelt-fuxf6rklarat}{%
\subsection{Enkelt förklarat}\label{enkelt-fuxf6rklarat}}

\textbf{Du får:} + Använda koden för icke-kommersiella ändamål, under
förutsättning att du anger källan. + Använda koden i ett projekt och
dela det under samma licens.

\textbf{Du får inte:} + Kopiera koden eller artiklarna för spridning. +
Använda koden eller artiklar i kommersiella syften av andra
utbildningscenter/företag än Campus Mölndal.

\hypertarget{mit-licens-fuxf6r-just-the-docs}{%
\section{MIT-licens för Just the
docs}\label{mit-licens-fuxf6r-just-the-docs}}

Upphovsrätt (c) 2023 Marcus Medina, Campus Mölndal

Härmed beviljas tillstånd, kostnadsfritt, till varje person som erhåller
en kopia av denna programvara och tillhörande dokumentationsfiler
(''Programvaran''), att använda Programvaran utan begränsningar,
inklusive, utan begränsning, rättigheter att använda, kopiera, ändra,
sammanfoga, publicera, distribuera, underlicensiera och/eller sälja
kopior av Programvaran samt att tillåta personer som erhåller
Programvaran att göra detta, under förutsättning att följande villkor
uppfylls:

Ovanstående upphovsrättsnotis och denna licensnotis ska inkluderas i
alla kopior eller betydande delar av Programvaran.

PROGRAMVARAN TILLHANDAHÅLLS ``SOM DEN ÄR'', UTAN GARANTI AV NÅGOT SLAG,
UTTRYCKT ELLER UNDERFÖRSTÅDD, INKLUSIVE MEN INTE BEGRÄNSAT TILL
GARANTIER OM SÄLJBARHET, LÄMPLIGHET FÖR ETT SPECIFIKT ÄNDAMÅL OCH
OÖVERTRÄDELSE AV IMMATERIELLA RÄTTIGHETER. I INGA HÄNDELSER SKA
UPPHOVSRÄTTSHAVERNA ELLER UPPHOVSRÄTTENS INNEHAVARE ANSVARA FÖR SKADOR
ELLER ANDRA FORDRINGAR, OAVSETT OM DET GÄLLER AVTALSBROTT, SKADESTÅND
ELLER ANNAT, SOM UPPKOMMER AV, UTAN BEGRÄNSNING, I SAMBAND MED
PROGRAMVARAN ELLER ANVÄNDNINGEN AV ELLER ANDRA VERKSAMHETER MED
PROGRAMVARAN.

\begin{itemize}
\tightlist
\item
  \href{https://opensource.org/licenses/MIT}{MIT-licens} (engelska)
\end{itemize}

\hypertarget{verktyg}{%
\section{Verktyg}\label{verktyg}}

I denna artikel kommer vi att titta på några användbara verktyg för att
underlätta utvecklingen av Java-program. Genom att använda dessa verktyg
kan vi öka vår produktivitet och förbättra kvaliteten på vår kod. Det är
viktigt att välja verktyg som bäst passar våra behov och att utforska
och använda dem på rätt sätt.

\hypertarget{ide-integrated-development-environment}{%
\subsection{1. IDE (Integrated Development
Environment)}\label{ide-integrated-development-environment}}

En IDE är en programvara som ger oss en miljö för att skriva, kompilera
och testa våra Java-program. Här är några populära IDE:er för Java:

\begin{itemize}
\tightlist
\item
  \textbf{Eclipse}: En kraftfull och flexibel IDE som är mycket populär
  bland Java-utvecklare. Den har omfattande funktioner och stöd för att
  underlätta utvecklingsprocessen.
\item
  \textbf{IntelliJ IDEA}: En annan populär IDE som erbjuder avancerade
  funktioner och verktyg för att öka produktiviteten. Den har också stöd
  för andra programmeringsspråk.
\item
  \textbf{NetBeans}: En lättanvänd IDE som är utvecklad specifikt för
  Java-programmering. Den har en intuitiv användargränssnitt och
  erbjuder en rad användbara funktioner.
\end{itemize}

\hypertarget{byggverktyg}{%
\subsection{2. Byggverktyg}\label{byggverktyg}}

Byggverktyg hjälper oss att automatisera byggprocessen för våra
Java-projekt. De hanterar kompilering, beroenden, testning och
paketering av vår kod. Här är några populära byggverktyg för Java:

\begin{itemize}
\tightlist
\item
  \textbf{Maven}: Ett kraftfullt verktyg för att hantera
  projektberoenden och automatisera byggprocessen. Det använder en
  XML-baserad konfigurationsfil för att definiera projektets struktur
  och beroenden.
\item
  \textbf{Gradle}: Ett flexibelt byggverktyg som erbjuder en DSL
  (Domain-Specific Language) för att definiera byggskript. Det stöder
  också Maven-repositorier och har enkla och läsbara
  konfigurationsfiler.
\end{itemize}

\hypertarget{versionshanteringssystem}{%
\subsection{3.
Versionshanteringssystem}\label{versionshanteringssystem}}

Ett versionshanteringssystem hjälper oss att hantera ändringar i vår kod
och samarbeta med andra utvecklare. Det spårar ändringar, hanterar
konflikter och gör det möjligt för oss att återställa tidigare versioner
av vår kod. Här är några populära versionshanteringssystem för Java:

\begin{itemize}
\tightlist
\item
  \textbf{Git}: Ett distribuerat versionshanteringssystem som är mycket
  populärt och används av många utvecklare runt om i världen. Det
  erbjuder enkla och kraftfulla kommandon för att hantera ändringar och
  samarbeta med andra.
\item
  \textbf{Subversion}: Ett centralt versionshanteringssystem som används
  av många företag och organisationer. Det erbjuder enkla kommandon och
  en webbgränssnitt för att hantera ändringar och samarbeta med andra.
\end{itemize}

\hypertarget{enhetstestningsramverk}{%
\subsection{4. Enhetstestningsramverk}\label{enhetstestningsramverk}}

Enhetstestning är en viktig del av utvecklingsprocessen för att
säkerställa att vår kod fungerar som förväntat. Ett
enhetstestningsramverk hjälper oss att skriva och köra tester för vår
Java-kod. Här är några populära enhetstestningsramverk för Java:

\begin{itemize}
\tightlist
\item
  \textbf{JUnit}: Ett robust och populärt enhetstestningsramverk för
  Java. Det erbjuder enkla och kraftfulla funktioner för att skriva och
  köra tester.
\item
  \textbf{TestNG}: Ett flexibelt och utökat enhetstestningsramverk för
  Java. Det erbjuder avancerade funktioner som parallell körning av
  tester och konfiguration av testfall.
\end{itemize}

\hypertarget{felsuxf6kning-och-profileringsverktyg}{%
\subsection{5. Felsökning och
profileringsverktyg}\label{felsuxf6kning-och-profileringsverktyg}}

Felsökning och profilering är viktiga för att hitta och åtgärda problem
i vår Java-kod. Här är några användbara verktyg för felsökning och
profilering i Java:

\begin{itemize}
\tightlist
\item
  \textbf{Eclipse Memory Analyzer}: Ett verktyg för att analysera
  minnesdumpar och hitta minnesläckor och andra minnesrelaterade
  problem.
\item
  \textbf{VisualVM}: Ett verktyg för att övervaka och profilera
  Java-program i realtid. Det ger detaljerad information om
  minnesanvändning, CPU-användning och andra prestandarelaterade
  mätvärden.
\end{itemize}

\hypertarget{summering}{%
\subsection{Summering}\label{summering}}

Att använda rätt verktyg kan göra en stor skillnad i vår
utvecklingsprocess. I denna artikel har vi tittat på några användbara
verktyg för att underlätta utvecklingen av Java-program. Genom att välja
och använda dessa verktyg på rätt sätt kan vi öka vår produktivitet och
förbättra kvaliteten på vår kod.

\hypertarget{installation}{%
\section{Installation}\label{installation}}

Här finns en lista på program du kan komma att behöva under
utbildningens gång.

\emph{Filen editerades senast 2022-12-09}

\hypertarget{git-installation}{%
\section{Git installation}\label{git-installation}}

\href{https://git-scm.com/}{Installera Git}\{: .btn \}

\emph{Filen editerades senast 2022-12-09} \#\# Beskrivning

Git är ett versionhanteringssystem som används för att spara och hantera
källkod. Det är ett av de mest använda versionhanteringssystemen i
dagens samhälle. Det är ett gratis program som är öppen källkod. Den här
installationen är för Windows. För att installera Git på Mac eller
Linux, gå till
\href{https://git-scm.com/book/en/v2/Getting-Started-Installing-Git}{Git
installation}. Du behöver inte den om du bara arbetar med Visual Studio,
men ska du in på Console / Terminalen så behöver du det.

\hypertarget{visual-studio-installation}{%
\section{Visual Studio installation}\label{visual-studio-installation}}

\href{https://visualstudio.microsoft.com/}{Installera Visual Studio}\{:
.btn \}

\emph{Filen editerades senast 2022-12-09} \#\# Addons till Visual Studio
som jag rekommenderar

\begin{itemize}
\tightlist
\item
  \href{https://marketplace.visualstudio.com/items?itemName=SteveCadwallader.CodeMaidVS2022}{Codemaid}
\item
  \href{https://marketplace.visualstudio.com/items?itemName=josefpihrt.Roslynator2022}{Roslynator}
\item
  \href{https://marketplace.visualstudio.com/items?itemName=GitHub.copilotvs}{Github
  Copilot}
\item
  \href{https://marketplace.visualstudio.com/items?itemName=martinw.SmartPaster}{Smart
  paste}
\item
  \href{https://marketplace.visualstudio.com/items?itemName=MadsKristensen.AddNewFile}{Add
  new page}
\end{itemize}

\hypertarget{beskrivning}{%
\subsection{Beskrivning}\label{beskrivning}}

Visual Studio är ett IDE för Java och .NET. Det är främst den som
används när man arbetar med Java.

\hypertarget{vs-code-installation}{%
\section{VS Code installation}\label{vs-code-installation}}

\href{https://code.visualstudio.com/}{Installera Visual Studio Code}\{:
.btn \}

\emph{Filen editerades senast 2022-12-09} \#\# Addons till VS Code som
jag rekommenderar

\begin{itemize}
\tightlist
\item
  \href{https://marketplace.visualstudio.com/items?itemName=josefpihrt-vscode.roslynator}{Roslynator}
\item
  {[}Github Copilot{]}
\end{itemize}

\hypertarget{beskrivning}{%
\subsection{Beskrivning}\label{beskrivning}}

VS Code är ett IDE för Java och .NET. Det är ett gratis program som är
öppen källkod. Den är mer anpassaningsbar än Visual Studio och den har
inte lika många funktioner som Visual Studio. Den är mer anpassad för
webbutveckling än Visual Studio.

\hypertarget{intellijinstallation}{%
\section{IntelliJinstallation}\label{intellijinstallation}}

\href{https://www.jetbrains.com/idea/}{Installera JetBrains IntelliJ}\{:
.btn \}

\hypertarget{addons-till-intellij-som-jag-rekommenderar}{%
\subsection{Addons till IntelliJ som jag
rekommenderar}\label{addons-till-intellij-som-jag-rekommenderar}}

\begin{itemize}
\tightlist
\item
  \href{https://plugins.jetbrains.com/plugin/17718-github-copilot}{Github
  Copilot}
\end{itemize}

\hypertarget{beskrivning}{%
\subsection{Beskrivning}\label{beskrivning}}

JetBrains IntelliJ IDEA är ett IDE för Java och andra
programmeringsspråk. Den erbjuder både en gratisversion (Community
Edition) och en betald version (Ultimate Edition) med fler funktioner.
Det är ett mycket populärt val bland Java-utvecklare tack vare dess
omfattande funktioner och användarvänliga gränssnitt.

Här är den begärda artikeln, omformulerad för NetBeans:

\hypertarget{netbeans-installation}{%
\section{NetBeans installation}\label{netbeans-installation}}

\href{https://netbeans.apache.org/download/index.html}{Installera
NetBeans}\{: .btn \}

\hypertarget{beskrivning}{%
\subsection{Beskrivning}\label{beskrivning}}

Apache NetBeans är ett gratis och öppen källkods IDE för Java och andra
programmeringsspråk som PHP, JavaScript, och C/C++. NetBeans erbjuder en
komplett uppsättning funktioner för professionell utveckling, inklusive
debugger, profiler, refactoringverktyg, statisk analys, versionhantering
och mycket mer. Dess användarvänliga gränssnitt och omfattande
hjälpmaterial gör det till ett bra val för både nybörjare och erfarna
utvecklare.

\hypertarget{rider-installation}{%
\section{Rider installation}\label{rider-installation}}

\href{https://www.jetbrains.com/rider/}{Installera Jetbrains Rider}\{:
.btn \}

\emph{Filen editerades senast 2022-12-09} \#\# Addons till Rider som jag
rekommenderar

\begin{itemize}
\tightlist
\item
  \href{https://plugins.jetbrains.com/plugin/17718-github-copilot}{Github
  Copilot}
\end{itemize}

\hypertarget{beskrivning}{%
\subsection{Beskrivning}\label{beskrivning}}

Jetbrains Rider är ett IDE för Java och .NET. Det är ett gratis program
som är öppen källkod. Den är som Visual Studio Code och den har
Jetbrains Resharper inbyggt.

\hypertarget{fuxf6rslag}{%
\section{Förslag}\label{fuxf6rslag}}

Här finns förslag på andra program som kan vara bra att använda under
studietiden.

\emph{Filen editerades senast 2022-12-09}

\hypertarget{evernote}{%
\section{Evernote}\label{evernote}}

\href{https://evernote.com/intl/sv/download/}{Installera Evernote}\{:
.btn .btn-blue \}

Evernote är ett program som kan användas för att skriva anteckningar,
skapa listor och göra anteckningar om bilder. Det är enkelt att dela med
sig av anteckningar till andra. Det finns även en mobilapp som gör det
enkelt att skriva anteckningar på språng. Det är gratis att skapa ett
konto och använda programmet. Men vill man ha mer utrymme eller fler
funktioner så får man skaffa ett betalkonto.
\href{https://evernote.com/intl/sv/students}{Studenter kan få 40 \%
rabatt på ett årsabonnemang på Evernote Personal.}. Det finns även en
\href{https://www.evernote.com/Login.action?targetUrl=\%2Fclient\%2Fweb}{webbversion}
som du kan använda om du inte vill installera programmet på din dator.

\emph{Filen editerades senast 2022-12-09}

Visst, här är en artikel om Github Student Developer Pack:

\hypertarget{githubstudent}{%
\section{Githubstudent}\label{githubstudent}}

Github Student Developer Pack är ett gratispaket som ger studenter
tillgång till ett brett utbud av verktyg och tjänster för att hjälpa dem
att lära sig och utveckla sina kodningskunskaper. Paketet inkluderar:

\begin{itemize}
\tightlist
\item
  Github Pro-konto: Detta ger studenter tillgång till premiumfunktioner
  som privat lagring, 2FA och 100 GB lagring.
\item
  GitHub Codespaces: Detta ger studenter möjlighet att utveckla kod i
  molnet utan att behöva installera något programvara på sin dator.
\item
  GitHub Sponsors: Detta ger studenter möjlighet att få sponsring från
  företag och individer för deras kodningsprojekt.
\item
  GitHub Education Cloud: Detta ger studenter tillgång till en mängd
  olika utbildningsresurser, inklusive kurser, handledningar och
  exempelkod.
\end{itemize}

Github Student Developer Pack är ett utmärkt sätt för studenter att lära
sig och utveckla sina kodningskunskaper. Paketet är gratis och
innehåller en mängd olika verktyg och tjänster som kan hjälpa studenter
att bli bättre kodare.

\hypertarget{fuxf6rdelar-med-github-student-developer-pack}{%
\subsection{Fördelar med Github Student Developer
Pack}\label{fuxf6rdelar-med-github-student-developer-pack}}

Det finns många fördelar med att använda Github Student Developer Pack.
Några av fördelarna inkluderar:

\begin{itemize}
\tightlist
\item
  Gratis: Github Student Developer Pack är helt gratis, vilket gör det
  till ett mycket prisvärt erbjudande för studenter.
\item
  Bredt utbud av verktyg och tjänster: Github Student Developer Pack
  innehåller ett brett utbud av verktyg och tjänster som kan hjälpa
  studenter att lära sig och utveckla sina kodningskunskaper.
\item
  Enkelt att använda: Github Student Developer Pack är mycket enkelt att
  använda, vilket gör det till ett bra val för studenter som är nya på
  kodning.
\item
  Stor community: Github har en stor community av användare och
  utvecklare, vilket gör det till ett bra ställe att lära sig och få
  hjälp med kodning.
\item
  Gratis Github Copilot
\end{itemize}

\hypertarget{hur-man-registrerar-sig-fuxf6r-github-student-developer-pack}{%
\subsection{Hur man registrerar sig för Github Student Developer
Pack}\label{hur-man-registrerar-sig-fuxf6r-github-student-developer-pack}}

För att registrera dig för Github Student Developer Pack måste du:

Du får mailadress från skolan, använd den för registrering och koppla
sedan ditt privata konto till den du registrerat för Github med skolans
konto.

\begin{enumerate}
\def\labelenumi{\arabic{enumi}.}
\tightlist
\item
  Ha en giltig studentmejladress.
\item
  Besök Githubs webbplats och klicka på ``Registrera dig för
  studentkonto''.
\item
  Ange din studentmejladress och lösenord.
\item
  Bekräfta din registrering genom att klicka på länken i det
  e-postmeddelande som du får från Github.
\item
  En gång du har registrerat dig kan du börja använda Github Student
  Developer Pack.
\end{enumerate}

Om det inte fungerar med mailadressen behöver du använda Studieintyg
från skolan, det är bara att fråga om det så mailar vi över den.

Alla Campus Mölndal YH studerande har tillgång till Github Student
Developer Pack.

\hypertarget{google-keep}{%
\section{Google Keep}\label{google-keep}}

\href{https://keep.google.com/}{Installera Google Keep}\{: .btn \}

\emph{Filen editerades senast 2022-12-09} \#\# Beskrivning

Google Keep är ett program som är gjort av Google. Man använder det till
att spara anteckningar. Det är och öppen källkod. Det är ungefär som att
ha en tavla med post-it lappar.

\hypertarget{miro-board}{%
\section{Miro Board}\label{miro-board}}

\href{https://miro.com/}{Öppna Miro Board}\{: .btn \}

\hypertarget{beskrivning}{%
\subsection{Beskrivning}\label{beskrivning}}

Miro är en verktyg för att skapa och dela idéer. Det är ett enkelt sätt
att samarbeta med andra. Det är enkelt att skapa olika typer av
innehåll, som till exempel flödesscheman, wireframes, moodboards,
mindmaps, backlog och mycket mer.

Det finns också en mobilapp som gör det enkelt att skriva anteckningar
när man är på språng.

Miro har en presentationsläge som kan användas för att presentera
skoluppgifter eller projekt.

\hypertarget{notion}{%
\section{Notion}\label{notion}}

\href{https://www.notion.so/}{Öppna Notion}\{: .btn .btn-blue \}

Notion är en digital anteckningsbok. Du kan skriva anteckningar, skapa
listor och göra anteckningar om bilder. Det är enkelt att dela med sig
av anteckningar till andra. Det finns även en mobilapp som gör det
enkelt att skriva anteckningar på språng. Vill man ha mer funktionalitet
kan man prenumerera på deras tjänst, men den är
\href{https://www.notion.so/product/notion-for-education}{gratis som pro
för studenter}.

\emph{Filen editerades senast 2022-12-09}

Visst, här är en artikel om AI-hjälpmedel som är tillgängliga för
studerande:

\hypertarget{ai}{%
\section{ai}\label{ai}}

Artificiell intelligens (AI) är ett brett fält inom datavetenskap som
handlar om att skapa intelligenta agenter. Intelligenta agenter är
system som kan tänka, lära sig och agera självständigt. AI används inom
en mängd olika områden, bland annat medicin, finans, transport och
utbildning.

Det vi ser mest av nu är Generativ AI, alltså AI som kan skapa saker.
Det kan vara allt från att skapa musik, bilder, texter och mycket mer.

Den kan skapa ny information från befintlig information. Den kan göra
detta genom att lära sig mönster och strukturer i data. När AI har lärt
sig dessa mönster kan den använda dem för att skapa nya exempel som
liknar de ursprungliga exemplen.

Ett populärt exempel på AI som skapar ny information är GANs. GANs
består av två neurala nätverk som arbetar tillsammans. Ett nätverk
(generatorn) skapar nya data, medan det andra nätverket (diskriminatorn)
försöker avgöra om de nya datan är äkta eller inte. Genom denna process
lär sig generatorn att skapa allt mer realistiska data.

AI som skapar ny information är ett kraftfullt verktyg som har potential
att revolutionera många olika branscher. Det är fortfarande under
utveckling, men det har redan börjat användas i en mängd olika
tillämpningar.

\hypertarget{anvuxe4ndning}{%
\subsection{Användning}\label{anvuxe4ndning}}

Det finns många olika sätt att använda AI i utbildning. AI kan användas
för att:

\begin{itemize}
\tightlist
\item
  Ge personligt stöd till studerande
\item
  Automatisera uppgifter
\item
  Skapa nya läromedel
\item
  Öka studerandenas engagemang
\item
  Förbättra studerandenas lärande
\end{itemize}

\hypertarget{begruxe4nsningar}{%
\subsection{Begränsningar}\label{begruxe4nsningar}}

AI är fortfarande under utveckling, och det finns vissa begränsningar
med AI-hjälpmedel. Till exempel kan AI-hjälpmedel vara dyra, och de kan
inte alltid förstå alla studerandes behov. Det är viktigt att vara
medveten om dessa begränsningar när man använder AI-hjälpmedel i
utbildning.

\hypertarget{att-tuxe4nka-puxe5-som-studerande}{%
\subsection{Att tänka på som
studerande}\label{att-tuxe4nka-puxe5-som-studerande}}

Om du är studerande är det viktigt att tänka på följande när du använder
AI-hjälpmedel:

\begin{itemize}
\tightlist
\item
  Använd AI-hjälpmedel för att stödja ditt lärande, inte för att ersätta
  det.
\item
  Var medveten om begränsningarna med AI-hjälpmedel.
\item
  Använd AI-hjälpmedel för att lära dig på ett sätt som är engagerande
  och meningsfullt för dig.
\end{itemize}

AI är ett kraftfullt verktyg som kan hjälpa studerande att lära sig mer
effektivt. Genom att vara medveten om AI:s begränsningar och fördelar
kan studerande använda AI-hjälpmedel för att förbättra sitt lärande.

\hypertarget{jetbrains-tools}{%
\section{Jetbrains-tools}\label{jetbrains-tools}}

JetBrains är en utvecklingsplattform som erbjuder ett brett utbud av
verktyg för utvecklare. Dessa verktyg är utformade för att hjälpa
utvecklare att vara mer produktiva och effektiva, och de kan användas
för att utveckla programvara i en mängd olika språk och plattformar.

JetBrains-verktyg är tillgängliga för både studenter och företag.
Studenter kan få en gratis licens för JetBrains-verktyg genom att
registrera sig för en JetBrains Education-konto.

Här är några av de mest populära JetBrains-verktygen:

\begin{itemize}
\tightlist
\item
  IntelliJ IDEA: En fullfjädrad integrerad utvecklingsmiljö (IDE) för
  Java, Kotlin, Scala och andra språk.
\item
  PyCharm: En fullfjädrad IDE för Python.
\item
  WebStorm: En fullfjädrad IDE för HTML, CSS och JavaScript.
\item
  CLion: En fullfjädrad IDE för C och C++.
\item
  Rider: En fullfjädrad IDE för .NET.
\item
  DataGrip: En fullfjädrad IDE för SQL.
\item
  GoLand: En fullfjädrad IDE för Go.
\item
  RubyMine: En fullfjädrad IDE för Ruby.
\item
  AppCode: En fullfjädrad IDE för Swift.
\item
  PhpStorm: En fullfjädrad IDE för PHP.
\end{itemize}

JetBrains-verktyg erbjuder en mängd olika funktioner som kan hjälpa
utvecklare att vara mer produktiva och effektiva. Dessa funktioner
inkluderar:

\begin{itemize}
\tightlist
\item
  Syntaxfärgning och autoinmatning
\item
  Felkontroll och intellisense
\item
  Kodförbättring och omstrukturering
\item
  Modulhantering och projekthantering
\item
  Testning och debuggning
\item
  Dokumentation och dokumentation
\end{itemize}

JetBrains-verktyg är ett kraftfullt verktyg som kan hjälpa utvecklare
att vara mer produktiva och effektiva. De är tillgängliga för både
studenter och företag, och de kan användas för att utveckla programvara
i en mängd olika språk och plattformar.

\hypertarget{nyttan-av-jetbrains-verktyg}{%
\subsection{Nyttan av
JetBrains-verktyg}\label{nyttan-av-jetbrains-verktyg}}

Det finns många fördelar med att använda JetBrains-verktyg. Några av de
mest uppenbara fördelarna är:

\begin{itemize}
\tightlist
\item
  De kan hjälpa dig att skriva bättre kod snabbare.
\item
  De kan hjälpa dig att hitta och fixa fel i din kod.
\item
  De kan hjälpa dig att organisera och strukturera din kod.
\item
  De kan hjälpa dig att testa din kod.
\item
  De kan hjälpa dig att dokumentera din kod.
\end{itemize}

\hypertarget{begruxe4nsningar-av-jetbrains-verktyg}{%
\subsection{Begränsningar av
JetBrains-verktyg}\label{begruxe4nsningar-av-jetbrains-verktyg}}

Det finns några få begränsningar med JetBrains-verktyg. Några av de mest
uppenbara begränsningarna är:

\begin{itemize}
\tightlist
\item
  De kan vara dyra.
\item
  De kan vara komplexa att använda.
\item
  De kan inte användas för alla språk och plattformar.
\end{itemize}

\hypertarget{att-tuxe4nka-puxe5-som-studerande}{%
\subsection{Att tänka på som
studerande}\label{att-tuxe4nka-puxe5-som-studerande}}

Om du är en studerande som funderar på att använda JetBrains-verktyg är
det några saker du bör tänka på. För det första är JetBrains-verktyg
dyra. Du kan dock få en gratis licens genom att registrera dig för en
JetBrains Education-konto. För det andra kan JetBrains-verktyg vara
komplexa att använda. Om du är nybörjare kan du behöva ta en kurs eller
läsa en manual för att lära dig hur du använder dem. För det tredje kan
JetBrains-verktyg inte användas för alla språk och plattformar. Om du
utvecklar programvara i ett språk eller på en plattform som inte stöds
av JetBrains-verktyg kan du behöva använda ett annat verktyg.

\hypertarget{tips}{%
\subsection{Tips}\label{tips}}

Om du bestämmer dig för att använda JetBrains-verktyg är det några tips
som kan hjälpa dig att få ut det mesta av dem. För det första, ta dig
tid att lära dig hur du använder verktygen. Det finns många resurser
tillgängliga online och i böcker som kan hjälpa dig att lära dig hur du
använder JetBrains-verktyg. För det andra, använd verktygen för att
automatisera dina uppgifter. JetBrains-verktyg kan hjälpa dig att
automatisera många av de uppgifter som du gör när du utvecklar
programvara. För det tredje, använd verktygen för att förbättra din kod.
JetBrains-verktyg kan hjälpa dig att hitta och fixa fel i din kod, och
de kan också hjälpa dig att förbättra kvaliteten på din kod.

\hypertarget{microsoft-tools}{%
\section{Microsoft-tools}\label{microsoft-tools}}

Microsoft erbjuder två olika IDE:er (integrerade utvecklingsmiljöer) för
programutveckling: Visual Studio Community och Visual Studio Code.

\begin{itemize}
\tightlist
\item
  \textbf{Visual Studio Community} är ett fullfjädrad IDE som främst
  används för .NET-utveckling, men det kan också användas för att
  utveckla i andra språk, till exempel Python och C++.
\item
  \textbf{Visual Studio Code} är en lättviktig texteditor som kan
  anpassas för att användas med en mängd olika språk, inklusive .NET,
  Python, JavaScript, TypeScript, C++, C\#, Java, PHP, Go och Rust.
\item
  SQL Server Management Studio (SSMS) är en IDE för SQL Server. Den kan
  användas för att utveckla och administrera SQL Server-databaser.
\end{itemize}

Båda IDE:erna är gratis att använda och är tillgängliga för Windows,
macOS och Linux.

\hypertarget{fuxf6rdelar-med-visual-studio-community-och-visual-studio-code}{%
\subsection{Fördelar med Visual Studio Community och Visual Studio
Code}\label{fuxf6rdelar-med-visual-studio-community-och-visual-studio-code}}

\begin{itemize}
\tightlist
\item
  Båda IDE:erna erbjuder en mängd olika funktioner som kan hjälpa
  utvecklare att vara mer produktiva och effektiva, till exempel
  syntaxfärgning, autoinmatning, felkontroll, kodförbättring och
  omstrukturering.
\item
  Båda IDE:erna har en stor community av användare och utvecklare som
  kan ge stöd och hjälp.
\item
  Båda IDE:erna är uppdaterade med de senaste funktionerna och
  buggfixarna.
\item
  SQL Server Management Studio (SSMS) hjälper till hanteringen av
  scheman, tabeller, index, vyer, lagrade procedurer, funktioner,
  triggers, användare, roller, säkerhet, säkerhetskopior,
  återställningar, replikering, loggar, jobb, schemaläggning,
  övervakning, underhållsplaner, databasdiagram mm.
\end{itemize}

\hypertarget{begruxe4nsningar-av-visual-studio-community-och-visual-studio-code}{%
\subsection{Begränsningar av Visual Studio Community och Visual Studio
Code}\label{begruxe4nsningar-av-visual-studio-community-och-visual-studio-code}}

\begin{itemize}
\tightlist
\item
  Visual Studio Community är inte lika kraftfullt som Visual Studio
  Professional eller Enterprise, och det saknas vissa funktioner som är
  tillgängliga i de dyrare versionerna.
\item
  Visual Studio Code är en lättviktig texteditor, och det saknar vissa
  funktioner som är tillgängliga i fullfjädrade IDE:er, till exempel
  intellisense och kodförbättring.
\item
  SQL Server Management Studio (SSMS) är inte tillgängligt för macOS
  eller Linux.
\item
  SQL Server äter upp mycket minne och kan vara långsamt. Så det är
  viktigt att ha en bra dator för att kunna köra det. SSMS är dock
  snäll.
\end{itemize}

\hypertarget{att-tuxe4nka-puxe5-som-studerande}{%
\subsection{Att tänka på som
studerande}\label{att-tuxe4nka-puxe5-som-studerande}}

Om du är en studerande som funderar på att använda Visual Studio
Community eller Visual Studio Code är det några saker du bör tänka på.

\begin{itemize}
\tightlist
\item
  Visual Studio Community är ett bra val för studenter som utvecklar i
  .NET.
\item
  Visual Studio Code är ett bra val för studenter som utvecklar i andra
  språk, till exempel Python, JavaScript eller TypeScript.
\item
  Det är bra att lära sig att använda båda IDE:erna, så att du kan välja
  det bästa verktyget för jobbet.
\item
  Det är bra att lära sig att skriva kod utan hjälpmedel, så att du inte
  blir beroende av IDE:erna.
\end{itemize}

\hypertarget{tips}{%
\subsection{Tips}\label{tips}}

Om du bestämmer dig för att använda Visual Studio Community eller Visual
Studio Code är det några tips som kan hjälpa dig att få ut det mesta av
dem.

\begin{itemize}
\tightlist
\item
  Ta dig tid att lära dig hur du använder verktygen. Det finns många
  resurser tillgängliga online och i böcker.
\item
  Använd verktygen för att automatisera dina uppgifter. Visual Studio
  Community och Visual Studio Code kan hjälpa dig att automatisera många
  av de uppgifter som du gör när du utvecklar programvara.
\item
  Använd verktygen för att förbättra din kod. Visual Studio Community
  och Visual Studio Code kan hjälpa dig att hitta och fixa fel i din
  kod, och de kan också hjälpa dig att förbättra kvaliteten på din kod.
\end{itemize}

\hypertarget{github-copilot}{%
\section{Github Copilot}\label{github-copilot}}

Github Copilot är en AI-assistent för programmerare som är utvecklad av
OpenAI och GitHub. Den är baserad på GPT-3, OpenAIs stora språkmodell,
och kan hjälpa programmerare att skriva kod snabbare och mer effektivt.

\hypertarget{anvuxe4ndning}{%
\subsection{Användning}\label{anvuxe4ndning}}

Github Copilot kan användas för att skriva alla typer av kod, inklusive
Python, Java, C++, och JavaScript. Den kan användas för att skriva ny
kod, eller för att förbättra befintlig kod.

Github Copilot fungerar genom att generera förslag på kod baserat på den
kod som du skriver. Du kan sedan acceptera eller avslå förslagen, eller
redigera dem efter behov.

\hypertarget{exempel-puxe5-prompts}{%
\subsection{Exempel på prompts}\label{exempel-puxe5-prompts}}

Här är några exempel på prompts som du kan använda med Github Copilot:

\begin{itemize}
\tightlist
\item
  // Skriv en funktion som returnerar summan av två tal.
\item
  // Skriv en klass som representerar en lista.
\item
  // Skriv en funktion som skriver ut alla tal från 1 till 10.
\end{itemize}

Efter att ha skrivit in prompten trycker du ENTER och väntar på att den
ska bli klar

\hypertarget{exempel}{%
\subsubsection{Exempel}\label{exempel}}

\begin{verbatim}
// Skapa en metod som rekursivt räknar ut Fibonacci-talen för ett givet tal.
\end{verbatim}

\textbf{Tryck enter}

\hypertarget{cb2}{}
\begin{Shaded}
\begin{Highlighting}[]
\KeywordTok{public} \DataTypeTok{static} \DataTypeTok{int} \FunctionTok{fibonacci}\OperatorTok{(}\DataTypeTok{int}\NormalTok{ n}\OperatorTok{)} \OperatorTok{\{}
    \ControlFlowTok{if} \OperatorTok{(}\NormalTok{n }\OperatorTok{\textless{}=} \DecValTok{1}\OperatorTok{)} \OperatorTok{\{}
        \ControlFlowTok{return}\NormalTok{ n}\OperatorTok{;}
    \OperatorTok{\}}
    \ControlFlowTok{return} \FunctionTok{fibonacci}\OperatorTok{(}\NormalTok{n }\OperatorTok{{-}} \DecValTok{1}\OperatorTok{)} \OperatorTok{+} \FunctionTok{fibonacci}\OperatorTok{(}\NormalTok{n }\OperatorTok{{-}} \DecValTok{2}\OperatorTok{);}
\OperatorTok{\}}
\end{Highlighting}
\end{Shaded}

Glöm inte att kommentera vart du har fått koden ifrån. I detta fall är
det från Github Copilot.

\hypertarget{cb3}{}
\begin{Shaded}
\begin{Highlighting}[]
\CommentTok{// Genererad av Github Copilot}
\end{Highlighting}
\end{Shaded}

\hypertarget{nyttan-av-det}{%
\subsection{Nyttan av det}\label{nyttan-av-det}}

Github Copilot kan hjälpa dig att skriva kod snabbare och mer effektivt.
Den kan också hjälpa dig att lära dig nya språk och ramverk, eftersom
den kan generera kod för dig att läsa och lära dig av. Det är viktigt
att du förstår den kod som den genererar, så att du kan lära dig av den.

\hypertarget{begruxe4nsningar}{%
\subsection{Begränsningar}\label{begruxe4nsningar}}

Github Copilot är fortfarande under utveckling, så den har vissa
begränsningar. Till exempel kan den inte förstå all kod, och den kan
ibland generera felaktig kod.

Github Copilot kostar per månad, så du bör inte använda den om du inte
har råd med den.

\hypertarget{att-tuxe4nka-puxe5-som-studerande}{%
\subsection{Att tänka på som
studerande}\label{att-tuxe4nka-puxe5-som-studerande}}

Om du är studerande är det viktigt att tänka på följande när du använder
Github Copilot:

\begin{itemize}
\tightlist
\item
  Ange alltid att du har använt Github Copilot när du lämnar in kod.
\item
  Kontrollera alltid koden som Github Copilot genererar noggrant för att
  se till att den är korrekt.
\item
  Använd inte Github Copilot för att fuska.
\item
  Github Copilot är gratis för studerande.
\end{itemize}

\hypertarget{plagiering-fusk}{%
\subsubsection{Plagiering / Fusk}\label{plagiering-fusk}}

Plagiering är att kopiera någon annans arbete utan att ange källan. Fusk
är att försöka få ett högre betyg än du har förtjänat genom att fuska.

Om du kopierar en klasskamrats kod utan att ange källan, så har du
plagierat. Om du kopierar en klasskamrats kod och anger källan, så har
du inte plagierat.

Om du fuska, så kommer du att få en lägre betyg eller bli utesluten från
kursen.

Det är viktigt att vara ärlig och att inte fuska. Om du har frågor om
plagiering eller fusk, så prata med din lärare.

\hypertarget{tabnine}{%
\section{Tabnine}\label{tabnine}}

Tabnine är en AI-baserad kodförslagstjänst som hjälper utvecklare att
skriva bättre kod snabbare. Den fungerar genom att analysera din kod och
sedan föreslå förbättringar, till exempel nya metoder, funktioner och
variabler. Tabnine är tillgänglig för en mängd olika språk, inklusive
Python, Java, C++ och JavaScript.

\hypertarget{anvuxe4ndning}{%
\subsection{Användning}\label{anvuxe4ndning}}

Tabnine kan användas på en mängd olika sätt. Du kan använda den för att:

\begin{itemize}
\tightlist
\item
  Skriva ny kod
\item
  Förbättra befintlig kod
\item
  Hitta och fixa fel
\item
  Lära dig nya kodningstekniker
\end{itemize}

\hypertarget{begruxe4nsningar}{%
\subsection{Begränsningar}\label{begruxe4nsningar}}

Tabnine är fortfarande under utveckling, så det finns några
begränsningar. Till exempel kan den inte förstå alla typer av kod, och
den kan ibland generera felaktig kod.

\hypertarget{att-tuxe4nka-puxe5-som-studerande}{%
\subsection{Att tänka på som
studerande}\label{att-tuxe4nka-puxe5-som-studerande}}

Om du är en studerande som funderar på att använda Tabnine är det några
saker du bör tänka på.

\begin{itemize}
\tightlist
\item
  Tabnine är ett verktyg som kan hjälpa dig att skriva bättre kod, men
  det är inte ett substitut för att lära dig att kod. Du bör fortfarande
  lära dig de grundläggande begreppen i kodning och hur du skriver
  korrekt kod.
\item
  Tabnine kan generera felaktig kod, så du bör alltid kontrollera koden
  som den genererar innan du använder den.
\item
  Du bör inte använda Tabnine för att fuska. Om du använder Tabnine för
  att generera kod som du inte förstår, kommer du inte att lära dig
  något och du kommer inte att kunna lösa problem som kräver att du
  förstår koden.
\item
  Tabnine är gratis men begränsad, den bättre versionen kostar, så du
  bör inte använda den om du inte har råd med den.
\end{itemize}

\hypertarget{tips}{%
\subsection{Tips}\label{tips}}

Om du bestämmer dig för att använda Tabnine är det några tips som kan
hjälpa dig att få ut det mesta av det.

\begin{itemize}
\tightlist
\item
  Lär dig hur man använder Tabnines inställningar för att anpassa det
  till dina behov.
\item
  Använd Tabnine för att få feedback på din kod, men använd inte den som
  en ersättning för att lära dig att kod.
\item
  Kontrollera alltid koden som Tabnine genererar innan du använder den.
\end{itemize}

\hypertarget{chatgpt}{%
\section{Chatgpt}\label{chatgpt}}

\hypertarget{vad-uxe4r-chatgpt}{%
\subsection{Vad är ChatGPT?}\label{vad-uxe4r-chatgpt}}

ChatGPT är en avancerad språkmodell utvecklad av OpenAI. Den bygger på
den banbrytande GPT-3.5-arkitekturen och är utformad för att kunna
generera text som svar på användargivna prompts. Modellen är tränad på
en enorm mängd textdata från olika källor och kan svara på frågor, ge
förklaringar och till och med skapa dialoger. ChatGPT har blivit en
användbar resurs för elever, forskare, författare och många andra som
behöver hjälp med att generera text av hög kvalitet.

\hypertarget{exempel-puxe5-prompts}{%
\subsection{Exempel på prompts}\label{exempel-puxe5-prompts}}

Prompts är de instruktioner eller frågor du ger till ChatGPT för att få
ett svar eller en textgenerering. Genom att formulera väldefinierade och
tydliga prompts kan du få bättre resultat från modellen. Här är några
exempel på prompts som du kan använda:

\begin{enumerate}
\def\labelenumi{\arabic{enumi}.}
\tightlist
\item
  ``Kan du förklara konceptet {[}lägg till begrepp eller ämne här{]} på
  ett enkelt sätt?''
\item
  ``Jag behöver hjälp med att förstå {[}specificera ämnet{]}. Kan du ge
  mig en grundläggande förklaring?''
\item
  ``Kan du ge mig några exempel på {[}specificera ämnet{]}?''
\item
  ``Vad är skillnaden mellan {[}specificera två begrepp{]}?''
\item
  ``Kan du ge mig några tips för att förbättra mina skrivfärdigheter?''
\item
  ``Jag har ett fel i min kod som jag inte kan hitta, kan du förklara
  vad som är fel?'' (klistra in koden i raden under)
\item
  ``Jag ska skapa ett projekt om {[}specificera ämnet{]}, kan du ge mig
  några idéer?''
\item
  ``Jag ska skapa ett projekt om {[}specificera ämnet{]}, kan du hjälpa
  mig att planera, steg för steg?''
\item
  ``Jag ska skapa ett projekt om {[}specificera ämnet{]}, kan du hjälpa
  mig att göra en utför backlog?''
\item
  ``Jag fick detta felmeddelande när jag körde min kod: {[}klistra in
  felmeddelandet här{]}. Vad betyder det?''
\end{enumerate}

Självklart kan du lägga till följdfrågor eller förtydliganden för att få
mer specifika svar.

\begin{enumerate}
\def\labelenumi{\arabic{enumi}.}
\tightlist
\item
  ``Kan du göra det mer utförligt?''
\item
  ``Kan du ge mig ett exempel?''
\item
  ``Kan du ge mig ett annat exempel?''
\item
  ``Kan du ge mig ett exempel som är mer relevant för {[}specificera
  ämnet{]}?''
\end{enumerate}

Dessa prompts kan användas som en startpunkt för att interagera med
ChatGPT och få information eller stöd för dina studier. Kom ihåg att
vara så specifik som möjligt i dina prompts för att få relevanta och
användbara svar.

Jag hoppas att dessa exempel på prompts hjälper dig att använda ChatGPT
mer effektivt som en studerande. Låt mig veta om det finns något annat
du behöver hjälp med!

\hypertarget{anvuxe4ndning}{%
\subsection{Användning}\label{anvuxe4ndning}}

ChatGPT kan användas på olika sätt för att hjälpa och underlätta
skrivprocessen. Här är några exempel på hur du kan dra nytta av ChatGPT:

\hypertarget{fuxe5-inspiration-till-nya-projekt}{%
\subsubsection{1. Få inspiration till nya
projekt:}\label{fuxe5-inspiration-till-nya-projekt}}

Du kan använda ChatGPT för att generera idéer och få inspiration när du
står inför skrivblockering eller behöver utforska nya ämnen. Genom att
ge en prompt eller en beskrivning kan du få förslag och riktlinjer för
att komma igång med ditt skrivande.

\hypertarget{anvuxe4nd-chatgpt-som-en-extra-luxe4rare}{%
\subsubsection{2. Använd ChatGPT som en extra
lärare:}\label{anvuxe4nd-chatgpt-som-en-extra-luxe4rare}}

Om du stöter på svåra begrepp eller inte förstår ett ämne kan ChatGPT
vara till hjälp. Du kan ställa frågor eller ge en kort förklaring till
modellen, och den kan generera en mer utförlig förklaring eller exempel
för att hjälpa till med din förståelse. Tänk på att även om ChatGPT kan
ge användbara svar, är det alltid bra att bekräfta informationen genom
andra källor.

\hypertarget{hantera-kodgenerering}{%
\subsubsection{3. Hantera kodgenerering:}\label{hantera-kodgenerering}}

ChatGPT kan också hjälpa till med kodgenerering genom att ge exempel
eller lösningsförslag för specifika programmeringsproblem. Det kan vara
användbart när du fastnar eller behöver hjälp med att implementera vissa
funktioner. Om du väljer att använda kod som genererats av ChatGPT, kom
ihåg att lägga till en kommentar som tydligt anger att koden är
genererad av modellen.

\hypertarget{begruxe4nsningar}{%
\subsection{Begränsningar}\label{begruxe4nsningar}}

Det är viktigt att vara medveten om begränsningarna med ChatGPT och att
använda det som ett verktyg med viss försiktighet:

\begin{itemize}
\item
  \textbf{Brist på kontextförståelse}: ChatGPT saknar kontextförståelse
  över olika frågor. Det innebär att det kan generera svar som kan verka
  rimliga men faktiskt vara felaktiga eller missvisande. Var alltid
  medveten om detta och verifiera information från andra tillförlitliga
  källor.
\item
  \textbf{Risk för generering av opålitlig information}: Eftersom
  modellen tränas på en mängd olika källor kan den även generera svar
  som är felaktiga, opålitliga eller kontroversiella. Var kritisk när du
  bedömer svaren och använd din egen bedömning för att avgöra om
  informationen är korrekt och pålitlig.
\item
  \textbf{Brist på mänskligt omdöme}: ChatGPT saknar förmågan att utöva
  mänskligt omdöme och kan därför inte bedöma eller anpassa sig till
  etiska eller moraliska aspekter av skrivandet. Var medveten om detta
  och använd din egen bedömning för att avgöra vad som är lämpligt och
  passande att använda från modellen.
\end{itemize}

\hypertarget{att-tuxe4nka-puxe5-som-studerande}{%
\subsection{Att tänka på som
studerande}\label{att-tuxe4nka-puxe5-som-studerande}}

När du använder ChatGPT som elev är det viktigt att vara medveten om
akademiska integritetsprinciper och undvika fusk eller plagiering. Här
är några riktlinjer att följa:

\begin{itemize}
\item
  \textbf{Använd ChatGPT för att få inspiration till nya projekt}:
  Modellen kan hjälpa dig att generera idéer och ge riktlinjer för ditt
  eget arbete. Se till att använda informationen som en utgångspunkt och
  utveckla den vidare på ett originellt sätt.
\item
  \textbf{Använd ChatGPT som en extra lärare}: Du kan använda modellen
  för att få förklaringar på saker du inte förstår, men det är viktigt
  att alltid komplettera med egen forskning och kontrollera
  informationen från andra tillförlitliga källor.
\item
  \textbf{Hantera kodgenerering med försiktighet}: Om du använder
  ChatGPT för att generera kod, kom ihåg att alltid ange att koden är
  genererad av modellen. Att kopiera en hel uppgift eller en
  klasskamrats kod utan att ange källan är inte tillåtet och anses vara
  plagiering eller fusk. Det är viktigt att visa att du har förstått
  uppgiften genom att bearbeta och anpassa den genererade koden på ett
  meningsfullt sätt.
\item
  \textbf{Var ärlig och respektera akademiska riktlinjer}: Att vara
  ärlig och respektera akademiska integritetsprinciper är avgörande för
  din personliga och professionella utveckling. Var noga med att följa
  dina lärares riktlinjer och ta ansvar för ditt eget arbete.
\end{itemize}

Kom ihåg att ChatGPT är ett verktyg som kan vara till hjälp, men det
ersätter inte din egen kreativitet, ansträngning och kritiska tänkande.
Genom att använda det på rätt sätt kan du dra nytta av dess potential
och förbättra din skrivförmåga.

\hypertarget{extra-prompts}{%
\subsection{Extra prompts}\label{extra-prompts}}

\textbf{Personlighet:}

``System prompt: Du ska nu agera som \{valfri personlighet\}'' + En
lärare i programmering som ska förklara för nybörjare + Darth Vader +
Elon Musk + Hund + Katt + Taylor Swift + Yoda + En svartsjuk flickvän /
pojkvän + En romantisk flickvän / pojkvän + Skåning / Göteborgare /
Stockholmare / Smålänning

\textbf{Förklaringar:} + ``Förklara \{valfritt ämne\} på ett enkelt
sätt'' + ``Förklara \{valfritt ämne\} ELI5'' + ``Förklara \{valfritt
ämne\} för en nybörjare''

\textbf{Prompt engineering:}

\begin{itemize}
\tightlist
\item
  ``Kan du hjälpa mig att skapa en Persona som passar mitt sätt att lära
  mig programmering, så att jag kan lära mig mer effektivt? Ställ frågor
  om mig och mina intressen, och skapa en personlighet som passar
  mig.''.
\end{itemize}

Efter att ha skrivit det kommer ChatGPT att ställa en massa frågor som
du ska svara på, ju mer detaljerat desto bättre. När du är klar med det,
så kommer den att skapa en personlighet som passar dig, och du kan
använda den för att ställa frågor till den. Skriv sedan
\textgreater\textgreater\textgreater{} System prompt: Du ska nu agera
som \{den personligheten\}. och klistra in hela personligheten den
skapade \textless\textless\textless{}

``\textgreater\textgreater\textgreater{} \textless\textless\textless''
är för att den ska fatta att allt hänger samman

Om du gillar personligheten, kopiera allt från
\textgreater\textgreater\textgreater{} till
\textless\textless\textless{} och spara det i notepadden så kan du
återanvända det i framtiden.

Du kan skapa personligheter till andra saker, det är upp till din
fantasi.

\hypertarget{luxe4nkar}{%
\subsection{Länkar}\label{luxe4nkar}}

\href{https://writesonic.com/blog/chatgpt-prompts/\#best-chatgpt-prompts-for-education}{215+
Chat prompts for ChatGPT}
\href{https://mpost.io/100-best-chatgpt-prompts-to-unleash-ais-potential/}{100
best ChatGPT prompts}
\href{https://www.euronews.com/next/2023/06/25/getting-the-most-out-of-chatgpt-these-are-the-most-useful-prompts-to-try-now}{Getting
the most out of ChatGPT}
\href{https://www.greataiprompts.com/chat-gpt/best-chat-gpt-prompts/}{120
best prompts for ChatGPT}

\hypertarget{bard}{%
\section{Bard}\label{bard}}

Bard är en stor språkmodell, även känd som en konversations-AI eller
chatbot, tränad för att vara informativ och omfattande. Den är tränad på
en enorm mängd textdata och kan kommunicera och generera
människoliknande text som svar på ett brett spektrum av uppmaningar och
frågor. Till exempel kan den ge sammanfattningar av faktaämnen eller
skapa berättelser.

\hypertarget{vad-uxe4r-bard}{%
\subsection{Vad är Bard?}\label{vad-uxe4r-bard}}

Bard är en stor språkmodell från Google AI, tränad på en enorm mängd
textdata. Den kan generera text, översätta språk, skriva olika typer av
kreativt innehåll och svara på dina frågor på ett informativt sätt. Bard
är fortfarande under utveckling, men den har lärt sig att utföra många
typer av uppgifter, inklusive

\begin{itemize}
\tightlist
\item
  Bard kommer att göra sitt bästa för att följa dina instruktioner och
  slutföra dina förfrågningar eftertänksamt.
\item
  Bard kommer att använda sin kunskap för att svara på dina frågor på
  ett omfattande och informativt sätt, även om de är öppna, utmanande
  eller konstiga.
\item
  Bard kommer att generera olika kreativa textformat, som dikter, kod,
  skript, musikstycken, e-post, brev, etc. Den kommer att göra mitt
  bästa för att uppfylla alla dina krav.
\end{itemize}

\hypertarget{anvuxe4ndning}{%
\subsection{Användning}\label{anvuxe4ndning}}

Bard kan användas för en mängd olika uppgifter, inklusive

\begin{itemize}
\tightlist
\item
  Att få hjälp med dina studier. Bard kan hjälpa dig att lära dig nya
  saker genom att ge dig sammanfattningar av faktaämnen eller genom att
  skapa berättelser som är relaterade till det du studerar.
\item
  Att skriva kreativt innehåll. Bard kan hjälpa dig att skriva dikter,
  kod, skript, musikstycken, e-post, brev, etc.
\item
  Att översätta språk. Bard kan översätta text från ett språk till ett
  annat.
\item
  Att svara på dina frågor. Bard kan svara på dina frågor på ett
  informativt sätt, även om de är öppna, utmanande eller konstiga.
\end{itemize}

\hypertarget{exempel-puxe5-prompts}{%
\subsection{Exempel på prompts}\label{exempel-puxe5-prompts}}

Bard kan svara på ett brett spektrum av uppmaningar och frågor. Här är
några exempel:

\begin{itemize}
\tightlist
\item
  Ge mig en sammanfattning av den franska revolutionen.
\item
  Skriv en dikt om kärlek.
\item
  Översätt den här meningen från engelska till spanska.
\item
  Varför är himlen blå?
\item
  Vad är meningen med livet?
\end{itemize}

\hypertarget{begruxe4nsningar}{%
\subsection{Begränsningar}\label{begruxe4nsningar}}

Bard är fortfarande under utveckling, så det finns vissa begränsningar
för vad den kan göra. Till exempel kan den inte alltid ge exakta eller
fullständiga svar på öppna, utmanande eller konstiga frågor. Dessutom är
Bard inte alltid kapabel att generera kreativt innehåll som är av hög
kvalitet.

\hypertarget{att-tuxe4nka-puxe5-som-studerande}{%
\subsection{Att tänka på som
studerande}\label{att-tuxe4nka-puxe5-som-studerande}}

Om du är studerande är det några saker du bör tänka på när du använder
Bard. För det första är det viktigt att komma ihåg att Bard inte är en
mänsklig lärare. Den kan inte ge dig samma nivå av personlig
uppmärksamhet och stöd som en mänsklig lärare kan. För det andra är det
viktigt att inte förlita dig för mycket på Bard. Det är viktigt att du
fortfarande lär dig och förstår det material du studerar. Slutligen är
det viktigt att använda Bard på ett etiskt sätt. Du bör inte använda
Bard för att fuska eller för att kopiera annans arbete.

\hypertarget{plagiering-fusk}{%
\subsubsection{Plagiering / Fusk}\label{plagiering-fusk}}

Plagiering är att kopiera annans arbete utan att ge kredit. Fusk är att
använda otillåtna hjälpmedel för att lösa en uppgift. Det är inte
tillåtet att fuska eller plagiera. Om du kopierar annans arbete utan att
ge kredit kommer du att bli disciplinerad. Om du använder otillåtna
hjälpmedel för att lösa en uppgift kommer du också att bli
disciplinerad.

\hypertarget{att-anvuxe4nda-bard-fuxf6r-att-fuxe5-inspiration-till-nya-projekt}{%
\subsubsection{Att använda Bard för att få inspiration till nya
projekt}\label{att-anvuxe4nda-bard-fuxf6r-att-fuxe5-inspiration-till-nya-projekt}}

Bard kan användas för att få inspiration till nya projekt. Till exempel

\begin{itemize}
\tightlist
\item
  \textbf{Fråga Bard om idéer.} Om du står inför en skrivstopp, eller om
  du letar efter nya idéer för ett projekt, kan du fråga Bard om hjälp.
  Bard kan komma med en mängd olika idéer, allt från berättelser till
  kod till presentationer.
\item
  \textbf{Be Bard om feedback.} När du har en idé kan du be Bard om
  feedback. Bard kan hjälpa dig att identifiera styrkor och svagheter i
  din idé, och den kan ge förslag på hur du kan förbättra den.
\item
  \textbf{Använd Bard som en brainstorming-partner.} Om du arbetar med
  ett team kan du använda Bard som en brainstorming-partner. Bard kan
  komma med nya idéer, och den kan hjälpa dig att lösa problem.
\end{itemize}

Bard är ett kraftfullt verktyg som kan hjälpa dig att få inspiration
till nya projekt. Genom att använda Bard kan du utöka dina idéer och
skapa något unikt och intressant.

Här är några ytterligare tips för att använda Bard för att få
inspiration till nya projekt:

\begin{itemize}
\tightlist
\item
  Var öppen för nya idéer. Bard kan komma med idéer som du aldrig skulle
  ha tänkt på själv.
\item
  Var inte rädd att experimentera. Bard kan hjälpa dig att testa nya
  idéer och se vad som fungerar.
\item
  Ha kul! Att få inspiration ska vara roligt. Om du inte har roligt
  kommer du inte att vara lika kreativ.
\end{itemize}

\hypertarget{att-anvuxe4nda-bard-som-en-extra-luxe4rare}{%
\subsubsection{Att använda Bard som en extra
lärare}\label{att-anvuxe4nda-bard-som-en-extra-luxe4rare}}

Bard kan användas som en extra lärare. Till exempel

\begin{itemize}
\item
  \textbf{Be Bard om hjälp med dina studier.} Om du stöter på svåra
  begrepp eller inte förstår ett ämne kan du be Bard om hjälp. Bard kan
  ge dig en sammanfattning av faktaämnen, och den kan ge dig exempel på
  hur du kan använda dina kunskaper.
\item
  ``Kan du förklara konceptet {[}lägg till begrepp eller ämne här{]} på
  ett enkelt sätt?''
\item
  ``Jag behöver hjälp med att förstå {[}specificera ämnet{]}. Kan du ge
  mig en grundläggande förklaring?''
\item
  ``Kan du ge mig några exempel på {[}specificera ämnet{]}?''
\item
  ``Vad är skillnaden mellan {[}specificera två begrepp{]}?''
\item
  ``Kan du ge mig några tips för att förbättra mina skrivfärdigheter?''
\item
  ``Jag har ett fel i min kod som jag inte kan hitta, kan du förklara
  vad som är fel?'' (klistra in koden i raden under)
\item
  ``Jag ska skapa ett projekt om {[}specificera ämnet{]}, kan du ge mig
  några idéer?''
\item
  ``Jag ska skapa ett projekt om {[}specificera ämnet{]}, kan du hjälpa
  mig att planera, steg för steg?''
\item
  ``Jag ska skapa ett projekt om {[}specificera ämnet{]}, kan du hjälpa
  mig att göra en utför backlog?''
\item
  ``Jag fick detta felmeddelande när jag körde min kod: {[}klistra in
  felmeddelandet här{]}. Vad betyder det?''
\end{itemize}

Självklart kan du lägga till följdfrågor eller förtydliganden för att få
mer specifika svar.

\begin{itemize}
\tightlist
\item
  ``Kan du göra det mer utförligt?''
\item
  ``Kan du ge mig ett exempel?''
\item
  ``Kan du ge mig ett annat exempel?''
\item
  ``Kan du ge mig ett exempel som är mer relevant för {[}specificera
  ämnet{]}?''
\end{itemize}

Dessa prompts kan användas som en startpunkt för att interagera med Bard
och få information eller stöd för dina studier. Kom ihåg att vara så
specifik som möjligt i dina prompts för att få relevanta och användbara
svar.

\hypertarget{tips}{%
\subsection{Tips}\label{tips}}

Här är några tips för att använda Bard mer effektivt som en studerande:

\begin{itemize}
\tightlist
\item
  Var specifik. Ju mer specifik du är, desto bättre svar kommer du att
  få.
\item
  Var tydlig. Ju tydligare du är, desto bättre svar kommer du att få.
\item
  Var tålmodig. Ibland kan det ta ett tag för Bard att generera ett
  svar.
\item
  Var kreativ. Bard kan hjälpa dig att tänka utanför boxen.
\item
  Var inte rädd för att experimentera. Bard kan hjälpa dig att testa nya
  idéer och se vad som fungerar.
\item
  Var inte rädd för att misslyckas. Bard kan hjälpa dig att lära dig av
  dina misstag.
\item
  Var inte rädd för att fråga. Bard kan hjälpa dig att lära dig mer om
  ett ämne.
\item
  Var inte rädd för att be om hjälp. Bard kan hjälpa dig att lära dig
  mer om ett ämne.
\item
  Var inte rädd för att be om feedback. Bard kan hjälpa dig att lära dig
  mer om ett ämne.
\item
  Var inte rädd för att be om förslag. Bard kan hjälpa dig att lära dig
  mer om ett ämne.
\end{itemize}

\textbf{Personligheter}

Bard kan anta olika personligheter. Här är några exempel:

\begin{itemize}
\tightlist
\item
  ``Du ska nu agera som en lärare för en studerande som behöver hjälp
  med att förstå ett ämne. Tänk på att förklara ELI5 och använda enkla
  ord.''
\item
  ``Du ska nu agera som en politiker, var så diffus som möjligt och
  använd många ord för att säga ingenting.''
\item
  ``Du ska nu agera som en poet, var så poetisk som möjligt och använd
  många metaforer.''
\item
  ``Du ska nu agera som en programmerare, var så teknisk som möjligt och
  använd många tekniska termer.''
\item
  ``Du ska nu agera som en diplomat, var så diplomatisk som möjligt och
  använd många ord för att säga ingenting.''
\item
  ``Du ska nu agera som en författare, var så kreativ som möjligt och
  använd många metaforer.''
\item
  ``Du ska nu agera som en bilförsäljare av begagnade bilar, var så
  övertygande som möjligt och använd många adjektiv.''
\item
  ``Du ska nu agera som en förälskad tonåring, var så romantisk som
  möjligt och använd många metaforer.''
\item
  ``Du ska nu agera som en kock, var så kreativ som möjligt och använd
  många metaforer.''
\item
  ``Du ska nu agera som en advokat, var så övertygande som möjligt och
  använd många adjektiv.''
\end{itemize}

\textbf{Känslor}

Bard kan anta olika känslor. Här är några exempel:

\begin{itemize}
\tightlist
\item
  ``Du ska nu agera som en glad person, var så glad som möjligt och
  använd många utropstecken.''
\item
  ``Du ska nu agera som en ledsen person, var så ledsen som möjligt och
  använd många utropstecken.''
\item
  ``Du ska nu agera som en arg person, var så arg som möjligt och använd
  många utropstecken.''
\end{itemize}

\textbf{Tidsperioder}

Bard kan anta olika tidsperioder. Här är några exempel:

\begin{itemize}
\tightlist
\item
  ``Du ska nu agera som en person från medeltiden, var så medeltida som
  möjligt och använd många medeltida termer.''
\item
  ``Du ska nu agera som en person från framtiden, var så futuristisk som
  möjligt och använd många futuristiska termer.''
\item
  ``Du ska nu agera som en person från 1800-talet, var så 1800-tal som
  möjligt och använd många 1800-talstermer.''
\item
  ``Du ska nu agera som en person från 1900-talet, var så 1900-tal som
  möjligt och använd många 1900-talstermer.''
\item
  ``Du ska nu agera som en person från 1980-talet, var så 2000-tal som
  möjligt och använd många 1980-talstermer.''
\item
  ``Du ska nu agera som en person från 1960-talet, var så 2000-tal som
  möjligt och använd många 1960-talstermer.''
\end{itemize}

\emph{Det är bra att tänka på att Bard är ganska ny så man får påminna
den om sin roll några gånger innan den går in i det helt}

\hypertarget{if}{%
\section{If}\label{if}}

En introduktion till ämnet If på språket `Svenska' och kodspråk `Java'.

\hypertarget{introduktion}{%
\subsection{Introduktion}\label{introduktion}}

If är en kontrollstruktur som används inom Java för att ställa logiska
frågor. Svaren på dessa frågor kan vara antingen sant eller falskt.
Baserat på svaret kommer Java att utföra olika handlingar. Om svaret är
sant kommer koden efter if att exekveras. Om svaret är falskt kommer
Java att hoppa över koden efter if och fortsätta med resten av
programmet.

\hypertarget{exempel}{%
\subsection{Exempel}\label{exempel}}

Här är ett exempel på hur man använder if i Java:

\hypertarget{cb1}{}
\begin{Shaded}
\begin{Highlighting}[]
\DataTypeTok{int}\NormalTok{ ålder }\OperatorTok{=} \DecValTok{18}\OperatorTok{;}
\ControlFlowTok{if} \OperatorTok{(}\NormalTok{ålder }\OperatorTok{\textgreater{}=} \DecValTok{18}\OperatorTok{)} \OperatorTok{\{}
    \BuiltInTok{System}\OperatorTok{.}\FunctionTok{out}\OperatorTok{.}\FunctionTok{println}\OperatorTok{(}\StringTok{"Du är myndig"}\OperatorTok{);}
\OperatorTok{\}}
\end{Highlighting}
\end{Shaded}

I exemplet ovan kontrollerar vi om åldern är större eller lika med 18.
Om det är sant skrivs meddelandet ``Du är myndig'' ut i konsolfönstret.
Om åldern är mindre än 18 kommer koden efter if att hoppas över.

\hypertarget{fuxf6rdelar}{%
\subsection{Fördelar}\label{fuxf6rdelar}}

\begin{itemize}
\tightlist
\item
  If-satsen ger oss möjlighet att kontrollera olika villkor och anpassa
  koden baserat på dessa villkor.
\item
  Genom att använda if-satsen kan vi skapa mer flexibla och dynamiska
  program som kan hantera olika scenarier.
\item
  If-satsen är grundläggande för att kunna implementera logik och
  beslutsfattande i våra program.
\end{itemize}

\hypertarget{begruxe4nsningar}{%
\subsection{Begränsningar}\label{begruxe4nsningar}}

\begin{itemize}
\tightlist
\item
  If-satser kan ibland leda till komplex och svårhanterlig kod om de
  används för mycket eller på ett ogenomtänkt sätt.
\item
  Det är viktigt att vara noggrann med att utvärdera och testa olika
  villkor för att undvika logiska fel eller buggar i koden.
\item
  En felaktigt formulerad eller felaktigt placerad if-sats kan leda till
  att koden inte fungerar som förväntat.
\end{itemize}

\hypertarget{anvuxe4ndningsomruxe5den}{%
\subsection{Användningsområden}\label{anvuxe4ndningsomruxe5den}}

If-satser används i många olika situationer och scenarier inom
programmering. Här är några exempel på användningsområden för if-satser
i Java:

\begin{enumerate}
\def\labelenumi{\arabic{enumi}.}
\tightlist
\item
  Användarens validering: Om vi vill kontrollera om en användare har
  angett giltig indata kan vi använda en if-sats för att validera input
  och vidta lämpliga åtgärder baserat på resultatet.
\item
  Styra programflödet: If-satser kan användas för att styra hur
  programmet ska bete sig baserat på olika villkor. Till exempel kan vi
  använda en if-sats för att avgöra om ett visst block av kod ska
  utföras eller inte.---
\item
  Loopar: If-satser kan användas inuti loopar för att kontrollera när
  loopen ska avslutas eller fortsätta. Detta ger oss möjlighet att skapa
  mer flexibla och dynamiska loopar.
\item
  Felsökning och hantering av undantag: If-satser kan användas för att
  hantera olika undantagssituationer och felsöka problem i koden. Genom
  att kontrollera olika villkor kan vi vidta lämpliga åtgärder för att
  hantera fel och undantag.
\end{enumerate}

\hypertarget{tldr}{%
\subsection{TL;DR}\label{tldr}}

If är ett kommando inom Java som används för att ställa logiska frågor
och utföra olika handlingar baserat på svaren. Det ger oss möjlighet att
kontrollera villkor och anpassa programbeteendet. If-satser används för
att validera användarinput, styra programflödet, hantera loopar och
hantera undantagssituationer.

\hypertarget{obligatorisk-dad-joke}{%
\subsection{Obligatorisk dad joke}\label{obligatorisk-dad-joke}}

Varför var ``if'' så bra på att lösa problem? För att den alltid hade en
``else'' i ärmen!

\hypertarget{else}{%
\section{Else}\label{else}}

Innehållsförteckning

\{: .text-delta \}

\begin{enumerate}
\def\labelenumi{\arabic{enumi}.}
\tightlist
\item
  TOC \{:toc\}
\end{enumerate}

\hypertarget{beskrivning}{%
\subsection{Beskrivning}\label{beskrivning}}

``Annars'' är ett nyckelord som används i programmeringsspråket Java för
att skapa en alternativ väg i en kodblockstruktur. Det ger möjlighet att
utföra en annan sekvens av kod om en ``om''-sats utvärderas som falsk.
Med andra ord kan ``annars'' användas för att hantera fall där
``om''-villkoret inte är uppfyllt, och programmet kan utföra en annan
uppsättning instruktioner istället. När en ``om''-sats utvärderas som
sann, utförs de kodinstruktioner som ligger inuti ``om''-blocket. Om
villkoret i ``om''-satsen däremot utvärderas som falskt, kommer
programmet att hoppa över kodblocket inom ``om'' och istället utföra
kodblocket inom ``annars''.

För att använda ``annars'' måste det alltid följas av en ``om''. Det är
också viktigt att notera att ``annars'' alltid kommer i slutet av en
``om''-sats och kan inte användas utanför den.

\hypertarget{exempel}{%
\subsection{Exempel}\label{exempel}}

Här är ett exempel som visar hur ``annars'' kan användas för att
kontrollera om en person är myndig eller inte baserat på deras ålder:

\hypertarget{cb1}{}
\begin{Shaded}
\begin{Highlighting}[]
\DataTypeTok{int}\NormalTok{ ålder }\OperatorTok{=} \DecValTok{18}\OperatorTok{;}
\ControlFlowTok{if} \OperatorTok{(}\NormalTok{ålder }\OperatorTok{\textgreater{}=} \DecValTok{18}\OperatorTok{)} \OperatorTok{\{}
    \BuiltInTok{System}\OperatorTok{.}\FunctionTok{out}\OperatorTok{.}\FunctionTok{println}\OperatorTok{(}\StringTok{"Du är myndig"}\OperatorTok{);}
\OperatorTok{\}} \ControlFlowTok{else} \OperatorTok{\{}
    \BuiltInTok{System}\OperatorTok{.}\FunctionTok{out}\OperatorTok{.}\FunctionTok{println}\OperatorTok{(}\StringTok{"Du är inte myndig"}\OperatorTok{);}
\OperatorTok{\}}
\end{Highlighting}
\end{Shaded}

I detta exempel är ``age'' satt till 18. Eftersom 18 är större än eller
lika med 18, utvärderas ``if''-villkoret som sant, och programmet
skriver ut ``Du är myndig''. Om vi ändrar värdet på ``age'' till 17
kommer ``if''-villkoret att utvärderas som falskt, och programmet
skriver ut ``Du är inte myndig''. Med hjälp av ``else'' kan vi låta
programmet välja olika vägar att följa beroende på villkoren.

\hypertarget{sammanfattning}{%
\subsection{Sammanfattning}\label{sammanfattning}}

``Else'' är ett viktigt nyckelord inom Java som ger möjlighet att skapa
en alternativ väg i kodblocket baserat på ett ``if''-villkor. Om
``if''-satsen utvärderas som sann utförs de kodinstruktioner som ligger
inuti ``if''-blocket, och om villkoret utvärderas som falskt utförs
istället kodinstruktionerna inom ``else''-blocket. Det ger programmerare
möjlighet att hantera olika fall och göra beslut i sina program.

\hypertarget{termer-och-fuxf6rklaringar}{%
\subsection{Termer och förklaringar}\label{termer-och-fuxf6rklaringar}}

\begin{longtable}[]{@{}
  >{\raggedright\arraybackslash}p{(\columnwidth - 8\tabcolsep) * \real{0.0300}}
  >{\raggedright\arraybackslash}p{(\columnwidth - 8\tabcolsep) * \real{0.6400}}
  >{\raggedright\arraybackslash}p{(\columnwidth - 8\tabcolsep) * \real{0.0100}}
  >{\raggedright\arraybackslash}p{(\columnwidth - 8\tabcolsep) * \real{0.0400}}
  >{\raggedright\arraybackslash}p{(\columnwidth - 8\tabcolsep) * \real{0.2500}}@{}}
\toprule\noalign{}
\begin{minipage}[b]{\linewidth}\raggedright
Term
\end{minipage} & \begin{minipage}[b]{\linewidth}\raggedright
Förklaring
\end{minipage} & \begin{minipage}[b]{\linewidth}\raggedright
\end{minipage} & \begin{minipage}[b]{\linewidth}\raggedright
\end{minipage} & \begin{minipage}[b]{\linewidth}\raggedright
\end{minipage} \\
\midrule\noalign{}
\endhead
\bottomrule\noalign{}
\endlastfoot
Else & Ett nyckelord i Java som används för att skapa en alternativ väg
i kodblocket när ``if''-villkoret är falskt. & & & \\
If & Ett nyckelord i Java som används för att skapa en villkorsbaserad
kodblockstruktur. Om villkoret är sant, utförs kodinstruktionerna inom
``if''-blocket & & Code block & En grupp av kodinstruktioner som är
grupperade tillsammans. \\
Condition & En logisk fråga som kan utvärderas som antingen sann eller
falsk. & & & \\
\end{longtable}

\hypertarget{obligatorisk-dad-joke}{%
\subsection{Obligatorisk dad-joke}\label{obligatorisk-dad-joke}}

Varför gillar programmerare att använda ``else''?

För att de inte gillar att vara i ``if-nite''! 😄

\hypertarget{else-if}{%
\section{Else if}\label{else-if}}

Else if är en viktig konstruktion inom programmering som erbjuder
flexibilitet och tydlighet när det gäller att hantera olika villkor och
beslut. Det kan vara användbart i olika scenarier där vi behöver göra
flera olika val baserat på olika förutsättningar.

\hypertarget{anvuxe4ndningsomruxe5den}{%
\subsection{Användningsområden}\label{anvuxe4ndningsomruxe5den}}

Else if kan tillämpas i olika situationer, inklusive:

\begin{enumerate}
\def\labelenumi{\arabic{enumi}.}
\tightlist
\item
  \textbf{Betygsberäkning}: Om vi har ett betygssystem där vi behöver
  tilldela olika betyg baserat på en students poäng, kan vi använda Else
  if för att göra olika beslut baserat på poängintervallet.
\item
  \textbf{Validering}: Om vi behöver validera olika typer av data, som
  användarinput eller inmatade värden, kan vi använda Else if för att
  göra olika valideringsbeslut baserat på typen av data eller
  innehållet.
\item
  \textbf{Menyval}: Om vi har en meny med flera alternativ kan vi
  använda Else if för att hantera olika val baserat på användarens
  inmatning.
\item
  \textbf{Datum- och tidshantering}: Vi kan använda Else if för att
  hantera olika datum- och tidsscenarier, t.ex. att kontrollera om det
  är morgon, eftermiddag eller kväll och vidta olika åtgärder baserat på
  det.
\end{enumerate}

\hypertarget{kodexempel}{%
\subsection{Kodexempel}\label{kodexempel}}

Här är några kodexempel som visar användningen av Else if i Java:

\hypertarget{betygsberuxe4kning}{%
\subsubsection{Betygsberäkning}\label{betygsberuxe4kning}}

\hypertarget{cb1}{}
\begin{Shaded}
\begin{Highlighting}[]
\DataTypeTok{int}\NormalTok{ poäng }\OperatorTok{=} \DecValTok{75}\OperatorTok{;}
\ControlFlowTok{if} \OperatorTok{(}\NormalTok{poäng }\OperatorTok{\textgreater{}=} \DecValTok{90}\OperatorTok{)} \OperatorTok{\{}
    \BuiltInTok{System}\OperatorTok{.}\FunctionTok{out}\OperatorTok{.}\FunctionTok{println}\OperatorTok{(}\StringTok{"A"}\OperatorTok{);}
\OperatorTok{\}} \ControlFlowTok{else} \ControlFlowTok{if} \OperatorTok{(}\NormalTok{poäng }\OperatorTok{\textgreater{}=} \DecValTok{80}\OperatorTok{)} \OperatorTok{\{}
    \BuiltInTok{System}\OperatorTok{.}\FunctionTok{out}\OperatorTok{.}\FunctionTok{println}\OperatorTok{(}\StringTok{"B"}\OperatorTok{);}
\OperatorTok{\}} \ControlFlowTok{else} \ControlFlowTok{if} \OperatorTok{(}\NormalTok{poäng }\OperatorTok{\textgreater{}=} \DecValTok{70}\OperatorTok{)} \OperatorTok{\{}
    \BuiltInTok{System}\OperatorTok{.}\FunctionTok{out}\OperatorTok{.}\FunctionTok{println}\OperatorTok{(}\StringTok{"C"}\OperatorTok{);}
\OperatorTok{\}} \ControlFlowTok{else} \ControlFlowTok{if} \OperatorTok{(}\NormalTok{poäng }\OperatorTok{\textgreater{}=} \DecValTok{60}\OperatorTok{)} \OperatorTok{\{}
    \BuiltInTok{System}\OperatorTok{.}\FunctionTok{out}\OperatorTok{.}\FunctionTok{println}\OperatorTok{(}\StringTok{"D"}\OperatorTok{);}
\OperatorTok{\}} \ControlFlowTok{else} \OperatorTok{\{}
    \BuiltInTok{System}\OperatorTok{.}\FunctionTok{out}\OperatorTok{.}\FunctionTok{println}\OperatorTok{(}\StringTok{"F"}\OperatorTok{);}
\OperatorTok{\}}
\end{Highlighting}
\end{Shaded}

I detta exempel utvärderar vi en elevs poäng och skriver ut det
motsvarande betyget beroende på poängintervallet.

\hypertarget{validering}{%
\subsubsection{Validering}\label{validering}}

Här utvärderar vi input från användaren

\hypertarget{cb2}{}
\begin{Shaded}
\begin{Highlighting}[]
\KeywordTok{import} \ImportTok{java}\OperatorTok{.}\ImportTok{util}\OperatorTok{.}\ImportTok{Scanner}\OperatorTok{;}

\KeywordTok{public} \KeywordTok{class}\NormalTok{ Main }\OperatorTok{\{}
    \KeywordTok{public} \DataTypeTok{static} \DataTypeTok{void} \FunctionTok{main}\OperatorTok{(}\BuiltInTok{String}\OperatorTok{[]}\NormalTok{ args}\OperatorTok{)} \OperatorTok{\{}
    \BuiltInTok{Scanner}\NormalTok{ scanner }\OperatorTok{=} \KeywordTok{new} \BuiltInTok{Scanner}\OperatorTok{(}\BuiltInTok{System}\OperatorTok{.}\FunctionTok{in}\OperatorTok{);}

    \BuiltInTok{System}\OperatorTok{.}\FunctionTok{out}\OperatorTok{.}\FunctionTok{println}\OperatorTok{(}\StringTok{"Vänligen ange din ålder:"}\OperatorTok{);}
    \DataTypeTok{int}\NormalTok{ age }\OperatorTok{=}\NormalTok{ scanner}\OperatorTok{.}\FunctionTok{nextInt}\OperatorTok{();}

    \ControlFlowTok{if} \OperatorTok{(}\NormalTok{age }\OperatorTok{\textless{}} \DecValTok{0}\OperatorTok{)} \OperatorTok{\{}
        \BuiltInTok{System}\OperatorTok{.}\FunctionTok{out}\OperatorTok{.}\FunctionTok{println}\OperatorTok{(}\StringTok{"Åldern kan inte vara negativ."}\OperatorTok{);}
    \OperatorTok{\}} \ControlFlowTok{else} \ControlFlowTok{if} \OperatorTok{(}\NormalTok{age }\OperatorTok{\textless{}} \DecValTok{3}\OperatorTok{)} \OperatorTok{\{}
        \BuiltInTok{System}\OperatorTok{.}\FunctionTok{out}\OperatorTok{.}\FunctionTok{println}\OperatorTok{(}\StringTok{"Barnrumpa, du är för söt!"}\OperatorTok{);}
    \OperatorTok{\}} \ControlFlowTok{else} \ControlFlowTok{if} \OperatorTok{(}\NormalTok{age }\OperatorTok{\textless{}} \DecValTok{18}\OperatorTok{)} \OperatorTok{\{}
        \BuiltInTok{System}\OperatorTok{.}\FunctionTok{out}\OperatorTok{.}\FunctionTok{println}\OperatorTok{(}\StringTok{"Ung och full av energi!"}\OperatorTok{);}
    \OperatorTok{\}} \ControlFlowTok{else} \ControlFlowTok{if} \OperatorTok{(}\NormalTok{age }\OperatorTok{\textless{}} \DecValTok{40}\OperatorTok{)} \OperatorTok{\{}
        \BuiltInTok{System}\OperatorTok{.}\FunctionTok{out}\OperatorTok{.}\FunctionTok{println}\OperatorTok{(}\StringTok{"Medelålders och i sin bästa form!"}\OperatorTok{);}
    \OperatorTok{\}} \ControlFlowTok{else} \ControlFlowTok{if} \OperatorTok{(}\NormalTok{age }\OperatorTok{\textless{}} \DecValTok{60}\OperatorTok{)} \OperatorTok{\{}
        \BuiltInTok{System}\OperatorTok{.}\FunctionTok{out}\OperatorTok{.}\FunctionTok{println}\OperatorTok{(}\StringTok{"Gamling, men fortfarande aktiv!"}\OperatorTok{);}
    \OperatorTok{\}} \ControlFlowTok{else} \ControlFlowTok{if} \OperatorTok{(}\NormalTok{age }\OperatorTok{\textless{}} \DecValTok{100}\OperatorTok{)} \OperatorTok{\{}
        \BuiltInTok{System}\OperatorTok{.}\FunctionTok{out}\OperatorTok{.}\FunctionTok{println}\OperatorTok{(}\StringTok{"Jisses, du är gammal!"}\OperatorTok{);}
    \OperatorTok{\}} \ControlFlowTok{else} \OperatorTok{\{}
        \BuiltInTok{System}\OperatorTok{.}\FunctionTok{out}\OperatorTok{.}\FunctionTok{println}\OperatorTok{(}\StringTok{"Du är uråldrig! Respekt!"}\OperatorTok{);}
    \OperatorTok{\}}

\NormalTok{    scanner}\OperatorTok{.}\FunctionTok{close}\OperatorTok{();}
    \OperatorTok{\}}
\OperatorTok{\}}
\end{Highlighting}
\end{Shaded}

\hypertarget{menyval}{%
\subsubsection{Menyval}\label{menyval}}

\hypertarget{cb3}{}
\begin{Shaded}
\begin{Highlighting}[]
\BuiltInTok{System}\OperatorTok{.}\FunctionTok{out}\OperatorTok{.}\FunctionTok{println}\OperatorTok{(}\StringTok{"Välj ett alternativ:"}\OperatorTok{);}
\BuiltInTok{System}\OperatorTok{.}\FunctionTok{out}\OperatorTok{.}\FunctionTok{println}\OperatorTok{(}\StringTok{"1. Visa saldo"}\OperatorTok{);}
\BuiltInTok{System}\OperatorTok{.}\FunctionTok{out}\OperatorTok{.}\FunctionTok{println}\OperatorTok{(}\StringTok{"2. Gör en insättning"}\OperatorTok{);}
\BuiltInTok{System}\OperatorTok{.}\FunctionTok{out}\OperatorTok{.}\FunctionTok{println}\OperatorTok{(}\StringTok{"3. Gör ett uttag"}\OperatorTok{);}
\DataTypeTok{int}\NormalTok{ val }\OperatorTok{=} \BuiltInTok{Integer}\OperatorTok{.}\FunctionTok{parseInt}\OperatorTok{(}\BuiltInTok{System}\OperatorTok{.}\FunctionTok{console}\OperatorTok{().}\FunctionTok{readLine}\OperatorTok{());}
\ControlFlowTok{if} \OperatorTok{(}\NormalTok{val }\OperatorTok{==} \DecValTok{1}\OperatorTok{)} \OperatorTok{\{}
    \CommentTok{// Visa saldo}
\OperatorTok{\}} \ControlFlowTok{else} \ControlFlowTok{if} \OperatorTok{(}\NormalTok{val }\OperatorTok{==} \DecValTok{2}\OperatorTok{)} \OperatorTok{\{}
    \CommentTok{// Gör en insättning}
\OperatorTok{\}} \ControlFlowTok{else} \ControlFlowTok{if} \OperatorTok{(}\NormalTok{val }\OperatorTok{==} \DecValTok{3}\OperatorTok{)} \OperatorTok{\{}
    \CommentTok{// Gör ett uttag}
\OperatorTok{\}} \ControlFlowTok{else} \OperatorTok{\{}
    \BuiltInTok{System}\OperatorTok{.}\FunctionTok{out}\OperatorTok{.}\FunctionTok{println}\OperatorTok{(}\StringTok{"Ogiltigt val"}\OperatorTok{);}
\OperatorTok{\}}
\end{Highlighting}
\end{Shaded}

I detta exempel kan användaren välja ett alternativ från en meny genom
att ange ett nummer. Beroende på användarens val kan olika handlingar
utföras.

\hypertarget{datum--och-tidshantering}{%
\subsubsection{Datum- och
tidshantering}\label{datum--och-tidshantering}}

\hypertarget{cb4}{}
\begin{Shaded}
\begin{Highlighting}[]
\KeywordTok{import} \ImportTok{java}\OperatorTok{.}\ImportTok{time}\OperatorTok{.}\ImportTok{LocalTime}\OperatorTok{;}

\NormalTok{LocalTime nu }\OperatorTok{=}\NormalTok{ LocalTime}\OperatorTok{.}\FunctionTok{now}\OperatorTok{();}
\ControlFlowTok{if} \OperatorTok{(}\NormalTok{nu}\OperatorTok{.}\FunctionTok{getHour}\OperatorTok{()} \OperatorTok{\textless{}} \DecValTok{12}\OperatorTok{)} \OperatorTok{\{}
    \BuiltInTok{System}\OperatorTok{.}\FunctionTok{out}\OperatorTok{.}\FunctionTok{println}\OperatorTok{(}\StringTok{"God morgon!"}\OperatorTok{);}
\OperatorTok{\}} \ControlFlowTok{else} \ControlFlowTok{if} \OperatorTok{(}\NormalTok{nu}\OperatorTok{.}\FunctionTok{getHour}\OperatorTok{()} \OperatorTok{\textless{}} \DecValTok{18}\OperatorTok{)} \OperatorTok{\{}
    \BuiltInTok{System}\OperatorTok{.}\FunctionTok{out}\OperatorTok{.}\FunctionTok{println}\OperatorTok{(}\StringTok{"God eftermiddag!"}\OperatorTok{);}
\OperatorTok{\}} \ControlFlowTok{else} \OperatorTok{\{}
    \BuiltInTok{System}\OperatorTok{.}\FunctionTok{out}\OperatorTok{.}\FunctionTok{println}\OperatorTok{(}\StringTok{"God kväll!"}\OperatorTok{);}
\OperatorTok{\}}
\end{Highlighting}
\end{Shaded}

I detta exempel används Else if för att kontrollera aktuell tid och
skriva ut lämpliga hälsningar baserat på tiden på dygnet.

\hypertarget{sammanfattning}{%
\subsection{Sammanfattning}\label{sammanfattning}}

Else if är en viktig konstruktion inom programmering som erbjuder
flexibilitet och tydlighet när det gäller att hantera olika villkor och
beslut. Det kan vara användbart i olika scenarier där vi behöver göra
flera olika val baserat på olika förutsättningar. Det är dock viktigt
att vara medveten om begränsningarna och komplexiteten som kan uppstå
när man använder Else if.

Obligatorisk dad-joke:

Varför gillar programmerare att använda Else if?

För att de har så många val i livet och villkorlig kärlek är bara en av
dem! 😄

\hypertarget{ternary-if}{%
\section{Ternary if}\label{ternary-if}}

Ternary if, även känd som conditional operator, är en kompakt syntax i
programmeringsspråket Java som tillåter oss att uttrycka en enkel
if-sats på en rad. Det ger oss möjlighet att utvärdera ett villkor och
välja en av två uttryck beroende på om villkoret är sant eller falskt.
Syntaxen för ternary if består av tre delar: villkor, frågetecken och
uttryck för sant och falskt. Det är ett kraftfullt verktyg som kan göra
koden mer koncis och läsbar.

\hypertarget{exempel}{%
\subsection{Exempel}\label{exempel}}

Här är ett exempel som visar hur man använder ternary if i Java:

\hypertarget{cb1}{}
\begin{Shaded}
\begin{Highlighting}[]
\DataTypeTok{int}\NormalTok{ age }\OperatorTok{=} \DecValTok{18}\OperatorTok{;}
\BuiltInTok{String}\NormalTok{ result }\OperatorTok{=}\NormalTok{ age }\OperatorTok{\textgreater{}=} \DecValTok{18} \OperatorTok{?} \StringTok{"Du är myndig"} \OperatorTok{:} \StringTok{"Du är inte myndig"}\OperatorTok{;}
\BuiltInTok{System}\OperatorTok{.}\FunctionTok{out}\OperatorTok{.}\FunctionTok{println}\OperatorTok{(}\NormalTok{result}\OperatorTok{);}
\end{Highlighting}
\end{Shaded}

I detta exempel tilldelas strängen ``Du är myndig'' till variabeln
``result'' om värdet av variabeln ``age'' är större eller lika med 18.
Annars tilldelas strängen ``Du är inte myndig''. Sedan skrivs värdet av
``result'' ut till konsolen. Det är viktigt att notera att ternary if är
ett uttryck och kan användas inuti andra uttryck eller tilldelningar.
Det ger oss möjlighet att göra kompakt och läsbar kod för att hantera
enkla villkor. En annan användning av ternary if är när vi redan har en
boolean-variabel och vi behöver välja mellan två värden baserat på dess
värde. Här är ett exempel:

\hypertarget{cb2}{}
\begin{Shaded}
\begin{Highlighting}[]
\DataTypeTok{boolean}\NormalTok{ catIsCute }\OperatorTok{=} \KeywordTok{true}\OperatorTok{;}
\BuiltInTok{String}\NormalTok{ result }\OperatorTok{=}\NormalTok{ catIsCute }\OperatorTok{?} \StringTok{"Katten är söt \textless{}3"} \OperatorTok{:} \StringTok{"Katten är inte söt :("}\OperatorTok{;}
\end{Highlighting}
\end{Shaded}

I det här exemplet tilldelas strängen ``Katten är söt \textless3'' till
variabeln ``result'' om värdet av ``catIsCute'' är sant (true). Annars
tilldelas strängen ``Katten är inte söt :(''.

Ternary if är ett verktyg som kan göra koden mer läsbar och koncis i
situationer där vi behöver göra en enkel villkorskontroll och tilldela
olika värden beroende på resultatet.

\hypertarget{slutsats}{%
\subsection{Slutsats}\label{slutsats}}

Ternary if, eller conditional operator, är ett användbart verktyg inom
programmering som gör det möjligt för oss att uttrycka en enkel if-sats
på en rad. Det hjälper oss att göra koden mer koncis och läsbar genom
att välja mellan två uttryck beroende på om ett villkor är sant eller
falskt. Genom att använda ternary if kan vi undvika att skriva en längre
if-sats när vi bara behöver hantera enkla villkor.

Ternary if använder följande syntax:

\hypertarget{cb3}{}
\begin{Shaded}
\begin{Highlighting}[]
\NormalTok{villkor }\OperatorTok{?}\NormalTok{ uttryck om sant }\OperatorTok{:}\NormalTok{ uttryck om falskt}
\end{Highlighting}
\end{Shaded}

Där ``villkor'' är det uttryck som utvärderas, ``?'' är frågetecknet som
markerar början på ternary if, ``uttryck om sant'' är det värde eller
uttryck som tilldelas om villkoret är sant, och ``uttryck om falskt'' är
det värde eller uttryck som tilldelas om villkoret är falskt.

Det är viktigt att notera att ternary if bara är lämplig för enkla
villkor och enkla uttryck. Om du behöver hantera mer komplexa villkor
eller flera uttryck kan det vara bättre att använda en vanlig if-sats.

\hypertarget{mer-luxe4sning}{%
\subsection{Mer läsning}\label{mer-luxe4sning}}

\begin{itemize}
\tightlist
\item
  \href{https://docs.oracle.com/en/java/javase/14/docs/api/java.base/java/lang/ConditionalOperator.html}{Java
  Conditional Operator (?:)}
\item
  \href{https://docs.oracle.com/en/java/javase/14/docs/api/java.base/java/lang/IfStatement.html}{Java
  If-else statement}
\end{itemize}

\hypertarget{termer}{%
\subsection{Termer}\label{termer}}

\begin{itemize}
\tightlist
\item
  Ternary if: En kompakt syntax i Java som tillåter oss att uttrycka en
  enkel if-sats på en rad.
\item
  Conditional operator: En annan term för ternary if, som beskriver dess
  användning för att utvärdera och välja mellan två uttryck baserat på
  ett villkor.
\end{itemize}

\hypertarget{obligatorisk-pappaskuxe4mt}{%
\subsection{Obligatorisk pappaskämt}\label{obligatorisk-pappaskuxe4mt}}

Varför är ternary if alltid så bekymrad?

För att den alltid vill vara säker på att det är \texttt{true} eller
\texttt{false}, inga ``may-be''!

\hypertarget{variabler}{%
\section{Variabler}\label{variabler}}

En variabel är en behållare som används för att lagra data. I Java måste
alla variabler deklareras innan de kan användas. Detta innebär att du
måste ange vilken typ av data som variabeln kommer att lagra.

graph LR A{[}Variabel{]} --\textgreater{} B{[}Deklarera{]} B
--\textgreater{} C{[}Initiera{]}

\hypertarget{exempel}{%
\subsection{Exempel}\label{exempel}}

\hypertarget{cb1}{}
\begin{Shaded}
\begin{Highlighting}[]
\CommentTok{// Deklarera en variabel av typen int.}
\DataTypeTok{int}\NormalTok{ number }\OperatorTok{=} \DecValTok{1138}\OperatorTok{;}

\CommentTok{// Deklarera en variabel av typen String.}
\BuiltInTok{String}\NormalTok{ name }\OperatorTok{=} \StringTok{"Obi Wan Kenobi"}\OperatorTok{;}

\CommentTok{// Deklarera en variabel av typen boolean.}
\DataTypeTok{boolean}\NormalTok{ isJedi }\OperatorTok{=} \KeywordTok{true}\OperatorTok{;}

\CommentTok{// Deklarera en variabel av typen double.}
\DataTypeTok{double}\NormalTok{ pi }\OperatorTok{=} \FloatTok{3.14}\OperatorTok{;}

\CommentTok{// Deklarera en variabel av typen char.}
\DataTypeTok{char}\NormalTok{ letter }\OperatorTok{=} \CharTok{\textquotesingle{}J\textquotesingle{}}\OperatorTok{;}
\end{Highlighting}
\end{Shaded}

I detta exempel har vi deklarerat fem variabler. Den första är en int,
vilket innebär att den kan lagra heltal. Den andra är en sträng
(String), vilket innebär att den kan lagra text. Den tredje är en
boolean, vilket innebär att den kan lagra sant eller falskt. Den fjärde
är en double, vilket innebär att den kan lagra decimaltal. Den femte är
en char, vilket innebär att den kan lagra ett tecken.

\hypertarget{typer}{%
\section{Typer}\label{typer}}

En typ är en klassificering av data som används för att bestämma vilken
typ av värde en variabel kan lagra. I Java finns det olika typer av
typer, inklusive primitiva typer och referenstyper.

\hypertarget{primitiva-typer}{%
\subsection{Primitiva typer}\label{primitiva-typer}}

Här är en lista över några vanliga primitiva typer i Java tillsammans
med deras beskrivning, minsta värde, högsta värde och storlek i minnet:

\begin{longtable}[]{@{}
  >{\raggedright\arraybackslash}p{(\columnwidth - 8\tabcolsep) * \real{0.0700}}
  >{\raggedright\arraybackslash}p{(\columnwidth - 8\tabcolsep) * \real{0.4000}}
  >{\raggedright\arraybackslash}p{(\columnwidth - 8\tabcolsep) * \real{0.2100}}
  >{\raggedright\arraybackslash}p{(\columnwidth - 8\tabcolsep) * \real{0.2200}}
  >{\raggedright\arraybackslash}p{(\columnwidth - 8\tabcolsep) * \real{0.0700}}@{}}
\toprule\noalign{}
\begin{minipage}[b]{\linewidth}\raggedright
Typ
\end{minipage} & \begin{minipage}[b]{\linewidth}\raggedright
Beskrivning
\end{minipage} & \begin{minipage}[b]{\linewidth}\raggedright
Minsta
\end{minipage} & \begin{minipage}[b]{\linewidth}\raggedright
Högsta
\end{minipage} & \begin{minipage}[b]{\linewidth}\raggedright
Storlek
\end{minipage} \\
\midrule\noalign{}
\endhead
\bottomrule\noalign{}
\endlastfoot
boolean & Booleskt värde & false & true & 1 byte \\
char & 16-bitars Unicode-tecken & `000' & `\,' & 2 byte \\
byte & 8-bitars heltal & -128 & 127 & 1 byte \\
short & 16-bitars heltal & -32768 & 32767 & 2 byte \\
int & 32-bitars heltal & -2147483648 & 2147483647 & 4 byte \\
long & 64-bitars heltal & -9223372036854775808 & 9223372036854775807 & 8
byte \\
float & 32-bitars flyttal med enkels precision & 1.4E-45 & 3.4028235E+38
& 4 byte \\
double & 64-bitars flyttal med dubbel precision & 4.9E-324 &
1.79769313486232E+308 & 8 byte \\
\end{longtable}

Observera att storlekarna som anges ovan är standardstorlekar och kan
variera beroende på implementationen och plattformen.

\hypertarget{referenstyper}{%
\subsection{Referenstyper}\label{referenstyper}}

Utöver de primitiva typerna har Java också referenstyper, som är typer
som lagrar referenser till objekt. Här är några exempel på vanliga
referenstyper i Java:

\begin{itemize}
\tightlist
\item
  \texttt{String}: En klass för att hantera textsträngar.
\item
  \texttt{Array}: En typ för att skapa och hantera arrayer av andra
  typer.
\item
  \texttt{Class}: En klass som representerar en typ vid körningstid.
\item
  \texttt{Interface}: En typ som definierar en samling metoder som en
  klass kan implementera.
\item
  \texttt{List}: En typ för att skapa och hantera listor av objekt.
\item
  \texttt{Map}: En typ för att skapa och hantera mappningar av
  nyckel-värde-par.
\end{itemize}

Referenstyper lagrar inte själva värdet utan pekar på objektet i minnet
där värdet finns lagrat. De ger oss möjligheten att hantera mer komplexa
datastrukturer och skapa hierarkier av objekt.

\hypertarget{termer}{%
\subsection{Termer}\label{termer}}

\begin{longtable}[]{@{}
  >{\raggedright\arraybackslash}p{(\columnwidth - 2\tabcolsep) * \real{0.1200}}
  >{\raggedright\arraybackslash}p{(\columnwidth - 2\tabcolsep) * \real{0.8700}}@{}}
\toprule\noalign{}
\begin{minipage}[b]{\linewidth}\raggedright
Term
\end{minipage} & \begin{minipage}[b]{\linewidth}\raggedright
Förklaring
\end{minipage} \\
\midrule\noalign{}
\endhead
\bottomrule\noalign{}
\endlastfoot
Array & En typ för att skapa och hantera arrayer av andra typer. \\
Boolesk & En typ med två möjliga värden: \texttt{true} och
\texttt{false}. \\
Char & En 16-bitars Unicode-tecken. \\
Double & En 64-bitars flyttal med dubbel precision. \\
Float & En 32-bitars flyttal med enkels precision. \\
Int & En 32-bitars heltalstyp. \\
Interface & En typ som definierar en samling metoder som en klass kan
implementera. \\
Klass & En mall för att skapa objekt. \\
List & En typ för att skapa och hantera listor av objekt. \\
Long & En 64-bitars heltalstyp. \\
Map & En typ för att skapa och hantera mappningar av
nyckel-värde-par. \\
Referenstyp & En typ som lagrar en referens till ett objekt. \\
Short & En 16-bitars heltalstyp. \\
Stack & En del av minnet som används för att lagra metoder och lokala
variabler. \\
String & En klass för att hantera textsträngar. \\
Variabel & En behållare som används för att lagra data. \\
Värdestyp & En typ som lagrar ett värde. \\
Wrapperklass & En klass som används för att ``linda in'' primitiva typer
och ge dem extra funktionalitet. \\
Wrapperobjekt & Ett objekt som innehåller en referens till en primitiv
typ. \\
\end{longtable}

\hypertarget{tldr-sammanfattning}{%
\subsection{TL;DR-sammanfattning}\label{tldr-sammanfattning}}

I Java finns det olika typer av typer, inklusive primitiva typer och
referenstyper. Primitiva typer inkluderar boolean, char, byte, short,
int, long, float och double. Dessa typer lagrar värden direkt.
Referenstyper inkluderar String, Array, Class, Interface, List och Map.
Dessa typer lagrar referenser till objekt och ger oss möjligheten att
hantera mer komplexa datastrukturer.

\hypertarget{datastrukturer}{%
\section{Datastrukturer}\label{datastrukturer}}

Innehållsförteckning

\{: .text-delta \} 1. TOC \{:toc\}

\hypertarget{beskrivning}{%
\subsection{Beskrivning}\label{beskrivning}}

Datastrukturer är organiserade sätt att lagra och hantera data i en
dator. Det finns olika typer av datastrukturer, såsom arrayer, listor,
träd, grafer och hashtabeller, var och en med sina egna unika egenskaper
och användningsområden. Dessa datastrukturer kan användas för att
effektivt söka, sortera, lagra och hämta data, och är en viktig del av
programmering och datavetenskap. Datastrukturer är lika mångsidiga som
superhjältar i DC- och Marvel-världen. Precis som hjältarna har de olika
förmågor och användningsområden som gör dem unika och passande för olika
situationer. Låt oss ta en titt på några av dessa datastrukturer och hur
de kan användas i programmering.

\hypertarget{array}{%
\subsection{Array}\label{array}}

En array är en samling av element av samma datatyp som lagras i minnet i
en sekventiell ordning. Varje element i arrayen kan nås med hjälp av ett
index. Indexet representerar positionen för varje element i arrayen, där
det första elementet har index 0, det andra elementet har index 1 och så
vidare. Arrayer används för att lagra och hantera en samling av värden
av samma typ, som till exempel en lista med tal eller strängar. Exempel:

\hypertarget{cb1}{}
\begin{Shaded}
\begin{Highlighting}[]
\BuiltInTok{String}\OperatorTok{[]}\NormalTok{ avengers }\OperatorTok{=} \KeywordTok{new} \BuiltInTok{String}\OperatorTok{[}\DecValTok{3}\OperatorTok{];}
\NormalTok{avengers}\OperatorTok{[}\DecValTok{0}\OperatorTok{]} \OperatorTok{=} \StringTok{"Iron Man"}\OperatorTok{;}
\NormalTok{avengers}\OperatorTok{[}\DecValTok{1}\OperatorTok{]} \OperatorTok{=} \StringTok{"Captain America"}\OperatorTok{;}
\NormalTok{avengers}\OperatorTok{[}\DecValTok{2}\OperatorTok{]} \OperatorTok{=} \StringTok{"Thor"}\OperatorTok{;}
\BuiltInTok{System}\OperatorTok{.}\FunctionTok{out}\OperatorTok{.}\FunctionTok{println}\OperatorTok{(}\NormalTok{avengers}\OperatorTok{[}\DecValTok{0}\OperatorTok{]);}  \CommentTok{// Output: Iron Man}
\end{Highlighting}
\end{Shaded}

\hypertarget{list}{%
\subsection{List}\label{list}}

En lista är en dynamisk samling av element av samma datatyp. Till
skillnad från en array kan storleken på en lista ändras dynamiskt när
element läggs till eller tas bort. Listor tillhandahåller också en mängd
användbara metoder för att söka, sortera och manipulera elementen i
listan. Listor används när du behöver en samling av element som kan
ändras i storlek under körningstid.

\hypertarget{cb2}{}
\begin{Shaded}
\begin{Highlighting}[]
\BuiltInTok{List}\OperatorTok{\textless{}}\DataTypeTok{int}\OperatorTok{\textgreater{}}\NormalTok{ numbers }\OperatorTok{=} \KeywordTok{new} \BuiltInTok{ArrayList}\OperatorTok{\textless{}\textgreater{}(}\BuiltInTok{Arrays}\OperatorTok{.}\FunctionTok{asList}\OperatorTok{(}\DecValTok{1}\OperatorTok{,} \DecValTok{2}\OperatorTok{,} \DecValTok{3}\OperatorTok{,} \DecValTok{4}\OperatorTok{,} \DecValTok{5}\OperatorTok{));}
\NormalTok{numbers}\OperatorTok{.}\FunctionTok{add}\OperatorTok{(}\DecValTok{6}\OperatorTok{);}
\BuiltInTok{System}\OperatorTok{.}\FunctionTok{out}\OperatorTok{.}\FunctionTok{println}\OperatorTok{(}\NormalTok{numbers}\OperatorTok{.}\FunctionTok{size}\OperatorTok{());}  \CommentTok{// Output: 6}
\end{Highlighting}
\end{Shaded}

\hypertarget{stack}{%
\subsection{Stack}\label{stack}}

En stack är en datastruktur som använder principen LIFO (Last In, First
Out), vilket innebär att det senast tillagda elementet är det första som
tas bort. Du kan lägga till element på toppen av stacken och ta bort
element från toppen. Stackar används när du behöver hantera data i
omvänd ordning av deras inmatningssekvens. Exempel:

\hypertarget{cb3}{}
\begin{Shaded}
\begin{Highlighting}[]
\BuiltInTok{Stack}\OperatorTok{\textless{}}\BuiltInTok{String}\OperatorTok{\textgreater{}}\NormalTok{ superheroes }\OperatorTok{=} \KeywordTok{new} \BuiltInTok{Stack}\OperatorTok{\textless{}\textgreater{}();}
\NormalTok{superheroes}\OperatorTok{.}\FunctionTok{push}\OperatorTok{(}\StringTok{"Spider{-}Man"}\OperatorTok{);}
\NormalTok{superheroes}\OperatorTok{.}\FunctionTok{push}\OperatorTok{(}\StringTok{"Iron Man"}\OperatorTok{);}
\NormalTok{superheroes}\OperatorTok{.}\FunctionTok{push}\OperatorTok{(}\StringTok{"Hulk"}\OperatorTok{);}
\BuiltInTok{System}\OperatorTok{.}\FunctionTok{out}\OperatorTok{.}\FunctionTok{println}\OperatorTok{(}\NormalTok{superheroes}\OperatorTok{.}\FunctionTok{pop}\OperatorTok{());}  \CommentTok{// Output: Hulk}
\end{Highlighting}
\end{Shaded}

\hypertarget{queue}{%
\subsection{Queue}\label{queue}}

En kö är en datastruktur som använder principen FIFO (First In, First
Out), vilket innebär att det första tillagda elementet är det första som
tas bort. Du kan lägga till element i slutet av kön och ta bort element
från början. Kön används när du behöver hantera data i samma ordning som
deras inmatningssekvens.

\hypertarget{cb4}{}
\begin{Shaded}
\begin{Highlighting}[]
\BuiltInTok{Queue}\OperatorTok{\textless{}}\BuiltInTok{String}\OperatorTok{\textgreater{}}\NormalTok{ heroes }\OperatorTok{=} \KeywordTok{new} \BuiltInTok{LinkedList}\OperatorTok{\textless{}\textgreater{}();}
\NormalTok{heroes}\OperatorTok{.}\FunctionTok{add}\OperatorTok{(}\StringTok{"Captain America"}\OperatorTok{);}
\NormalTok{heroes}\OperatorTok{.}\FunctionTok{add}\OperatorTok{(}\StringTok{"Black Widow"}\OperatorTok{);}
\NormalTok{heroes}\OperatorTok{.}\FunctionTok{add}\OperatorTok{(}\StringTok{"Hawkeye"}\OperatorTok{);}
\BuiltInTok{System}\OperatorTok{.}\FunctionTok{out}\OperatorTok{.}\FunctionTok{println}\OperatorTok{(}\NormalTok{heroes}\OperatorTok{.}\FunctionTok{poll}\OperatorTok{());}  \CommentTok{// Output: Captain America}
\end{Highlighting}
\end{Shaded}

\hypertarget{dictionary}{%
\subsection{Dictionary}\label{dictionary}}

Ett dictionary är en datastruktur som lagrar data som en samling av
nyckel-värde-par. Varje nyckel måste vara unik och kan användas för att
hämta det tillhörande värdet. Dictionaries används när du behöver snabb
åtkomst till värden baserat på en given nyckel.

\hypertarget{cb5}{}
\begin{Shaded}
\begin{Highlighting}[]
\BuiltInTok{HashMap}\OperatorTok{\textless{}}\BuiltInTok{String}\OperatorTok{,} \BuiltInTok{Integer}\OperatorTok{\textgreater{}}\NormalTok{ heroStats }\OperatorTok{=} \KeywordTok{new} \BuiltInTok{HashMap}\OperatorTok{\textless{}\textgreater{}();}
\NormalTok{heroStats}\OperatorTok{.}\FunctionTok{put}\OperatorTok{(}\StringTok{"Iron Man"}\OperatorTok{,} \DecValTok{100}\OperatorTok{);}
\NormalTok{heroStats}\OperatorTok{.}\FunctionTok{put}\OperatorTok{(}\StringTok{"Captain America"}\OperatorTok{,} \DecValTok{90}\OperatorTok{);}
\NormalTok{heroStats}\OperatorTok{.}\FunctionTok{put}\OperatorTok{(}\StringTok{"Thor"}\OperatorTok{,} \DecValTok{95}\OperatorTok{);}
\BuiltInTok{System}\OperatorTok{.}\FunctionTok{out}\OperatorTok{.}\FunctionTok{println}\OperatorTok{(}\NormalTok{heroStats}\OperatorTok{.}\FunctionTok{get}\OperatorTok{(}\StringTok{"Iron Man"}\OperatorTok{));}  \CommentTok{// Output: 100}
\end{Highlighting}
\end{Shaded}

\hypertarget{hashset}{%
\subsection{HashSet}\label{hashset}}

En HashSet är en datastruktur som lagrar unika element utan någon
specifik ordning. HashSet använder en hashfunktion för att beräkna
hashvärdet för varje element, vilket möjliggör snabb insättning,
borttagning och sökning av element. HashSet används när du behöver lagra
en samling av unika element och inte behöver behålla någon specifik
ordning.

\hypertarget{cb6}{}
\begin{Shaded}
\begin{Highlighting}[]
\KeywordTok{import} \ImportTok{java}\OperatorTok{.}\ImportTok{util}\OperatorTok{.}\ImportTok{HashSet}\OperatorTok{;}
\KeywordTok{public} \KeywordTok{class}\NormalTok{ Main }\OperatorTok{\{}
\KeywordTok{public} \DataTypeTok{static} \DataTypeTok{void} \FunctionTok{main}\OperatorTok{(}\BuiltInTok{String}\OperatorTok{[]}\NormalTok{ args}\OperatorTok{)} \OperatorTok{\{}
\BuiltInTok{HashSet}\OperatorTok{\textless{}}\BuiltInTok{String}\OperatorTok{\textgreater{}}\NormalTok{ villains }\OperatorTok{=} \KeywordTok{new} \BuiltInTok{HashSet}\OperatorTok{\textless{}}\BuiltInTok{String}\OperatorTok{\textgreater{}();}
\NormalTok{villains}\OperatorTok{.}\FunctionTok{add}\OperatorTok{(}\StringTok{"Joker"}\OperatorTok{);}
\NormalTok{villains}\OperatorTok{.}\FunctionTok{add}\OperatorTok{(}\StringTok{"Loki"}\OperatorTok{);}
\NormalTok{villains}\OperatorTok{.}\FunctionTok{add}\OperatorTok{(}\StringTok{"Green Goblin"}\OperatorTok{);}
\BuiltInTok{System}\OperatorTok{.}\FunctionTok{out}\OperatorTok{.}\FunctionTok{println}\OperatorTok{(}\NormalTok{villains}\OperatorTok{.}\FunctionTok{size}\OperatorTok{());}  \CommentTok{// Output: 3}
\OperatorTok{\}}
\OperatorTok{\}}
\end{Highlighting}
\end{Shaded}

I exemplen ovan har vi använt namnen på några av de populära
superhjältarna och skurkarna från DC och Marvel för att illustrera
användningen av olika datastrukturer. Precis som superhjältar och
skurkar kan vara olika och har olika förmågor, har också
datastrukturerna olika egenskaper och användningsområden.

\hypertarget{uxe4nnu-fler-datastrukturer}{%
\subsection{Ännu fler
datastrukturer}\label{uxe4nnu-fler-datastrukturer}}

Exemplen ovan är inte alla datastrukturer som finns. Artikeln fokuserar
på några vanliga datastrukturer, men det finns fler datastrukturer som
kan vara användbara inom programmering och datavetenskap. Här är några
ytterligare exempel på datastrukturer som kan vara värda att undersöka:

\begin{itemize}
\tightlist
\item
  Länkad lista: En datastruktur där elementen är länkade tillsammans i
  en kedja.
\item
  Träd: En hierarkisk datastruktur där varje element har en överordnad
  och noll eller flera underordnade element.
\item
  Graf: En datastruktur som består av noder (vertices) och kanter
  (edges) som förbinder dessa noder.
\item
  Heap: En speciell typ av trädstruktur där varje nod har ett värde
  större eller mindre än eller lika med sina underordnade noder.
\item
  Trie: En speciell typ av trädstruktur som används för att effektivt
  lagra och söka efter ord och mönster.
\item
  Hash-tabell: En datastruktur som använder en hashfunktion för att
  snabbt lagra och hämta värden baserat på en given nyckel.
\item
  Graf: En datastruktur som består av noder (vertices) och kanter
  (edges) som förbinder dessa noder. Grafer används för att representera
  relationer och nätverk mellan olika objekt.---
\item
  Heap: En speciell typ av trädstruktur där varje nod har ett värde
  större eller mindre än eller lika med sina underordnade noder. Heapar
  används ofta för att implementera prioritetsköer och
  sorteringsalgoritmer som heap sort.
\item
  Trie: En speciell typ av trädstruktur som används för att effektivt
  lagra och söka efter ord och mönster. Tries används ofta i
  applikationer som textkomprimering, stavningskontroll och
  autokomplettering.
\item
  Hash Table: En datastruktur som använder en hashfunktion för att
  snabbt lagra och hämta värden baserat på en given nyckel.
  Hash-tabeller är användbara för att implementera snabbt åtkomstbara
  datastrukturer som uppslagslistor och uppslagsböcker.
\item
  Set: En datastruktur som lagrar unika element utan någon specifik
  ordning. Set kan användas för att utföra olika mängdoperationer som
  union, snitt och differens.
\item
  Stack: En datastruktur som följer principen ``last in, first out''
  (LIFO), vilket innebär att det senast tillagda elementet är det första
  som tas bort. Stackar används ofta för att hantera
  återuppringningsinformation i funktioner eller för att utvärdera
  uttryck i matematiska uttryck.
\item
  Queue: En datastruktur som följer principen ``first in, first out''
  (FIFO), vilket innebär att det första tillagda elementet är det första
  som tas bort. Köer används ofta för att hantera processer och trådar
  eller för att implementera buffertar i kommunikationssystem. Dessa är
  bara några exempel på datastrukturer som du kan lära dig mer om. Varje
  datastruktur har sina egna unika egenskaper och användningsområden,
  och det kan vara värdefullt att utforska dem för att utöka din kunskap
  om datastrukturer och deras tillämpningar.
\end{itemize}

\hypertarget{referenser}{%
\subsection{Referenser}\label{referenser}}

\begin{itemize}
\tightlist
\item
  \href{https://en.wikipedia.org/wiki/Data_structure}{Wikipedia}
\item
  \href{https://www.geeksforgeeks.org/data-structures/}{GeeksforGeeks}
\item
  \href{https://www.tutorialspoint.com/data_structures_algorithms/index.htm}{TutorialsPoint}
\item
  \href{https://www.programiz.com/dsa}{Programiz}
\item
  \href{https://www.khanacademy.org/computing/computer-science/algorithms}{Khan
  Academy}
\item
  \href{https://www.codecademy.com/learn/learn-data-structures}{Codecademy}
\item
  \href{https://www.youtube.com/playlistlist=PLu0W_9lII9agICnT8t4iYVSZ3eykIAOME}{CodeWithHarry}
\end{itemize}

\hypertarget{array-uxf6vningar}{%
\section{Array övningar}\label{array-uxf6vningar}}

Här finns några övningar för att träna på att använda arrayer i Java.

Innehållsförteckning

\{: .text-delta \} 1. Innehållsförteckning \{:toc\}

\hypertarget{skapa-en-array}{%
\subsection{1. Skapa en array}\label{skapa-en-array}}

Skapa en array med namnet \texttt{numbers} som innehåller följande tal:
1, 2, 3, 4, 5.

Lösning

\hypertarget{cb1}{}
\begin{Shaded}
\begin{Highlighting}[]
\DataTypeTok{int}\OperatorTok{[]}\NormalTok{ numbers }\OperatorTok{=} \OperatorTok{\{}\DecValTok{1}\OperatorTok{,} \DecValTok{2}\OperatorTok{,} \DecValTok{3}\OperatorTok{,} \DecValTok{4}\OperatorTok{,} \DecValTok{5}\OperatorTok{\};}
\end{Highlighting}
\end{Shaded}

\hypertarget{tilldela-vuxe4rden-till-en-array}{%
\subsection{2. Tilldela värden till en
array}\label{tilldela-vuxe4rden-till-en-array}}

Skapa en array med namnet \texttt{fruits} som kan hålla 3 fruktnamn.
Tilldela värdena ``äpple'', ``banan'' och ``apelsin'' till arrayen.

Lösning

\hypertarget{cb2}{}
\begin{Shaded}
\begin{Highlighting}[]
\BuiltInTok{String}\OperatorTok{[]}\NormalTok{ fruits }\OperatorTok{=} \OperatorTok{\{}\StringTok{"äpple"}\OperatorTok{,} \StringTok{"banan"}\OperatorTok{,} \StringTok{"apelsin"}\OperatorTok{\};}
\end{Highlighting}
\end{Shaded}

\hypertarget{huxe4mta-vuxe4rden-fruxe5n-en-array}{%
\subsection{3. Hämta värden från en
array}\label{huxe4mta-vuxe4rden-fruxe5n-en-array}}

Använd indexering för att hämta det andra värdet från arrayen
\texttt{numbers}. int andraTalet = numbers{[}1{]};

\hypertarget{uppdatera-vuxe4rden-i-en-array}{%
\subsection{4. Uppdatera värden i en
array}\label{uppdatera-vuxe4rden-i-en-array}}

Uppdatera det första värdet i arrayen \texttt{fruits} till ``päron''.

Lösning

\hypertarget{cb3}{}
\begin{Shaded}
\begin{Highlighting}[]
\NormalTok{fruits}\OperatorTok{[}\DecValTok{0}\OperatorTok{]} \OperatorTok{=} \StringTok{"päron"}\OperatorTok{;}
\end{Highlighting}
\end{Shaded}

\hypertarget{loopa-igenom-en-array}{%
\subsection{5. Loopa igenom en array}\label{loopa-igenom-en-array}}

Använd en \texttt{for}-loop för att skriva ut varje frukt i arrayen
\texttt{fruits}.

Lösning

\hypertarget{cb4}{}
\begin{Shaded}
\begin{Highlighting}[]
\ControlFlowTok{for} \OperatorTok{(}\DataTypeTok{int}\NormalTok{ i }\OperatorTok{=} \DecValTok{0}\OperatorTok{;}\NormalTok{ i }\OperatorTok{\textless{}}\NormalTok{ fruits}\OperatorTok{.}\FunctionTok{length}\OperatorTok{;}\NormalTok{ i}\OperatorTok{++)} \OperatorTok{\{}
    \BuiltInTok{System}\OperatorTok{.}\FunctionTok{out}\OperatorTok{.}\FunctionTok{println}\OperatorTok{(}\NormalTok{fruits}\OperatorTok{[}\NormalTok{i}\OperatorTok{]);}
\OperatorTok{\}}
\end{Highlighting}
\end{Shaded}

\hypertarget{array-av-objekt}{%
\subsection{6. Array av objekt}\label{array-av-objekt}}

Skapa en array med namnet \texttt{characters} som kan hålla objekt av
typen \texttt{Character}. Skapa två \texttt{Character}-objekt med namnen
``Iron Man'' och ``Captain America'' och placera dem i arrayen.

Lösning

\hypertarget{cb5}{}
\begin{Shaded}
\begin{Highlighting}[]
\BuiltInTok{Character}\OperatorTok{[]}\NormalTok{ characters }\OperatorTok{=} \KeywordTok{new} \BuiltInTok{Character}\OperatorTok{[}\DecValTok{2}\OperatorTok{];}
\NormalTok{characters}\OperatorTok{[}\DecValTok{0}\OperatorTok{]} \OperatorTok{=} \KeywordTok{new} \BuiltInTok{Character}\OperatorTok{(}\StringTok{"Iron Man"}\OperatorTok{);}
\NormalTok{characters}\OperatorTok{[}\DecValTok{1}\OperatorTok{]} \OperatorTok{=} \KeywordTok{new} \BuiltInTok{Character}\OperatorTok{(}\StringTok{"Captain America"}\OperatorTok{);}
\end{Highlighting}
\end{Shaded}

\hypertarget{tvuxe5dimensionell-array}{%
\subsection{7. Tvådimensionell array}\label{tvuxe5dimensionell-array}}

Skapa en tvådimensionell array med namnet \texttt{matrix} som har 3
rader och 3 kolumner. Fyll arrayen med värdena 1, 2, 3, 4, 5, 6, 7, 8,
9.

Lösning

\hypertarget{cb6}{}
\begin{Shaded}
\begin{Highlighting}[]
\DataTypeTok{int}\OperatorTok{[][]}\NormalTok{ matrix }\OperatorTok{=} \OperatorTok{\{}
    \OperatorTok{\{}\DecValTok{1}\OperatorTok{,} \DecValTok{2}\OperatorTok{,} \DecValTok{3}\OperatorTok{\},}
    \OperatorTok{\{}\DecValTok{4}\OperatorTok{,} \DecValTok{5}\OperatorTok{,} \DecValTok{6}\OperatorTok{\},}
    \OperatorTok{\{}\DecValTok{7}\OperatorTok{,} \DecValTok{8}\OperatorTok{,} \DecValTok{9}\OperatorTok{\}\};}
\end{Highlighting}
\end{Shaded}

\hypertarget{suxf6k-efter-ett-vuxe4rde-i-en-array}{%
\subsection{8. Sök efter ett värde i en
array}\label{suxf6k-efter-ett-vuxe4rde-i-en-array}}

Skriv en metod med namnet \texttt{containsValue} som tar emot en array
av heltal och ett heltal att söka efter. Metoden ska returnera
\texttt{true} om värdet finns i arrayen, annars \texttt{false}.

Lösning

\hypertarget{cb7}{}
\begin{Shaded}
\begin{Highlighting}[]
\KeywordTok{public} \DataTypeTok{static} \DataTypeTok{boolean} \FunctionTok{containsValue}\OperatorTok{(}\DataTypeTok{int}\OperatorTok{[]}\NormalTok{ array}\OperatorTok{,} \DataTypeTok{int}\NormalTok{ value}\OperatorTok{)} \OperatorTok{\{}
    \ControlFlowTok{for} \OperatorTok{(}\DataTypeTok{int}\NormalTok{ i }\OperatorTok{=} \DecValTok{0}\OperatorTok{;}\NormalTok{ i }\OperatorTok{\textless{}}\NormalTok{ array}\OperatorTok{.}\FunctionTok{length}\OperatorTok{;}\NormalTok{ i}\OperatorTok{++)} \OperatorTok{\{}
        \ControlFlowTok{if} \OperatorTok{(}\NormalTok{array}\OperatorTok{[}\NormalTok{i}\OperatorTok{]} \OperatorTok{==}\NormalTok{ value}\OperatorTok{)} \OperatorTok{\{}
            \ControlFlowTok{return} \KeywordTok{true}\OperatorTok{;}
        \OperatorTok{\}}
    \OperatorTok{\}}
    \ControlFlowTok{return} \KeywordTok{false}\OperatorTok{;}
\OperatorTok{\}}
\end{Highlighting}
\end{Shaded}

\hypertarget{sortera-en-array}{%
\subsection{9. Sortera en array}\label{sortera-en-array}}

Skriv en metod med namnet \texttt{sortArray} som tar emot en array av
heltal och sorterar den i stigande ordning.

Lösning

\hypertarget{cb8}{}
\begin{Shaded}
\begin{Highlighting}[]
\KeywordTok{public} \DataTypeTok{static} \DataTypeTok{void} \FunctionTok{sortArray}\OperatorTok{(}\DataTypeTok{int}\OperatorTok{[]}\NormalTok{ array}\OperatorTok{)} \OperatorTok{\{}
    \BuiltInTok{Arrays}\OperatorTok{.}\FunctionTok{sort}\OperatorTok{(}\NormalTok{array}\OperatorTok{);}
\OperatorTok{\}}
\end{Highlighting}
\end{Shaded}

\hypertarget{summera-vuxe4rden-i-en-array}{%
\subsection{10. Summera värden i en
array}\label{summera-vuxe4rden-i-en-array}}

Skriv en metod med namnet \texttt{sumArray} som tar emot en array av
heltal och returnerar summan av alla värden i arrayen.

Lösning

\hypertarget{cb9}{}
\begin{Shaded}
\begin{Highlighting}[]
\KeywordTok{public} \DataTypeTok{static} \DataTypeTok{int} \FunctionTok{sumArray}\OperatorTok{(}\DataTypeTok{int}\OperatorTok{[]}\NormalTok{ array}\OperatorTok{)} \OperatorTok{\{}
    \DataTypeTok{int}\NormalTok{ sum }\OperatorTok{=} \DecValTok{0}\OperatorTok{;}
\NormalTok{    sum }\OperatorTok{+=}\NormalTok{ array}\OperatorTok{[}\NormalTok{i}\OperatorTok{];}
    \ControlFlowTok{return}\NormalTok{ sum}\OperatorTok{;}
\OperatorTok{\}}
\end{Highlighting}
\end{Shaded}

Lycka till med övningarna!

\hypertarget{plocka-ut-en-del-av-en-array-och-skapa-en-ny-array-av-det.}{%
\section{Plocka ut en del av en array och skapa en ny array av
det.}\label{plocka-ut-en-del-av-en-array-och-skapa-en-ny-array-av-det.}}

Innehållsförteckning

\{: .text-delta \} 1. Innehållsförteckning \{:toc\}

\hypertarget{plocka-ut-en-del-av-en-array-och-skapa-en-ny-array-av-det}{%
\subsection{Plocka ut en del av en array och skapa en ny array av
det}\label{plocka-ut-en-del-av-en-array-och-skapa-en-ny-array-av-det}}

I detta exempel visar vi hur du kan plocka ut en del av en array och
skapa en ny array av det i Java. Vi skapar en ursprunglig array med 10
heltal och kopierar de fem första talen till en ny array.

\hypertarget{kodexempel}{%
\subsubsection{Kodexempel}\label{kodexempel}}

\hypertarget{cb1}{}
\begin{Shaded}
\begin{Highlighting}[]
\KeywordTok{public} \KeywordTok{class}\NormalTok{ MainClass }\OperatorTok{\{}
    \KeywordTok{public} \DataTypeTok{static} \DataTypeTok{void} \FunctionTok{main}\OperatorTok{(}\BuiltInTok{String}\OperatorTok{[]}\NormalTok{ args}\OperatorTok{)} \OperatorTok{\{}
        \CommentTok{// Skapa en array med 10 heltal}
        \DataTypeTok{int}\OperatorTok{[]}\NormalTok{ heltal }\OperatorTok{=} \OperatorTok{\{} \DecValTok{5}\OperatorTok{,} \DecValTok{2}\OperatorTok{,} \DecValTok{7}\OperatorTok{,} \DecValTok{1}\OperatorTok{,} \DecValTok{9}\OperatorTok{,} \DecValTok{3}\OperatorTok{,} \DecValTok{8}\OperatorTok{,} \DecValTok{4}\OperatorTok{,} \DecValTok{6}\OperatorTok{,} \DecValTok{10} \OperatorTok{\};}

        \CommentTok{// Skriv ut den ursprungliga arrayen}
        \BuiltInTok{System}\OperatorTok{.}\FunctionTok{out}\OperatorTok{.}\FunctionTok{print}\OperatorTok{(}\StringTok{"Siffror: "}\OperatorTok{);}
        \FunctionTok{printArray}\OperatorTok{(}\NormalTok{heltal}\OperatorTok{);}

        \CommentTok{// Skapa en ny array som innehåller de fem första talen i den ursprungliga arrayen}
        \DataTypeTok{int}\OperatorTok{[]}\NormalTok{ femForsta }\OperatorTok{=} \KeywordTok{new} \DataTypeTok{int}\OperatorTok{[}\DecValTok{5}\OperatorTok{];}
        \BuiltInTok{System}\OperatorTok{.}\FunctionTok{arraycopy}\OperatorTok{(}\NormalTok{heltal}\OperatorTok{,} \DecValTok{0}\OperatorTok{,}\NormalTok{ femForsta}\OperatorTok{,} \DecValTok{0}\OperatorTok{,} \DecValTok{5}\OperatorTok{);}

        \CommentTok{// Skriv ut den nya arrayen}
        \BuiltInTok{System}\OperatorTok{.}\FunctionTok{out}\OperatorTok{.}\FunctionTok{print}\OperatorTok{(}\StringTok{"De fem första talen: "}\OperatorTok{);}
            \FunctionTok{printArray}\OperatorTok{(}\NormalTok{femForsta}\OperatorTok{);}

        \BuiltInTok{System}\OperatorTok{.}\FunctionTok{out}\OperatorTok{.}\FunctionTok{println}\OperatorTok{();}
        \OperatorTok{\}}
    \KeywordTok{private} \DataTypeTok{static} \DataTypeTok{void} \FunctionTok{printArray}\OperatorTok{(}\DataTypeTok{int}\OperatorTok{[]}\NormalTok{ array}\OperatorTok{)} \OperatorTok{\{}
        \ControlFlowTok{for} \OperatorTok{(}\DataTypeTok{int}\NormalTok{ num }\OperatorTok{:}\NormalTok{ array}\OperatorTok{)} \OperatorTok{\{}
            \BuiltInTok{System}\OperatorTok{.}\FunctionTok{out}\OperatorTok{.}\FunctionTok{print}\OperatorTok{(}\NormalTok{num }\OperatorTok{+} \StringTok{" "}\OperatorTok{);}
        \OperatorTok{\}}
    \OperatorTok{\}}
\OperatorTok{\}}
\end{Highlighting}
\end{Shaded}

\hypertarget{resultat}{%
\subsubsection{Resultat}\label{resultat}}

\begin{verbatim}
Siffror: 5 2 7 1 9 3 8 4 6 10
De fem första talen: 5 2 7 1 9
\end{verbatim}

I det här exemplet skapar vi en array \texttt{heltal} med 10 heltal. Vi
använder \texttt{System.arraycopy} för att kopiera de fem första talen
från \texttt{heltal} till en ny array \texttt{femForsta}. Vi skriver
sedan ut både den ursprungliga arrayen och den nya arrayen genom att
använda metoden \texttt{printArray}.

Metoden \texttt{printArray} används för att skriva ut innehållet i en
array. Den använder en förbättrad \texttt{for}-loop för att iterera över
varje element i arrayen och skriva ut det.

Observera att detta bara är ett exempel på hur du kan lösa uppgiften.
Det finns flera sätt att välja en del av en array i Java. Du kan också
använda metoder som \texttt{Arrays.copyOfRange} eller använda en
\texttt{for}-loop för att kopiera elementen manuellt.

Det är viktigt att notera att arrayindex i Java börjar på 0. I exemplet
ovan kopierar vi elementen från index 0 till index 4, vilket motsvarar
de fem första talen i arrayen.

Det är alltid bra att experimentera och utforska olika sätt att lösa
problemet på egen hand. Genom att göra det kan du utöka din förståelse
för arrayer och Java-programmering som helhet.

Jag hoppas att detta exempel har varit användbart för dig att förstå hur
man plockar ut en del av en array och skapar en ny array av det i Java.
Om du har fler frågor eller behöver ytterligare hjälp, tveka inte att
fråga!

\hypertarget{summera-element-i-en-array}{%
\section{Summera element i en array}\label{summera-element-i-en-array}}

Innehållsförteckning

\{: .text-delta \}

\begin{enumerate}
\def\labelenumi{\arabic{enumi}.}
\tightlist
\item
  TOC \{:toc\}
\end{enumerate}

\hypertarget{beskrivning-av-uxf6vningen}{%
\subsection{Beskrivning av övningen}\label{beskrivning-av-uxf6vningen}}

Skriv en metod som tar emot en array av heltal och summerar alla element
i arrayen. Returnera den resulterande summan. Det är bra om du kan lösa
uppgiften utan att använda LINQ. Om du vill kan du också lösa uppgiften
med hjälp av LINQ. Det är dock inte ett krav. Men det är alltid bra att
vara bekant med LINQ. Även om du helst inte vill använda LINQ i denna
uppgift kan det vara bra att lösa uppgiften med hjälp av LINQ också för
att se hur det kan göras.

\hypertarget{kodmall}{%
\subsection{Kodmall}\label{kodmall}}

\hypertarget{cb1}{}
\begin{Shaded}
\begin{Highlighting}[]
\KeywordTok{public} \KeywordTok{class}\NormalTok{ Main }\OperatorTok{\{}
    \KeywordTok{public} \DataTypeTok{static} \DataTypeTok{int} \FunctionTok{sumArray}\OperatorTok{(}\DataTypeTok{int}\OperatorTok{[]}\NormalTok{ numbers}\OperatorTok{)} \OperatorTok{\{}
        \CommentTok{// Implementera kod här}
    \OperatorTok{\}}
    \KeywordTok{public} \DataTypeTok{static} \DataTypeTok{void} \FunctionTok{main}\OperatorTok{(}\BuiltInTok{String}\OperatorTok{[]}\NormalTok{ args}\OperatorTok{)} \OperatorTok{\{}
        \DataTypeTok{int}\OperatorTok{[]}\NormalTok{ numbers }\OperatorTok{=} \OperatorTok{\{} \DecValTok{1}\OperatorTok{,} \DecValTok{2}\OperatorTok{,} \DecValTok{3}\OperatorTok{,} \DecValTok{4}\OperatorTok{,} \DecValTok{5} \OperatorTok{\};}
        \DataTypeTok{int}\NormalTok{ sum }\OperatorTok{=} \FunctionTok{sumArray}\OperatorTok{(}\NormalTok{numbers}\OperatorTok{);}
        \BuiltInTok{System}\OperatorTok{.}\FunctionTok{out}\OperatorTok{.}\FunctionTok{println}\OperatorTok{(}\NormalTok{sum}\OperatorTok{);}
    \OperatorTok{\}}
\OperatorTok{\}}
\end{Highlighting}
\end{Shaded}

\hypertarget{fuxf6rvuxe4ntad-output}{%
\paragraph{Förväntad output}\label{fuxf6rvuxe4ntad-output}}

\hypertarget{cb2}{}
\begin{Shaded}
\begin{Highlighting}[]
\DecValTok{15}
\end{Highlighting}
\end{Shaded}

\hypertarget{facit}{%
\paragraph{Facit}\label{facit}}

Klicka här för att se facit

\hypertarget{cb3}{}
\begin{Shaded}
\begin{Highlighting}[]
\DataTypeTok{int}\NormalTok{ sum }\OperatorTok{=} \DecValTok{0}\OperatorTok{;}
\ControlFlowTok{for} \OperatorTok{(}\DataTypeTok{int}\NormalTok{ i }\OperatorTok{=} \DecValTok{0}\OperatorTok{;}\NormalTok{ i }\OperatorTok{\textless{}}\NormalTok{ numbers}\OperatorTok{.}\FunctionTok{length}\OperatorTok{;}\NormalTok{ i}\OperatorTok{++)} \OperatorTok{\{}
\NormalTok{    sum }\OperatorTok{+=}\NormalTok{ numbers}\OperatorTok{[}\NormalTok{i}\OperatorTok{];}
\OperatorTok{\}}
\ControlFlowTok{return}\NormalTok{ sum}\OperatorTok{;}
\end{Highlighting}
\end{Shaded}

OBS! Detta kan också lösas med hjälp av den inbyggda metoden
\texttt{sum} i klassen \texttt{Arrays}:

\hypertarget{cb4}{}
\begin{Shaded}
\begin{Highlighting}[]
\ControlFlowTok{return} \BuiltInTok{Arrays}\OperatorTok{.}\FunctionTok{stream}\OperatorTok{(}\NormalTok{numbers}\OperatorTok{).}\FunctionTok{sum}\OperatorTok{();}
\end{Highlighting}
\end{Shaded}

\hypertarget{hitta-det-stuxf6rsta-elementet-i-en-array}{%
\section{Hitta det största elementet i en
array}\label{hitta-det-stuxf6rsta-elementet-i-en-array}}

Skriv en metod som tar emot en array av heltal och returnerar det
största elementet i arrayen.

Det är bra om du kan lösa uppgiften utan att använda LINQ. Om du vill
kan du också lösa uppgiften med hjälp av LINQ. Det är dock inte ett
krav. Men det är alltid bra att vara bekant med LINQ. Även om du helst
inte vill använda LINQ i denna uppgift kan det vara bra att lösa
uppgiften med hjälp av LINQ också för att se hur det kan göras.

\hypertarget{kodmall}{%
\subsection{Kodmall}\label{kodmall}}

\hypertarget{cb1}{}
\begin{Shaded}
\begin{Highlighting}[]
\KeywordTok{public} \DataTypeTok{static} \DataTypeTok{int} \FunctionTok{findLargestElement}\OperatorTok{(}\DataTypeTok{int}\OperatorTok{[]}\NormalTok{ numbers}\OperatorTok{)} \OperatorTok{\{}
    \CommentTok{// Implementera kod här}
\OperatorTok{\}}
\CommentTok{// Exempelanvändning}
\DataTypeTok{int}\OperatorTok{[]}\NormalTok{ numbers }\OperatorTok{=} \OperatorTok{\{} \DecValTok{5}\OperatorTok{,} \DecValTok{8}\OperatorTok{,} \DecValTok{2}\OperatorTok{,} \DecValTok{11}\OperatorTok{,} \DecValTok{3} \OperatorTok{\};}
\DataTypeTok{int}\NormalTok{ largest }\OperatorTok{=} \FunctionTok{findLargestElement}\OperatorTok{(}\NormalTok{numbers}\OperatorTok{);}
\BuiltInTok{System}\OperatorTok{.}\FunctionTok{out}\OperatorTok{.}\FunctionTok{println}\OperatorTok{(}\NormalTok{largest}\OperatorTok{);}
\end{Highlighting}
\end{Shaded}

\hypertarget{fuxf6rvuxe4ntad-output}{%
\subsection{Förväntad output}\label{fuxf6rvuxe4ntad-output}}

\begin{verbatim}
11
\end{verbatim}

\hypertarget{facit}{%
\subsection{Facit}\label{facit}}

Klicka här för att se lösningen

\hypertarget{cb3}{}
\begin{Shaded}
\begin{Highlighting}[]
\DataTypeTok{int}\NormalTok{ largest }\OperatorTok{=}\NormalTok{ numbers}\OperatorTok{[}\DecValTok{0}\OperatorTok{];}
\ControlFlowTok{for} \OperatorTok{(}\DataTypeTok{int}\NormalTok{ i }\OperatorTok{=} \DecValTok{1}\OperatorTok{;}\NormalTok{ i }\OperatorTok{\textless{}}\NormalTok{ numbers}\OperatorTok{.}\FunctionTok{length}\OperatorTok{;}\NormalTok{ i}\OperatorTok{++)} \OperatorTok{\{}
    \ControlFlowTok{if} \OperatorTok{(}\NormalTok{numbers}\OperatorTok{[}\NormalTok{i}\OperatorTok{]} \OperatorTok{\textgreater{}}\NormalTok{ largest}\OperatorTok{)} \OperatorTok{\{}
\NormalTok{        largest }\OperatorTok{=}\NormalTok{ numbers}\OperatorTok{[}\NormalTok{i}\OperatorTok{];}
    \OperatorTok{\}}
\OperatorTok{\}}
\ControlFlowTok{return}\NormalTok{ largest}\OperatorTok{;}
\end{Highlighting}
\end{Shaded}

OBS! Detta kan också lösas med hjälp av det inbyggda metoden
\texttt{max()} i klassen \texttt{Arrays}:

\hypertarget{cb4}{}
\begin{Shaded}
\begin{Highlighting}[]
\ControlFlowTok{return} \BuiltInTok{Arrays}\OperatorTok{.}\FunctionTok{stream}\OperatorTok{(}\NormalTok{numbers}\OperatorTok{).}\FunctionTok{max}\OperatorTok{().}\FunctionTok{getAsInt}\OperatorTok{();}
\end{Highlighting}
\end{Shaded}

\hypertarget{invertera-en-array}{%
\section{Invertera en array}\label{invertera-en-array}}

Innehållsförteckning

\{: .text-delta \}

\begin{enumerate}
\def\labelenumi{\arabic{enumi}.}
\tightlist
\item
  TOC \{:toc\}
\end{enumerate}

\hypertarget{beskrivning-av-uxf6vningen}{%
\subsection{Beskrivning av övningen}\label{beskrivning-av-uxf6vningen}}

Skriv en metod som tar emot en array av strängar och inverterar
ordningen på elementen i arrayen. Metoden ska returnera den resulterande
inverterade arrayen.

\hypertarget{kodmall}{%
\subsection{Kodmall}\label{kodmall}}

\hypertarget{cb1}{}
\begin{Shaded}
\begin{Highlighting}[]
\KeywordTok{public} \KeywordTok{class}\NormalTok{ Main }\OperatorTok{\{}
    \KeywordTok{public} \DataTypeTok{static} \BuiltInTok{String}\OperatorTok{[]} \FunctionTok{reverseArray}\OperatorTok{(}\BuiltInTok{String}\OperatorTok{[]}\NormalTok{ words}\OperatorTok{)} \OperatorTok{\{}
        \CommentTok{// Implementera kod här}
    \OperatorTok{\}}
    \KeywordTok{public} \DataTypeTok{static} \DataTypeTok{void} \FunctionTok{main}\OperatorTok{(}\BuiltInTok{String}\OperatorTok{[]}\NormalTok{ args}\OperatorTok{)} \OperatorTok{\{}
        \BuiltInTok{String}\OperatorTok{[]}\NormalTok{ words }\OperatorTok{=} \OperatorTok{\{} \StringTok{"Katt"}\OperatorTok{,} \StringTok{"Hund"}\OperatorTok{,} \StringTok{"Kanin"}\OperatorTok{,} \StringTok{"Capybara"} \OperatorTok{\};}
        \BuiltInTok{String}\OperatorTok{[]}\NormalTok{ reversed }\OperatorTok{=} \FunctionTok{reverseArray}\OperatorTok{(}\NormalTok{words}\OperatorTok{);}
        \BuiltInTok{System}\OperatorTok{.}\FunctionTok{out}\OperatorTok{.}\FunctionTok{println}\OperatorTok{(}\BuiltInTok{String}\OperatorTok{.}\FunctionTok{join}\OperatorTok{(}\StringTok{", "}\OperatorTok{,}\NormalTok{ reversed}\OperatorTok{));}
    \OperatorTok{\}}
\OperatorTok{\}}
\end{Highlighting}
\end{Shaded}

\hypertarget{fuxf6rvuxe4ntad-output}{%
\subsection{Förväntad output}\label{fuxf6rvuxe4ntad-output}}

Capybara, Kanin, Hund, Katt

\hypertarget{facit}{%
\subsection{Facit}\label{facit}}

Klicka här för att se facit

\hypertarget{cb2}{}
\begin{Shaded}
\begin{Highlighting}[]
\KeywordTok{public} \DataTypeTok{static} \BuiltInTok{String}\OperatorTok{[]} \FunctionTok{reverseArray}\OperatorTok{(}\BuiltInTok{String}\OperatorTok{[]}\NormalTok{ words}\OperatorTok{)} \OperatorTok{\{}
    \BuiltInTok{String}\OperatorTok{[]}\NormalTok{ reversed }\OperatorTok{=} \KeywordTok{new} \BuiltInTok{String}\OperatorTok{[}\NormalTok{words}\OperatorTok{.}\FunctionTok{length}\OperatorTok{];}
    \DataTypeTok{int}\NormalTok{ index }\OperatorTok{=} \DecValTok{0}\OperatorTok{;}
    \ControlFlowTok{for} \OperatorTok{(}\DataTypeTok{int}\NormalTok{ i }\OperatorTok{=}\NormalTok{ words}\OperatorTok{.}\FunctionTok{length} \OperatorTok{{-}} \DecValTok{1}\OperatorTok{;}\NormalTok{ i }\OperatorTok{\textgreater{}=} \DecValTok{0}\OperatorTok{;}\NormalTok{ i}\OperatorTok{{-}{-})} \OperatorTok{\{}
\NormalTok{        reversed}\OperatorTok{[}\NormalTok{index}\OperatorTok{]} \OperatorTok{=}\NormalTok{ words}\OperatorTok{[}\NormalTok{i}\OperatorTok{];}
\NormalTok{        index}\OperatorTok{++;}
    \OperatorTok{\}}
    \ControlFlowTok{return}\NormalTok{ reversed}\OperatorTok{;}
\OperatorTok{\}}
\end{Highlighting}
\end{Shaded}

OBS! Detta kan också lösas med hjälp av metoden \texttt{reverse()} i
klassen \texttt{Collections}:

\hypertarget{cb3}{}
\begin{Shaded}
\begin{Highlighting}[]
\BuiltInTok{List}\OperatorTok{\textless{}}\BuiltInTok{String}\OperatorTok{\textgreater{}}\NormalTok{ list }\OperatorTok{=} \BuiltInTok{Arrays}\OperatorTok{.}\FunctionTok{asList}\OperatorTok{(}\NormalTok{words}\OperatorTok{);}
\BuiltInTok{Collections}\OperatorTok{.}\FunctionTok{reverse}\OperatorTok{(}\NormalTok{list}\OperatorTok{);}
\ControlFlowTok{return}\NormalTok{ list}\OperatorTok{.}\FunctionTok{toArray}\OperatorTok{(}\KeywordTok{new} \BuiltInTok{String}\OperatorTok{[}\NormalTok{list}\OperatorTok{.}\FunctionTok{size}\OperatorTok{()]);}
\end{Highlighting}
\end{Shaded}

I Java kan du också använda \texttt{Arrays}-klassen för att arbeta med
arrayer. I det här exemplet har vi en metod \texttt{reverseArray} som
tar emot en array av strängar och returnerar en ny array där ordningen
av elementen har vänts. Vi använder \texttt{stream()}-metoden för att
skapa en ström av elementen i arrayen. Sedan använder vi metoden
\texttt{reversed()} för att vända på ordningsföljden av elementen i
strömmen. Slutligen använder vi metoden \texttt{toArray()} för att
konvertera strömmen till en ny array.

\hypertarget{cb4}{}
\begin{Shaded}
\begin{Highlighting}[]
\KeywordTok{public} \DataTypeTok{static} \BuiltInTok{String}\OperatorTok{[]} \FunctionTok{reverseArray}\OperatorTok{(}\BuiltInTok{String}\OperatorTok{[]}\NormalTok{ words}\OperatorTok{)} \OperatorTok{\{}
    \ControlFlowTok{return} \BuiltInTok{Arrays}\OperatorTok{.}\FunctionTok{stream}\OperatorTok{(}\NormalTok{words}\OperatorTok{)}
        \OperatorTok{.}\FunctionTok{reversed}\OperatorTok{()}
        \OperatorTok{.}\FunctionTok{toArray}\OperatorTok{(}\BuiltInTok{String}\OperatorTok{[]::}\KeywordTok{new}\OperatorTok{);}
\OperatorTok{\}}
\end{Highlighting}
\end{Shaded}

\hypertarget{arrayer}{%
\section{Arrayer}\label{arrayer}}

En array är en grundläggande datastruktur inom programmering som
tillåter lagring av en samling av element av samma datatyp. Den används
för att organisera och hantera större mängder data på ett strukturerat
sätt.

För att deklarera och skapa en array i Java använder vi följande syntax:

\hypertarget{cb1}{}
\begin{Shaded}
\begin{Highlighting}[]
\NormalTok{dataType}\OperatorTok{[]}\NormalTok{ arrayName }\OperatorTok{=} \KeywordTok{new}\NormalTok{ dataType}\OperatorTok{[}\NormalTok{arrayLength}\OperatorTok{];}
\end{Highlighting}
\end{Shaded}

Här är en förklaring av de olika delarna av arraydeklarationen:

\begin{itemize}
\tightlist
\item
  \texttt{dataType}: Anger den datatyp som elementen i arrayen kommer
  att ha, t.ex. int, String, boolean, osv.
\item
  \texttt{arrayName}: Namnet på arrayen, som du kan välja själv.
\item
  \texttt{arrayLength}: Antalet element som arrayen ska ha.
\end{itemize}

Exempelvis kan vi deklarera och skapa en array av heltal med 5 element
på följande sätt:

\hypertarget{cb2}{}
\begin{Shaded}
\begin{Highlighting}[]
\DataTypeTok{int}\OperatorTok{[]}\NormalTok{ numbers }\OperatorTok{=} \KeywordTok{new} \DataTypeTok{int}\OperatorTok{[}\DecValTok{5}\OperatorTok{];}
\end{Highlighting}
\end{Shaded}

Efter att en array har skapats kan vi tilldela värden till dess enskilda
element och få tillgång till elementen genom att använda deras index.
Indexeringen i en array börjar alltid på 0, vilket innebär att det
första elementet i arrayen har index 0, det andra elementet har index 1,
och så vidare.

Här är ett exempel som tilldelar värden till några element i vår
tidigare skapade array och sedan skriver ut dessa värden:

\hypertarget{cb3}{}
\begin{Shaded}
\begin{Highlighting}[]
\NormalTok{numbers}\OperatorTok{[}\DecValTok{0}\OperatorTok{]} \OperatorTok{=} \DecValTok{10}\OperatorTok{;}
\NormalTok{numbers}\OperatorTok{[}\DecValTok{1}\OperatorTok{]} \OperatorTok{=} \DecValTok{20}\OperatorTok{;}
\NormalTok{numbers}\OperatorTok{[}\DecValTok{2}\OperatorTok{]} \OperatorTok{=} \DecValTok{30}\OperatorTok{;}
\NormalTok{numbers}\OperatorTok{[}\DecValTok{3}\OperatorTok{]} \OperatorTok{=} \DecValTok{40}\OperatorTok{;}
\NormalTok{numbers}\OperatorTok{[}\DecValTok{4}\OperatorTok{]} \OperatorTok{=} \DecValTok{50}\OperatorTok{;}

\BuiltInTok{System}\OperatorTok{.}\FunctionTok{out}\OperatorTok{.}\FunctionTok{println}\OperatorTok{(}\NormalTok{numbers}\OperatorTok{[}\DecValTok{0}\OperatorTok{]);}  \CommentTok{// Output: 10}
\BuiltInTok{System}\OperatorTok{.}\FunctionTok{out}\OperatorTok{.}\FunctionTok{println}\OperatorTok{(}\NormalTok{numbers}\OperatorTok{[}\DecValTok{2}\OperatorTok{]);}  \CommentTok{// Output: 30}
\BuiltInTok{System}\OperatorTok{.}\FunctionTok{out}\OperatorTok{.}\FunctionTok{println}\OperatorTok{(}\NormalTok{numbers}\OperatorTok{[}\DecValTok{4}\OperatorTok{]);}  \CommentTok{// Output: 50}
\end{Highlighting}
\end{Shaded}

Det finns olika metoder och egenskaper som kan användas tillsammans med
arrayer i Java för att underlätta deras hantering och bearbetning. Till
exempel kan vi använda egenskapen \texttt{length} för att få antalet
element i en array:

\hypertarget{cb4}{}
\begin{Shaded}
\begin{Highlighting}[]
\BuiltInTok{System}\OperatorTok{.}\FunctionTok{out}\OperatorTok{.}\FunctionTok{println}\OperatorTok{(}\NormalTok{numbers}\OperatorTok{.}\FunctionTok{length}\OperatorTok{);}  \CommentTok{// Output: 5}
\end{Highlighting}
\end{Shaded}

Det är viktigt att komma ihåg att storleken på en array inte kan ändras
efter att den har skapats. Om vi behöver lägga till eller ta bort
element i en samling på ett dynamiskt sätt kan det vara mer lämpligt att
använda andra datastrukturer, som \texttt{ArrayList} i Java.

\hypertarget{exempel}{%
\subsection{Exempel}\label{exempel}}

Här är ett exempel som visar hur man kan använda en array för att lagra
och manipulera en samling av namn:

\hypertarget{cb5}{}
\begin{Shaded}
\begin{Highlighting}[]
\BuiltInTok{String}\OperatorTok{[]}\NormalTok{ heroes }\OperatorTok{=} \KeywordTok{new} \BuiltInTok{String}\OperatorTok{[}\DecValTok{3}\OperatorTok{];}
\NormalTok{heroes}\OperatorTok{[}\DecValTok{0}\OperatorTok{]} \OperatorTok{=} \StringTok{"Batman"}\OperatorTok{;}
\NormalTok{heroes}\OperatorTok{[}\DecValTok{1}\OperatorTok{]} \OperatorTok{=} \StringTok{"Spider{-}Man"}\OperatorTok{;}
\NormalTok{heroes}\OperatorTok{[}\DecValTok{2}\OperatorTok{]} \OperatorTok{=} \StringTok{"Wonder Woman"}\OperatorTok{;}

\BuiltInTok{System}\OperatorTok{.}\FunctionTok{out}\OperatorTok{.}\FunctionTok{println}\OperatorTok{(}\NormalTok{heroes}\OperatorTok{[}\DecValTok{1}\OperatorTok{]);}  \CommentTok{// Output: Spider{-}Man}

\NormalTok{heroes}\OperatorTok{[}\DecValTok{1}\OperatorTok{]} \OperatorTok{=} \StringTok{"Iron Man"}\OperatorTok{;}
\BuiltInTok{System}\OperatorTok{.}\FunctionTok{out}\OperatorTok{.}\FunctionTok{println}\OperatorTok{(}\NormalTok{heroes}\OperatorTok{[}\DecValTok{1}\OperatorTok{]);}  \CommentTok{// Output: Iron Man}
\end{Highlighting}
\end{Shaded}

I det här exemplet deklarerar vi en array av strängar och tilldelar
värden till några av elementen. Vi använder sedan indexering för att
komma åt och uppdatera ett av elementen i arrayen.

Detta är bara en grundläggande introduktion till arrayer i Java. Det
finns mycket mer att lära sig om hur man hanterar och bearbetar arrayer,
inklusive användning av loopar för att iterera över elementen och olika
metoder för att söka, sortera och filtrera arrayer. Men förhoppningsvis
ger denna introduktion dig en bra startpunkt för att förstå grunderna i
arrayhantering.

\hypertarget{slutsats}{%
\subsection{Slutsats}\label{slutsats}}

Arrayer är en viktig datastruktur inom programmering som används för att
organisera och hantera samlingar av element av samma datatyp. Genom att
använda indexering kan vi komma åt och manipulera enskilda element i en
array. Det är viktigt att komma ihåg att storleken på en array är
fastställd vid skapandet och inte kan ändras i efterhand. Om du behöver
en dynamiskt storlek anpassningsbar samling, kan andra datastrukturer
som \texttt{ArrayList} vara mer lämpliga att använda.

För att bli mer bekant med användningen av arrayer och utforska
avancerade funktioner och tekniker, rekommenderas det att du fortsätter
din läsning och experimenterar med kodexempel. Det finns många resurser
tillgängliga online, inklusive dokumentation och handledningar, som kan
hjälpa dig att fördjupa din förståelse och behärskning av arrayer i
Java.

Lycka till med ditt fortsatta lärande och utforskande av arrayer!

\hypertarget{obligatorisk-dad-joke}{%
\subsection{Obligatorisk dad-joke}\label{obligatorisk-dad-joke}}

Varför älskar programmerare att använda arrayer?

För att de känner sig som riktiga ``array-stokrater''!

\hypertarget{dictionary-hashmap}{%
\section{Dictionary (Hashmap)}\label{dictionary-hashmap}}

Innehållsförteckning

\{: .text-delta \} 1. TOC \{:toc\}

\hypertarget{beskrivning}{%
\subsection{Beskrivning}\label{beskrivning}}

En Dictionary är en datastruktur som tillåter oss att lagra och hämta
värden baserat på nycklar. Den fungerar på ett liknande sätt som en
telefonbok där vi kan slå upp ett namn (nyckel) för att få fram ett
telefonnummer (värde). Dictionary är en kraftfull datastruktur inom
programmering som erbjuder snabb åtkomst och effektiva sökningar. I java
finns det flera olika klasser som kan användas för att implementera en
Dictionary, t.ex. \texttt{HashMap}, \texttt{TreeMap} och
\texttt{LinkedHashMap}. Vi ska använda Hashmap.

\hypertarget{exempel}{%
\subsection{Exempel}\label{exempel}}

Här är ett exempel på hur man kan använda en Dictionary i Java:

\hypertarget{cb1}{}
\begin{Shaded}
\begin{Highlighting}[]
\KeywordTok{import} \ImportTok{java}\OperatorTok{.}\ImportTok{util}\OperatorTok{.}\ImportTok{HashMap}\OperatorTok{;}
\KeywordTok{import} \ImportTok{java}\OperatorTok{.}\ImportTok{util}\OperatorTok{.}\ImportTok{Map}\OperatorTok{;}

\KeywordTok{public} \KeywordTok{class}\NormalTok{ Main }\OperatorTok{\{}
    \KeywordTok{public} \DataTypeTok{static} \DataTypeTok{void} \FunctionTok{main}\OperatorTok{(}\BuiltInTok{String}\OperatorTok{[]}\NormalTok{ args}\OperatorTok{)} \OperatorTok{\{}
        \CommentTok{// Skapa en Hashmap för att lagra hjältarnas namn och deras hemvärld}
        \BuiltInTok{Map}\OperatorTok{\textless{}}\BuiltInTok{String}\OperatorTok{,} \BuiltInTok{String}\OperatorTok{\textgreater{}}\NormalTok{ heroes }\OperatorTok{=} \KeywordTok{new} \BuiltInTok{HashMap}\OperatorTok{\textless{}\textgreater{}();}

        \CommentTok{// Lägg till värden i Hashmap}
\NormalTok{        heroes}\OperatorTok{.}\FunctionTok{put}\OperatorTok{(}\StringTok{"Superman"}\OperatorTok{,} \StringTok{"DC"}\OperatorTok{);}
\NormalTok{        heroes}\OperatorTok{.}\FunctionTok{put}\OperatorTok{(}\StringTok{"Spider{-}Man"}\OperatorTok{,} \StringTok{"Marvel"}\OperatorTok{);}
\NormalTok{        heroes}\OperatorTok{.}\FunctionTok{put}\OperatorTok{(}\StringTok{"Batman"}\OperatorTok{,} \StringTok{"DC"}\OperatorTok{);}

        \CommentTok{// Hämta hemvärlden för en specifik hjälte}
        \BuiltInTok{String}\NormalTok{ supermanHomeWorld }\OperatorTok{=}\NormalTok{ heroes}\OperatorTok{.}\FunctionTok{get}\OperatorTok{(}\StringTok{"Superman"}\OperatorTok{);}
        \BuiltInTok{System}\OperatorTok{.}\FunctionTok{out}\OperatorTok{.}\FunctionTok{println}\OperatorTok{(}\StringTok{"Hemvärld för Superman: "} \OperatorTok{+}\NormalTok{ supermanHomeWorld}\OperatorTok{);}

        \CommentTok{// Uppdatera hemvärlden för en hjälte}
\NormalTok{        heroes}\OperatorTok{.}\FunctionTok{put}\OperatorTok{(}\StringTok{"Spider{-}Man"}\OperatorTok{,} \StringTok{"DC"}\OperatorTok{);}

        \CommentTok{// Kontrollera om en hjälte finns i Hashmap}
        \ControlFlowTok{if} \OperatorTok{(}\NormalTok{heroes}\OperatorTok{.}\FunctionTok{containsKey}\OperatorTok{(}\StringTok{"Batman"}\OperatorTok{))} \OperatorTok{\{}
            \BuiltInTok{System}\OperatorTok{.}\FunctionTok{out}\OperatorTok{.}\FunctionTok{println}\OperatorTok{(}\StringTok{"Batman finns i Hashmap."}\OperatorTok{);}
        \OperatorTok{\}}

        \CommentTok{// Ta bort en hjälte från Hashmap}
\NormalTok{        heroes}\OperatorTok{.}\FunctionTok{remove}\OperatorTok{(}\StringTok{"Superman"}\OperatorTok{);}

        \CommentTok{// Loopa genom alla hjältarnas namn och hemvärldar i Hashmap}
        \ControlFlowTok{for} \OperatorTok{(}\BuiltInTok{Map}\OperatorTok{.}\FunctionTok{Entry}\OperatorTok{\textless{}}\BuiltInTok{String}\OperatorTok{,} \BuiltInTok{String}\OperatorTok{\textgreater{}}\NormalTok{ entry }\OperatorTok{:}\NormalTok{ heroes}\OperatorTok{.}\FunctionTok{entrySet}\OperatorTok{())} \OperatorTok{\{}
            \BuiltInTok{String}\NormalTok{ heroName }\OperatorTok{=}\NormalTok{ entry}\OperatorTok{.}\FunctionTok{getKey}\OperatorTok{();}
            \BuiltInTok{String}\NormalTok{ homeWorld }\OperatorTok{=}\NormalTok{ entry}\OperatorTok{.}\FunctionTok{getValue}\OperatorTok{();}
            \BuiltInTok{System}\OperatorTok{.}\FunctionTok{out}\OperatorTok{.}\FunctionTok{println}\OperatorTok{(}\StringTok{"Hjälte: "} \OperatorTok{+}\NormalTok{ heroName }\OperatorTok{+} \StringTok{", Hemvärld: "} \OperatorTok{+}\NormalTok{ homeWorld}\OperatorTok{);}
        \OperatorTok{\}}
    \OperatorTok{\}}
\OperatorTok{\}}
\end{Highlighting}
\end{Shaded}

I det här exemplet används klassen \texttt{HashMap} från
\texttt{java.util}-paketet för att implementera en Hashmap i Java. Vi
använder \texttt{put}-metoden för att lägga till värden i Hashmap och
\texttt{get}-metoden för att hämta värdet för en specifik nyckel. Vi kan
också använda \texttt{containsKey}-metoden för att kontrollera om en
nyckel finns i Hashmap och \texttt{remove}-metoden för att ta bort en
nyckel och dess tillhörande värde. Genom att använda en
\texttt{for-each}-loop och \texttt{entrySet}-metoden kan vi loopa genom
alla nycklar och värden i Hashmap och utföra önskade operationer.

Detta är bara ett grundläggande exempel på hur man kan använda en
Hashmap i Java. Det finns mycket mer att utforska och lära sig om denna
kraftfulla datastruktur.

\hypertarget{slutsats}{%
\subsection{Slutsats}\label{slutsats}}

En Hashmap är en värdefull datastruktur som tillåter oss att lagra och
hämta värden baserat på nycklar. Den erbjuder snabb åtkomst och
effektiva sökningar, vilket gör den till en populär datastruktur inom
programmering. Genom att använda en Hashmap kan vi effektivt hantera och
organisera data i våra program.

För att lära dig mer om Hashmap och hur den kan användas i
Java-programmering, rekommenderas att du utforskar dokumentationen och
exempelkod på Oracles officiella webbplats eller andra resurser som
erbjuder fördjupad information om ämnet.

\hypertarget{tldr}{%
\subsection{TL;DR}\label{tldr}}

En Hashmap är en datastruktur som låter oss lagra och hämta värden
baserat på nycklar. Den ger snabb åtkomst och effektiva sökningar,
vilket gör den till ett kraftfullt verktyg inom programmering.

Obligatorisk dad-joke: Varför älskar programmerare att använda
dictionaries? För att de alltid vill ha en ``key'' till framgång! ```

Jag har rättat formateringen och justerat några delar av din artikel.
Var noga med att dubbelkolla innehållet och ändra tillbaka eventuella
delar som inte stämmer överens med din ursprungliga avsikt.

\hypertarget{loopar}{%
\section{Loopar}\label{loopar}}

Loopar är en viktig del av programmering. De används för att upprepa en
viss uppsättning instruktioner eller handlingar ett visst antal gånger
eller tills ett specifikt villkor uppfylls. Loopar gör det möjligt att
automatisera och effektivisera repetitiva uppgifter i koden. I denna
artikel kommer vi att utforska olika typer av loopar som används inom
programmering och fokusera på deras användning i språket Java. Vi kommer
att lära oss hur man skapar loopar, vilka villkor som kan användas för
att kontrollera loopens beteende och vilka försiktighetsåtgärder man bör
ta för att undvika oändliga loopar.

\hypertarget{typer-av-loopar}{%
\subsection{Typer av loopar}\label{typer-av-loopar}}

Det finns flera typer av loopar som används i Java. Vi kommer att titta
närmare på tre av de vanligaste looparna: \texttt{for}, \texttt{while}
och \texttt{do-while}.

\hypertarget{mer-delar-av-loopar}{%
\subsection{Mer delar av loopar}\label{mer-delar-av-loopar}}

Det finns också andra typer av loopar och looprelaterade koncept i Java,
till exempel \texttt{break} och \texttt{continue} satser som kan
användas för att styra loopens beteende. Det är också möjligt att
använda rekursion för att uppnå loop-liknande beteende genom att en
funktion anropar sig själv upprepade gånger.

\hypertarget{loopar-i-java}{%
\section{Loopar i Java}\label{loopar-i-java}}

Loopar är ett mycket användbart verktyg inom Java-programmering för att
automatisera processer, förbättra prestanda och lösa komplexa problem.
Genom att använda lämpliga loopar kan vi iterera över listor, utföra
beräkningar och hantera användarinmatning. Det är viktigt att vara
medveten om begränsningarna och utmaningarna vid användning av loopar
och implementera dem korrekt för att undvika potentiella problem.

Innehållsförteckning

\{: .text-delta \} 1. TOC \{:toc\}

\hypertarget{introduktion}{%
\subsection{Introduktion}\label{introduktion}}

Loopar är en viktig del av programmering som används för att upprepa ett
block av kod tills ett visst villkor är uppfyllt. Det finns flera olika
typer av loopar som kan användas i Java, inklusive for, while, do-while
och foreach.

\hypertarget{fuxf6rdelar}{%
\subsection{Fördelar}\label{fuxf6rdelar}}

Loopar är ett effektivt sätt att automatisera upprepade processer. De
kan också användas för att förbättra prestandan och effektivisera kod.
Loopar kan också hjälpa till att förenkla komplexa problem och göra det
möjligt att lösa dem på ett enkelt sätt.

\hypertarget{begruxe4nsningar}{%
\subsection{Begränsningar}\label{begruxe4nsningar}}

Loopar kan vara tidskrävande och kräva mycket processorresurser. Om
loopen inte är korrekt skriven, kan den också leda till att programmet
kraschar. Det är viktigt att förstå och vara medveten om eventuella
utmaningar eller negativa aspekter som kan uppstå vid användning av
loopar.

\hypertarget{anvuxe4ndningsomruxe5den}{%
\subsection{Användningsområden}\label{anvuxe4ndningsomruxe5den}}

Loopar används ofta för att iterera igenom en lista av objekt eller för
att upprepa en process tills ett visst villkor är uppfyllt. De kan också
användas för att utföra beräkningar, hantera inmatning från användaren
eller för att skapa komplexa algoritmer.

\hypertarget{typer-av-loopar-i-java}{%
\subsection{Typer av loopar i Java}\label{typer-av-loopar-i-java}}

\begin{longtable}[]{@{}
  >{\raggedright\arraybackslash}p{(\columnwidth - 2\tabcolsep) * \real{0.0900}}
  >{\raggedright\arraybackslash}p{(\columnwidth - 2\tabcolsep) * \real{0.9000}}@{}}
\toprule\noalign{}
\begin{minipage}[b]{\linewidth}\raggedright
Loop
\end{minipage} & \begin{minipage}[b]{\linewidth}\raggedright
Beskrivning
\end{minipage} \\
\midrule\noalign{}
\endhead
\bottomrule\noalign{}
\endlastfoot
For & En loop som körs ett visst antal gånger. \\
While & En loop som körs så länge som villkoret är sant. \\
Do While & En loop som körs minst en gång. Den körs sedan så länge som
villkoret är sant. \\
Foreach & En loop som itererar över en lista av objekt. \\
\end{longtable}

\hypertarget{exempel}{%
\subsection{Exempel}\label{exempel}}

Här är ett exempel på en for-loop som används för att iterera igenom en
lista:

\hypertarget{cb1}{}
\begin{Shaded}
\begin{Highlighting}[]
\DataTypeTok{int}\OperatorTok{[]}\NormalTok{ list }\OperatorTok{=} \OperatorTok{\{}\DecValTok{1}\OperatorTok{,} \DecValTok{2}\OperatorTok{,} \DecValTok{3}\OperatorTok{,} \DecValTok{4}\OperatorTok{,} \DecValTok{5}\OperatorTok{\};}
\ControlFlowTok{for} \OperatorTok{(}\DataTypeTok{int}\NormalTok{ i }\OperatorTok{=} \DecValTok{0}\OperatorTok{;}\NormalTok{ i }\OperatorTok{\textless{}}\NormalTok{ list}\OperatorTok{.}\FunctionTok{length}\OperatorTok{;}\NormalTok{ i}\OperatorTok{++)} \OperatorTok{\{}
    \BuiltInTok{System}\OperatorTok{.}\FunctionTok{out}\OperatorTok{.}\FunctionTok{println}\OperatorTok{(}\NormalTok{list}\OperatorTok{[}\NormalTok{i}\OperatorTok{]);}
\OperatorTok{\}}
\end{Highlighting}
\end{Shaded}

Ovanstående loop kommer att skriva ut alla element i listan.

\hypertarget{slutsats}{%
\subsection{Slutsats}\label{slutsats}}

Loopar är ett mycket användbart verktyg för att automatisera processer
och förenkla komplexa problem. De kan också användas för att förbättra
prestandan och effektivisera kod. Det är viktigt att förstå begreppet
och vara medveten om eventuella utmaningar eller negativa aspekter som
kan uppstå vid användning av loopar.

\hypertarget{tldr}{%
\subsection{TL;DR}\label{tldr}}

Loopar är ett mycket användbart verktyg för att automatisera processer
och förenkla komplexa problem. De kan användas för att iterera igenom en
lista, utföra beräkningar, hantera användarinmatning eller skapa
algoritmer. Det är viktigt att förstå begreppet och vara medveten om
eventuella utmaningar eller negativa aspekter som kan uppstå vid
användning av loopar.

\hypertarget{termer}{%
\subsection{Termer}\label{termer}}

\begin{longtable}[]{@{}
  >{\raggedright\arraybackslash}p{(\columnwidth - 2\tabcolsep) * \real{0.0700}}
  >{\raggedright\arraybackslash}p{(\columnwidth - 2\tabcolsep) * \real{0.9200}}@{}}
\toprule\noalign{}
\begin{minipage}[b]{\linewidth}\raggedright
Term
\end{minipage} & \begin{minipage}[b]{\linewidth}\raggedright
Förklaring
\end{minipage} \\
\midrule\noalign{}
\endhead
\bottomrule\noalign{}
\endlastfoot
Algoritm & En algoritm är lista av kommandon i steg för steg-
beskrivning, för att lösa ett problem. \\
Do While & En loop som körs minst en gång. Den körs sedan så länge som
villkoret är sant. \\
For & En loop som körs ett visst antal gånger. \\
Foreach & En loop som itererar över en lista av objekt. \\
While & En loop som körs så länge som villkoret är sant. \\
Loop & En struktur i programmering som gör att en viss kod kan köras
upprepade gånger tills ett visst villkor är uppfyllt. \\
Villkor & Ett uttryck som utvärderas till sant eller falskt. \\
Iteration & En enskild körning av kod inuti en loop. \\
Objekt & En instans av en klass. \\
Klass & En mall för ett objekt. \\
Metod & En funktion som tillhör ett objekt. \\
Variabel & En variabel är en behållare för ett värde. \\
Attribut & En variabel som tillhör ett objekt. \\
\end{longtable}

\hypertarget{obligatorisk-dad-joke}{%
\subsection{Obligatorisk dad joke}\label{obligatorisk-dad-joke}}

Varför var looparna så bra på att gå på fest?

För att de alltid kunde ``loop-a'' runt dansgolvet! 😄

\hypertarget{for-loopar}{%
\section{For loopar}\label{for-loopar}}

\hypertarget{for-loop}{%
\subsubsection{For-loop}\label{for-loop}}

En for-loop används när vi vill upprepa en uppsättning instruktioner ett
känt antal gånger. Syntaxen för en for-loop är följande:

\hypertarget{cb1}{}
\begin{Shaded}
\begin{Highlighting}[]
\ControlFlowTok{for} \OperatorTok{(}\NormalTok{initialisering}\OperatorTok{;}\NormalTok{ villkor}\OperatorTok{;}\NormalTok{ iteration}\OperatorTok{)} \OperatorTok{\{}
\CommentTok{// Kod som ska upprepas}
\OperatorTok{\}}
\end{Highlighting}
\end{Shaded}

Här är en förklaring av de olika delarna i en for-loop:

\begin{itemize}
\tightlist
\item
  \texttt{initialisering}: Här initieras en räknare eller en variabel
  som används för att kontrollera antalet iterationer. Detta görs
  vanligtvis genom att tilldela ett startvärde till räknaren.
\item
  \texttt{villkor}: Här definieras villkoret som kontrollerar om loopen
  ska fortsätta att upprepas eller avslutas. Om villkoret är sant,
  fortsätter loopen att köras. Om villkoret är falskt, avslutas loopen
  och programmet fortsätter med koden efter loopen.
\item
  \texttt{iteration}: Här specificeras hur räknaren eller variabeln ska
  ändras vid varje iteration. Vanligtvis ökar eller minskar man värdet
  på räknaren med ett visst steg.
\end{itemize}

Här är ett exempel på en for-loop som skriver ut talen 1 till 5:

\hypertarget{cb2}{}
\begin{Shaded}
\begin{Highlighting}[]
\ControlFlowTok{for} \OperatorTok{(}\DataTypeTok{int}\NormalTok{ i }\OperatorTok{=} \DecValTok{1}\OperatorTok{;}\NormalTok{ i }\OperatorTok{\textless{}=} \DecValTok{5}\OperatorTok{;}\NormalTok{ i}\OperatorTok{++)} \OperatorTok{\{}
    \BuiltInTok{System}\OperatorTok{.}\FunctionTok{out}\OperatorTok{.}\FunctionTok{println}\OperatorTok{(}\NormalTok{i}\OperatorTok{);}
\OperatorTok{\}}
\end{Highlighting}
\end{Shaded}

I exemplet ovan initieras räknaren (variabeln) \texttt{i} till 1. Sedan
kontrolleras villkoret \texttt{i\ \textless{}=\ 5}. Om villkoret är
sant, skrivs värdet på räknaren ut och räknaren ökar med ett steg. Detta
upprepas tills villkoret är falskt, dvs. när räknaren når värdet 6.

´´´text 1 2 3 4 5 ´´´

\hypertarget{fuxf6rdelar}{%
\subsection{Fördelar}\label{fuxf6rdelar}}

For-loopar har flera fördelar jämfört med andra typer av loopar:

\begin{itemize}
\tightlist
\item
  De är enkla att använda och förstå.
\item
  De är effektiva eftersom de inte kräver några extra variabler eller
  villkor för att kontrollera antalet iterationer.
\item
  De är säkrare än andra typer av loopar eftersom de inte kräver någon
  manuell hantering av indexer eller iterationsspecifikationer.
\end{itemize}

\hypertarget{nackdelar}{%
\subsection{Nackdelar}\label{nackdelar}}

For-loopar har också några nackdelar jämfört med andra typer av loopar:

\begin{itemize}
\tightlist
\item
  De är inte lika flexibla som andra typer av loopar eftersom de kräver
  att du känner till antalet iterationer i förväg.
\item
  De är inte lika läsbara som andra typer av loopar eftersom de kräver
  att du känner till antalet iterationer i förväg.
\item
  Man kan få problem om man lägger till eller tar bort element i
  samlingen under iterationen.
\end{itemize}

\hypertarget{foreach-loop}{%
\section{Foreach-loop}\label{foreach-loop}}

En foreach-loop används för att iterera över en samling av objekt eller
värden och utföra en handling för varje element i samlingen. Syntaxen
för en foreach-loop är följande:

\hypertarget{cb1}{}
\begin{Shaded}
\begin{Highlighting}[]
\ControlFlowTok{for} \OperatorTok{(}\DataTypeTok{var}\NormalTok{ element }\OperatorTok{:}\NormalTok{ samling}\OperatorTok{)} \OperatorTok{\{}
\CommentTok{// Kod som ska utföras för varje element}
\OperatorTok{\}}
\end{Highlighting}
\end{Shaded}

Här är en förklaring av de olika delarna i en foreach-loop:

\begin{itemize}
\tightlist
\item
  \texttt{element}: En variabel som används för att representera varje
  element i samlingen när loopen itererar över den.
\item
  \texttt{samling}: Den samling av objekt eller värden som ska itereras
  över, t.ex. en lista, ett fält eller en array.
\end{itemize}

Här är ett exempel på en foreach-loop som skriver ut varje element i en
lista:

\hypertarget{cb2}{}
\begin{Shaded}
\begin{Highlighting}[]
\BuiltInTok{List}\OperatorTok{\textless{}}\BuiltInTok{Integer}\OperatorTok{\textgreater{}}\NormalTok{ numbers }\OperatorTok{=} \KeywordTok{new} \BuiltInTok{ArrayList}\OperatorTok{\textless{}\textgreater{}(}\BuiltInTok{Arrays}\OperatorTok{.}\FunctionTok{asList}\OperatorTok{(}\DecValTok{1}\OperatorTok{,} \DecValTok{2}\OperatorTok{,} \DecValTok{3}\OperatorTok{,} \DecValTok{4}\OperatorTok{,} \DecValTok{5}\OperatorTok{));}

\ControlFlowTok{for} \OperatorTok{(}\DataTypeTok{int}\NormalTok{ number }\OperatorTok{:}\NormalTok{ numbers}\OperatorTok{)} \OperatorTok{\{}
    \BuiltInTok{System}\OperatorTok{.}\FunctionTok{out}\OperatorTok{.}\FunctionTok{println}\OperatorTok{(}\NormalTok{number}\OperatorTok{);}
\OperatorTok{\}}
\end{Highlighting}
\end{Shaded}

Resultatet av koden ovan kommer att vara:

\begin{Shaded}
\begin{Highlighting}[]
\NormalTok{1}
\NormalTok{2}
\NormalTok{3}
\NormalTok{4}
\NormalTok{5}
\end{Highlighting}
\end{Shaded}

Foreach-loopar är särskilt användbara när du vill arbeta med varje
element i en samling utan att behöva oroa dig för att hantera indexer
och iterationsspecifikationer manuellt.

\hypertarget{fuxf6rdelar}{%
\subsection{Fördelar}\label{fuxf6rdelar}}

Foreach-loopar har flera fördelar jämfört med andra typer av loopar:

\begin{itemize}
\tightlist
\item
  De är enkla att använda och förstå.
\item
  De är säkrare än andra typer av loopar eftersom de inte kräver någon
  manuell hantering av indexer eller iterationsspecifikationer.
\item
  De är mer läsbara än andra typer av loopar eftersom de inte kräver
  någon manuell hantering av indexer eller iterationsspecifikationer.
\end{itemize}

\hypertarget{nackdelar}{%
\subsection{Nackdelar}\label{nackdelar}}

Foreach-loopar har också några nackdelar:

\begin{itemize}
\tightlist
\item
  De kan inte användas för att iterera över en samling baklänges.
\item
  Om du har värden i samlingen kan du inte änra dem medan du itererar
  över dem. Du måste kopiera värdena till en annan samling först.
\end{itemize}

\hypertarget{nuxe4stlade-loopar}{%
\section{Nästlade loopar}\label{nuxe4stlade-loopar}}

Nästlade loopar är loopar som finns inuti en annan loop.

Detta innebär att den inre loopen upprepas lika många gånger som den
yttre loopen bestämmer.

Om tänker dig att du ska säga ``A B C'' tre gånger.

Du säger då ``ABC ABC ABC''

I kodform skulle detta bli

\hypertarget{cb1}{}
\begin{Shaded}
\begin{Highlighting}[]
\DataTypeTok{char}\OperatorTok{[]}\NormalTok{ letters }\OperatorTok{=} \OperatorTok{\{}\CharTok{\textquotesingle{}A\textquotesingle{}}\OperatorTok{,} \CharTok{\textquotesingle{}B\textquotesingle{}}\OperatorTok{,} \CharTok{\textquotesingle{}C\textquotesingle{}}\OperatorTok{\};}

\ControlFlowTok{for}\OperatorTok{(}\DataTypeTok{int}\NormalTok{ i }\OperatorTok{=} \DecValTok{0}\OperatorTok{;}\NormalTok{ i }\OperatorTok{\textless{}} \DecValTok{3}\OperatorTok{;}\NormalTok{ i}\OperatorTok{++)} \OperatorTok{\{}
    \FunctionTok{foreach}\OperatorTok{(}\DataTypeTok{char}\NormalTok{ letter in letters}\OperatorTok{)} \OperatorTok{\{}
        \BuiltInTok{System}\OperatorTok{.}\FunctionTok{out}\OperatorTok{.}\FunctionTok{print}\OperatorTok{(}\NormalTok{letter}\OperatorTok{);}
    \OperatorTok{\}}
\NormalTok{    system}\OperatorTok{.}\FunctionTok{out}\OperatorTok{.}\FunctionTok{print}\OperatorTok{(}\CharTok{\textquotesingle{} \textquotesingle{}}\OperatorTok{);}
\OperatorTok{\}}
\end{Highlighting}
\end{Shaded}

med resultatet

\begin{Shaded}
\begin{Highlighting}[]
\NormalTok{ABC ABC ABC}
\end{Highlighting}
\end{Shaded}

De används när vi behöver upprepa en uppsättning instruktioner i flera
dimensioner. Till exempel kan vi ha en loop som representerar rader och
en inre loop som representerar kolumner för att bearbeta en matris. Här
är ett exempel på en nästlad loop:

\hypertarget{cb3}{}
\begin{Shaded}
\begin{Highlighting}[]
\ControlFlowTok{for} \OperatorTok{(}\DataTypeTok{int}\NormalTok{ row }\OperatorTok{=} \DecValTok{1}\OperatorTok{;}\NormalTok{ row }\OperatorTok{\textless{}=} \DecValTok{3}\OperatorTok{;}\NormalTok{ row}\OperatorTok{++)} \OperatorTok{\{}
    \ControlFlowTok{for} \OperatorTok{(}\DataTypeTok{int}\NormalTok{ col }\OperatorTok{=} \DecValTok{1}\OperatorTok{;}\NormalTok{ col }\OperatorTok{\textless{}=} \DecValTok{3}\OperatorTok{;}\NormalTok{ col}\OperatorTok{++)} \OperatorTok{\{}
        \BuiltInTok{System}\OperatorTok{.}\FunctionTok{out}\OperatorTok{.}\FunctionTok{println}\OperatorTok{(}\StringTok{"Row: "} \OperatorTok{+}\NormalTok{ row }\OperatorTok{+} \StringTok{", Column: "} \OperatorTok{+}\NormalTok{ col}\OperatorTok{);}
    \OperatorTok{\}}
\OperatorTok{\}}
\end{Highlighting}
\end{Shaded}

med följande resultat:

\hypertarget{cb4}{}
\begin{Shaded}
\begin{Highlighting}[]
\NormalTok{Row}\OperatorTok{:} \DecValTok{1}\OperatorTok{,}\NormalTok{ Column}\OperatorTok{:} \DecValTok{1}
\NormalTok{Row}\OperatorTok{:} \DecValTok{1}\OperatorTok{,}\NormalTok{ Column}\OperatorTok{:} \DecValTok{2}
\NormalTok{Row}\OperatorTok{:} \DecValTok{1}\OperatorTok{,}\NormalTok{ Column}\OperatorTok{:} \DecValTok{3}
\NormalTok{Row}\OperatorTok{:} \DecValTok{2}\OperatorTok{,}\NormalTok{ Column}\OperatorTok{:} \DecValTok{1}
\NormalTok{Row}\OperatorTok{:} \DecValTok{2}\OperatorTok{,}\NormalTok{ Column}\OperatorTok{:} \DecValTok{2}
\NormalTok{Row}\OperatorTok{:} \DecValTok{2}\OperatorTok{,}\NormalTok{ Column}\OperatorTok{:} \DecValTok{3}
\NormalTok{Row}\OperatorTok{:} \DecValTok{3}\OperatorTok{,}\NormalTok{ Column}\OperatorTok{:} \DecValTok{1}
\NormalTok{Row}\OperatorTok{:} \DecValTok{3}\OperatorTok{,}\NormalTok{ Column}\OperatorTok{:} \DecValTok{2}
\NormalTok{Row}\OperatorTok{:} \DecValTok{3}\OperatorTok{,}\NormalTok{ Column}\OperatorTok{:} \DecValTok{3}
\end{Highlighting}
\end{Shaded}

När nästlade loopar används är det viktigt att vara medveten om
prestanda och hur många iterationer som utförs. Nästlade loopar kan öka
tidskomplexiteten för ett program och leda till längre exekveringstider
om de inte används effektivt.

Loopar är kraftfulla verktyg som hjälper till att automatisera
upprepningsuppgifter och utföra operationer på en samling av objekt
eller värden. Genom att välja rätt typ av loop och använda den på rätt
sätt kan vi skapa mer effektiv och strukturerad kod.

\hypertarget{while}{%
\section{While}\label{while}}

While-loopar är en typ av loop inom programmering som används för att
upprepa en kodsekvens så länge ett visst villkor är sant. Det är en
grundläggande kontrollstruktur som tillåter oss att utföra repeterade
handlingar utan att behöva skriva samma kod flera gånger. I en
while-loop utvärderas villkoret före varje iteration, och om det är sant
körs loopen och kodblocket inuti utförs. När villkoret blir falskt
avslutas loopen och programmet fortsätter med den efterföljande koden.

\hypertarget{fuxf6rdelar}{%
\subsection{Fördelar}\label{fuxf6rdelar}}

While-loopar erbjuder flera fördelar och används i olika situationer:

\begin{enumerate}
\def\labelenumi{\arabic{enumi}.}
\item
  \textbf{Upprepa kod}: En while-loop tillåter oss att upprepa en
  kodsekvens ett godtyckligt antal gånger baserat på ett villkor. Det är
  särskilt användbart när vi behöver utföra samma handlingar eller
  beräkningar igen och igen tills ett visst kriterium är uppfyllt.
\item
  \textbf{Enkelhet}: While-loopar är relativt enkla att förstå och
  implementera. De har en tydlig struktur och kräver inte avancerade
  koncept eller komplicerade syntaxregler.
\item
  \textbf{Effektivitet}: Genom att använda en while-loop kan vi undvika
  att upprepa samma kod manuellt. Istället kan vi använda en loop för
  att automatisera upprepningen och minska kodupprepning, vilket sparar
  tid och minskar risken för fel.
\item
  \textbf{Anpassningsbarhet}: Med en while-loop kan vi enkelt ändra
  villkoret som avgör när loopen ska avslutas. Det ger oss flexibilitet
  att anpassa loopen efter olika scenarier och ändra beteendet med hjälp
  av villkor.
\item
  \textbf{Möjlighet till oändlig loop}: Genom att använda en lämpligt
  konstruerad while-loop kan vi skapa en oändlig loop, vilket innebär
  att loopen fortsätter att köras tills ett avbrytningsvillkor uppfylls.
  Detta kan vara användbart i vissa situationer där vi vill att en
  kodsekvens ska köras kontinuerligt tills vi explicit bryter ut ur
  loopen.
\end{enumerate}

\hypertarget{begruxe4nsningar}{%
\subsection{Begränsningar}\label{begruxe4nsningar}}

Trots sina fördelar har while-loopar vissa begränsningar att vara
medveten om:

\begin{enumerate}
\def\labelenumi{\arabic{enumi}.}
\item
  \textbf{Risk för oändlig loop}: Om villkoret i en while-loop alltid är
  sant kan loopen köra för evigt och leda till att programmet fryser
  eller inte svarar. Det är viktigt att vara försiktig när man skapar
  while-loopar och se till att det finns ett sätt att bryta ut ur loopen
  när det behövs.
\item
  \textbf{Uppdatering av villkor}: För att undvika oändliga loopar eller
  loopar som aldrig körs måste vi se till att uppdatera villkoret i
  while-loopen på rätt sätt. Annars kan loopen antingen inte köras alls
  eller köras för evigt utan att uppfylla det avsedda syftet.
\item
  \textbf{Risk för buggar}: Om vi inte är noggranna när vi skriver
  villkoret i en while-loop kan det leda till buggar och felaktigt
  beteende i vårt program. Det är viktigt att noggrant testa och
  validera villkoret för att undvika sådana problem.
\end{enumerate}

\hypertarget{anvuxe4ndningsomruxe5den}{%
\subsection{Användningsområden}\label{anvuxe4ndningsomruxe5den}}

While-loopar kan användas i en rad olika situationer och scenarier,
inklusive:

\begin{enumerate}
\def\labelenumi{\arabic{enumi}.}
\item
  \textbf{Inmatningsvalidering}: Vi kan använda en while-loop för att
  validera användarinput och be användaren att ange korrekta värden
  tills det görs.
\item
  \textbf{Iteration över en samling}: Med en while-loop kan vi iterera
  över en samling av objekt och utföra en viss åtgärd för varje objekt
  tills ett visst villkor är uppfyllt.
\item
  \textbf{Upprepning av en kodsekvens}: Vi kan använda en while-loop för
  att upprepa en kodsekvens tills ett visst villkor är sant. Det kan
  vara användbart när vi vill utföra en uppgift eller en beräkning ett
  visst antal gånger eller tills ett specifikt kriterium är uppfyllt.
\item
  \textbf{Simulering av händelser}: Med en while-loop kan vi simulera
  händelser och uppdatera tillståndet på olika delar av vårt program
  tills ett visst villkor är uppfyllt.
\end{enumerate}

\hypertarget{exempel}{%
\subsection{Exempel}\label{exempel}}

Här är ett exempel som visar hur man använder en while-loop i Java:

\hypertarget{cb1}{}
\begin{Shaded}
\begin{Highlighting}[]
\KeywordTok{public} \KeywordTok{class}\NormalTok{ WhileExample }\OperatorTok{\{}
    \KeywordTok{public} \DataTypeTok{static} \DataTypeTok{void} \FunctionTok{main}\OperatorTok{(}\BuiltInTok{String}\OperatorTok{[]}\NormalTok{ args}\OperatorTok{)} \OperatorTok{\{}
        \DataTypeTok{int}\NormalTok{ count }\OperatorTok{=} \DecValTok{0}\OperatorTok{;}
        \ControlFlowTok{while} \OperatorTok{(}\NormalTok{count }\OperatorTok{\textless{}} \DecValTok{5}\OperatorTok{)} \OperatorTok{\{}
            \BuiltInTok{System}\OperatorTok{.}\FunctionTok{out}\OperatorTok{.}\FunctionTok{println}\OperatorTok{(}\StringTok{"Count: "} \OperatorTok{+}\NormalTok{ count}\OperatorTok{);}
\NormalTok{            count}\OperatorTok{++;}
        \OperatorTok{\}}
    \OperatorTok{\}}
\OperatorTok{\}}
\end{Highlighting}
\end{Shaded}

I detta exempel använder vi en while-loop för att skriva ut värdet på
variabeln \texttt{count} tills \texttt{count} når värdet 5. För varje
iteration ökar vi värdet på \texttt{count} med 1 och skriver sedan ut
värdet. Loopen kommer att fortsätta tills \texttt{count} når värdet 5.

Output: Count: 0 Count: 1 Count: 2 Count: 3 Count: 4

Detta är bara ett exempel på hur man kan använda en while-loop. Det
finns många fler möjligheter och variationer beroende på dina specifika
behov och scenarier.

\hypertarget{sammanfattning}{%
\subsection{Sammanfattning}\label{sammanfattning}}

While-loopar är en viktig kontrollstruktur inom programmering som låter
oss upprepa en kodsekvens så länge ett villkor är sant. De erbjuder
enkelhet, effektivitet och anpassningsbarhet i våra program. Det är
viktigt att vara medveten om riskerna med oändliga loopar och att se
till att uppdatera villkoret på rätt sätt för att undvika buggar. Med
rätt användning kan while-loopar vara kraftfulla verktyg för att hantera
upprepade uppgifter och logik i våra program.

Fortsätt utforska och experimentera med while-loopar och andra loopar
inom programmering för att bygga robusta och effektiva program.

\hypertarget{do-while}{%
\section{Do While}\label{do-while}}

Do While är en loop som körs minst en gång. Den körs sedan så länge som
villkoret är sant.

Innehållsförteckning

\{: .text-delta \} 1. TOC \{:toc\}

\hypertarget{beskrivning}{%
\subsection{Beskrivning}\label{beskrivning}}

Do While är en loopstruktur i Java som liknar While-loopen. Skillnaden
är att i en Do While-loop kontrolleras villkoret efter varje iteration,
vilket innebär att loopen alltid körs minst en gång. Detta är användbart
när du vill att en viss kod ska utföras minst en gång innan villkoret
kontrolleras.

\hypertarget{exempel}{%
\subsection{Exempel}\label{exempel}}

För att förstå hur en Do While-loop fungerar, låt oss titta på följande
kodexempel:

\hypertarget{cb1}{}
\begin{Shaded}
\begin{Highlighting}[]
\DataTypeTok{int}\NormalTok{ i }\OperatorTok{=} \DecValTok{0}\OperatorTok{;}
\ControlFlowTok{do} \OperatorTok{\{}
    \BuiltInTok{System}\OperatorTok{.}\FunctionTok{out}\OperatorTok{.}\FunctionTok{println}\OperatorTok{(}\NormalTok{i}\OperatorTok{);}
\NormalTok{    i}\OperatorTok{++;}
\OperatorTok{\}} \ControlFlowTok{while} \OperatorTok{(}\NormalTok{i }\OperatorTok{\textless{}} \DecValTok{10}\OperatorTok{);}
\end{Highlighting}
\end{Shaded}

I detta exempel deklarerar vi en variabel \texttt{i} och tilldelar den
värdet 0. Sedan har vi en Do While-loop som kontrollerar om \texttt{i}
är mindre än 10. Inuti loopen skriver den ut värdet av \texttt{i} och
ökar sedan värdet med 1. Loopen fortsätter att köras så länge som
\texttt{i} är mindre än 10.

Med följande output:

\begin{Shaded}
\begin{Highlighting}[]
\NormalTok{0}
\NormalTok{1}
\NormalTok{2}
\NormalTok{3}
\NormalTok{4}
\NormalTok{5}
\NormalTok{6}
\NormalTok{7}
\NormalTok{8}
\NormalTok{9}
\end{Highlighting}
\end{Shaded}

Medan villkoret i Do While-loopen är sant, körs koden inuti loopen. När
villkoret blir falskt, avslutas loopen och programmet fortsätter med
resten av koden efter loopen.

I detta specifika exempel kommer loopen att köra 10 gånger och skriva ut
värdena från 0 till 9.

\hypertarget{sammanfattning}{%
\subsection{Sammanfattning}\label{sammanfattning}}

Do While är en loopstruktur i Java som körs minst en gång och sedan
fortsätter att köra så länge som villkoret är sant. Det är användbart
när du behöver utföra en viss kod minst en gång innan villkoret
kontrolleras.

\hypertarget{termer}{%
\subsection{Termer}\label{termer}}

\begin{longtable}[]{@{}
  >{\raggedright\arraybackslash}p{(\columnwidth - 2\tabcolsep) * \real{0.1100}}
  >{\raggedright\arraybackslash}p{(\columnwidth - 2\tabcolsep) * \real{0.8800}}@{}}
\toprule\noalign{}
\begin{minipage}[b]{\linewidth}\raggedright
Term
\end{minipage} & \begin{minipage}[b]{\linewidth}\raggedright
Förklaring
\end{minipage} \\
\midrule\noalign{}
\endhead
\bottomrule\noalign{}
\endlastfoot
Do While-loop & En loopstruktur i Java som körs minst en gång och sedan
fortsätter att köra så länge som villkoret är sant. \\
Loop & En struktur i programmering som gör att en viss kod kan köras
upprepade gånger tills ett visst villkor är uppfyllt. \\
Iteration & En enskild körning av kod inuti en loop. \\
Villkor & Ett uttryck som avgör om en loop ska fortsätta köras eller
inte. \\
\end{longtable}

\hypertarget{rekursiva-loopar}{%
\section{Rekursiva loopar}\label{rekursiva-loopar}}

Rekursion är en teknik där en funktion kallar sig själv för att utföra
en uppgift. Detta kan vara användbart när du har en uppgift som kan
delas upp i mindre liknande deluppgifter. Ett vanligt exempel på
rekursion är beräkningen av det n:te Fibonacci-talet. Här är en rekursiv
funktion som beräknar Fibonacci-talet:

\hypertarget{cb1}{}
\begin{Shaded}
\begin{Highlighting}[]
\DataTypeTok{int} \FunctionTok{Fibonacci}\OperatorTok{(}\DataTypeTok{int}\NormalTok{ n}\OperatorTok{)} \OperatorTok{\{}
    \ControlFlowTok{if} \OperatorTok{(}\NormalTok{n }\OperatorTok{\textless{}=} \DecValTok{1}\OperatorTok{)}
        \ControlFlowTok{return}\NormalTok{ n}\OperatorTok{;}
    \ControlFlowTok{else}
        \ControlFlowTok{return} \FunctionTok{Fibonacci}\OperatorTok{(}\NormalTok{n }\OperatorTok{{-}} \DecValTok{1}\OperatorTok{)} \OperatorTok{+} \FunctionTok{Fibonacci}\OperatorTok{(}\NormalTok{n }\OperatorTok{{-}} \DecValTok{2}\OperatorTok{);}
\end{Highlighting}
\end{Shaded}

Du kan sedan anropa funktionen med önskat värde för att få det
motsvarande Fibonacci-talet:

\hypertarget{cb2}{}
\begin{Shaded}
\begin{Highlighting}[]
\DataTypeTok{int}\NormalTok{ result }\OperatorTok{=} \FunctionTok{Fibonacci}\OperatorTok{(}\DecValTok{6}\OperatorTok{);}
\BuiltInTok{System}\OperatorTok{.}\FunctionTok{out}\OperatorTok{.}\FunctionTok{println}\OperatorTok{(}\NormalTok{result}\OperatorTok{);} \CommentTok{// Resultatet blir 8}
\end{Highlighting}
\end{Shaded}

Bästa sättet att förstå det är att köra det i debugläge och stega ett
steg i taget. Skriv om koden såhär:

\hypertarget{cb3}{}
\begin{Shaded}
\begin{Highlighting}[]
\DataTypeTok{int} \FunctionTok{Fibonacci}\OperatorTok{(}\DataTypeTok{int}\NormalTok{ n}\OperatorTok{)} \OperatorTok{\{}
    \ControlFlowTok{if} \OperatorTok{(}\NormalTok{n }\OperatorTok{\textless{}=} \DecValTok{1}\OperatorTok{)}
        \ControlFlowTok{return}\NormalTok{ n}\OperatorTok{;}
    \ControlFlowTok{else}
        \ControlFlowTok{return}
            \FunctionTok{Fibonacci}\OperatorTok{(}\NormalTok{n }\OperatorTok{{-}} \DecValTok{1}\OperatorTok{)} \OperatorTok{+}
            \FunctionTok{Fibonacci}\OperatorTok{(}\NormalTok{n }\OperatorTok{{-}} \DecValTok{2}\OperatorTok{);}
\OperatorTok{\}}
\end{Highlighting}
\end{Shaded}

Vi försöker förklara flödet

\begin{enumerate}
\def\labelenumi{\arabic{enumi}.}
\tightlist
\item
  Först anropas metoden med värdet 6.
\item
  Eftersom 6 inte är mindre än eller lika med 1, går vi vidare till
  else-blocket.
\item
  Metoden anropar sig själv två gånger med de två föregående talen i
  Fibonacci-sekvensen, det vill säga \texttt{Fibonacci(6\ -\ 1)} och
  \texttt{Fibonacci(6\ -\ 2)}.
\item
  Nu fortsätter vi att evaluera dessa anrop separat.
\item
  Första anropet \texttt{Fibonacci(6\ -\ 1)} utvärderas till
  \texttt{Fibonacci(5)}.
\item
  För att beräkna \texttt{Fibonacci(5)} måste vi fortsätta att anropa
  metoden med mindre värden tills vi når basfallet.
\item
  Detta resulterar i en serie anrop \texttt{Fibonacci(4)},
  \texttt{Fibonacci(3)}, \texttt{Fibonacci(2)}, och
  \texttt{Fibonacci(1)}.
\item
  När vi når basfallet \texttt{Fibonacci(1)} returneras värdet 1.
\item
  Därefter fortsätter vi att evaluera anropen på vägen tillbaka.
  \texttt{Fibonacci(2)} returneras också med värdet 1.
\item
  Nu kan vi beräkna \texttt{Fibonacci(3)} genom att addera
  \texttt{Fibonacci(2)} och \texttt{Fibonacci(1)}, vilket ger 2.
\item
  Vi fortsätter att evaluera anropen och beräkna värdena för
  \texttt{Fibonacci(4)} och \texttt{Fibonacci(5)}.
\item
  Slutligen, när vi har beräknat \texttt{Fibonacci(6)} genom att addera
  \texttt{Fibonacci(5)} och \texttt{Fibonacci(4)}, returneras värdet 8.
\end{enumerate}

Så resultatet av att anropa \texttt{Fibonacci(6)} är 8 enligt
definitionen av Fibonacci-sekvensen.

Det är värt att notera att den här implementationen är rekursiv och kan
vara ineffektiv för stora värden på \texttt{n} på grund av att den utför
många redundanta beräkningar. Det finns alternativa iterativa eller
memoiserade implementationer som kan förbättra prestandan för
Fibonacci-sekvensen.

Sätt en breakpoint på raden \texttt{return\ n;} och kör programmet.
Stega igenom koden och se vad som händer.

\hypertarget{rekursion-i-java}{%
\subsection{Rekursion i Java}\label{rekursion-i-java}}

Rekursion är en cool teknik som används inom programmering för att lösa
problem genom att bryta ner dem i mindre delproblem. Det innebär att en
funktion kan anropa sig själv för att lösa ett problem. Det är dock
viktigt att ha en stoppvillkor för att undvika oändlig rekursion. Här är
ett exempel på rekursion i Java:

\hypertarget{cb4}{}
\begin{Shaded}
\begin{Highlighting}[]
\KeywordTok{public} \KeywordTok{class}\NormalTok{ RekursionExempel }\OperatorTok{\{}
    \KeywordTok{public} \DataTypeTok{static} \DataTypeTok{void} \FunctionTok{main}\OperatorTok{(}\BuiltInTok{String}\OperatorTok{[]}\NormalTok{ args}\OperatorTok{)} \OperatorTok{\{}

        \DataTypeTok{int}\NormalTok{ resultat }\OperatorTok{=} \FunctionTok{summera}\OperatorTok{(}\DecValTok{5}\OperatorTok{);}
        \BuiltInTok{System}\OperatorTok{.}\FunctionTok{out}\OperatorTok{.}\FunctionTok{println}\OperatorTok{(}\StringTok{"Summan är: "} \OperatorTok{+}\NormalTok{ resultat}\OperatorTok{);}
    \OperatorTok{\}}

    \KeywordTok{public} \DataTypeTok{static} \DataTypeTok{int} \FunctionTok{summera}\OperatorTok{(}\DataTypeTok{int}\NormalTok{ n}\OperatorTok{)} \OperatorTok{\{}
        \ControlFlowTok{if} \OperatorTok{(}\NormalTok{n }\OperatorTok{==} \DecValTok{1}\OperatorTok{)} \OperatorTok{\{}
            \ControlFlowTok{return} \DecValTok{1}\OperatorTok{;}
        \OperatorTok{\}} \ControlFlowTok{else} \OperatorTok{\{}
            \ControlFlowTok{return}\NormalTok{ n }\OperatorTok{+} \FunctionTok{summera}\OperatorTok{(}\NormalTok{n }\OperatorTok{{-}} \DecValTok{1}\OperatorTok{);}
    \OperatorTok{\}}
\OperatorTok{\}}
\end{Highlighting}
\end{Shaded}

I det här exemplet använder vi en rekursiv funktion \texttt{summera} för
att beräkna summan av alla heltal från 1 till det givna talet
\texttt{n}. Funktionen har en stoppvillkor som säger att om \texttt{n}
är lika med 1, returnera 1. Annars anropar funktionen sig själv med
\texttt{n-1} och lägger till \texttt{n} i resultatet. När vi kör
programmet kommer det att skriva ut summan av alla heltal från 1 till 5,
vilket är 15.

Rekursion kan vara svårt att förstå i början, men med övning och
erfarenhet kan du behärska denna kraftfulla teknik. Fortsätt att
utforska och experimentera med rekursion för att utveckla dina
programmeringskunskaper!

Om du vill ha mer inspiration kan du titta på några av följande filmer:

\begin{itemize}
\tightlist
\item
  \href{https://www.imdb.com/title/tt0945513/}{Source code}
\item
  \href{https://www.imdb.com/title/tt5308322/}{Happy Death day}
\item
  \href{https://www.imdb.com/title/tt8155288/}{Happy Death day 2U}
\item
  \href{https://www.imdb.com/title/tt0107048/}{Groundhog day}
\item
  \href{https://www.imdb.com/title/tt1631867/}{Edge of tomorrow}
\end{itemize}

\hypertarget{obligatorisk-dad-joke}{%
\subsection{Obligatorisk dad-joke}\label{obligatorisk-dad-joke}}

Varför gillar programmerare att använda rekursion i skämt?

För att förstå rekursion i skämt måste du först förstå rekursion i
skämt.

\hypertarget{evig-loop}{%
\section{Evig loop}\label{evig-loop}}

En oändlig loop är en loop som inte har något villkor för att avslutas
och kommer att köra kontinuerligt tills programmet avbryts. Det kan vara
användbart i vissa situationer, till exempel när du skapar ett program
som körs i en slinga och kontinuerligt utför några uppgifter i
bakgrunden. Här är ett exempel på en oändlig loop:

\hypertarget{cb1}{}
\begin{Shaded}
\begin{Highlighting}[]
\ControlFlowTok{while} \OperatorTok{(}\KeywordTok{true}\OperatorTok{)} \OperatorTok{\{}
\CommentTok{// Kod som körs kontinuerligt}
\OperatorTok{\}}
\end{Highlighting}
\end{Shaded}

Observera att en oändlig loop kan leda till att programmet hänger sig om
det inte finns något sätt att avbryta loopen manuellt.

Detta gör man när man testar eller när man har en brytning i loopen som
gör att den avslutas.

\hypertarget{fuxf6rdelar}{%
\subsection{Fördelar}\label{fuxf6rdelar}}

Oändliga loopar har flera fördelar jämfört med andra typer av loopar:

\begin{itemize}
\tightlist
\item
  De är enkla att använda och förstå.
\item
  De är effektiva eftersom de inte kräver några extra variabler eller
  villkor för att kontrollera antalet iterationer.
\item
  De är säkrare än andra typer av loopar eftersom de inte kräver någon
  manuell hantering av indexer eller iterationsspecifikationer.
\end{itemize}

\hypertarget{nackdelar}{%
\subsection{Nackdelar}\label{nackdelar}}

Oändliga loopar har flera nackdelar jämfört med andra typer av loopar:

\begin{itemize}
\tightlist
\item
  De kan vara svåra att felsöka eftersom de inte har något villkor för
  att avsluta loopen.
\item
  De kan leda till att programmet hänger sig om det inte finns något
  sätt att avbryta loopen manuellt.
\item
  De kan vara svåra att felsöka eftersom de inte har något villkor för
  att avsluta loopen.
\end{itemize}

\hypertarget{exempel}{%
\subsection{Exempel}\label{exempel}}

Här är ett exempel på en oändlig loop som skriver ut talen 1 till 5:

\hypertarget{cb2}{}
\begin{Shaded}
\begin{Highlighting}[]
\ControlFlowTok{while} \OperatorTok{(}\KeywordTok{true}\OperatorTok{)} \OperatorTok{\{}
    \BuiltInTok{System}\OperatorTok{.}\FunctionTok{out}\OperatorTok{.}\FunctionTok{println}\OperatorTok{(}\StringTok{"1"}\OperatorTok{);}
    \BuiltInTok{System}\OperatorTok{.}\FunctionTok{out}\OperatorTok{.}\FunctionTok{println}\OperatorTok{(}\StringTok{"2"}\OperatorTok{);}
    \BuiltInTok{System}\OperatorTok{.}\FunctionTok{out}\OperatorTok{.}\FunctionTok{println}\OperatorTok{(}\StringTok{"3"}\OperatorTok{);}
    \BuiltInTok{System}\OperatorTok{.}\FunctionTok{out}\OperatorTok{.}\FunctionTok{println}\OperatorTok{(}\StringTok{"4"}\OperatorTok{);}
    \BuiltInTok{System}\OperatorTok{.}\FunctionTok{out}\OperatorTok{.}\FunctionTok{println}\OperatorTok{(}\StringTok{"5"}\OperatorTok{);}
\OperatorTok{\}}
\end{Highlighting}
\end{Shaded}

Resultatet av koden ovan kommer att vara:

\begin{Shaded}
\begin{Highlighting}[]
\NormalTok{1}
\NormalTok{2}
\NormalTok{3}
\NormalTok{4}
\NormalTok{5}
\NormalTok{1}
\NormalTok{2}
\NormalTok{3}
\NormalTok{osv...}
\end{Highlighting}
\end{Shaded}

\hypertarget{exempel-puxe5-ouxe4ndlig-loop}{%
\subsection{Exempel på oändlig
loop}\label{exempel-puxe5-ouxe4ndlig-loop}}

Här är ett exempel på en meny som körs i en oändlig loop. Fördelen med
detta är att vi inte behöver veta hur många gånger användaren vill
upprepa loopen, utan vi kan låta användaren avsluta loopen när den är
klar.

\hypertarget{cb4}{}
\begin{Shaded}
\begin{Highlighting}[]
\ControlFlowTok{while} \OperatorTok{(}\KeywordTok{true}\OperatorTok{)} \OperatorTok{\{}
    \BuiltInTok{System}\OperatorTok{.}\FunctionTok{out}\OperatorTok{.}\FunctionTok{println}\OperatorTok{(}\StringTok{"Välj ett alternativ:"}\OperatorTok{);}
    \BuiltInTok{System}\OperatorTok{.}\FunctionTok{out}\OperatorTok{.}\FunctionTok{println}\OperatorTok{(}\StringTok{"1. Visa alla produkter"}\OperatorTok{);}
    \BuiltInTok{System}\OperatorTok{.}\FunctionTok{out}\OperatorTok{.}\FunctionTok{println}\OperatorTok{(}\StringTok{"2. Lägg till en produkt"}\OperatorTok{);}
    \BuiltInTok{System}\OperatorTok{.}\FunctionTok{out}\OperatorTok{.}\FunctionTok{println}\OperatorTok{(}\StringTok{"3. Ta bort en produkt"}\OperatorTok{);}
    \BuiltInTok{System}\OperatorTok{.}\FunctionTok{out}\OperatorTok{.}\FunctionTok{println}\OperatorTok{(}\StringTok{"4. Avsluta programmet"}\OperatorTok{);}

    \DataTypeTok{int}\NormalTok{ choice }\OperatorTok{=}\NormalTok{ scanner}\OperatorTok{.}\FunctionTok{nextInt}\OperatorTok{();}

    \CommentTok{// Observera att switch break inte bryter loopen i sig}
    \CommentTok{// same name different shit... eller vad man ska säga}
\NormalTok{    bool breakLoop}\OperatorTok{=}\KeywordTok{false}\OperatorTok{;}
    \ControlFlowTok{switch} \OperatorTok{(}\NormalTok{choice}\OperatorTok{)} \OperatorTok{\{}
        \ControlFlowTok{case} \DecValTok{1}\OperatorTok{:}
            \CommentTok{// Visa alla produkter}
            \ControlFlowTok{break}\OperatorTok{;}
        \ControlFlowTok{case} \DecValTok{2}\OperatorTok{:}
            \CommentTok{// Lägg till en produkt}
            \ControlFlowTok{break}\OperatorTok{;}
        \ControlFlowTok{case} \DecValTok{3}\OperatorTok{:}
            \CommentTok{// Ta bort en produkt}
            \ControlFlowTok{break}\OperatorTok{;}
        \ControlFlowTok{case} \DecValTok{4}\OperatorTok{:}
            \CommentTok{// Avsluta programmet}
\NormalTok{            breakLoop}\OperatorTok{=}\KeywordTok{true}\OperatorTok{;}
            \ControlFlowTok{break}\OperatorTok{;}
        \KeywordTok{default}\OperatorTok{:}
            \BuiltInTok{System}\OperatorTok{.}\FunctionTok{out}\OperatorTok{.}\FunctionTok{println}\OperatorTok{(}\StringTok{"Felaktigt val"}\OperatorTok{);}
            \ControlFlowTok{break}\OperatorTok{;}
    \OperatorTok{\}}

    \ControlFlowTok{if} \OperatorTok{(}\NormalTok{breakLoop}\OperatorTok{)} \ControlFlowTok{break}\OperatorTok{;} \CommentTok{// bryter loopen}
\OperatorTok{\}}
\end{Highlighting}
\end{Shaded}

\hypertarget{loop-kontroller}{%
\section{Loop kontroller}\label{loop-kontroller}}

För att få mer kontroll på loopar kan vi använda oss av loop kontroller.
Dessa är \texttt{break}, \texttt{continue} och \texttt{return}.

\hypertarget{break}{%
\subsection{Break}\label{break}}

\texttt{break} används för att avbryta en loop. Detta kan vara
användbart om vi vill avbryta en loop innan den är klar.

\hypertarget{cb1}{}
\begin{Shaded}
\begin{Highlighting}[]
\DataTypeTok{int}\NormalTok{ max }\OperatorTok{=} \DecValTok{5}\OperatorTok{;}
\ControlFlowTok{for}\OperatorTok{(}\DataTypeTok{int}\NormalTok{ i }\OperatorTok{=} \DecValTok{0}\OperatorTok{;}\NormalTok{ i }\OperatorTok{\textless{}} \DecValTok{10}\OperatorTok{;}\NormalTok{ i}\OperatorTok{++)} \OperatorTok{\{}
    \ControlFlowTok{if}\OperatorTok{(}\NormalTok{i }\OperatorTok{\textgreater{}=}\NormalTok{ max}\OperatorTok{)} \OperatorTok{\{}
        \ControlFlowTok{break}\OperatorTok{;}
    \OperatorTok{\}}
    \BuiltInTok{System}\OperatorTok{.}\FunctionTok{out}\OperatorTok{.}\FunctionTok{println}\OperatorTok{(}\NormalTok{i}\OperatorTok{);}
\OperatorTok{\}}
\end{Highlighting}
\end{Shaded}

Resultatet av koden ovan kommer att vara:

\begin{Shaded}
\begin{Highlighting}[]
\NormalTok{0}
\NormalTok{1}
\NormalTok{2}
\NormalTok{3}
\NormalTok{4}
\end{Highlighting}
\end{Shaded}

Loopen är en loop som ska gå upp till tio, men vi styr den med
\texttt{break} så att den bara går upp till fem som vi angett som
maxvärde.

\hypertarget{continue}{%
\subsection{Continue}\label{continue}}

\texttt{continue} används för att hoppa över en iteration i en loop.
Detta kan vara användbart om vi vill hoppa över en iteration i en loop.

\hypertarget{cb3}{}
\begin{Shaded}
\begin{Highlighting}[]
\DataTypeTok{int}\NormalTok{ mod}\OperatorTok{=}\DecValTok{3}\OperatorTok{;}
\ControlFlowTok{for}\OperatorTok{(}\DataTypeTok{int}\NormalTok{ i }\OperatorTok{=} \DecValTok{0}\OperatorTok{;}\NormalTok{ i }\OperatorTok{\textless{}} \DecValTok{10}\OperatorTok{;}\NormalTok{ i}\OperatorTok{++)} \OperatorTok{\{}
    \ControlFlowTok{if}\OperatorTok{(}\NormalTok{i }\OperatorTok{\%}\NormalTok{ mod }\OperatorTok{==} \DecValTok{0}\OperatorTok{)} \OperatorTok{\{}
        \ControlFlowTok{continue}\OperatorTok{;}
    \OperatorTok{\}}
    \BuiltInTok{System}\OperatorTok{.}\FunctionTok{out}\OperatorTok{.}\FunctionTok{println}\OperatorTok{(}\NormalTok{i}\OperatorTok{);}
\OperatorTok{\}}
\end{Highlighting}
\end{Shaded}

Resultatet av koden ovan kommer att vara:

\begin{Shaded}
\begin{Highlighting}[]
\NormalTok{1}
\NormalTok{2}
\NormalTok{4}
\NormalTok{5}
\NormalTok{7}
\NormalTok{8}
\end{Highlighting}
\end{Shaded}

Med mod variabeln kan vi styra hur många iterationer som ska hoppas
över. I exemplet ovan hoppar vi över var tredje iteration. Continue
används ofta tillsammans med en if-sats. Den gör att loopen hoppar över
allt följande kod i loopen och går direkt till nästa iteration.

\hypertarget{return}{%
\subsection{Return}\label{return}}

\texttt{return} används för att avsluta en metod. Detta kan vara
användbart om vi vill avsluta en metod innan den är klar.

\hypertarget{cb5}{}
\begin{Shaded}
\begin{Highlighting}[]
\KeywordTok{public} \DataTypeTok{static} \DataTypeTok{void} \FunctionTok{main}\OperatorTok{(}\BuiltInTok{String}\OperatorTok{[]}\NormalTok{ args}\OperatorTok{)} \OperatorTok{\{}
    \DataTypeTok{int}\NormalTok{ max }\OperatorTok{=} \DecValTok{5}\OperatorTok{;}
    \ControlFlowTok{for}\OperatorTok{(}\DataTypeTok{int}\NormalTok{ i }\OperatorTok{=} \DecValTok{0}\OperatorTok{;}\NormalTok{ i }\OperatorTok{\textless{}} \DecValTok{10}\OperatorTok{;}\NormalTok{ i}\OperatorTok{++)} \OperatorTok{\{}
        \ControlFlowTok{if}\OperatorTok{(}\NormalTok{i }\OperatorTok{\textgreater{}=}\NormalTok{ max}\OperatorTok{)} \OperatorTok{\{}
            \ControlFlowTok{return}\OperatorTok{;}
        \OperatorTok{\}}
        \BuiltInTok{System}\OperatorTok{.}\FunctionTok{out}\OperatorTok{.}\FunctionTok{println}\OperatorTok{(}\NormalTok{i}\OperatorTok{);}
    \OperatorTok{\}}
\OperatorTok{\}}
\end{Highlighting}
\end{Shaded}

Resultatet av koden ovan kommer att vara:

\begin{Shaded}
\begin{Highlighting}[]
\NormalTok{0}
\NormalTok{1}
\NormalTok{2}
\NormalTok{3}
\NormalTok{4}
\end{Highlighting}
\end{Shaded}

Hela metoden avslutas när vi når \texttt{return} och vi kommer inte till
\texttt{System.out.println(i);} som ligger efter.

\hypertarget{yield}{%
\subsection{yield}\label{yield}}

\texttt{yield} används för att avsluta en iterator. Detta kan vara
användbart om vi vill avsluta en iterator innan den är klar.

\hypertarget{cb7}{}
\begin{Shaded}
\begin{Highlighting}[]
\KeywordTok{public} \DataTypeTok{static} \DataTypeTok{void} \FunctionTok{main}\OperatorTok{(}\BuiltInTok{String}\OperatorTok{[]}\NormalTok{ args}\OperatorTok{)} \OperatorTok{\{}
    \DataTypeTok{int}\NormalTok{ max }\OperatorTok{=} \DecValTok{5}\OperatorTok{;}
    \ControlFlowTok{for}\OperatorTok{(}\DataTypeTok{int}\NormalTok{ i }\OperatorTok{=} \DecValTok{0}\OperatorTok{;}\NormalTok{ i }\OperatorTok{\textless{}} \DecValTok{10}\OperatorTok{;}\NormalTok{ i}\OperatorTok{++)} \OperatorTok{\{}
        \ControlFlowTok{if}\OperatorTok{(}\NormalTok{i }\OperatorTok{\textgreater{}=}\NormalTok{ max}\OperatorTok{)} \OperatorTok{\{}
\NormalTok{            yield}\OperatorTok{;}
        \OperatorTok{\}}
        \BuiltInTok{System}\OperatorTok{.}\FunctionTok{out}\OperatorTok{.}\FunctionTok{println}\OperatorTok{(}\NormalTok{i}\OperatorTok{);}
    \OperatorTok{\}}
\OperatorTok{\}}
\end{Highlighting}
\end{Shaded}

Resultatet av koden ovan kommer att vara:

\begin{Shaded}
\begin{Highlighting}[]
\NormalTok{0}
\NormalTok{1}
\NormalTok{2}
\NormalTok{3}
\NormalTok{4}
\end{Highlighting}
\end{Shaded}

Hela iteratorn avslutas när vi når \texttt{yield} och vi kommer inte
till \texttt{System.out.println(i);} som ligger efter. Vad är då
skillnaden mellan detta och break? Jo, break avslutar bara loopen, medan
yield avslutar hela iteratorn. Detta är användbart om vi vill avsluta en
iterator innan den är klar.

Vi kollar på ett bättre exempel:

\hypertarget{cb9}{}
\begin{Shaded}
\begin{Highlighting}[]
\CommentTok{// kalla på en metod som yieldar}

\KeywordTok{public} \DataTypeTok{static} \DataTypeTok{void} \FunctionTok{main}\OperatorTok{(}\BuiltInTok{String}\OperatorTok{[]}\NormalTok{ args}\OperatorTok{)} \OperatorTok{\{}
    \BuiltInTok{System}\OperatorTok{.}\FunctionTok{out}\OperatorTok{.}\FunctionTok{println}\OperatorTok{(}\StringTok{"Start"}\OperatorTok{);}
    \BuiltInTok{System}\OperatorTok{.}\FunctionTok{out}\OperatorTok{.}\FunctionTok{println}\OperatorTok{(}\FunctionTok{GetNumber}\OperatorTok{());}
    \BuiltInTok{System}\OperatorTok{.}\FunctionTok{out}\OperatorTok{.}\FunctionTok{println}\OperatorTok{(}\FunctionTok{GetNumber}\OperatorTok{());}
    \BuiltInTok{System}\OperatorTok{.}\FunctionTok{out}\OperatorTok{.}\FunctionTok{println}\OperatorTok{(}\FunctionTok{GetNumber}\OperatorTok{());}
    \BuiltInTok{System}\OperatorTok{.}\FunctionTok{out}\OperatorTok{.}\FunctionTok{println}\OperatorTok{(}\FunctionTok{GetNumber}\OperatorTok{());}
    \BuiltInTok{System}\OperatorTok{.}\FunctionTok{out}\OperatorTok{.}\FunctionTok{println}\OperatorTok{(}\StringTok{"End"}\OperatorTok{);}
\OperatorTok{\}}

\KeywordTok{public} \DataTypeTok{static} \DataTypeTok{int} \FunctionTok{GetNumber}\OperatorTok{()} \OperatorTok{\{}
    \ControlFlowTok{for}\OperatorTok{(}\DataTypeTok{int}\NormalTok{ i }\OperatorTok{=} \DecValTok{0}\OperatorTok{;}\NormalTok{ i }\OperatorTok{\textless{}} \DecValTok{10}\OperatorTok{;}\NormalTok{ i}\OperatorTok{++)} \OperatorTok{\{}
        \ControlFlowTok{if}\OperatorTok{(}\NormalTok{i }\OperatorTok{\textgreater{}=} \DecValTok{5}\OperatorTok{)} \OperatorTok{\{}
\NormalTok{            yield}\OperatorTok{;}
        \OperatorTok{\}}
        \BuiltInTok{System}\OperatorTok{.}\FunctionTok{out}\OperatorTok{.}\FunctionTok{println}\OperatorTok{(}\NormalTok{i}\OperatorTok{);}
    \OperatorTok{\}}
    \ControlFlowTok{return} \DecValTok{0}\OperatorTok{;}
\OperatorTok{\}}
\end{Highlighting}
\end{Shaded}

Resultatet av koden ovan kommer att vara:

\begin{Shaded}
\begin{Highlighting}[]
\NormalTok{Start}
\NormalTok{0}
\NormalTok{1}
\NormalTok{2}
\NormalTok{3}
\NormalTok{4}
\NormalTok{End}
\end{Highlighting}
\end{Shaded}

Vad händer om vi byter ut yield mot break?

\hypertarget{cb11}{}
\begin{Shaded}
\begin{Highlighting}[]
\CommentTok{// kalla på en metod som yieldar}

\KeywordTok{public} \DataTypeTok{static} \DataTypeTok{void} \FunctionTok{main}\OperatorTok{(}\BuiltInTok{String}\OperatorTok{[]}\NormalTok{ args}\OperatorTok{)} \OperatorTok{\{}
    \BuiltInTok{System}\OperatorTok{.}\FunctionTok{out}\OperatorTok{.}\FunctionTok{println}\OperatorTok{(}\StringTok{"Start"}\OperatorTok{);}
    \BuiltInTok{System}\OperatorTok{.}\FunctionTok{out}\OperatorTok{.}\FunctionTok{println}\OperatorTok{(}\FunctionTok{GetNumber}\OperatorTok{());}
    \BuiltInTok{System}\OperatorTok{.}\FunctionTok{out}\OperatorTok{.}\FunctionTok{println}\OperatorTok{(}\FunctionTok{GetNumber}\OperatorTok{());}
    \BuiltInTok{System}\OperatorTok{.}\FunctionTok{out}\OperatorTok{.}\FunctionTok{println}\OperatorTok{(}\StringTok{"End"}\OperatorTok{);}
\OperatorTok{\}}

\KeywordTok{public} \DataTypeTok{static} \DataTypeTok{int} \FunctionTok{GetNumber}\OperatorTok{()} \OperatorTok{\{}
    \ControlFlowTok{for}\OperatorTok{(}\DataTypeTok{int}\NormalTok{ i }\OperatorTok{=} \DecValTok{0}\OperatorTok{;}\NormalTok{ i }\OperatorTok{\textless{}} \DecValTok{10}\OperatorTok{;}\NormalTok{ i}\OperatorTok{++)} \OperatorTok{\{}
        \ControlFlowTok{if}\OperatorTok{(}\NormalTok{i }\OperatorTok{\textgreater{}=} \DecValTok{5}\OperatorTok{)} \OperatorTok{\{}
            \ControlFlowTok{break}\OperatorTok{;}
        \OperatorTok{\}}
        \BuiltInTok{System}\OperatorTok{.}\FunctionTok{out}\OperatorTok{.}\FunctionTok{println}\OperatorTok{(}\NormalTok{i}\OperatorTok{);}
    \OperatorTok{\}}
    \ControlFlowTok{return} \DecValTok{0}\OperatorTok{;}
\OperatorTok{\}}
\end{Highlighting}
\end{Shaded}

Resultatet av koden ovan kommer att vara:

\begin{Shaded}
\begin{Highlighting}[]
\NormalTok{Start}
\NormalTok{0}
\NormalTok{1}
\NormalTok{2}
\NormalTok{3}
\NormalTok{4}
\NormalTok{0}
\NormalTok{1}
\NormalTok{2}
\NormalTok{3}
\NormalTok{4}
\NormalTok{End}
\end{Highlighting}
\end{Shaded}

Yield ger oss alltså en möjlighet att avsluta en iterator innan den är
klar. Detta är användbart vi vill anropa en metod men bara få en del av
resultatet åt gången. Visst är det coolt.

Tänk dig att du har en lista med namn men vill bara arbeta med en åt
gången utan att behöva kapsla in din kod i en jätteloop. Då kan du
använda yield.

Det är ett sätt att slippa ifrån loopar, vilket man gärna gör i
funktionell programmering.

\hypertarget{objektorienterad-programmering-oop}{%
\section{Objektorienterad programmering
(OOP)}\label{objektorienterad-programmering-oop}}

I denna avsnitt kommer vi att titta på ämnet Objektorienterad
programmering (OOP).

\hypertarget{nuxe4r-du-har-luxe4st-detta-kommer-du-att-kunna}{%
\subsection{När du har läst detta kommer du att
kunna}\label{nuxe4r-du-har-luxe4st-detta-kommer-du-att-kunna}}

\begin{itemize}
\tightlist
\item
  Förstå och förklara vad Objektorienterad programmering (OOP) är och
  dess relevans inom programmering.
\item
  Diskutera fördelar och begränsningar med Objektorienterad
  programmering (OOP).
\item
  Identifiera olika användningsområden där Objektorienterad
  programmering (OOP) kan tillämpas.
\item
  Förstå och tolka ett kodexempel som använder Objektorienterad
  programmering (OOP).
\item
  Känna dig inspirerad att utforska mer om ämnet och lära dig mer!
\end{itemize}

\hypertarget{innehuxe5llsfuxf6rteckning}{%
\subsection{Innehållsförteckning}\label{innehuxe5llsfuxf6rteckning}}

\begin{itemize}
\tightlist
\item
  \protect\hyperlink{introduktion}{Introduktion}
\item
  \protect\hyperlink{fuxf6rdelar}{Fördelar}
\item
  \protect\hyperlink{begruxe4nsningar}{Begränsningar}
\item
  \protect\hyperlink{anvuxe4ndningsomruxe5den}{Användningsområden}
\item
  \protect\hyperlink{exempel}{Exempel}
\item
  \protect\hyperlink{slutsats}{Slutsats}
\item
  \protect\hyperlink{tldr}{TL;DR}
\end{itemize}

\hypertarget{introduktion}{%
\subsection{Introduktion}\label{introduktion}}

Välkommen till världen av Objektorienterad programmering (OOP)! Det är
en spännande grej inom programmering som handlar om att organisera kod
på ett sätt som liknar hur vi tänker och interagerar med saker i den
verkliga världen. Vi skapar objekt som har egenskaper och beteenden,
precis som riktiga grejer vi kan pilla på och få dem att göra saker!

\hypertarget{fuxf6rdelar}{%
\subsection{Fördelar}\label{fuxf6rdelar}}

Så, varför är OOP så coolt? Jo, det finns faktiskt några riktigt häftiga
fördelar med det. Kolla in:

\begin{itemize}
\tightlist
\item
  \textbf{Modularitet}: Med OOP kan vi dela upp vår kod i små,
  självständiga objekt. Det gör det enklare att bygga och underhålla vår
  kod och ger oss möjlighet att återanvända objekten på olika ställen.
  Precis som att bygga med lego!---
\item
  \textbf{Återanvändbarhet}: Vi kan använda samma objekt om och om igen
  i olika delar av vår kod eller till och med i olika projekt. Det
  sparar tid och minskar mängden krånglig kod. Det är bra för oss och
  bra för vår planet!
\item
  \textbf{Lätt att förstå och ändra}: Om vi behöver göra ändringar i en
  del av koden påverkar det inte resten av systemet. Vi kan fixa buggar
  eller göra förbättringar utan att oroa oss för att allt annat ska
  sluta fungera. Det är som att bygga ett hus med olika moduler som kan
  bytas ut eller uppgraderas utan att påverka resten av huset.
\item
  \textbf{Abstraktion}: Genom att använda abstraktion kan vi förenkla
  komplexa system genom att bara visa de viktigaste detaljerna för
  användaren. Det är som att köra en bil, vi behöver inte veta hur
  motorn fungerar för att kunna köra bilen.
\item
  \textbf{Flexibilitet}: OOP ger oss möjlighet att skapa hierarkier av
  objekt och använda koncept som arv och polymorfism. Det ger oss
  flexibilitet att skapa olika typer av objekt och hantera dem på ett
  enkelt sätt. Det är som att ha en verktygslåda med olika verktyg som
  vi kan använda för att lösa olika problem.
\end{itemize}

\hypertarget{begruxe4nsningar}{%
\subsection{Begränsningar}\label{begruxe4nsningar}}

Nu ska vi vara ärliga här, OOP har också sina begränsningar och
utmaningar. Det är viktigt att vara medveten om dem:

\begin{itemize}
\tightlist
\item
  \textbf{Inlärningskurva}: OOP kan vara lite knepigt att lära sig i
  början. Det finns en del koncept och termer att förstå, men det är
  värt det! När du väl har lärt dig grunderna kommer du att kunna skapa
  fantastiska saker!---
\item
  \textbf{Prestandaöverväganden}: Ibland kan OOP leda till viss
  prestandaförlust eftersom det kan finnas några ``bakom kulisserna''
  -operationer som påverkar kodens hastighet. Detta är dock vanligtvis
  inget att oroa sig för i de flesta fall.
\item
  \textbf{Designkomplexitet}: Om OOP-koncept inte används på rätt sätt
  kan det leda till överdriven designkomplexitet och göra koden svårare
  att förstå och underhålla. Det är viktigt att använda OOP på ett sätt
  som gör koden enklare och inte mer komplicerad.
\end{itemize}

\hypertarget{tilluxe4mpningsomruxe5den}{%
\subsection{Tillämpningsområden}\label{tilluxe4mpningsomruxe5den}}

OOP är inte bara ett fancy koncept som bara fungerar i laboratoriemiljö.
Det har faktiskt många praktiska tillämpningar! Här är några exempel:

\begin{itemize}
\tightlist
\item
  \textbf{Applikationsutveckling}: Om vi bygger applikationer med
  komplexa datastrukturer och beteenden kan OOP vara till stor hjälp.
  Det hjälper oss att organisera koden och göra den lättare att hantera.
\item
  \textbf{Spelutveckling}: Har du någonsin velat bygga ditt eget spel?
  OOP är vägen att gå! Det hjälper oss att modellera spelobjekt, hantera
  spellogik och skapa interaktiva spelupplevelser.
\item
  \textbf{Webbutveckling}: Inom webbutveckling kan OOP hjälpa oss att
  skapa återanvändbara och skalbara komponenter. Det hjälper oss också
  att implementera designmönster som gör vår kod renare och mer
  effektiv.
\item
  \textbf{Simuleringar}: OOP används ofta inom simuleringar för att
  modellera och interagera med simulerade enheter och beteenden. Det kan
  vara allt från vetenskapliga simuleringar till spelutveckling.
\item
  \textbf{Databashantering}: Till och med databashantering kan dra nytta
  av OOP. Vi kan abstrahera databasåtkomst och använda objekt för att
  enkelt kommunicera med vår databas.
\end{itemize}

\hypertarget{exempel}{%
\subsection{Exempel}\label{exempel}}

Nu ska vi titta på ett konkret exempel för att se OOP i aktion! Vi ska
skapa en klass som representerar en bil i Java:

\hypertarget{cb1}{}
\begin{Shaded}
\begin{Highlighting}[]
\KeywordTok{public} \KeywordTok{class}\NormalTok{ Bil }\OperatorTok{\{}
\KeywordTok{private} \BuiltInTok{String}\NormalTok{ märke}\OperatorTok{;}
\KeywordTok{private} \BuiltInTok{String}\NormalTok{ modell}\OperatorTok{;}
\KeywordTok{private} \DataTypeTok{int}\NormalTok{ årsmodell}\OperatorTok{;}
\KeywordTok{public} \FunctionTok{Bil}\OperatorTok{(}\BuiltInTok{String}\NormalTok{ märke}\OperatorTok{,} \BuiltInTok{String}\NormalTok{ modell}\OperatorTok{,} \DataTypeTok{int}\NormalTok{ årsmodell}\OperatorTok{)} \OperatorTok{\{}
\KeywordTok{this}\OperatorTok{.}\FunctionTok{m}\NormalTok{ärke }\OperatorTok{=}\NormalTok{ märke}\OperatorTok{;}
\KeywordTok{this}\OperatorTok{.}\FunctionTok{modell} \OperatorTok{=}\NormalTok{ modell}\OperatorTok{;}
\KeywordTok{this}\OperatorTok{.}\NormalTok{årsmodell }\OperatorTok{=}\NormalTok{ årsmodell}\OperatorTok{;}
\OperatorTok{\}}
\KeywordTok{public} \BuiltInTok{String} \FunctionTok{getMärke}\OperatorTok{()} \OperatorTok{\{}
\ControlFlowTok{return}\NormalTok{ märke}\OperatorTok{;}
\KeywordTok{public} \DataTypeTok{void} \FunctionTok{setMärke}\OperatorTok{(}\BuiltInTok{String}\NormalTok{ märke}\OperatorTok{)} \OperatorTok{\{}
\KeywordTok{public} \BuiltInTok{String} \FunctionTok{getModell}\OperatorTok{()} \OperatorTok{\{}
\ControlFlowTok{return}\NormalTok{ modell}\OperatorTok{;}
\KeywordTok{public} \DataTypeTok{void} \FunctionTok{setModell}\OperatorTok{(}\BuiltInTok{String}\NormalTok{ modell}\OperatorTok{)} \OperatorTok{\{}
\KeywordTok{public} \DataTypeTok{int} \FunctionTok{getÅrsmodell}\OperatorTok{()} \OperatorTok{\{}
\ControlFlowTok{return}\NormalTok{ årsmodell}\OperatorTok{;}
\KeywordTok{public} \DataTypeTok{void} \FunctionTok{setÅrsmodell}\OperatorTok{(}\DataTypeTok{int}\NormalTok{ årsmodell}\OperatorTok{)} \OperatorTok{\{}
\KeywordTok{public} \DataTypeTok{void} \FunctionTok{köra}\OperatorTok{()} \OperatorTok{\{}
\BuiltInTok{System}\OperatorTok{.}\FunctionTok{out}\OperatorTok{.}\FunctionTok{println}\OperatorTok{(}\StringTok{"Bilen kör."}\OperatorTok{);}
\KeywordTok{public} \DataTypeTok{void} \FunctionTok{stanna}\OperatorTok{()} \OperatorTok{\{}
\BuiltInTok{System}\OperatorTok{.}\FunctionTok{out}\OperatorTok{.}\FunctionTok{println}\OperatorTok{(}\StringTok{"Bilen stannar."}\OperatorTok{);}
\OperatorTok{\}}
\end{Highlighting}
\end{Shaded}

I det här exemplet har vi en klass som heter ``Bil''. Den har tre
privata instansvariabler: ``märke'', ``modell'' och ``årsmodell''. Vi
har också en konstruktor som tar emot dessa tre variabler och tilldelar
dem till de motsvarande instansvariablerna.

Vi har också getter- och setter-metoder för varje instansvariabel för
att kunna få och sätta värden på dem. Dessutom har vi två metoder,
``köra'' och ``stanna'', som bara skriver ut meddelanden till konsolen.

Nu kan vi skapa objekt av klassen ``Bil'' och använda dess metoder:

\hypertarget{cb2}{}
\begin{Shaded}
\begin{Highlighting}[]
\KeywordTok{public} \KeywordTok{class}\NormalTok{ Main }\OperatorTok{\{}
\KeywordTok{public} \DataTypeTok{static} \DataTypeTok{void} \FunctionTok{main}\OperatorTok{(}\BuiltInTok{String}\OperatorTok{[]}\NormalTok{ args}\OperatorTok{)} \OperatorTok{\{}
\NormalTok{Bil bil }\OperatorTok{=} \KeywordTok{new} \FunctionTok{Bil}\OperatorTok{(}\StringTok{"Volvo"}\OperatorTok{,} \StringTok{"V70"}\OperatorTok{,} \DecValTok{2022}\OperatorTok{);}
\BuiltInTok{System}\OperatorTok{.}\FunctionTok{out}\OperatorTok{.}\FunctionTok{println}\OperatorTok{(}\StringTok{"Märke: "} \OperatorTok{+}\NormalTok{ bil}\OperatorTok{.}\FunctionTok{getM}\NormalTok{ärke}\OperatorTok{());}
\BuiltInTok{System}\OperatorTok{.}\FunctionTok{out}\OperatorTok{.}\FunctionTok{println}\OperatorTok{(}\StringTok{"Modell: "} \OperatorTok{+}\NormalTok{ bil}\OperatorTok{.}\FunctionTok{getModell}\OperatorTok{());}
\BuiltInTok{System}\OperatorTok{.}\FunctionTok{out}\OperatorTok{.}\FunctionTok{println}\OperatorTok{(}\StringTok{"Årsmodell: "} \OperatorTok{+}\NormalTok{ bil}\OperatorTok{.}\FunctionTok{get}\NormalTok{Årsmodell}\OperatorTok{());}
\NormalTok{bil}\OperatorTok{.}\FunctionTok{k}\NormalTok{öra}\OperatorTok{();}
\NormalTok{bil}\OperatorTok{.}\FunctionTok{stanna}\OperatorTok{();}
\OperatorTok{\}}
\OperatorTok{\}}
\end{Highlighting}
\end{Shaded}

I detta exempel skapar vi en instans av klassen ``Bil'' med hjälp av
konstruktorn och tilldelar den till variabeln ``bil''. Sedan använder vi
getter-metoderna för att få värdena på instansvariablerna och skriver ut
dem till konsolen. Slutligen anropar vi metoder för att köra och stanna
bilen.

\hypertarget{cb3}{}
\begin{Shaded}
\begin{Highlighting}[]
\KeywordTok{public} \KeywordTok{class}\NormalTok{ Car }\OperatorTok{\{}
\KeywordTok{private} \BuiltInTok{String}\NormalTok{ brand}\OperatorTok{;}
\KeywordTok{private} \BuiltInTok{String}\NormalTok{ color}\OperatorTok{;}
\KeywordTok{private} \DataTypeTok{int}\NormalTok{ speed}\OperatorTok{;}
\KeywordTok{public} \FunctionTok{Car}\OperatorTok{(}\BuiltInTok{String}\NormalTok{ brand}\OperatorTok{,} \BuiltInTok{String}\NormalTok{ color}\OperatorTok{)} \OperatorTok{\{}
\KeywordTok{this}\OperatorTok{.}\FunctionTok{brand} \OperatorTok{=}\NormalTok{ brand}\OperatorTok{;}
\KeywordTok{this}\OperatorTok{.}\FunctionTok{color} \OperatorTok{=}\NormalTok{ color}\OperatorTok{;}
\KeywordTok{this}\OperatorTok{.}\FunctionTok{speed} \OperatorTok{=} \DecValTok{0}\OperatorTok{;}
\KeywordTok{public} \DataTypeTok{void} \FunctionTok{accelerate}\OperatorTok{(}\DataTypeTok{int}\NormalTok{ value}\OperatorTok{)} \OperatorTok{\{}
\NormalTok{speed }\OperatorTok{+=}\NormalTok{ value}\OperatorTok{;}
\OperatorTok{\}}
\KeywordTok{public} \DataTypeTok{void} \FunctionTok{brake}\OperatorTok{(}\DataTypeTok{int}\NormalTok{ value}\OperatorTok{)} \OperatorTok{\{}
\NormalTok{speed }\OperatorTok{{-}=}\NormalTok{ value}\OperatorTok{;}
\OperatorTok{\}}
\KeywordTok{public} \DataTypeTok{int} \FunctionTok{getSpeed}\OperatorTok{()} \OperatorTok{\{}
\ControlFlowTok{return}\NormalTok{ speed}\OperatorTok{;}
\OperatorTok{\}}

\KeywordTok{public} \KeywordTok{class}\NormalTok{ Main }\OperatorTok{\{}
\KeywordTok{public} \DataTypeTok{static} \DataTypeTok{void} \FunctionTok{main}\OperatorTok{(}\BuiltInTok{String}\OperatorTok{[]}\NormalTok{ args}\OperatorTok{)} \OperatorTok{\{}
\NormalTok{Car myCar }\OperatorTok{=} \KeywordTok{new} \FunctionTok{Car}\OperatorTok{(}\StringTok{"Volvo"}\OperatorTok{,} \StringTok{"blå"}\OperatorTok{);}
\NormalTok{myCar}\OperatorTok{.}\FunctionTok{accelerate}\OperatorTok{(}\DecValTok{20}\OperatorTok{);}
\BuiltInTok{System}\OperatorTok{.}\FunctionTok{out}\OperatorTok{.}\FunctionTok{println}\OperatorTok{(}\NormalTok{myCar}\OperatorTok{.}\FunctionTok{getSpeed}\OperatorTok{());} \CommentTok{// Output: 20}
\NormalTok{myCar}\OperatorTok{.}\FunctionTok{brake}\OperatorTok{(}\DecValTok{10}\OperatorTok{);}
\BuiltInTok{System}\OperatorTok{.}\FunctionTok{out}\OperatorTok{.}\FunctionTok{println}\OperatorTok{(}\NormalTok{myCar}\OperatorTok{.}\FunctionTok{getSpeed}\OperatorTok{());} \CommentTok{// Output: 10}
\OperatorTok{\}}
\OperatorTok{\}}
\end{Highlighting}
\end{Shaded}

I detta exempel skapar vi en bilklass (\texttt{Car}) med egenskaper och
metoder för att accelerera, bromsa och hämta hastigheten. I
\texttt{main}-metoden skapar vi en instans av \texttt{Car} och testar
dess metoder.

\hypertarget{slutsats}{%
\subsection{Slutsats}\label{slutsats}}

All right! Du har precis fått en introduktion till Objektorienterad
programmering (OOP). Det är en viktig del av programmeringsvärlden och
något som du definitivt bör utforska mer! Med OOP kan du organisera din
kod på ett sätt som gör den lättare att förstå, underhålla och
återanvända. Tänk på fördelarna och användningsområdena vi har
diskuterat och låt det inspirera dig att ta dig an OOP och skapa
fantastiska saker! Lycka till!

\hypertarget{tldr}{%
\subsection{TL;DR}\label{tldr}}

OOP är en programmeringsmetodik som hjälper oss att organisera vår kod
genom att skapa objekt med egenskaper och beteenden. Det ger oss
modularitet, återanvändbarhet och lättare underhåll. OOP kan användas
inom applikationsutveckling, spelutveckling, webbutveckling,
simuleringar och databashantering. Ta en titt på vårt exempel med en
bilklass och kom igång med OOP! Nu är det dags att dyka djupare in i
världen av OOP och skapa magi med din kod!

\hypertarget{klasser-och-objekt-inom-programmering-med-java}{%
\section{Klasser och Objekt inom programmering med
Java}\label{klasser-och-objekt-inom-programmering-med-java}}

En artikel som utforskar ämnet Klasser och Objekt inom programmering med
Java.

Innehållsförteckning

\{: .text-delta \}

\begin{enumerate}
\def\labelenumi{\arabic{enumi}.}
\tightlist
\item
  TOC \{:toc\}
\end{enumerate}

\hypertarget{tldr}{%
\subsection{TL;DR}\label{tldr}}

Klasser och Objekt är centrala koncept inom objektorienterad
programmering med Java. Klasser fungerar som mallar för att skapa objekt
med definierade egenskaper och beteenden. Genom att använda Klasser och
Objekt kan vi skapa modulära och återanvändbara program, hantera
komplexitet och modellera relationer mellan olika entiteter. Det finns
fördelar som modularitet och abstraktion, men också begränsningar som
inlärningskurva och designkomplexitet. Klasser och Objekt kan tillämpas
inom olika områden som applikationsutveckling, spelutveckling och
simuleringar.

\hypertarget{nuxe4r-du-luxe4st-detta-ska-du-kunna}{%
\subsection{När du läst detta ska du
kunna}\label{nuxe4r-du-luxe4st-detta-ska-du-kunna}}

\begin{itemize}
\tightlist
\item
  Förstå och förklara vad Klasser och Objekt är och deras relevans inom
  programmering med Java.
\item
  Diskutera fördelar och begränsningar med Klasser och Objekt.
\item
  Identifiera olika användningsområden där Klasser och Objekt kan
  tillämpas.
\item
  Förstå och tolka ett kodexempel som använder Klasser och Objekt i
  Java.
\item
  Sammanfatta viktiga insikter och rekommendationer för vidare läsning.
\end{itemize}

\hypertarget{introduktion}{%
\subsection{Introduktion}\label{introduktion}}

Välkommen till denna artikel som kommer att utforska konceptet Klasser
och Objekt inom programmering med Java. Klasser och Objekt är
grundläggande byggstenar inom objektorienterad programmering (OOP) och
spelar en central roll för att skapa modulära och strukturerade program.
Genom att förstå och behärska dessa koncept kan du skapa kraftfulla och
flexibla program som är lättare att underhålla och återanvända.

\hypertarget{vad-uxe4r-klasser-och-objekt}{%
\subsection{Vad är Klasser och
Objekt?}\label{vad-uxe4r-klasser-och-objekt}}

Klasser och Objekt är centrala koncept inom objektorienterad
programmering med Java som hjälper oss att organisera och strukturera
vår kod på ett mer logiskt sätt.

En \textbf{klass} kan betraktas som en mall eller ritning som definierar
hur ett objekt av den klassen ska se ut och agera. Klassen innehåller
definitioner av attribut (också kallade egenskaper) och metoder som
beskriver objektets egenskaper och beteenden.

Ett \textbf{objekt} är en specifik förekomst av en klass. Det kan ses
som en konkret entitet med egenskaper som är definierade av klassen och
som kan utföra de metoder som klassen tillhandahåller.

En analogi som ofta används för att förklara klasser och objekt är att
klassen är som en ritning för att skapa ett hus och objektet är själva
huset som skapas enligt ritningen. Ritningen definierar husets
egenskaper (antalet rum, färg, storlek, etc.) och metoder (öppna dörren,
tända lampan, etc.), medan själva huset är den specifika instansen som
skapas baserat på ritningen.

\hypertarget{klassers-namn}{%
\subsubsection{Klassers namn}\label{klassers-namn}}

Klasser kallas för olika namn beroende på hur de används. Här är den
uppdaterade tabellen med basklass, subklass och några namn för olika
designmönster:

\begin{longtable}[]{@{}
  >{\raggedright\arraybackslash}p{(\columnwidth - 6\tabcolsep) * \real{0.0700}}
  >{\raggedright\arraybackslash}p{(\columnwidth - 6\tabcolsep) * \real{0.7900}}
  >{\raggedright\arraybackslash}p{(\columnwidth - 6\tabcolsep) * \real{0.1200}}
  >{\raggedright\arraybackslash}p{(\columnwidth - 6\tabcolsep) * \real{0.0100}}@{}}
\toprule\noalign{}
\begin{minipage}[b]{\linewidth}\raggedright
Term
\end{minipage} & \begin{minipage}[b]{\linewidth}\raggedright
Förklaring
\end{minipage} & \begin{minipage}[b]{\linewidth}\raggedright
Andra namn
\end{minipage} & \begin{minipage}[b]{\linewidth}\raggedright
\end{minipage} \\
\midrule\noalign{}
\endhead
\bottomrule\noalign{}
\endlastfoot
Klass & En klass är som en ritning för att bygga något. Det berättar
vilka delar och funktioner som något ska ha när det skapas. & & \\
Abstrakt klass & En abstrakt klass är som en mall där man bara får en
idé om hur något ska se ut och fungera, men inte exakt hur det ska
göras. & & \\
Gränssnitt & Ett gränssnitt är som en överenskommelse där man bestämmer
vilka saker man behöver kunna göra för att passa in i en viss grupp. &
Interface & \\
Konkret klass & En konkret klass är som en färdig produkt som kan
användas direkt, precis som en leksak som man kan leka med direkt när
man köper den. & & \\
Basklass & En basklass är som en övergripande klass som innehåller
gemensam funktionalitet och egenskaper som är ärvt av flera subklasser.
& Superklass & \\
Subklass & En subklass är som en specialiserad klass som ärver
funktionalitet och egenskaper från en basklass och kan lägga till eller
ändra dem. & Underklass & \\
Instans & En instans är som en specifik sak som man skapar baserat på
ritningen eller mallen i en klass. & Objekt & \\
Modell & En modell är som en beskrivning av hur något ser ut och
fungerar, t.ex. hur man ska spara information i en databas. & & \\
Typ & En typ är som en kategori som används för att organisera och
hantera olika sorters saker, som olika typer av mätningar eller enheter.
& & \\
Objekt & Ett objekt är som en grej som man kan röra vid och använda. Det
är det konkreta resultatet av att använda en ritning eller mall. & & \\
Entitet & En entitet är som en specifik grej som man vill hålla reda på
och spara information om, som en person eller en produkt i en databas. &
& \\
Helper & En helper är som en bästa vän som alltid är där för att hjälpa
till med små uppgifter och göra saker lite enklare för dig. &
Hjälpklass, Verktygsklass & \\
Utility & En utility är som en verktygslåda full med användbara saker
som hjälper dig att lösa specifika problem eller göra svåra saker
enklare. & & \\
Komponent & En komponent är som en del av något större, som en pusselbit
som passar in i en större bild och har en specifik funktion. & & \\
Enhetsobjekt & Ett enhetsobjekt är som en representation av en specifik
enhet eller apparat som kan utföra vissa uppgifter eller ha viss funkt &
ionalitet. & \\
Verktyg & Ett verktyg är som en hjälpande hand som underlättar och
effektiviserar olika uppgifter eller processer. & & \\
Modul & En modul är som en självständig del av ett större system, som
kan kopplas in och användas för att utföra specifika uppgifter. & & \\
Singleton & Singleton är ett designmönster som används för att se till
att endast en instans av en klass skapas och att den kan nås globalt. &
& \\
Fabrik & Fabrik är ett designmönster som används för att skapa objekt
utan att avslöja den konkreta implementationen och istället använda en
gemensam gränssnitt. & Factory & \\
Byggare & Byggare är ett designmönster som används för att skapa objekt
stegvis och möjliggör olika sätt att bygga upp komplexa objekt. &
Builder & \\
Strategi & Strategi är ett designmönster som används för att välja och
använda en av flera möjliga algoritmer eller beteenden vid körningstid.
& Strategy & \\
Observatör & Observatör är ett designmönster som används för att
etablera en publikera/prenumerera-mekanism där flera objekt kan lyssna
på händelser från ett annat objekt. & Observer & \\
\end{longtable}

Som du ser har klasser många namn, man kan säga att det är en klass för
varje tillfälle. Men inget att oroa sig för, du kommer att lära dig mer
om dem allt eftersom du fortsätter att lära dig om programmering. Ju mer
man använder klasser, desto bättre lär man sig namnen.

\hypertarget{fuxf6rdelar}{%
\subsection{Fördelar}\label{fuxf6rdelar}}

Användningen av Klasser och Objekt erbjuder flera fördelar inom
programmering med Java:

\begin{enumerate}
\def\labelenumi{\arabic{enumi}.}
\item
  \textbf{Modularitet och återanvändbarhet}: Klasser möjliggör modulär
  kod genom att separera olika delar av programmet i olika klasser.
  Detta gör det enklare att hantera och underhålla koden samt möjliggör
  återanvändning av kod genom att skapa nya objekt baserat på en
  befintlig klass.
\item
  \textbf{Abstraktion och hantering av komplexitet}: Genom att använda
  klasser kan vi abstrahera bort detaljer och fokusera på de väsentliga
  egenskaperna och beteendena hos objektet. Detta hjälper till att
  hantera komplexitet och gör koden mer läsbar och underhållbar.
\item
  \textbf{Kapsling och informationsskydd}: Klasser möjliggör att vi kan
  begränsa åtkomsten till objektets egenskaper och metoder. Detta
  främjar informationsskydd och hjälper till att undvika oavsiktliga
  ändringar av objektets tillstånd.
\item
  \textbf{Hantering av relationer mellan objekt}: Genom att använda
  klasser och objekt kan vi definiera och hantera relationer mellan
  olika objekt. Detta möjliggör att vi kan skapa mer realistiska och
  flexibla modeller av den verkliga världen i våra program.
\end{enumerate}

\hypertarget{begruxe4nsningar}{%
\subsection{Begränsningar}\label{begruxe4nsningar}}

Det finns också vissa begränsningar eller utmaningar med att använda
Klasser och Objekt i Java:

\begin{enumerate}
\def\labelenumi{\arabic{enumi}.}
\item
  \textbf{Inlärningskurva}: Konceptet med Klasser och Objekt kan vara
  svårt att förstå i början, särskilt för nybörjare inom programmering.
  Att förstå sambandet mellan klasser och objekt och att behärska olika
  OOP-koncept kan ta tid och övning.
\item
  \textbf{Prestandaöverhead}: Objektorienterad kod kan vara lite mer
  resurskrävande än procedurorienterad kod på grund av det extra
  lagringsutrymmet som krävs för att hålla reda på objektens tillstånd
  och beteenden. Detta kan vara en faktor att ta hänsyn till när det
  krävs maximal prestanda i en applikation.
\item
  \textbf{Designkomplexitet}: Att designa och planera klasser och objekt
  kan vara en utmaning. Att hitta rätt abstraktionsnivå, definiera
  korrekta relationer och undvika överkomplicerade hierarkier kan vara
  svårt och kräver erfarenhet och bra designprinciper.
\end{enumerate}

\hypertarget{anvuxe4ndningsomruxe5den}{%
\subsection{Användningsområden}\label{anvuxe4ndningsomruxe5den}}

Klasser och Objekt kan tillämpas i en mängd olika scenarier inom
programmering med Java. Här är några exempel på användningsområden:

\begin{enumerate}
\def\labelenumi{\arabic{enumi}.}
\item
  \textbf{Applikationsutveckling}: Klasser och Objekt används i stor
  utsträckning vid utveckling av applikationer, oavsett om det är
  webbapplikationer, mobilappar eller skrivbordsprogram. Genom att
  strukturera koden med hjälp av klasser och objekt blir programmet mer
  organiserat och lättare att underhålla.
\item
  \textbf{Spelutveckling}: Inom spelutveckling används Klasser och
  Objekt för att skapa olika spelobjekt, karaktärer, världar och mycket
  mer. Genom att använda objektorienterad programmering kan
  spelutvecklare skapa komplexa och interaktiva spelvärldar.
\item
  \textbf{Simuleringar}: Simuleringsprogram och modelleringsverktyg kan
  dra nytta av Klasser och Objekt för att representera och simulera
  olika entiteter och processer. Genom att använda objekt för att
  modellera olika aspekter av systemet kan simuleringar bli mer
  realistiska och flexibla.
\item
  \textbf{Databashantering}: Vid databashantering används ofta
  objektorienterade koncept för att modellera och hantera data. Objekt
  kan representera tabeller, rader och kolumner i en databas och
  möjliggöra en mer flexibel och hanterbar databasstruktur.
\end{enumerate}

\hypertarget{exempelkod}{%
\subsection{Exempelkod}\label{exempelkod}}

För att illustrera användningen av Klasser och Objekt, låt oss tänka oss
att vi bygger ett program för att hantera en biblioteksdatabase. Vi kan
använda Klasser och Objekt för att representera olika entiteter i
biblioteket, till exempel böcker och medlemmar. Här är ett kodexempel i
Java:

\hypertarget{cb1}{}
\begin{Shaded}
\begin{Highlighting}[]
\KeywordTok{public} \KeywordTok{class} \BuiltInTok{Book} \OperatorTok{\{}
\KeywordTok{private} \BuiltInTok{String}\NormalTok{ title}\OperatorTok{;}
\KeywordTok{private} \BuiltInTok{String}\NormalTok{ author}\OperatorTok{;}
\KeywordTok{private} \DataTypeTok{int}\NormalTok{ year}\OperatorTok{;}

\KeywordTok{public} \BuiltInTok{Book}\OperatorTok{(}\BuiltInTok{String}\NormalTok{ title}\OperatorTok{,} \BuiltInTok{String}\NormalTok{ author}\OperatorTok{,} \DataTypeTok{int}\NormalTok{ year}\OperatorTok{)} \OperatorTok{\{}
\KeywordTok{this}\OperatorTok{.}\FunctionTok{title} \OperatorTok{=}\NormalTok{ title}\OperatorTok{;}
\KeywordTok{this}\OperatorTok{.}\FunctionTok{author} \OperatorTok{=}\NormalTok{ author}\OperatorTok{;}
\KeywordTok{this}\OperatorTok{.}\FunctionTok{year} \OperatorTok{=}\NormalTok{ year}\OperatorTok{;}
\OperatorTok{\}}

\KeywordTok{public} \DataTypeTok{void} \FunctionTok{displayInfo}\OperatorTok{()} \OperatorTok{\{}
\BuiltInTok{System}\OperatorTok{.}\FunctionTok{out}\OperatorTok{.}\FunctionTok{println}\OperatorTok{(}\StringTok{"Title: "} \OperatorTok{+}\NormalTok{ title}\OperatorTok{);}
\BuiltInTok{System}\OperatorTok{.}\FunctionTok{out}\OperatorTok{.}\FunctionTok{println}\OperatorTok{(}\StringTok{"Author: "} \OperatorTok{+}\NormalTok{ author}\OperatorTok{);}
\BuiltInTok{System}\OperatorTok{.}\FunctionTok{out}\OperatorTok{.}\FunctionTok{println}\OperatorTok{(}\StringTok{"Year: "} \OperatorTok{+}\NormalTok{ year}\OperatorTok{);}
\OperatorTok{\}}
\OperatorTok{\}}
\end{Highlighting}
\end{Shaded}

Nu har vi skapat en klass för böcker som innehåller information om
titel, författare och år. Vi kan sedan skapa ett objekt av klassen och
använda dess metoder för att visa informationen om boken:

\hypertarget{cb2}{}
\begin{Shaded}
\begin{Highlighting}[]
\KeywordTok{public} \KeywordTok{class} \BuiltInTok{Member} \OperatorTok{\{}
\KeywordTok{private} \BuiltInTok{String}\NormalTok{ name}\OperatorTok{;}
\KeywordTok{private} \DataTypeTok{int}\NormalTok{ age}\OperatorTok{;}
\KeywordTok{private} \BuiltInTok{List}\OperatorTok{\textless{}}\BuiltInTok{Book}\OperatorTok{\textgreater{}}\NormalTok{ borrowedBooks}\OperatorTok{;}

\KeywordTok{public} \BuiltInTok{Member}\OperatorTok{(}\BuiltInTok{String}\NormalTok{ name}\OperatorTok{,} \DataTypeTok{int}\NormalTok{ age}\OperatorTok{)} \OperatorTok{\{}
\KeywordTok{this}\OperatorTok{.}\FunctionTok{name} \OperatorTok{=}\NormalTok{ name}\OperatorTok{;}
\KeywordTok{this}\OperatorTok{.}\FunctionTok{age} \OperatorTok{=}\NormalTok{ age}\OperatorTok{;}
\KeywordTok{this}\OperatorTok{.}\FunctionTok{borrowedBooks} \OperatorTok{=} \KeywordTok{new} \BuiltInTok{ArrayList}\OperatorTok{\textless{}\textgreater{}();}
\OperatorTok{\}}

\KeywordTok{public} \DataTypeTok{void} \FunctionTok{borrowBook}\OperatorTok{(}\BuiltInTok{Book}\NormalTok{ book}\OperatorTok{)} \OperatorTok{\{}
\NormalTok{borrowedBooks}\OperatorTok{.}\FunctionTok{add}\OperatorTok{(}\NormalTok{book}\OperatorTok{);}
\OperatorTok{\}}

\KeywordTok{public} \DataTypeTok{void} \FunctionTok{displayInfo}\OperatorTok{()} \OperatorTok{\{}
\BuiltInTok{System}\OperatorTok{.}\FunctionTok{out}\OperatorTok{.}\FunctionTok{println}\OperatorTok{(}\StringTok{"Name: "} \OperatorTok{+}\NormalTok{ name}\OperatorTok{);}
\BuiltInTok{System}\OperatorTok{.}\FunctionTok{out}\OperatorTok{.}\FunctionTok{println}\OperatorTok{(}\StringTok{"Age: "} \OperatorTok{+}\NormalTok{ age}\OperatorTok{);}
\BuiltInTok{System}\OperatorTok{.}\FunctionTok{out}\OperatorTok{.}\FunctionTok{println}\OperatorTok{(}\StringTok{"Borrowed Books:"}\OperatorTok{);}
\ControlFlowTok{for} \OperatorTok{(}\BuiltInTok{Book}\NormalTok{ book }\OperatorTok{:}\NormalTok{ borrowedBooks}\OperatorTok{)} \OperatorTok{\{}
\BuiltInTok{System}\OperatorTok{.}\FunctionTok{out}\OperatorTok{.}\FunctionTok{println}\OperatorTok{(}\StringTok{"{-} "} \OperatorTok{+}\NormalTok{ book}\OperatorTok{.}\FunctionTok{getTitle}\OperatorTok{());}
\OperatorTok{\}}
\OperatorTok{\}}
\OperatorTok{\}}
\end{Highlighting}
\end{Shaded}

Här har vi skapat en klass för medlemmar som innehåller information om
namn, ålder och en lista över lånade böcker. Vi kan sedan skapa ett
objekt av klassen och använda dess metoder för att visa informationen om
medlemmen:

Nu kan vi använda klasserna för att skapa ett bibliotekssystem:

\hypertarget{cb3}{}
\begin{Shaded}
\begin{Highlighting}[]
\CommentTok{// Exempel på användning av Klasser och Objekt}
\BuiltInTok{Book}\NormalTok{ book }\OperatorTok{=} \KeywordTok{new} \BuiltInTok{Book}\OperatorTok{(}\StringTok{"The Catcher in the Rye"}\OperatorTok{,} \StringTok{"J.D. Salinger"}\OperatorTok{,} \DecValTok{1951}\OperatorTok{);}
\BuiltInTok{Member}\NormalTok{ member }\OperatorTok{=} \KeywordTok{new} \BuiltInTok{Member}\OperatorTok{(}\StringTok{"John Doe"}\OperatorTok{,} \DecValTok{30}\OperatorTok{);}
\NormalTok{member}\OperatorTok{.}\FunctionTok{borrowBook}\OperatorTok{(}\NormalTok{book}\OperatorTok{);}

\NormalTok{book}\OperatorTok{.}\FunctionTok{displayInfo}\OperatorTok{();}
\NormalTok{member}\OperatorTok{.}\FunctionTok{displayInfo}\OperatorTok{();}
\end{Highlighting}
\end{Shaded}

Vi kan skapa ett objekt av varje klass och använda deras metoder för att
visa informationen om boken och medlemmen.

I exemplet skapar vi en bok, sedan skapar vi en person som heter John
Doe och är 30 år gammal. Vi lånar sedan boken till John Doe och visar
informationen om boken och medlemmen.

När vi har båda klasserna instansierade kan vi nu tala om för boken att
visa sin information och för medlemmen att visa sin information. Vi kan
också lägga till en bok till medlemmens lånade böcker och sedan visa
informationen om medlemmen igen för att se att boken har lagts till.

\hypertarget{output}{%
\subsubsection{Output}\label{output}}

\begin{verbatim}
Title: The Catcher in the Rye
Author: J.D. Salinger
Year: 1951
Name: John Doe
Age: 30
Borrowed Books:
- The Catcher in the Rye
\end{verbatim}

Så kan vi använda klasserna, \texttt{Book} och \texttt{Member}, som
representerar en bok och en medlem i biblioteket. Vi kan skapa objekt av
varje klass och använda deras metoder för att visa informationen om
boken och medlemmen.

\hypertarget{slutsats}{%
\subsection{Slutsats}\label{slutsats}}

Klasser och Objekt är grundläggande koncept inom objektorienterad
programmering med Java som hjälper oss att organisera och strukturera
vår kod på ett mer logiskt och effektivt sätt. Genom att använda Klasser
och Objekt kan vi skapa modulära och återanvändbara program, hantera
komplexitet och modellera relationer mellan olika entiteter.

Det är viktigt att förstå både fördelarna och begränsningarna med att
använda Klasser och Objekt och att kunna identifiera olika
användningsområden där dessa koncept kan tillämpas. Genom att behärska
Klasser och Objekt kan du utveckla välstrukturerad och lättunderhållen
kod.

För vidare läsning och fördjupning rekommenderas att studera
OOP-principer, designmönster och mer avancerade koncept som arv och
polymorfism.

\hypertarget{obligatorisk-dad-joke}{%
\subsection{Obligatorisk dad joke}\label{obligatorisk-dad-joke}}

Varför ville objektet inte gå på festen?\\
För att det inte hade någon klass!

\hypertarget{klasskomposition}{%
\section{Klasskomposition}\label{klasskomposition}}

Klasskomposition är en teknik inom objektorienterad programmering (OOP)
där mindre klasser kombineras för att bygga större och mer komplexa
klasser. Genom att använda klasskomposition kan vi skapa mer modulära,
flexibla och underhållbara lösningar.

\hypertarget{fuxf6rdelar}{%
\subsection{Fördelar}\label{fuxf6rdelar}}

Att använda klasskomposition i din kod kan erbjuda flera fördelar:

\begin{itemize}
\item
  \textbf{Återanvändbarhet}: Genom att kombinera mindre och
  specialiserade klasser kan vi skapa återanvändbar kod. De mindre
  klasserna kan användas i olika sammanhang och kan även användas i
  andra större klasser om det behövs.
\item
  \textbf{Flexibilitet}: Genom att använda klasskomposition blir koden
  mer flexibel. Vi kan lägga till eller ta bort komponenter genom att
  lägga till eller ta bort de mindre klasserna som ingår i den större
  klassen.
\item
  \textbf{Tydlighet och underhållbarhet}: Genom att dela upp
  funktionaliteten i mindre och självständiga klasser blir koden lättare
  att förstå och underhålla. Varje klass har sitt eget ansvarsområde och
  kan ändras separat utan att påverka resten av systemet.
\end{itemize}

\hypertarget{begruxe4nsningar}{%
\subsection{Begränsningar}\label{begruxe4nsningar}}

Det finns också vissa begränsningar att tänka på när man använder
klasskomposition:

\begin{itemize}
\item
  \textbf{Komplexitet}: Klasskomposition kan öka komplexiteten i koden
  om den inte används på ett genomtänkt sätt. Det är viktigt att
  noggrant planera och designa klassrelationerna för att undvika onödig
  komplexitet.
\item
  \textbf{Beroenden}: Om en klass är beroende av en annan klass genom
  klasskomposition, kan ändringar i den andra klassen påverka
  funktionaliteten hos den första klassen. Det är viktigt att vara
  medveten om beroendena mellan klasserna och att hantera dem på ett
  lämpligt sätt.
\item
  \textbf{Prestanda}: Klasskomposition kan påverka prestandan hos
  programmet eftersom det kan innebära fler objekt och därmed mer
  minnesanvändning och bearbetningstid. Det är viktigt att överväga
  prestandaaspekterna när man använder klasskomposition och att optimera
  koden vid behov.
\end{itemize}

\hypertarget{exempel-puxe5-anvuxe4ndning}{%
\subsection{Exempel på användning}\label{exempel-puxe5-anvuxe4ndning}}

Klasskomposition kan tillämpas i en mängd olika scenarier. Här är några
exempel på användningsområden där klasskomposition är användbart:

\begin{itemize}
\item
  \textbf{Bil och motor}: En bilklass kan innehålla en instans av en
  motorklass för att hantera bilens motorfunktionalitet.
\item
  \textbf{Dator och processor}: En datorklass kan innehålla en instans
  av en processor för att hantera beräkningsfunktionalitet.
\item
  \textbf{Bank och konton}: En bankklass kan innehålla flera instanser
  av en kontoklass för att hantera olika bankkonton.
\item
  \textbf{Spel och karaktärer}: Ett spel kan ha en spelklass som
  innehåller flera instanser av en karaktärsklass för att hantera
  spelkaraktärer.
\end{itemize}

\hypertarget{kodexempel}{%
\subsection{Kodexempel}\label{kodexempel}}

Låt oss nu titta på ett kodexempel för att förstå hur klasskomposition
kan tillämpas i praktiken. I det här exemplet har vi två klasser:
\texttt{Person} och \texttt{Address}. Klassen \texttt{Person}
representerar en person och innehåller en instans av klassen
\texttt{Address} för att hantera adressrelaterad information.

\hypertarget{cb1}{}
\begin{Shaded}
\begin{Highlighting}[]
\KeywordTok{public} \KeywordTok{class}\NormalTok{ Address }\OperatorTok{\{}
\KeywordTok{private} \BuiltInTok{String}\NormalTok{ street}\OperatorTok{;}
\KeywordTok{private} \BuiltInTok{String}\NormalTok{ city}\OperatorTok{;}
\KeywordTok{private} \BuiltInTok{String}\NormalTok{ state}\OperatorTok{;}
\KeywordTok{private} \BuiltInTok{String}\NormalTok{ zipCode}\OperatorTok{;}

\KeywordTok{public} \FunctionTok{Address}\OperatorTok{(}\BuiltInTok{String}\NormalTok{ street}\OperatorTok{,} \BuiltInTok{String}\NormalTok{ city}\OperatorTok{,} \BuiltInTok{String}\NormalTok{ state}\OperatorTok{,} \BuiltInTok{String}\NormalTok{ zipCode}\OperatorTok{)} \OperatorTok{\{}
\KeywordTok{this}\OperatorTok{.}\FunctionTok{street} \OperatorTok{=}\NormalTok{ street}\OperatorTok{;}
\KeywordTok{this}\OperatorTok{.}\FunctionTok{city} \OperatorTok{=}\NormalTok{ city}\OperatorTok{;}
\KeywordTok{this}\OperatorTok{.}\FunctionTok{state} \OperatorTok{=}\NormalTok{ state}\OperatorTok{;}
\KeywordTok{this}\OperatorTok{.}\FunctionTok{zipCode} \OperatorTok{=}\NormalTok{ zipCode}\OperatorTok{;}
\OperatorTok{\}}

\CommentTok{// Getters and setters}

\KeywordTok{public} \BuiltInTok{String} \FunctionTok{getStreet}\OperatorTok{()} \OperatorTok{\{}\ControlFlowTok{return}\NormalTok{ street}\OperatorTok{;\}}
\KeywordTok{public} \BuiltInTok{String} \FunctionTok{getCity}\OperatorTok{()} \OperatorTok{\{}\ControlFlowTok{return}\NormalTok{ city}\OperatorTok{;\}}
\KeywordTok{public} \BuiltInTok{String} \FunctionTok{getState}\OperatorTok{()} \OperatorTok{\{}\ControlFlowTok{return}\NormalTok{ state}\OperatorTok{;\}}
\KeywordTok{public} \BuiltInTok{String} \FunctionTok{getZipCode}\OperatorTok{()} \OperatorTok{\{}\ControlFlowTok{return}\NormalTok{ zipCode}\OperatorTok{;\}}
\KeywordTok{public} \DataTypeTok{void} \FunctionTok{setStreet}\OperatorTok{(}\BuiltInTok{String}\NormalTok{ street}\OperatorTok{)} \OperatorTok{\{}\KeywordTok{this}\OperatorTok{.}\FunctionTok{street} \OperatorTok{=}\NormalTok{ street}\OperatorTok{;\}}
\KeywordTok{public} \DataTypeTok{void} \FunctionTok{setCity}\OperatorTok{(}\BuiltInTok{String}\NormalTok{ city}\OperatorTok{)} \OperatorTok{\{}\KeywordTok{this}\OperatorTok{.}\FunctionTok{city} \OperatorTok{=}\NormalTok{ city}\OperatorTok{;\}}
\KeywordTok{public} \DataTypeTok{void} \FunctionTok{setState}\OperatorTok{(}\BuiltInTok{String}\NormalTok{ state}\OperatorTok{)} \OperatorTok{\{}\KeywordTok{this}\OperatorTok{.}\FunctionTok{state} \OperatorTok{=}\NormalTok{ state}\OperatorTok{;\}}
\KeywordTok{public} \DataTypeTok{void} \FunctionTok{setZipCode}\OperatorTok{(}\BuiltInTok{String}\NormalTok{ zipCode}\OperatorTok{)} \OperatorTok{\{}\KeywordTok{this}\OperatorTok{.}\FunctionTok{zipCode} \OperatorTok{=}\NormalTok{ zipCode}\OperatorTok{;\}}

\AttributeTok{@Override}
\KeywordTok{public} \BuiltInTok{String} \FunctionTok{toString}\OperatorTok{()} \OperatorTok{\{}
\ControlFlowTok{return} \StringTok{"Address\{"} \OperatorTok{+}
\StringTok{"street=\textquotesingle{}"} \OperatorTok{+}\NormalTok{ street }\OperatorTok{+} \CharTok{\textquotesingle{}\textbackslash{}\textquotesingle{}\textquotesingle{}} \OperatorTok{+}
\StringTok{", city=\textquotesingle{}"} \OperatorTok{+}\NormalTok{ city }\OperatorTok{+} \CharTok{\textquotesingle{}\textbackslash{}\textquotesingle{}\textquotesingle{}} \OperatorTok{+}
\StringTok{", state=\textquotesingle{}"} \OperatorTok{+}\NormalTok{ state }\OperatorTok{+} \CharTok{\textquotesingle{}\textbackslash{}\textquotesingle{}\textquotesingle{}} \OperatorTok{+}
\StringTok{", zipCode=\textquotesingle{}"} \OperatorTok{+}\NormalTok{ zipCode }\OperatorTok{+} \CharTok{\textquotesingle{}\textbackslash{}\textquotesingle{}\textquotesingle{}} \OperatorTok{+}
\CharTok{\textquotesingle{}\}\textquotesingle{}}\OperatorTok{;}
\OperatorTok{\}}
\OperatorTok{\}}
\end{Highlighting}
\end{Shaded}

Nu har vi en klass som reprenterar en adress. Den innehåller attribut
för gatuadress, stad, län och postnummer. Klassen har också en
konstruktor som tar emot dessa attribut och en metod \texttt{toString()}
som returnerar en strängrepresentation av adressen.

\hypertarget{cb2}{}
\begin{Shaded}
\begin{Highlighting}[]
\KeywordTok{public} \KeywordTok{class}\NormalTok{ Person }\OperatorTok{\{}
\KeywordTok{private} \BuiltInTok{String}\NormalTok{ name}\OperatorTok{;}
\KeywordTok{private} \DataTypeTok{int}\NormalTok{ age}\OperatorTok{;}
\KeywordTok{private}\NormalTok{ Address address}\OperatorTok{;}

\KeywordTok{public} \FunctionTok{Person}\OperatorTok{(}\BuiltInTok{String}\NormalTok{ name}\OperatorTok{,} \DataTypeTok{int}\NormalTok{ age}\OperatorTok{,}\NormalTok{ Address address}\OperatorTok{)} \OperatorTok{\{}
\KeywordTok{this}\OperatorTok{.}\FunctionTok{name} \OperatorTok{=}\NormalTok{ name}\OperatorTok{;}
\KeywordTok{this}\OperatorTok{.}\FunctionTok{age} \OperatorTok{=}\NormalTok{ age}\OperatorTok{;}
\KeywordTok{this}\OperatorTok{.}\FunctionTok{address} \OperatorTok{=}\NormalTok{ address}\OperatorTok{;}
\OperatorTok{\}}

\KeywordTok{public} \BuiltInTok{String} \FunctionTok{getAddressDetails}\OperatorTok{()} \OperatorTok{\{}
\ControlFlowTok{return} \StringTok{"Address: "} \OperatorTok{+}\NormalTok{ address}\OperatorTok{.}\FunctionTok{getStreet}\OperatorTok{()} \OperatorTok{+} \StringTok{", "} \OperatorTok{+}\NormalTok{ address}\OperatorTok{.}\FunctionTok{getCity}\OperatorTok{()} \OperatorTok{+} \StringTok{", "} \OperatorTok{+}\NormalTok{ address}\OperatorTok{.}\FunctionTok{getState}\OperatorTok{()} \OperatorTok{+} \StringTok{", "} \OperatorTok{+}\NormalTok{ address}\OperatorTok{.}\FunctionTok{getZipCode}\OperatorTok{();}
\OperatorTok{\}}
\OperatorTok{\}}
\end{Highlighting}
\end{Shaded}

Här har vi en klass som representerar en person. Den innehåller attribut
för namn, ålder och en instans av klassen \texttt{Address} för att
hantera adressuppgifter. Klassen har också en metod
\texttt{getAddressDetails()} som returnerar en strängrepresentation av
personens adress.

\hypertarget{cb3}{}
\begin{Shaded}
\begin{Highlighting}[]
\CommentTok{// Användning av klasskomposition}
\NormalTok{Address address }\OperatorTok{=} \KeywordTok{new} \FunctionTok{Address}\OperatorTok{(}\StringTok{"123 Main St"}\OperatorTok{,} \StringTok{"Cityville"}\OperatorTok{,} \StringTok{"Stateville"}\OperatorTok{,} \StringTok{"12345"}\OperatorTok{);}
\NormalTok{Person person }\OperatorTok{=} \KeywordTok{new} \FunctionTok{Person}\OperatorTok{(}\StringTok{"John Doe"}\OperatorTok{,} \DecValTok{25}\OperatorTok{,}\NormalTok{ address}\OperatorTok{);}
\BuiltInTok{System}\OperatorTok{.}\FunctionTok{out}\OperatorTok{.}\FunctionTok{println}\OperatorTok{(}\NormalTok{person}\OperatorTok{.}\FunctionTok{getAddressDetails}\OperatorTok{());}
\end{Highlighting}
\end{Shaded}

I det här exemplet har vi en \texttt{Address}-klass som innehåller
attribut för gatuadress, stad, delstat och postnummer. Klassen
\texttt{Person} innehåller attribut för namn, ålder och en instans av
\texttt{Address}-klassen för att hantera personens adressuppgifter.
Klassen \texttt{Person} har också en metod \texttt{getAddressDetails()}
som returnerar en strängrepresentation av personens adress.

För att använda dessa klasser kan vi skapa en instans av \texttt{Person}
och tilldela en instans av \texttt{Address} till dess
\texttt{address}-attribut. Vi kan sedan använda
\texttt{getAddressDetails()}-metoden för att få personens
adressuppgifter.

\hypertarget{cb4}{}
\begin{Shaded}
\begin{Highlighting}[]
\NormalTok{Address address }\OperatorTok{=} \KeywordTok{new} \FunctionTok{Address}\OperatorTok{(}\StringTok{"123 Main St"}\OperatorTok{,} \StringTok{"Cityville"}\OperatorTok{,} \StringTok{"Stateville"}\OperatorTok{,} \StringTok{"12345"}\OperatorTok{);}
\NormalTok{Person person }\OperatorTok{=} \KeywordTok{new} \FunctionTok{Person}\OperatorTok{(}\StringTok{"John Doe"}\OperatorTok{,} \DecValTok{25}\OperatorTok{,}\NormalTok{ address}\OperatorTok{);}
\BuiltInTok{System}\OperatorTok{.}\FunctionTok{out}\OperatorTok{.}\FunctionTok{println}\OperatorTok{(}\NormalTok{person}\OperatorTok{.}\FunctionTok{getAddressDetails}\OperatorTok{());}
\end{Highlighting}
\end{Shaded}

\hypertarget{output}{%
\subsubsection{Output}\label{output}}

\begin{verbatim}
Address: 123 Main St, Cityville, Stateville, 12345
\end{verbatim}

I det här kodexemplet skapar vi en instans av \texttt{Address} med
adressuppgifterna ``123 Main St'', ``Cityville'', ``Stateville'' och
``12345''. Vi skapar sedan en instans av \texttt{Person}

med namnet ``John Doe'', ålder 25 och den tidigare skapade
adressinstansen. Slutligen skriver vi ut personens adressuppgifter med
hjälp av \texttt{getAddressDetails()}-metoden.

\hypertarget{slutsats}{%
\subsection{Slutsats}\label{slutsats}}

Klasskomposition är en kraftfull teknik inom objektorienterad
programmering som låter oss skapa mer flexibla, återanvändbara och
underhållbara lösningar. Genom att kombinera mindre och specialiserade
klasser kan vi bygga större och mer komplexa system utan att offra
tydlighet och prestanda. Det är viktigt att förstå fördelarna och
begränsningarna med klasskomposition för att kunna använda det på bästa
sätt.

För vidare läsning och fördjupning i ämnet rekommenderas att utforska
designprinciper och mönster inom objektorienterad programmering, såsom
``SOLID-principerna'' och ``komposition över arv''. Genom att utveckla
din kunskap om dessa koncept kan du bli en skickligare programmerare och
bygga mer effektiva och flexibla program.

Absolut! Här är hela artikeln om inkapsling i Java:

\hypertarget{inkapsling}{%
\section{Inkapsling}\label{inkapsling}}

Inkapsling är en viktig princip inom objektorienterad programmering som
handlar om att kombinera data och metoder inom en klass och kontrollera
åtkomsten till dem. Det främjar säkerhet, moduläritet, återanvändbarhet
och kodunderhåll. Inkapsling kan tillämpas i dataklasser, API-design och
användas tillsammans med arv och polymorfism.

Innehållsförteckning

\{: .text-delta \} 1. Innehållsförteckning \{:toc\}

\hypertarget{vad-uxe4r-inkapsling}{%
\subsection{Vad är inkapsling?}\label{vad-uxe4r-inkapsling}}

Inkapsling innebär att kombinera data och metoder inom en klass och
kontrollera åtkomsten till dem. Genom att använda inkapsling kan vi
definiera vilka medlemmar som är tillgängliga utanför klassen och hur de
kan manipuleras. Detta möjliggör en tydlig separation mellan
implementationen av en klass och dess användning i andra delar av
programmet.

Inkapsling har flera fördelar inom programmering:

\begin{enumerate}
\def\labelenumi{\arabic{enumi}.}
\item
  \textbf{Moduläritet}: Inkapsling möjliggör att klasser kan vara
  självständiga enheter med en tydlig gränssnitt för att interagera med
  andra delar av programmet. Detta underlättar utveckling, underhåll och
  återanvändbarhet av kod, eftersom det gör det lättare att isolera och
  ändra specifika delar av koden utan att påverka resten av programmet.
\item
  \textbf{Säkerhet}: Genom att använda inkapsling kan vi begränsa
  åtkomsten till viss data och funktionalitet. Detta gör det möjligt att
  skydda känslig information och förhindra oavsiktlig manipulation av
  data. Genom att använda privata medlemmar kan vi säkerställa att vissa
  operationer bara kan utföras inom klassen eller av utvalda metoder.
\item
  \textbf{Kodunderhåll}: Inkapsling främjar en bättre struktur och
  organisation av koden. Genom att gruppera relaterade data och
  funktioner inom en klass blir det lättare att förstå och ändra koden.
  Detta underlättar underhåll och felsökning, eftersom det minimerar
  risken för biverkningar och felaktiga ändringar i andra delar av
  programmet.
\end{enumerate}

\hypertarget{begruxe4nsningar-med-inkapsling}{%
\subsection{Begränsningar med
inkapsling}\label{begruxe4nsningar-med-inkapsling}}

Även om inkapsling har många fördelar finns det också vissa
begränsningar att vara medveten om:

\begin{enumerate}
\def\labelenumi{\arabic{enumi}.}
\item
  \textbf{Komplexitet}: Inkapsling kan introducera en viss komplexitet i
  koden, särskilt när det gäller att hantera olika åtkomstnivåer och
  beroenden mellan klasser. Det är viktigt att noga planera och
  organisera klasser för att undvika överdriven komplexitet och onödiga
  beroenden.
\item
  \textbf{Prestanda}: Vissa åtkomstmodifierare, som \texttt{private},
  kan medföra en viss prestandaförlust när det gäller att komma åt data
  och utföra operationer. Detta beror på att åtkomst till privata
  medlemmar kräver ytterligare overhead i form av metodanrop eller
  reflektion. Det är dock viktigt att notera att prestandaförlusten
  oftast är försumbar och att inkapslingens fördelar oftast överväger
  den.
\end{enumerate}

\hypertarget{anvuxe4ndningsomruxe5den-fuxf6r-inkapsling-i-java}{%
\subsection{Användningsområden för inkapsling i
Java}\label{anvuxe4ndningsomruxe5den-fuxf6r-inkapsling-i-java}}

Inkapsling kan tillämpas på olika sätt inom Java-programmering. Här är
några vanliga användningsområden:

\begin{enumerate}
\def\labelenumi{\arabic{enumi}.}
\item
  \textbf{Klasser och objekt}: Inkapsling används i stor utsträckning
  för att definiera klasser och objekt i Java. Genom att använda
  åtkomstmodifierare, som \texttt{public}, \texttt{private} och
  \texttt{protected}, kan vi definiera vilka medlemmar som är
  tillgängliga utanför klassen och hur de kan manipuleras.
\item
  \textbf{Egenskaper (Getters och Setters)}: Genom att använda getters
  och setters kan vi kontrollera åtkomsten till objektens data och
  möjliggöra läsning och skrivning av privata medlemmar på ett
  kontrollerat sätt. Genom att använda getters och setters kan vi
  implementera validering och logik för att säkerställa korrekt
  användning av data.
\item
  \textbf{Gränssnitt (Interfaces)}: Inkapsling används för att definiera
  gränssnitt i Java. Ett gränssnitt definierar en uppsättning metoder
  som en klass kan implementera. Genom att använda gränssnitt kan vi
  separera definitionen av en klass från dess implementation och
  möjliggöra en mer flexibel och modulär kodstruktur.
\item
  \textbf{Nedarvning (Inheritance)}: Inkapsling används också i samband
  med nedarvning för att definiera och kontrollera åtkomsten till
  medlemmar i en basklass. Genom att använda åtkomstmodifierare kan vi
  definiera vilka medlemmar som ärver och vilka som inte ärver till en
  underklass.
\end{enumerate}

\hypertarget{exempel-puxe5-inkapsling-i-java}{%
\subsection{Exempel på inkapsling i
Java}\label{exempel-puxe5-inkapsling-i-java}}

För att ge en bättre förståelse för inkapsling i Java kan vi titta på
ett exempel på en klass som använder inkapsling:

\hypertarget{cb1}{}
\begin{Shaded}
\begin{Highlighting}[]
\KeywordTok{public} \KeywordTok{class}\NormalTok{ BankAccount }\OperatorTok{\{}
\KeywordTok{private} \DataTypeTok{double}\NormalTok{ balance}\OperatorTok{;}

\KeywordTok{public} \FunctionTok{BankAccount}\OperatorTok{(}\DataTypeTok{double}\NormalTok{ initialBalance}\OperatorTok{)} \OperatorTok{\{}
\NormalTok{balance }\OperatorTok{=}\NormalTok{ initialBalance}\OperatorTok{;}
\OperatorTok{\}}

\KeywordTok{public} \DataTypeTok{void} \FunctionTok{deposit}\OperatorTok{(}\DataTypeTok{double}\NormalTok{ amount}\OperatorTok{)} \OperatorTok{\{}
\NormalTok{balance }\OperatorTok{+=}\NormalTok{ amount}\OperatorTok{;}
\OperatorTok{\}}

\KeywordTok{public} \DataTypeTok{void} \FunctionTok{withdraw}\OperatorTok{(}\DataTypeTok{double}\NormalTok{ amount}\OperatorTok{)} \OperatorTok{\{}
\ControlFlowTok{if} \OperatorTok{(}\NormalTok{amount }\OperatorTok{\textless{}=}\NormalTok{ balance}\OperatorTok{)} \OperatorTok{\{}
\NormalTok{balance }\OperatorTok{{-}=}\NormalTok{ amount}\OperatorTok{;}
\OperatorTok{\}} \ControlFlowTok{else} \OperatorTok{\{}
\BuiltInTok{System}\OperatorTok{.}\FunctionTok{out}\OperatorTok{.}\FunctionTok{println}\OperatorTok{(}\StringTok{"Otillräckliga medel."}\OperatorTok{);}
\OperatorTok{\}}
\OperatorTok{\}}

\KeywordTok{public} \DataTypeTok{double} \FunctionTok{getBalance}\OperatorTok{()} \OperatorTok{\{}
\ControlFlowTok{return}\NormalTok{ balance}\OperatorTok{;}
\OperatorTok{\}}
\OperatorTok{\}}

\KeywordTok{public} \KeywordTok{class}\NormalTok{ Main }\OperatorTok{\{}
\KeywordTok{public} \DataTypeTok{static} \DataTypeTok{void} \FunctionTok{main}\OperatorTok{(}\BuiltInTok{String}\OperatorTok{[]}\NormalTok{ args}\OperatorTok{)} \OperatorTok{\{}

\NormalTok{BankAccount account }\OperatorTok{=} \KeywordTok{new} \FunctionTok{BankAccount}\OperatorTok{(}\DecValTok{1000}\OperatorTok{);}
\BuiltInTok{System}\OperatorTok{.}\FunctionTok{out}\OperatorTok{.}\FunctionTok{println}\OperatorTok{(}\StringTok{"Nuvarande saldo: "} \OperatorTok{+}\NormalTok{ account}\OperatorTok{.}\FunctionTok{getBalance}\OperatorTok{());}
\NormalTok{account}\OperatorTok{.}\FunctionTok{deposit}\OperatorTok{(}\DecValTok{500}\OperatorTok{);}
\BuiltInTok{System}\OperatorTok{.}\FunctionTok{out}\OperatorTok{.}\FunctionTok{println}\OperatorTok{(}\StringTok{"Efter insättning: "} \OperatorTok{+}\NormalTok{ account}\OperatorTok{.}\FunctionTok{getBalance}\OperatorTok{());}
\NormalTok{account}\OperatorTok{.}\FunctionTok{withdraw}\OperatorTok{(}\DecValTok{200}\OperatorTok{);}
\BuiltInTok{System}\OperatorTok{.}\FunctionTok{out}\OperatorTok{.}\FunctionTok{println}\OperatorTok{(}\StringTok{"Efter uttag: "} \OperatorTok{+}\NormalTok{ account}\OperatorTok{.}\FunctionTok{getBalance}\OperatorTok{());}
\OperatorTok{\}}
\OperatorTok{\}}
\end{Highlighting}
\end{Shaded}

I detta exempel har vi en klass \texttt{BankAccount} som representerar
en bankräkning. Klassen har en privat medlemsvariabel \texttt{balance}
som håller reda på kontots saldo. Klassen har också metoder för att
sätta in pengar (\texttt{deposit}), ta ut pengar (\texttt{withdraw}) och
hämta saldot (\texttt{getBalance}).

I \texttt{Main}-klassen skapar vi en instans av \texttt{BankAccount} och
använder dess metoder för att utföra insättningar och uttag på
bankkontot.

\hypertarget{ett-annat-exempel}{%
\subsection{Ett annat exempel}\label{ett-annat-exempel}}

Vi kan också titta på ett annat exempel på inkapsling i Java:

\hypertarget{cb2}{}
\begin{Shaded}
\begin{Highlighting}[]
\KeywordTok{public} \KeywordTok{class}\NormalTok{ Superhero }\OperatorTok{\{}
\KeywordTok{private} \BuiltInTok{String}\NormalTok{ name}\OperatorTok{;}    \CommentTok{// Privat medlemsvariabel för superhjältens namn}
\KeywordTok{private} \BuiltInTok{String}\NormalTok{ powers}\OperatorTok{;}  \CommentTok{// Privat medlemsvariabel för superhjältens krafter}

\CommentTok{// Konstruktor för att skapa en instans av Superhero med namn och krafter}
\KeywordTok{public} \FunctionTok{Superhero}\OperatorTok{(}\BuiltInTok{String}\NormalTok{ name}\OperatorTok{,} \BuiltInTok{String}\NormalTok{ powers}\OperatorTok{)} \OperatorTok{\{}
\KeywordTok{this}\OperatorTok{.}\FunctionTok{name} \OperatorTok{=}\NormalTok{ name}\OperatorTok{;}
\KeywordTok{this}\OperatorTok{.}\FunctionTok{powers} \OperatorTok{=}\NormalTok{ powers}\OperatorTok{;}
\OperatorTok{\}}

\CommentTok{// Getter{-}metod för att hämta superhjältens namn}
\KeywordTok{public} \BuiltInTok{String} \FunctionTok{getName}\OperatorTok{()} \OperatorTok{\{}
\ControlFlowTok{return}\NormalTok{ name}\OperatorTok{;}
\OperatorTok{\}}

\CommentTok{// Getter{-}metod för att hämta superhjältens krafter}
\KeywordTok{public} \BuiltInTok{String} \FunctionTok{getPowers}\OperatorTok{()} \OperatorTok{\{}
\ControlFlowTok{return}\NormalTok{ powers}\OperatorTok{;}
\OperatorTok{\}}

\CommentTok{// Setter{-}metod för att sätta superhjältens krafter}
\KeywordTok{public} \DataTypeTok{void} \FunctionTok{setPowers}\OperatorTok{(}\BuiltInTok{String}\NormalTok{ powers}\OperatorTok{)} \OperatorTok{\{}
\KeywordTok{this}\OperatorTok{.}\FunctionTok{powers} \OperatorTok{=}\NormalTok{ powers}\OperatorTok{;}
\OperatorTok{\}}

\CommentTok{// Metod för att visa superhjältens namn och krafter}
\KeywordTok{public} \DataTypeTok{void} \FunctionTok{displaySuperhero}\OperatorTok{()} \OperatorTok{\{}
\BuiltInTok{System}\OperatorTok{.}\FunctionTok{out}\OperatorTok{.}\FunctionTok{println}\OperatorTok{(}\StringTok{"Namn: "} \OperatorTok{+}\NormalTok{ name}\OperatorTok{);}
\BuiltInTok{System}\OperatorTok{.}\FunctionTok{out}\OperatorTok{.}\FunctionTok{println}\OperatorTok{(}\StringTok{"Krafter: "} \OperatorTok{+}\NormalTok{ powers}\OperatorTok{);}
\OperatorTok{\}}
\OperatorTok{\}}

\KeywordTok{public} \KeywordTok{class}\NormalTok{ Main }\OperatorTok{\{}
\KeywordTok{public} \DataTypeTok{static} \DataTypeTok{void} \FunctionTok{main}\OperatorTok{(}\BuiltInTok{String}\OperatorTok{[]}\NormalTok{ args}\OperatorTok{)} \OperatorTok{\{}
\CommentTok{// Skapa en instans av Superhero med namnet "Spider{-}Man" och krafterna "Wall{-}crawling, superhuman strength"}
\NormalTok{Superhero superhero }\OperatorTok{=} \KeywordTok{new} \FunctionTok{Superhero}\OperatorTok{(}\StringTok{"Spider{-}Man"}\OperatorTok{,} \StringTok{"Wall{-}crawling, superhuman strength"}\OperatorTok{);}

\CommentTok{// Visa superhjältens namn och krafter med hjälp av displaySuperhero{-}metoden}
\NormalTok{superhero}\OperatorTok{.}\FunctionTok{displaySuperhero}\OperatorTok{();}

\CommentTok{// Uppdatera superhjältens krafter med hjälp av setPowers{-}metoden}
\NormalTok{superhero}\OperatorTok{.}\FunctionTok{setPowers}\OperatorTok{(}\StringTok{"Web{-}slinging, spider{-}sense"}\OperatorTok{);}

\CommentTok{// Visa superhjältens uppdaterade krafter med hjälp av getPowers{-}metoden}
\BuiltInTok{System}\OperatorTok{.}\FunctionTok{out}\OperatorTok{.}\FunctionTok{println}\OperatorTok{(}\StringTok{"Uppdaterade krafter: "} \OperatorTok{+}\NormalTok{ superhero}\OperatorTok{.}\FunctionTok{getPowers}\OperatorTok{());}
\OperatorTok{\}}
\OperatorTok{\}}
\end{Highlighting}
\end{Shaded}

I detta exempel har vi en klass \texttt{Superhero} som representerar en
superhjälte. Klassen har två privata medlemsvariabler, \texttt{name} och
\texttt{powers}, för att hålla superhjältens namn och krafter. Här är
förklaringar för varje del av koden:

\begin{itemize}
\tightlist
\item
  I konstruktorn \texttt{Superhero} tar vi emot namn och krafter som
  parametrar och tilldelar dem till de privata medlemsvariablerna med
  hjälp av \texttt{this}-referensen.
\item
  \texttt{getName} är en getter-metod som returnerar superhjältens namn.
\item
  \texttt{getPowers} är en getter-metod som returnerar superhjältens
  krafter.
\item
  \texttt{setPowers} är en setter-metod som tar emot en ny kraft som
  parameter och uppdaterar den privata medlemsvariabeln \texttt{powers}
  med den nya kraften.
\item
  \texttt{displaySuperhero} är en metod som visar superhjältens namn och
  krafter genom att skriva ut dem till konsolen.
\item
  I \texttt{Main}-klassens \texttt{main}-metod skapar vi en instans av
  \texttt{Superhero} med namnet ``Spider-Man'' och krafterna
  ``Wall-crawling, superhuman strength''.
\item
  Vi anropar \texttt{displaySuperhero}-metoden på den skapade instansen
  för att visa superhjältens namn och krafter.
\item
  Sedan använder vi \texttt{setPowers}-metoden för att uppdatera
  superhjältens krafter till ``Web-slinging, spider-sense''.
\item
  Till sist använder vi \texttt{getPowers}-metoden för att hämta och
  skriva ut de uppdaterade krafterna till konsolen.
\end{itemize}

På så sätt kan du skapa och manipulera en instans av
\texttt{Superhero}-klassen med hjälp av inkapsling i Java.

Jag hoppas att detta förtydligar exemplet för dig! Låt mig veta om det
finns något mer jag kan hjälpa dig med.

\hypertarget{sammanfattning}{%
\subsection{Sammanfattning}\label{sammanfattning}}

Inkapsling är en viktig princip inom objektorienterad programmering som
möjliggör att data och funktioner som hör samman hålls tillsammans inom
en enhet, kallad en klass. Genom att använda inkapsling kan vi definiera
vilka medlemmar som är tillgängliga utanför klassen och hur de kan
manipuleras. Detta främjar moduläritet, återanvändbarhet och säkerhet i
programkoden.

Inkapsling har flera fördelar, såsom moduläritet, säkerhet och
kodunderhåll. Det finns dock också vissa begränsningar, såsom ökad
komplexitet och eventuell prestandaförlust. Inkapsling kan tillämpas på
olika sätt inom Java-programmering, inklusive klasser, getters och
setters, gränssnitt och nedarvning.

Genom att förstå inkapslingens koncept och användningsområden kan du
skapa välstrukturerad och lättunderhållen kod i Java. Lycka till med
dina programmeringsprojekt!

\hypertarget{obligatorisk-dad-joke}{%
\subsection{Obligatorisk dad-joke}\label{obligatorisk-dad-joke}}

Varför älskar programmerare att använda inkapsling?

För att de inte vill läcka sina privata medlemmar!

// Arv är en princip inom OOP där en klass kan ärva egenskaper och
beteenden från en annan överordnad klass, vilket gör det möjligt att
återanvända kod och skapa hierarkier av klasser.

// Superklassen är den överordnade klassen som fungerar som en mall
eller ritning för de nedärvande klasserna. Subklasserna ärver och kan
utöka funktionaliteten från superklassen. // Exempel: // Superklassen
``Fordon'' kan ha egenskaper och beteenden som är gemensamma för alla
fordon, t.ex. hastighet och färg. // Subklasserna ``Bil'' och
``Motorcykel'' kan ärva dessa egenskaper och beteenden från ``Fordon''
och lägga till sina egna unika egenskaper och beteenden, t.ex. antal
dörrar och typ av motor. // På så sätt kan vi återanvända kod och
undvika att duplicera samma kod i flera klasser. Vi kan också
strukturera och organisera vår kod på ett modulärt sätt, vilket
förbättrar underhållbarheten och läsbarheten. // Arv möjliggör också
polymorfism, vilket innebär att en instans av en subklass kan användas
där en instans av superklassen förväntas. Detta gör det möjligt att
skriva generell kod som kan hantera olika typer av objekt. // En metod
som tar emot en parameter av typen ``Fordon'' kan användas för att
hantera både ``Bil'' och ``Motorcykel'' objekt, eftersom de båda är
subklasser till ``Fordon''. // Arv kan vara en kraftfull teknik för att
organisera och strukturera kod, men det kan också leda till problem om
det inte används på rätt sätt. Det är viktigt att tänka på följande
begränsningar och rekommendationer när du använder arv: // - Arv bör
endast användas när det finns en verklig ``är ett'' relation mellan
superklassen och subklassen. Om relationen bara är en ``har ett''
relation, kan det vara mer lämpligt att använda komposition istället för
arv. // - Arv kan leda till en hierarki av klasser som kan bli svår att
underhålla och förstå. Det är viktigt att noggrant planera och designa
hierarkin för att undvika överdriven komplexitet. // - Om en subklass
behöver ändra eller utöka funktionaliteten från superklassen, kan det
vara mer lämpligt att använda interfaces eller abstrakta klasser
istället för arv. Detta möjliggör flexibilitet och undviker problem med
klasshierarkin. // - Det är viktigt att använda rätt
namngivningskonventioner och vara konsekvent i hela klasshierarkin för
att underlätta förståelsen och läsbarheten. // - Det är också viktigt
att vara medveten om eventuell kodupprepning i klasshierarkin och
undvika det genom att extrahera gemensam funktionalitet till separata
metoder eller klasser. // - Slutligen är det alltid bra att dokumentera
klasshierarkin och dess relationer för att underlätta förståelsen och
underhållningen av koden. // Med dessa rekommendationer kan arv vara en
kraftfull teknik för att organisera och strukturera kod inom OOP. \# Arv

Arv kan tillämpas i olika situationer och användningsområden inom
programmering. Här är några exempel: 1. Skapa en hierarki av klasser där
varje klass representerar en mer specialiserad version av den
överordnade klassen. Till exempel kan vi ha en superklass ``Djur'' och
subklasser som ``Hund'', ``Katt'' och ``Fågel'' som ärver och utökar
funktionaliteten från ``Djur''. 2. Återanvända kod genom att extrahera
gemensam funktionalitet till en överordnad klass och låta subklasserna
ärva den. På så sätt undviker vi duplicering av kod och håller vår kod
baserad på principen ``Don't Repeat Yourself'' (DRY). 3. Implementera
polymorfism genom att använda superklasser och subklasser. Polymorfism
möjliggör att en instans av en subklass kan behandlas som en instans av
superklassen, vilket gör att vi kan skriva generell kod som fungerar med
olika typer av objekt. 4. Skapa en modellering av verkliga scenarier där
hierarkier och relationer mellan olika objekt är viktiga. Till exempel
kan vi modellera en hierarki av fordon där vi har en superklass
``Fordon'' och subklasser som ``Bil'', ``Motorcykel'' och ``Buss''. Det
finns många andra användningsområden för arv beroende på de specifika
kraven och designen av programmet. Det är viktigt att noggrant planera
och designa hierarkin för att få ut det bästa av arv och undvika problem
eller överdriven komplexitet. \#\# Kodexempel: Arv i Java

Låt oss nu titta på ett kodexempel som visar hur arv kan användas i
Java. Vi kommer att använda exemplet med hierarkin av fordon som nämndes
tidigare.

\begin{verbatim}
// Superklassen "Fordon"
class Fordon {
private int hastighet;
private String färg;
public Fordon(int hastighet, String färg) {
this.hastighet = hastighet;
this.färg = färg;
}
public int getHastighet() {
return hastighet;
public String getFärg() {
return färg;
}
// Subklassen "Bil" ärver från "Fordon"
class Bil extends Fordon {
private int antalDörrar;
public Bil(int hastighet, String färg, int antalDörrar) {
super(hastighet, färg);
this.antalDörrar = antalDörrar;
public int getAntalDörrar() {
return antalDörrar;
// Subklassen "Motorcykel" ärver från "Fordon"
class Motorcykel extends Fordon {
private String typAvMotor;
public Motorcykel(int hastighet, String färg, String typAvMotor) {
this.typAvMotor = typAvMotor;
public String getTypAvMotor() {
return typAvMotor;
// Genom arv kan en hierarki av klasser skapas där gemensamma egenskaper och beteenden
// placeras i överordnade klasser och specialiserade egenskaper och beteenden läggs till
// i de nedärvande klasserna. Detta möjliggör återanvändning av kod och effektivisering av
// utvecklingsprocessen.
// Fördelar
// Arv erbjuder flera fördelar inom programmering:
// 1. Kodåteranvändning: Genom att använda arv kan du återanvända kod från överordnade
//    klasser i de nedärvande klasserna. Detta minskar behovet av att skriva samma kod
//    flera gånger och förbättrar därmed kodens underhållbarhet och läsbarhet.
// 2. Modulär design: Arv bidrar till en modulär design av program. Genom att dela upp
//    funktionaliteten i olika klasser och använda arv kan du separera olika ansvarsområden
//    och skapa en hierarki av klasser som är enklare att förstå och hantera.
// 3. Kodens struktur: Arv kan förbättra kodens struktur och organisering. Genom att placera
//    gemensam funktionalitet i överordnade klasser och specialisera den i nedärvande klasser
//    blir koden mer lättläslig och intuitiv.
// 4. Utbytbarhet: Arv möjliggör att objekt av en nedärvande klass kan användas där objekt
//    av en överordnad klass förväntas. Detta skapar möjligheten att behandla objekt på ett
//    enhetligt sätt och gör koden mer flexibel och skalbar.
// Begränsningar
// Trots sina fördelar har arv vissa begränsningar och kompromisser:
// 1. Tätt kopplade klasser: Genom att använda arv skapas en tät koppling mellan överordnade
//    och nedärvande klasser. Om du gör ändringar i överordnade klasser kan det påverka alla
//    nedärvande klasser, vilket kan vara komplicerat att hantera och underhålla.
// 2. Brist på flexibilitet: Arv kan begränsa flexibiliteten i en kodbas. Om hierarkin av
//    klasser inte är korrekt utformad kan det bli svårt att lägga till eller ändra
//    funktionalitet på ett smidigt sätt.
// 3. Ökad komplexitet: När hierarkin av klasser blir djup och komplex kan det bli svårt att
//    förstå och hantera koden. Det är viktigt att noggrant planera och organisera
//    klasshierarkin för att undvika överflödig komplexitet.
// Användningsområden
// Arv kan tillämpas i olika scenarier och användningsområden inom programmering. Här är
// några exempel:
// 1. GUI-ramverk: I grafiska användargränssnittsramverk används arv för att skapa hierarkier
//    av användargränssnittskomponenter. Till exempel kan en överordnad klass "Komponent"
//    innehålla grundläggande egenskaper och beteenden, medan nedärvande klasser som "Knapp"
//    och "Textfält" specialiserar funktionaliteten.
// 2. Spelprogrammering: I spelutveckling kan arv användas för att skapa en hierarki av
//    spelobjekt. Till exempel kan en överordnad klass "Spelobjekt" innehålla gemensamma
//    egenskaper och beteenden, medan nedärvande klasser som "Fiende" och "Spelare"
//    specialiserar funktionaliteten för specifika spelkaraktärer.
\end{verbatim}

För att visa hur arv fungerar i Java kan vi använda följande kodexempel:
// En överordnad klass ``Karaktär'' som representerar en generell
spelkaraktär. class Karaktär \{ String namn; int hälsa; public
Karaktär(String namn, int hälsa) \{ this.namn = namn; this.hälsa =
hälsa; public void attack() \{ System.out.println(namn + ''
attackerar!{}``); // En nedärvande klass''Fiende'' som specialiserar
funktionaliteten för en fiende i spelet. class Fiende extends Karaktär
\{ int styrka; public Fiende(String namn, int hälsa, int styrka) \{
super(namn, hälsa); this.styrka = styrka; System.out.println(namn + ''
attackerar med styrka '' + styrka + ``!''); // En nedärvande klass
``Spelare'' som specialiserar funktionaliteten för spelaren i spelet.
class Spelare extends Karaktär \{ int nivå; public Spelare(String namn,
int hälsa, int nivå) \{ this.nivå = nivå; System.out.println(namn + ''
attackerar med nivå '' + nivå + ``!''); public class ArvExempel \{
public static void main(String{[}{]} args) \{ // Skapar en fiende och en
spelare. Fiende fiende = new Fiende(``Ondskans mästare'', 100, 10);
Spelare spelare = new Spelare(``Hjälten'', 100, 5); // Anropar
attack-metoden för fienden och spelaren. fiende.attack();
spelare.attack(); I detta kodexempel används arv för att skapa en
hierarki av spelkaraktärer. Den överordnade klassen ``Karaktär''
innehåller gemensamma egenskaper och beteenden för alla spelkaraktärer,
medan nedärvande klasser som ``Fiende'' och ``Spelare'' specialiserar
funktionaliteten för specifika spelkaraktärer. Båda nedärvande klasserna
har en egen implementation av attack-metoden som anropas när de
attackerar. När vi kör programmet skapar vi en instans av Fiende och en
instans av Spelare och anropar--- 3. \textbf{Databashanterare}: I en
databashanterare kan arv användas för att skapa en hierarki av
databasobjekt. Till exempel kan en överordnad klass ``DatabaseObject''
innehålla generella funktioner för att hantera databasoperationer, medan
nedärvande klasser som ``Table'' (Tabell) och ``Query'' (Fråga)
specialiserar funktionaliteten för specifika databasentiteter. Detta är
bara några exempel på användningsområden där arv kan tillämpas inom
programmering. Principen om arv kan vara användbar i olika typer av
program och system, oavsett om det är grafiska användargränssnitt, spel
eller databashantering. \#\# Exempelkod - Arv i en berättelse

För att bättre förstå arv kan vi titta på ett kodexempel som illustrerar
användningen av arv genom en berättelse. Anta att vi bygger ett spel där
vi har olika typer av karaktärer, inklusive fiender och hjältar. Vi kan
använda arv för att skapa en hierarki av karaktärsklasser. // Definiera
överordnad klass Karaktär public class Karaktär \{ public String Namn;
public int Hälsa; // Definiera nedärvande klass Fiende public class
Fiende extends Karaktär \{ public void Attackera() \{ // Implementera
attacklogik för fiender // Definiera nedärvande klass Hjälte public
class Hjälte extends Karaktär \{ public void Försvara() \{ //
Implementera försvarlogik för hjältar // Användning av arv i spellogik
Fiende fiende = new Fiende(); fiende.Namn = ``Ond skurk''; fiende.Hälsa
= 100; fiende.Attackera(); Hjälte hjälte = new Hjälte(); hjälte.Namn =
``Modig hjälte''; hjälte.Hälsa = 100; hjälte.Försvara(); I det här
exemplet skapas en överordnad klass \texttt{Karaktär} som har egenskaper
för \texttt{Namn} (Name) och \texttt{Hälsa} (Health). Därefter skapas
två nedärvande klasser, \texttt{Fiende} (Enemy) och \texttt{Hjälte}
(Hero), som båda ärver från \texttt{Karaktär} och lägger till sina egna
specifika metoder (\texttt{Attackera} och \texttt{Försvara}). I
spellogiken skapar vi en instans av \texttt{Fiende} och sätter dess
egenskaper (\texttt{Namn} och \texttt{Hälsa}) och använder sedan dess
specifika metod \texttt{Attackera}. På samma sätt skapar vi en instans
av \texttt{Hjälte}, sätter dess egenskaper och använder dess specifika
metod \texttt{Försvara}. Genom att använda arv kan vi enkelt skapa en
hierarki av karaktärsklasser och använda deras specifika metoder och
egenskaper baserat på deras roll i spelet.I detta kodexempel har vi en
överordnad klass \texttt{Karaktär} som innehåller gemensamma egenskaper
för både fiender och hjältar. Genom att ärva från \texttt{Karaktär} kan
vi definiera specialiserad funktionalitet för fiender och hjältar i
deras respektive nedärvande klasser \texttt{Fiende} och \texttt{Hjälte}.
Vi kan sedan skapa instanser av dessa klasser och använda deras unika
funktioner, som \texttt{Attackera()} för fiender och \texttt{Försvara()}
för hjältar. \#\#\# Utdata

(fiende.Attackera() skriver ut något här) (hjälte.Försvara() skriver ut
något här) Detta är bara ett enkelt exempel som visar hur arv kan
användas för att skapa hierarkier av klasser och dela funktionalitet
mellan dem. I praktiken kan arv vara mycket mer kraftfullt och komplext
i sina tillämpningar. \#\# Slutsats

Arv är en viktig princip inom objektorienterad programmering som
möjliggör återanvändning av kod och skapar hierarkier av klasser. Genom
att använda arv kan vi strukturera och organisera vår kod på ett
modulärt sätt, vilket förbättrar underhållbarheten och läsbarheten. I
denna artikel har vi utforskat arvets fördelar, inklusive
kodåteranvändning, modulär design och strukturerad kod. Vi har också
diskuterat dess begränsningar och utmaningar, såsom tät koppling och
ökad komplexitet. Arv kan tillämpas inom olika områden inom
programmering, från grafiska användargränssnittsramverk till
spelprogrammering och databashanterare. Genom att förstå och behärska
arv kan du utveckla effektivare och mer flexibla program. För att
fördjupa dina kunskaper rekommenderar vi att du fortsätter läsa om arv,
utforskar mer avancerade koncept som abstrakt arv och gränssnitt och
experimenterar med att använda arv i dina egna programmeringsprojekt.
\#\# TL;DRArv i Java:

Arv i Java är en princip som möjliggör att en klass kan ärva egenskaper
och beteenden från en annan klass. Detta möjliggör kodåteranvändning och
skapar en hierarki av klasser. För att implementera arv i Java används
nyckelordet ``extends''. En klass kan ärva från en annan klass genom att
använda extends-nyckelordet och ange namnet på den överordnade klassen.
Klassen som ärver kallas subklassen och den överordnade klassen kallas
superklassen. Här är ett exempel som visar hur arv fungerar i Java: //
Superklass class Djur \{ public void sägLjud() \{
System.out.println(``Djuret säger ljudet''); // Subklass som ärver från
superklassen Djur class Hund extends Djur \{ System.out.println(``Hunden
säger voff''); public class Main \{ Hund hund = new Hund();
hund.sägLjud(); // Output: Hunden säger voff I exemplet ovan har vi en
superklass ``Djur'' och en subklass ``Hund''. Subklassen ``Hund'' ärver
egenskaper och beteenden från superklassen ``Djur''. Subklassen kan
också överlagra metoder från superklassen för att ändra deras beteende.
I main-metoden skapar vi en instans av subklassen ``Hund'' och anropar
metoden ``sägLjud()''. Eftersom subklassen har överlagrat metoden,
kommer den att visa det specifika ljudet för hundar. Arv möjliggör också
flera nivåer av hierarkier. En subklass kan i sin tur fungera som
superklass för en annan subklass. På så sätt kan vi bygga mer komplexa
strukturer av klasser och utnyttja arvets fördelar. Det är viktigt att
notera att Java inte stöder flerfaldig arv, vilket innebär att en klass
inte kan ärva från flera superklasser samtidigt. Java använder istället
gränssnitt (interface) för att uppnå liknande funktionalitet. Gränssnitt
låter en klass implementera flera gränssnitt samtidigt, vilket ger
flexibilitet och kodåteranvändning.

\hypertarget{komposition-uxf6ver-arv}{%
\section{Komposition över arv}\label{komposition-uxf6ver-arv}}

\hypertarget{vad-uxe4r-det-fuxf6r-nuxe5got}{%
\subsection{Vad är det för något?}\label{vad-uxe4r-det-fuxf6r-nuxe5got}}

Komposition över arv är ett designmönster inom objektorienterad
programmering som förespråkar att ``sammansättning'' (komposition) bör
användas istället för ``arv'' (inheritance) för att uppnå
återanvändbarhet och flexibilitet i kod. Detta mönster är baserat på
principen om att klasser bör sammansättas av andra klasser snarare än
att ärva från dem.

\hypertarget{varfuxf6r-uxe4r-det-viktigt}{%
\subsection{Varför är det viktigt?}\label{varfuxf6r-uxe4r-det-viktigt}}

Komposition över arv är ett viktigt koncept inom objektorienterad
programmering eftersom det hjälper oss att skapa flexibla och
återanvändbara klasser. Genom att använda komposition kan vi enkelt
kombinera olika klasser för att skapa nya klasser med olika beteenden
och egenskaper. Detta gör det möjligt för oss att skapa modulära och
flexibla program som är lätta att underhålla och utöka.

\hypertarget{fuxf6rdelar}{%
\subsection{Fördelar}\label{fuxf6rdelar}}

Komposition över arv erbjuder flera fördelar och används i olika
situationer:

\begin{enumerate}
\def\labelenumi{\arabic{enumi}.}
\item
  \textbf{Flexibilitet}: Genom att använda komposition kan vi enkelt
  kombinera olika klasser för att skapa nya klasser med olika beteenden
  och egenskaper. Detta gör det möjligt för oss att skapa modulära och
  flexibla program som är lätta att underhålla och utöka.
\item
  \textbf{Återanvändbarhet}: Genom att använda komposition kan vi
  återanvända befintlig kod och undvika att duplicera kod. Detta gör det
  möjligt för oss att skapa modulära och flexibla program som är lätta
  att underhålla och utöka.
\item
  \textbf{Moduläritet}: Genom att använda komposition kan vi skapa
  modulära program som är lätta att underhålla och utöka. Detta gör det
  möjligt för oss att skapa modulära och flexibla program som är lätta
  att underhålla och utöka.
\item
  \textbf{Minskad beroende}: Genom att använda komposition kan vi minska
  beroendet mellan klasser och undvika att duplicera kod. Detta gör det
  möjligt för oss att skapa modulära och flexibla program som är lätta
  att underhålla och utöka.
\end{enumerate}

\hypertarget{begruxe4nsningar}{%
\subsection{Begränsningar}\label{begruxe4nsningar}}

Komposition över arv har vissa begränsningar att vara medveten om:

\begin{enumerate}
\def\labelenumi{\arabic{enumi}.}
\item
  \textbf{Ökad komplexitet}: Genom att använda komposition kan vi öka
  komplexiteten i vårt program. Detta kan göra det svårare att förstå
  och underhålla koden.
\item
  \textbf{Mer kod att skriva}: Genom att använda komposition kan vi
  behöva skriva mer kod för att uppnå samma resultat som med arv. Detta
  kan göra det svårare att förstå och underhålla koden.
\end{enumerate}

\hypertarget{enkel-fuxf6rklaring}{%
\subsection{Enkel förklaring}\label{enkel-fuxf6rklaring}}

Man skulle förenkla det hela med att säga att vi simulerar ``arv av
flera klasser'' och undvika att duplicera kod. Detta gör det möjligt för
oss att skapa modulära och flexibla program som är lätta att underhålla
och utöka.

\hypertarget{exempel-1}{%
\subsection{Exempel 1}\label{exempel-1}}

Känns det rörigt? Lugnt, vi kollar på ett exempel\ldots{} Vi har
faktiskt använt detta innan\ldots{}

kolla koden nedan

\hypertarget{cb1}{}
\begin{Shaded}
\begin{Highlighting}[]
\KeywordTok{class}\NormalTok{ Pet}
\OperatorTok{\{}
    \KeywordTok{private} \BuiltInTok{String}\NormalTok{ name}\OperatorTok{;}
    \KeywordTok{private} \DataTypeTok{int}\NormalTok{ age}\OperatorTok{;}
    \KeywordTok{private} \BuiltInTok{String}\NormalTok{ breed}\OperatorTok{;}

    \KeywordTok{public} \FunctionTok{Pet}\OperatorTok{(}\BuiltInTok{String}\NormalTok{ name}\OperatorTok{,} \DataTypeTok{int}\NormalTok{ age}\OperatorTok{,} \BuiltInTok{String}\NormalTok{ breed}\OperatorTok{)}
    \OperatorTok{\{}
        \KeywordTok{this}\OperatorTok{.}\FunctionTok{name} \OperatorTok{=}\NormalTok{ name}\OperatorTok{;}
        \KeywordTok{this}\OperatorTok{.}\FunctionTok{age} \OperatorTok{=}\NormalTok{ age}\OperatorTok{;}
        \KeywordTok{this}\OperatorTok{.}\FunctionTok{breed} \OperatorTok{=}\NormalTok{ breed}\OperatorTok{;}
    \OperatorTok{\}}

    \KeywordTok{public} \BuiltInTok{String} \FunctionTok{getName}\OperatorTok{()\{} \ControlFlowTok{return}\NormalTok{ name}\OperatorTok{;} \OperatorTok{\}}
    \KeywordTok{public} \DataTypeTok{int} \FunctionTok{getAge}\OperatorTok{()\{} \ControlFlowTok{return}\NormalTok{ age}\OperatorTok{;} \OperatorTok{\}}
    \KeywordTok{public} \BuiltInTok{String} \FunctionTok{getBreed}\OperatorTok{()\{} \ControlFlowTok{return}\NormalTok{ breed}\OperatorTok{;} \OperatorTok{\}}
    \KeywordTok{public} \DataTypeTok{void} \FunctionTok{setName}\OperatorTok{(}\BuiltInTok{String}\NormalTok{ name}\OperatorTok{)\{} \KeywordTok{this}\OperatorTok{.}\FunctionTok{name} \OperatorTok{=}\NormalTok{ name}\OperatorTok{;} \OperatorTok{\}}
    \KeywordTok{public} \DataTypeTok{void} \FunctionTok{setAge}\OperatorTok{(}\DataTypeTok{int}\NormalTok{ age}\OperatorTok{)\{} \KeywordTok{this}\OperatorTok{.}\FunctionTok{age} \OperatorTok{=}\NormalTok{ age}\OperatorTok{;} \OperatorTok{\}}
    \KeywordTok{public} \DataTypeTok{void} \FunctionTok{setBreed}\OperatorTok{(}\BuiltInTok{String}\NormalTok{ breed}\OperatorTok{)\{} \KeywordTok{this}\OperatorTok{.}\FunctionTok{breed} \OperatorTok{=}\NormalTok{ breed}\OperatorTok{;} \OperatorTok{\}}
\OperatorTok{\}}
\KeywordTok{class}\NormalTok{ Person}
\OperatorTok{\{}
    \KeywordTok{private} \BuiltInTok{String}\NormalTok{ name}\OperatorTok{;}
    \KeywordTok{private} \DataTypeTok{int}\NormalTok{ age}\OperatorTok{;}

    \KeywordTok{public} \FunctionTok{Person}\OperatorTok{(}\BuiltInTok{String}\NormalTok{ name}\OperatorTok{,} \DataTypeTok{int}\NormalTok{ age}\OperatorTok{)}
    \OperatorTok{\{}
        \KeywordTok{this}\OperatorTok{.}\FunctionTok{name} \OperatorTok{=}\NormalTok{ name}\OperatorTok{;}
        \KeywordTok{this}\OperatorTok{.}\FunctionTok{age} \OperatorTok{=}\NormalTok{ age}\OperatorTok{;}
    \OperatorTok{\}}

    \KeywordTok{public} \BuiltInTok{String} \FunctionTok{getName}\OperatorTok{()\{} \ControlFlowTok{return}\NormalTok{ name}\OperatorTok{;} \OperatorTok{\}}
    \KeywordTok{public} \DataTypeTok{int} \FunctionTok{getAge}\OperatorTok{()\{} \ControlFlowTok{return}\NormalTok{ age}\OperatorTok{;} \OperatorTok{\}}
    \KeywordTok{public} \DataTypeTok{void} \FunctionTok{setName}\OperatorTok{(}\BuiltInTok{String}\NormalTok{ name}\OperatorTok{)\{} \KeywordTok{this}\OperatorTok{.}\FunctionTok{name} \OperatorTok{=}\NormalTok{ name}\OperatorTok{;} \OperatorTok{\}}
    \KeywordTok{public} \DataTypeTok{void} \FunctionTok{setAge}\OperatorTok{(}\DataTypeTok{int}\NormalTok{ age}\OperatorTok{)\{} \KeywordTok{this}\OperatorTok{.}\FunctionTok{age} \OperatorTok{=}\NormalTok{ age}\OperatorTok{;} \OperatorTok{\}}
\OperatorTok{\}}

\KeywordTok{class}\NormalTok{ HumanPet}
\OperatorTok{\{}
    \KeywordTok{private}\NormalTok{ Person person}\OperatorTok{;}
    \KeywordTok{private}\NormalTok{ Pet pet}\OperatorTok{;}

    \KeywordTok{public} \FunctionTok{HumanPet}\OperatorTok{(}\NormalTok{Person person}\OperatorTok{,}\NormalTok{ Pet pet}\OperatorTok{)}
    \OperatorTok{\{}
        \KeywordTok{this}\OperatorTok{.}\FunctionTok{person} \OperatorTok{=}\NormalTok{ person}\OperatorTok{;}
        \KeywordTok{this}\OperatorTok{.}\FunctionTok{pet} \OperatorTok{=}\NormalTok{ pet}\OperatorTok{;}
    \OperatorTok{\}}

    \KeywordTok{public}\NormalTok{ Person }\FunctionTok{getPerson}\OperatorTok{()\{} \ControlFlowTok{return}\NormalTok{ person}\OperatorTok{;} \OperatorTok{\}}
    \KeywordTok{public}\NormalTok{ Pet }\FunctionTok{getPet}\OperatorTok{()\{} \ControlFlowTok{return}\NormalTok{ pet}\OperatorTok{;} \OperatorTok{\}}
    \KeywordTok{public} \DataTypeTok{void} \FunctionTok{setPerson}\OperatorTok{(}\NormalTok{Person person}\OperatorTok{)\{} \KeywordTok{this}\OperatorTok{.}\FunctionTok{person} \OperatorTok{=}\NormalTok{ person}\OperatorTok{;} \OperatorTok{\}}
    \KeywordTok{public} \DataTypeTok{void} \FunctionTok{setPet}\OperatorTok{(}\NormalTok{Pet pet}\OperatorTok{)\{} \KeywordTok{this}\OperatorTok{.}\FunctionTok{pet} \OperatorTok{=}\NormalTok{ pet}\OperatorTok{;} \OperatorTok{\}}
\OperatorTok{\}}
\end{Highlighting}
\end{Shaded}

Här hade vi kunnat ärva \texttt{Person} eller \texttt{Pet} men vi väljer
självklart kompostion över arv. Oavsett variant skulle det orsakat
problem. Så nu kan vi ärva \texttt{HumanPet} och få med oss både
\texttt{Person} och \texttt{Pet}\ldots{} om vi nu vill det.

\hypertarget{cb2}{}
\begin{Shaded}
\begin{Highlighting}[]
\KeywordTok{class}\NormalTok{ Family}
\OperatorTok{\{}
    \KeywordTok{private}\NormalTok{ HumanPet father}\OperatorTok{;}
    \KeywordTok{private}\NormalTok{ HumanPet mother}\OperatorTok{;}
    \KeywordTok{private}\NormalTok{ HumanPet}\OperatorTok{[]}\NormalTok{ child}\OperatorTok{;}
    \KeywordTok{private}\NormalTok{ Pet}\OperatorTok{[]}\NormalTok{ familyPets}\OperatorTok{;}

    \KeywordTok{public} \FunctionTok{Family}\OperatorTok{(}\NormalTok{HumanPet father}\OperatorTok{,}\NormalTok{ HumanPet mother}\OperatorTok{,}\NormalTok{ HumanPet}\OperatorTok{[]}\NormalTok{ child}\OperatorTok{,}\NormalTok{ Pet}\OperatorTok{[]}\NormalTok{ familyPets}\OperatorTok{)}
    \OperatorTok{\{}
        \KeywordTok{this}\OperatorTok{.}\FunctionTok{father} \OperatorTok{=}\NormalTok{ father}\OperatorTok{;}
        \KeywordTok{this}\OperatorTok{.}\FunctionTok{mother} \OperatorTok{=}\NormalTok{ mother}\OperatorTok{;}
        \KeywordTok{this}\OperatorTok{.}\FunctionTok{child} \OperatorTok{=}\NormalTok{ child}\OperatorTok{;}
        \KeywordTok{this}\OperatorTok{.}\FunctionTok{familyPets} \OperatorTok{=}\NormalTok{ familyPets}\OperatorTok{;}
    \OperatorTok{\}}

    \KeywordTok{public}\NormalTok{ HumanPet }\FunctionTok{getFather}\OperatorTok{()\{} \ControlFlowTok{return}\NormalTok{ father}\OperatorTok{;} \OperatorTok{\}}
    \KeywordTok{public}\NormalTok{ HumanPet }\FunctionTok{getMother}\OperatorTok{()\{} \ControlFlowTok{return}\NormalTok{ mother}\OperatorTok{;} \OperatorTok{\}}
    \KeywordTok{public}\NormalTok{ HumanPet }\FunctionTok{getChild}\OperatorTok{()\{} \ControlFlowTok{return}\NormalTok{ child}\OperatorTok{;} \OperatorTok{\}}
    \KeywordTok{public}\NormalTok{ Pet}\OperatorTok{[]} \FunctionTok{getFamilyPets}\OperatorTok{()\{} \ControlFlowTok{return}\NormalTok{ familyPets}\OperatorTok{;} \OperatorTok{\}}
    \KeywordTok{public} \DataTypeTok{void} \FunctionTok{setFather}\OperatorTok{(}\NormalTok{HumanPet father}\OperatorTok{)\{} \KeywordTok{this}\OperatorTok{.}\FunctionTok{father} \OperatorTok{=}\NormalTok{ father}\OperatorTok{;} \OperatorTok{\}}
    \KeywordTok{public} \DataTypeTok{void} \FunctionTok{setMother}\OperatorTok{(}\NormalTok{HumanPet mother}\OperatorTok{)\{} \KeywordTok{this}\OperatorTok{.}\FunctionTok{mother} \OperatorTok{=}\NormalTok{ mother}\OperatorTok{;} \OperatorTok{\}}
    \KeywordTok{public} \DataTypeTok{void} \FunctionTok{setChild}\OperatorTok{(}\NormalTok{HumanPet child}\OperatorTok{)\{} \KeywordTok{this}\OperatorTok{.}\FunctionTok{child} \OperatorTok{=}\NormalTok{ child}\OperatorTok{;} \OperatorTok{\}}
    \KeywordTok{public} \DataTypeTok{void} \FunctionTok{setFamilyPets}\OperatorTok{(}\NormalTok{Pet}\OperatorTok{[]}\NormalTok{ familyPets}\OperatorTok{)\{} \KeywordTok{this}\OperatorTok{.}\FunctionTok{familyPets} \OperatorTok{=}\NormalTok{ familyPets}\OperatorTok{;} \OperatorTok{\}}
\OperatorTok{\}}
\end{Highlighting}
\end{Shaded}

Nu kan vi skapa en familj där alla har var sitt husdjur.

\hypertarget{exempel-2}{%
\subsection{Exempel 2}\label{exempel-2}}

För att förstå hur man implementerar ``Komposition över Arv'' i
praktiken, låt oss titta på ett exempel med Java-kod. Anta att vi har
ett program som hanterar karaktärer i en Star Wars-värld. Vi har två
klasser: \texttt{Character} och \texttt{Weapon}.

\hypertarget{cb3}{}
\begin{Shaded}
\begin{Highlighting}[]
\KeywordTok{public} \KeywordTok{class} \BuiltInTok{Character} \OperatorTok{\{}
\KeywordTok{private} \BuiltInTok{String}\NormalTok{ name}\OperatorTok{;}
\KeywordTok{private}\NormalTok{ Weapon weapon}\OperatorTok{;}

\KeywordTok{public} \BuiltInTok{Character}\OperatorTok{(}\BuiltInTok{String}\NormalTok{ name}\OperatorTok{,}\NormalTok{ Weapon weapon}\OperatorTok{)} \OperatorTok{\{}
\KeywordTok{this}\OperatorTok{.}\FunctionTok{name} \OperatorTok{=}\NormalTok{ name}\OperatorTok{;}
\KeywordTok{this}\OperatorTok{.}\FunctionTok{weapon} \OperatorTok{=}\NormalTok{ weapon}\OperatorTok{;}
\OperatorTok{\}}

\KeywordTok{public} \DataTypeTok{void} \FunctionTok{attack}\OperatorTok{()} \OperatorTok{\{}
\ControlFlowTok{if} \OperatorTok{(}\NormalTok{weapon }\OperatorTok{!=} \KeywordTok{null}\OperatorTok{)} \OperatorTok{\{}
\NormalTok{weapon}\OperatorTok{.}\FunctionTok{use}\OperatorTok{();}
\OperatorTok{\}} \ControlFlowTok{else} \OperatorTok{\{}
\BuiltInTok{System}\OperatorTok{.}\FunctionTok{out}\OperatorTok{.}\FunctionTok{println}\OperatorTok{(}\NormalTok{name }\OperatorTok{+} \StringTok{" has no weapon."}\OperatorTok{);}
\OperatorTok{\}}
\OperatorTok{\}}
\OperatorTok{\}}
\end{Highlighting}
\end{Shaded}

Ovan skapade vi karaktärer, vapen och använde dem i vår applikation. Vi
skapade en instans av \texttt{Character} och tilldelade den ett vapen.
Sedan anropade vi \texttt{attack()}-metoden för att se hur karaktären
använder sitt tilldelade vapen.

Nu ska vi skapa vapen.

\hypertarget{cb4}{}
\begin{Shaded}
\begin{Highlighting}[]
\KeywordTok{public} \KeywordTok{interface}\NormalTok{ Weapon }\OperatorTok{\{}
\DataTypeTok{void} \FunctionTok{use}\OperatorTok{();}
\OperatorTok{\}}

\KeywordTok{public} \KeywordTok{class}\NormalTok{ Lightsaber }\KeywordTok{implements}\NormalTok{ Weapon }\OperatorTok{\{}
\AttributeTok{@Override}
\KeywordTok{public} \DataTypeTok{void} \FunctionTok{use}\OperatorTok{()} \OperatorTok{\{}
\BuiltInTok{System}\OperatorTok{.}\FunctionTok{out}\OperatorTok{.}\FunctionTok{println}\OperatorTok{(}\StringTok{"Swinging lightsaber!"}\OperatorTok{);}
\OperatorTok{\}}
\OperatorTok{\}}

\KeywordTok{public} \KeywordTok{class}\NormalTok{ Blaster }\KeywordTok{implements}\NormalTok{ Weapon }\OperatorTok{\{}
\AttributeTok{@Override}
\KeywordTok{public} \DataTypeTok{void} \FunctionTok{use}\OperatorTok{()} \OperatorTok{\{}
\BuiltInTok{System}\OperatorTok{.}\FunctionTok{out}\OperatorTok{.}\FunctionTok{println}\OperatorTok{(}\StringTok{"Shooting blaster!"}\OperatorTok{);}
\OperatorTok{\}}
\OperatorTok{\}}
\end{Highlighting}
\end{Shaded}

I det här exemplet representerar \texttt{Character} en karaktär i Star
Wars och \texttt{Weapon} representerar ett vapen som kan användas av
karaktären. Istället för att använda arv för att definiera olika typer
av karaktärer med specifika vapen, använder vi komposition för att
sammansätta karaktärer med vapen.

En karaktär har ett namn och en referens till ett vapen. I
\texttt{Character}-klassen finns en metod \texttt{attack()} som anropar
\texttt{use()}-metoden för det tilldelade vapnet. Om karaktären inte har
något vapen skrivs ett meddelande ut.

Vi har också en \texttt{Weapon}-gränssnitt som definierar en enda metod
\texttt{use()}. Både \texttt{Lightsaber}- och \texttt{Blaster}-klasserna
implementerar detta gränssnitt och tillhandahåller implementationer för
\texttt{use()}-metoden.

För att använda detta i vår applikation kan vi skapa instanser av
\texttt{Character} och \texttt{Weapon} och sätta ihop dem genom att
tilldela ett vapen till en karaktär.

\hypertarget{cb5}{}
\begin{Shaded}
\begin{Highlighting}[]
\KeywordTok{public} \KeywordTok{class}\NormalTok{ Main }\OperatorTok{\{}
\KeywordTok{public} \DataTypeTok{static} \DataTypeTok{void} \FunctionTok{main}\OperatorTok{(}\BuiltInTok{String}\OperatorTok{[]}\NormalTok{ args}\OperatorTok{)} \OperatorTok{\{}
\NormalTok{Weapon lightsaber }\OperatorTok{=} \KeywordTok{new} \FunctionTok{Lightsaber}\OperatorTok{();}
\BuiltInTok{Character}\NormalTok{ jedi }\OperatorTok{=} \KeywordTok{new} \BuiltInTok{Character}\OperatorTok{(}\StringTok{"Luke Skywalker"}\OperatorTok{,}\NormalTok{ lightsaber}\OperatorTok{);}
\NormalTok{jedi}\OperatorTok{.}\FunctionTok{attack}\OperatorTok{();} \CommentTok{// Output: Swinging lightsaber!}

\NormalTok{Weapon blaster }\OperatorTok{=} \KeywordTok{new} \FunctionTok{Blaster}\OperatorTok{();}
\BuiltInTok{Character}\NormalTok{ bountyHunter }\OperatorTok{=} \KeywordTok{new} \BuiltInTok{Character}\OperatorTok{(}\StringTok{"Boba Fett"}\OperatorTok{,}\NormalTok{ blaster}\OperatorTok{);}
\NormalTok{bountyHunter}\OperatorTok{.}\FunctionTok{attack}\OperatorTok{();} \CommentTok{// Output: Shooting blaster!}

\BuiltInTok{Character}\NormalTok{ padawan }\OperatorTok{=} \KeywordTok{new} \BuiltInTok{Character}\OperatorTok{(}\StringTok{"Ahsoka Tano"}\OperatorTok{,} \KeywordTok{null}\OperatorTok{);}
\NormalTok{padawan}\OperatorTok{.}\FunctionTok{attack}\OperatorTok{();} \CommentTok{// Output: Ahsoka Tano has no weapon.}
\OperatorTok{\}}
\OperatorTok{\}}
\end{Highlighting}
\end{Shaded}

I \texttt{Main}-klassen skapar vi instanser av \texttt{Lightsaber} och
\texttt{Blaster} som representerar vapen. Sedan skapar vi instanser av
\texttt{Character} och tilldelar respektive vapen till dem. Genom att
anropa \texttt{attack()}-metoden för varje karaktär kan vi se hur de
använder sina tilldelade vapen.

Detta exempel visar hur vi kan använda komposition för att sammansätta
klasser och uppnå flexibilitet och återanvändbarhet. Istället för att ha
en hierarki av klasser med olika kombinationer av vapen, kan vi enkelt
kombinera olika karaktärer och vapen genom att ändra vilka
objektinstanser som tilldelas till varje karaktär.

\hypertarget{sammanfattning}{%
\subsection{Sammanfattning}\label{sammanfattning}}

I denna artikel har vi utforskat konceptet ``Komposition över Arv'' och
dess betydelse inom programmering. Vi har sett fördelarna med att
använda komposition, såsom flexibilitet, moduläritet, återanvändbarhet
och minskat beroende. Vi har också identifierat begränsningar som ökad
komplexitet och mer kod att skriva. Genom ett kodexempel har vi sett hur
man kan implementera komposition i praktiken med hjälp av Java.

För vidare läsning rekommenderas att utforska olika designmönster och
principer inom objektorienterad programmering, såsom SOLID-principerna
och Design Patterns-boken av Erich Gamma m.fl.

\hypertarget{uxe5tkomstmoderator}{%
\section{Åtkomstmoderator}\label{uxe5tkomstmoderator}}

En åtkomstmodifikator är en modifierare som används för att definiera
åtkomsten till en klassmedlem. Detta innebär att du kan bestämma vilka
klassmedlemmar som ska vara tillgängliga för andra klasser och vilka som
ska vara privata för klassen.

Det finns fyra huvudsakliga åtkomstmodifikatorer i Java:

\begin{enumerate}
\def\labelenumi{\arabic{enumi}.}
\tightlist
\item
  \textbf{public}: Tillåter åtkomst till en klassmedlem från alla andra
  klasser.
\item
  \textbf{private}: Tillåter åtkomst till en klassmedlem endast från den
  klass där den är deklarerad.
\item
  \textbf{protected}: Tillåter åtkomst till en klassmedlem endast från
  den klass där den är deklarerad och alla dess underklasser.
\item
  \textbf{ingen åtkomstmodifikator}: Om ingen åtkomstmodifikator anges,
  är standardåtkomsten ``package-private''. Det innebär att
  klassmedlemmen är tillgänglig endast inom samma paket.
\item
  \textbf{final}: Tillåter inte att en klassmedlem ändras efter att den
  har deklarerats.
\item
  \textbf{package-private} (default): Tillåter åtkomst till en
  klassmedlem endast inom samma paket.
\item
  \textbf{static}: Tillåter åtkomst till en klassmedlem utan att skapa
  en instans av klassen.
\item
  \textbf{abstract}: Tillåter inte att en klassmedlem har en
  implementation.
\item
  \textbf{transient}: Tillåter inte att en klassmedlem sparas när ett
  objekt serialiseras.
\item
  \textbf{volatile}: Tillåter inte att en klassmedlem cachas av JVM.
\item
  \textbf{synchronized}: Tillåter inte att flera trådar samtidigt ändrar
  en klassmedlem.
\item
  \textbf{native}: Tillåter inte att en klassmedlem implementeras i
  Java.
\item
  \textbf{sealed}: Tillåter inte att en klassmedlem ärver från en annan
  klass.
\end{enumerate}

Exempel:

\hypertarget{cb1}{}
\begin{Shaded}
\begin{Highlighting}[]
\KeywordTok{public} \KeywordTok{class}\NormalTok{ Person }\OperatorTok{\{}
    \KeywordTok{public} \BuiltInTok{String}\NormalTok{ name}\OperatorTok{;} \CommentTok{// Tillgänglig från alla klasser}
    \KeywordTok{private} \DataTypeTok{int}\NormalTok{ age}\OperatorTok{;} \CommentTok{// Endast tillgänglig från Person{-}klassen}
    \KeywordTok{protected} \BuiltInTok{String}\NormalTok{ address}\OperatorTok{;} \CommentTok{// Tillgänglig från Person{-}klassen och dess underklasser}
    \BuiltInTok{String}\NormalTok{ phoneNumber}\OperatorTok{;} \CommentTok{// Endast tillgänglig inom samma paket som Person{-}klassen}
\OperatorTok{\}}
\end{Highlighting}
\end{Shaded}

I exemplet ovan är \texttt{name} tillgänglig från alla klasser,
\texttt{age} är endast tillgänglig från \texttt{Person}-klassen,
\texttt{address} är tillgänglig från \texttt{Person}-klassen och dess
underklasser, och \texttt{phoneNumber} är endast tillgänglig inom samma
paket som \texttt{Person}-klassen.

\hypertarget{private}{%
\section{Private}\label{private}}

Privat är en åtkomstmodifikator som gör att en klassmedlem endast är
tillgänglig för den klass där den är deklarerad.

Innehållsförteckning

\{: .text-delta \} 1. TOC \{:toc\}

\hypertarget{beskrivning}{%
\subsection{Beskrivning}\label{beskrivning}}

Vi kan använda privat för att göra en klassmedlem endast tillgänglig för
den klass där den är deklarerad. Detta är användbart om vi vill att en
klassmedlem endast ska vara tillgänglig för den klass där den är
deklarerad och inte för någon annan klass.

\hypertarget{exempel}{%
\subsection{Exempel}\label{exempel}}

Låt oss titta på ett exempel där vi använder privat för att göra en
klassmedlem endast tillgänglig för den klass där den är deklarerad:

\hypertarget{cb1}{}
\begin{Shaded}
\begin{Highlighting}[]
\KeywordTok{public} \KeywordTok{class}\NormalTok{ Person }\OperatorTok{\{}
    \KeywordTok{private} \BuiltInTok{String}\NormalTok{ name}\OperatorTok{;}
    \KeywordTok{private} \DataTypeTok{int}\NormalTok{ age}\OperatorTok{;}
\OperatorTok{\}}
\end{Highlighting}
\end{Shaded}

I detta exempel har vi en klass som heter Person. Vi har också två
egenskaper, name och age. Båda är privata, vilket innebär att de endast
är tillgängliga för klassen Person. Även om name och age borde vara
tillgängliga för alla klasser som länkas till denna, kommer de inte att
vara tillgängliga för någon annan klass på grund av att de är privata.

Detta innebär att vi inte kan använda dem i klassen Person själv heller.
Vi kan inte ens använda dem i en konstruktor i Person-klassen. Detta
beror på att privata medlemmar endast är tillgängliga för den klass där
de är deklarerade.

Jaja det är ju bara ett exempel. Det får räcka som Dad-Joke för den här
artikeln.

\begin{verbatim}
 ___       _______   ___
 |   |     |   _   | |   |
 |.  |     |.  |   | |.  |
 |.  |___  |.  |   | |.  |___
 |:  1   | |:  1   | |:  1   |
 |::.. . | |::.. . | |::.. . |
 `-------' `-------' `-------'
\end{verbatim}

\hypertarget{protected}{%
\section{Protected}\label{protected}}

Protected är en åtkomstmodifierare som gör att en klass, metod eller
egenskap är tillgänglig för klassen den är deklarerad i och alla klasser
som ärver från den.

Innehållsförteckning

\{: .text-delta \} 1. TOC \{:toc\}

\hypertarget{beskrivning}{%
\subsection{Beskrivning}\label{beskrivning}}

När vi arbetar med polymorfism kan det vara bra att göra en metod eller
egenskap tillgänglig för alla klasser som ärver från en klass. Detta gör
vi genom att använda protected. Protected är som private för alla
klasser, utom den som ärver. Klasser kan inte ärva privata medlemmar, så
detta är det bästa alternativet.

\hypertarget{exempel}{%
\subsection{Exempel}\label{exempel}}

\hypertarget{cb1}{}
\begin{Shaded}
\begin{Highlighting}[]
\CommentTok{// Klassen Person ärver inte från någon annan klass}
\CommentTok{// Eller ja... den ärver från Objekt, det gör alla klasser :)}
\KeywordTok{public} \KeywordTok{class}\NormalTok{ Person }\OperatorTok{\{}
    \KeywordTok{protected} \BuiltInTok{String}\NormalTok{ name}\OperatorTok{;}
\OperatorTok{\}}

\CommentTok{// Klassen Student ärver från Person och kan därför använda "name"}
\KeywordTok{public} \KeywordTok{class}\NormalTok{ Student }\KeywordTok{extends}\NormalTok{ Person }\OperatorTok{\{}
    \KeywordTok{public} \FunctionTok{Student}\OperatorTok{(}\BuiltInTok{String}\NormalTok{ name}\OperatorTok{)} \OperatorTok{\{}
        \KeywordTok{this}\OperatorTok{.}\FunctionTok{name} \OperatorTok{=}\NormalTok{ name}\OperatorTok{;}
    \OperatorTok{\}}

\CommentTok{// Klassen Student kan nu använda "name" och skriva ut det}
\KeywordTok{public} \DataTypeTok{void} \FunctionTok{printInfo}\OperatorTok{()} \OperatorTok{\{}
    \BuiltInTok{System}\OperatorTok{.}\FunctionTok{out}\OperatorTok{.}\FunctionTok{println}\OperatorTok{(}\StringTok{"Student Info: "} \OperatorTok{+}\NormalTok{ name}\OperatorTok{);}
\OperatorTok{\}}
\end{Highlighting}
\end{Shaded}

\hypertarget{fuxf6rklaring}{%
\subsection{Förklaring}\label{fuxf6rklaring}}

I exemplet ovan är ``name'' propertyn i klassen Person skyddad för alla
klasser förutom de som ärver från den, vilket gör att klassen Student
kan lägga till ett namn och skriva ut det, medan inga andra klasser kan
göra det.

\hypertarget{obligatorisk-dad-joke}{%
\subsection{Obligatorisk Dad-Joke}\label{obligatorisk-dad-joke}}

Varför var det nyckelordet ``protected'' så nervöst hela tiden?

För att det alltid ville hålla sitt privata liv\ldots{} väl,
``protected''! 😄

\hypertarget{public}{%
\section{Public}\label{public}}

Public is an access modifier that allows a class, method, or property to
be accessible to all classes.

Table of Contents

\{: .text-delta \} 1. TOC \{:toc\}

\hypertarget{beskrivning}{%
\subsection{Beskrivning}\label{beskrivning}}

Med public kan vi göra en klass, metod eller egenskap tillgänglig för
alla klasser. Detta är bra om vi vill dela kod mellan flera klasser i
olika projekt.

\hypertarget{exempel}{%
\subsection{Exempel}\label{exempel}}

\hypertarget{cb1}{}
\begin{Shaded}
\begin{Highlighting}[]
\KeywordTok{public} \KeywordTok{class}\NormalTok{ Person }\OperatorTok{\{}
    \KeywordTok{public} \BuiltInTok{String}\NormalTok{ name}\OperatorTok{;}
    \KeywordTok{public} \DataTypeTok{int}\NormalTok{ age}\OperatorTok{;}
\OperatorTok{\}}
\end{Highlighting}
\end{Shaded}

\hypertarget{fuxf6rklaring}{%
\subsection{Förklaring}\label{fuxf6rklaring}}

I exemplet ovan är klassen Person tillgänglig för alla klasser i alla
projekt. Detta gör att vi kan skapa en ny instans av Person i en annan
klass och lägga till ett namn och ålder.

\hypertarget{getters-och-setters}{%
\subsection{Getters och setters?}\label{getters-och-setters}}

Varför har man Getters och Setters istället för publika variabler?

\begin{longtable}[]{@{}
  >{\raggedright\arraybackslash}p{(\columnwidth - 4\tabcolsep) * \real{0.3300}}
  >{\raggedright\arraybackslash}p{(\columnwidth - 4\tabcolsep) * \real{0.3300}}
  >{\raggedright\arraybackslash}p{(\columnwidth - 4\tabcolsep) * \real{0.3300}}@{}}
\toprule\noalign{}
\begin{minipage}[b]{\linewidth}\raggedright
Namn
\end{minipage} & \begin{minipage}[b]{\linewidth}\raggedright
Fördel
\end{minipage} & \begin{minipage}[b]{\linewidth}\raggedright
Nackdel
\end{minipage} \\
\midrule\noalign{}
\endhead
\bottomrule\noalign{}
\endlastfoot
Getter & Du kan göra mer än bara returnera en variabel. Du kan till
exempel validera värdet innan du returnerar det. Du kan returnera
resultat av data i din klass istället för en specifik variabel. & Du
måste skriva mer kod. \\
Setter & Du kan göra mer än bara tilldela en variabel. Du kan till
exempel validera värdet innan du tilldelar det. Om värdet inte är bra
för din klass kan kasta ett felmeddelande eller bara vägra att spara
värdet & Du måste skriva mer kod. \\
Public & Du behöver inte skriva mer kod. & Du kan inte göra mer än att
tilldela eller returnera en variabel. \\
\end{longtable}

\hypertarget{exempel-1}{%
\subsection{Exempel}\label{exempel-1}}

\hypertarget{cb2}{}
\begin{Shaded}
\begin{Highlighting}[]
\CommentTok{// Setter som inte accepterar negativa tal}
\KeywordTok{public} \DataTypeTok{void} \FunctionTok{setAge}\OperatorTok{(}\DataTypeTok{int}\NormalTok{ age}\OperatorTok{)} \OperatorTok{\{}
    \ControlFlowTok{if} \OperatorTok{(}\NormalTok{age }\OperatorTok{\textless{}} \DecValTok{0}\OperatorTok{)} \OperatorTok{\{}
        \ControlFlowTok{throw} \KeywordTok{new} \BuiltInTok{IllegalArgumentException}\OperatorTok{(}\StringTok{"Age cannot be negative"}\OperatorTok{);}
    \OperatorTok{\}}
    \KeywordTok{this}\OperatorTok{.}\FunctionTok{age} \OperatorTok{=}\NormalTok{ age}\OperatorTok{;}
\OperatorTok{\}}

\CommentTok{// Getter som returnerar åldern på en person baserat på en annan variabel (Date)}
\KeywordTok{public} \DataTypeTok{int} \FunctionTok{getAge}\OperatorTok{(}\BuiltInTok{Date}\NormalTok{ date}\OperatorTok{)} \OperatorTok{\{}
    \ControlFlowTok{return}\NormalTok{ date}\OperatorTok{.}\FunctionTok{getYear}\OperatorTok{()} \OperatorTok{{-}} \KeywordTok{this}\OperatorTok{.}\FunctionTok{birthDate}\OperatorTok{.}\FunctionTok{getYear}\OperatorTok{();}
\OperatorTok{\}}
\end{Highlighting}
\end{Shaded}

\hypertarget{sammanfattning}{%
\subsection{Sammanfattning}\label{sammanfattning}}

Public är en åtkomstmodifikator som gör att en klass, metod eller
egenskap är tillgänglig för alla klasser. Detta är bra om vi vill dela
kod mellan flera klasser i olika projekt. Att dela metoder till externa
klasser är ett bra sätt för att få klasser att samarbeta.

\hypertarget{static}{%
\section{Static}\label{static}}

En static är en åtkomstmodifikator som gör att en klassmedlem tillhör
klassen och inte objektet.

Innehållsförteckning

\{: .text-delta \}

\begin{enumerate}
\def\labelenumi{\arabic{enumi}.}
\tightlist
\item
  TOC \{:toc\}
\end{enumerate}

\hypertarget{beskrivning}{%
\subsection{Beskrivning}\label{beskrivning}}

Vi kan använda static för att göra en klassmedlem tillhör klassen och
inte objektet. Detta är användbart om vi vill dela data mellan olika
klasser utan att behöva skapa flera instanser av en klass.

När koden kompileras kommer alla statiska medlemmar att tilldelas ett
gemensamt minnesutrymme som delas av alla instanser av klassen. Detta
gör att vi kan dela data mellan olika klasser utan att behöva skapa
flera instanser av en klass. Detta är användbart om vi vill dela data
mellan olika klasser utan att behöva skapa flera instanser av en klass.

mindmap root((Minnet)) Instanser av klasser Instansvariabler
Instansmetoder Statiskt Statiska variabler Statiska metoder

I detta diagram representerar ``Minne'' den allmänna minnesplatsen där
programmet körs.

\begin{itemize}
\tightlist
\item
  \texttt{Instanser\ av\ klasser} representerar alla instanser av
  klasser som skapas under körningen.
\item
  \texttt{Statiskt} representerar alla statiska variabler och metoder
  som skapas under körningen.
\end{itemize}

Observera att diagrammet bara visar en förenklad representation av hur
en statisk metod och instansvariabler lagras i minnet vid kompileringen.
Det kan finnas fler faktorer och detaljer som påverkar
minnesorganisationen beroende på den specifika implementationen och
miljön där programmet körs.

Hoppas att detta diagram hjälper till att illustrera hur en statisk
metod läggs i minnet! Låt mig veta om du har fler frågor.

\hypertarget{exempel}{%
\subsection{Exempel}\label{exempel}}

\hypertarget{cb1}{}
\begin{Shaded}
\begin{Highlighting}[]
\KeywordTok{public} \KeywordTok{class}\NormalTok{ UserSettings }\OperatorTok{\{}
    \KeywordTok{public} \DataTypeTok{static} \BuiltInTok{String}\NormalTok{ UserName}\OperatorTok{;}
    \KeywordTok{public} \DataTypeTok{static} \DataTypeTok{int}\NormalTok{ Password}\OperatorTok{;}
    \KeywordTok{public} \DataTypeTok{static} \DataTypeTok{boolean}\NormalTok{ DarkMode }\OperatorTok{=} \KeywordTok{true}\OperatorTok{;}
\OperatorTok{\}}
\end{Highlighting}
\end{Shaded}

I detta exempel har vi en klass som heter \texttt{UserSettings}. Vi har
också tre statiska egenskaper: \texttt{UserName}, \texttt{Password} och
\texttt{DarkMode}. Dessa egenskaper är tillgängliga för alla klasser i
samma projekt. Men om en annan klass i ett annat projekt länkar till
detta projekt kommer de inte att kunna använda
\texttt{UserSettings}-klassen.

Det är viktigt att notera att i Java kan vi inte ha statiska egenskaper
med automatisk implementering som i C\#. Istället måste vi använda
publika statiska variabler och få tillgång till dem direkt genom
klassnamnet.

För att använda \texttt{UserSettings}-klassen kan vi tilldela värden
till de statiska egenskaperna innan vi använder dem:

\hypertarget{cb2}{}
\begin{Shaded}
\begin{Highlighting}[]
\NormalTok{UserSettings}\OperatorTok{.}\FunctionTok{UserName} \OperatorTok{=} \StringTok{"JohnDoe"}\OperatorTok{;}
\NormalTok{UserSettings}\OperatorTok{.}\FunctionTok{Password} \OperatorTok{=} \DecValTok{123456}\OperatorTok{;}
\NormalTok{UserSettings}\OperatorTok{.}\FunctionTok{DarkMode} \OperatorTok{=} \KeywordTok{true}\OperatorTok{;}
\end{Highlighting}
\end{Shaded}

Och sedan kan vi använda dessa värden i andra delar av koden:

\hypertarget{cb3}{}
\begin{Shaded}
\begin{Highlighting}[]
\BuiltInTok{System}\OperatorTok{.}\FunctionTok{out}\OperatorTok{.}\FunctionTok{println}\OperatorTok{(}\StringTok{"Användarnamn: "} \OperatorTok{+}\NormalTok{ UserSettings}\OperatorTok{.}\FunctionTok{UserName}\OperatorTok{);}
\BuiltInTok{System}\OperatorTok{.}\FunctionTok{out}\OperatorTok{.}\FunctionTok{println}\OperatorTok{(}\StringTok{"Lösenord: "} \OperatorTok{+}\NormalTok{ UserSettings}\OperatorTok{.}\FunctionTok{Password}\OperatorTok{);}
\BuiltInTok{System}\OperatorTok{.}\FunctionTok{out}\OperatorTok{.}\FunctionTok{println}\OperatorTok{(}\StringTok{"Mörkt läge: "} \OperatorTok{+}\NormalTok{ UserSettings}\OperatorTok{.}\FunctionTok{DarkMode}\OperatorTok{);}
\end{Highlighting}
\end{Shaded}

Detta ger oss möjligheten att dela data mellan olika klasser utan att
behöva skapa flera instanser av \texttt{UserSettings}-klassen. Vi kan
också ändra värdena i \texttt{UserSettings}-klassen och ha tillgång till
de uppdaterade värdena överallt där vi använder klassen.

Detta är bara en av många användningar av \texttt{static} i Java. Det
kan också användas för att skapa statiska metoder och block, vilket vi
kommer att utforska i andra artiklar.

\hypertarget{obligatorisk-dad-joke}{%
\subsection{Obligatorisk Dad-joke}\label{obligatorisk-dad-joke}}

Varför var den statiska metoden så dålig på att socialisera?

För att den aldrig kunde få någon att ``dynamiskt'' intressera sig för
den! 😔

\hypertarget{polymorfism}{%
\section{Polymorfism}\label{polymorfism}}

Polymorfism kommer från grekiskans poly = många och morphe = form. Det
är en av de viktigaste egenskaperna i objektorienterad programmering.

Innehållsförteckning

\{: .text-delta \} 1. TOC \{:toc\}

\hypertarget{beskrivning}{%
\subsection{Beskrivning}\label{beskrivning}}

Polymorfism är en egenskap som gör att objekt som ärver från samma
familj av en klass eller ett interface kan hanteras som om de vore samma
typ. Detta ger oss möjligheten att skapa en lista av olika objekt från
samma familj och ändå kunna hantera dem på samma sätt. Polymorfism ger
oss många fördelar:

\begin{itemize}
\tightlist
\item
  Möjlighet att byta ut delar av programmet utan att behöva ändra resten
  av programmet.
\item
  Möjlighet att skapa en metod som kan hantera olika typer av objekt
  utan att behöva skapa en metod för varje typ av objekt.
\item
  Möjlighet att skapa plugins till programmet utan att behöva ändra i
  programmet, med mera.
\end{itemize}

Här är den korrigerade och förbättrade versionen av artikeln med
kommentarer:

\hypertarget{cb1}{}
\begin{Shaded}
\begin{Highlighting}[]
\FunctionTok{\# Interfaces}

\NormalTok{Gränssnitt är ett kraftfullt verktyg i Java som ger oss möjlighet att skapa flexibla och återanvändbara komponenter i våra program. Genom att använda gränssnitt kan vi implementera polymorfism och separera implementation och användning av komponenter.}

\FunctionTok{\#\# Beskrivning}

\NormalTok{Ett gränssnitt är en typ som definierar en uppsättning metoder, egenskaper och händelser som en klass kan implementera. En klass som implementerar ett gränssnitt måste implementera alla dess medlemmar. Därför säger man att gränssnitt är ett kontrakt.}

\FunctionTok{\#\# Exempel}

\NormalTok{Låt oss titta på ett exempel där vi skapar ett gränssnitt som heter }\InformationTok{\textasciigrave{}Animal\textasciigrave{}}\NormalTok{:}

\InformationTok{\textasciigrave{}\textasciigrave{}\textasciigrave{}java}
\KeywordTok{interface}\NormalTok{ Animal }\OperatorTok{\{}
\BuiltInTok{String} \FunctionTok{getName}\OperatorTok{();}
\DataTypeTok{void} \FunctionTok{eat}\OperatorTok{();}
\DataTypeTok{void} \FunctionTok{sleep}\OperatorTok{();}
\DataTypeTok{void} \FunctionTok{shit}\OperatorTok{();}
\OperatorTok{\}}
\end{Highlighting}
\end{Shaded}

I detta exempel har vi ett gränssnitt som heter \texttt{Animal}. Vi har
också tre metoder och en egenskap. Observera att gränssnitt inte kan ha
fält, men de kan ha egenskaper!

\hypertarget{cb2}{}
\begin{Shaded}
\begin{Highlighting}[]
\KeywordTok{class}\NormalTok{ Cat }\KeywordTok{implements}\NormalTok{ Animal }\OperatorTok{\{}
\KeywordTok{private} \BuiltInTok{String}\NormalTok{ name}\OperatorTok{;}

\KeywordTok{public} \FunctionTok{Cat}\OperatorTok{(}\BuiltInTok{String}\NormalTok{ name}\OperatorTok{)} \OperatorTok{\{}
\KeywordTok{this}\OperatorTok{.}\FunctionTok{name} \OperatorTok{=}\NormalTok{ name}\OperatorTok{;}
\OperatorTok{\}}

\KeywordTok{public} \BuiltInTok{String} \FunctionTok{getName}\OperatorTok{()} \OperatorTok{\{}
\ControlFlowTok{return}\NormalTok{ name}\OperatorTok{;}
\OperatorTok{\}}

\KeywordTok{public} \DataTypeTok{void} \FunctionTok{eat}\OperatorTok{()} \OperatorTok{\{}
\BuiltInTok{System}\OperatorTok{.}\FunctionTok{out}\OperatorTok{.}\FunctionTok{println}\OperatorTok{(}\NormalTok{name }\OperatorTok{+} \StringTok{" is eating."}\OperatorTok{);}
\OperatorTok{\}}

\KeywordTok{public} \DataTypeTok{void} \FunctionTok{sleep}\OperatorTok{()} \OperatorTok{\{}
\BuiltInTok{System}\OperatorTok{.}\FunctionTok{out}\OperatorTok{.}\FunctionTok{println}\OperatorTok{(}\NormalTok{name }\OperatorTok{+} \StringTok{" is sleeping."}\OperatorTok{);}
\OperatorTok{\}}

\KeywordTok{public} \DataTypeTok{void} \FunctionTok{shit}\OperatorTok{()} \OperatorTok{\{}
\BuiltInTok{System}\OperatorTok{.}\FunctionTok{out}\OperatorTok{.}\FunctionTok{println}\OperatorTok{(}\NormalTok{name }\OperatorTok{+} \StringTok{" is taking a shit."}\OperatorTok{);}
\OperatorTok{\}}
\OperatorTok{\}}
\end{Highlighting}
\end{Shaded}

I detta exempel har vi en klass som heter \texttt{Cat} och den
implementerar gränssnittet \texttt{Animal}. Vi har en privat egenskap
för namnet på katten och en konstruktor för att sätta namnet. Vi
implementerar också alla metoder från gränssnittet \texttt{Animal}.

Nu kan vi skapa en instans av katten och använda dess metoder:

\hypertarget{cb3}{}
\begin{Shaded}
\begin{Highlighting}[]
\KeywordTok{public} \KeywordTok{class}\NormalTok{ Main }\OperatorTok{\{}
\KeywordTok{public} \DataTypeTok{static} \DataTypeTok{void} \FunctionTok{main}\OperatorTok{(}\BuiltInTok{String}\OperatorTok{[]}\NormalTok{ args}\OperatorTok{)} \OperatorTok{\{}
\NormalTok{Cat cat }\OperatorTok{=} \KeywordTok{new} \FunctionTok{Cat}\OperatorTok{(}\StringTok{"Whiskers"}\OperatorTok{);}
\NormalTok{cat}\OperatorTok{.}\FunctionTok{eat}\OperatorTok{();}
\NormalTok{cat}\OperatorTok{.}\FunctionTok{sleep}\OperatorTok{();}
\NormalTok{cat}\OperatorTok{.}\FunctionTok{shit}\OperatorTok{();}
\OperatorTok{\}}
\OperatorTok{\}}
\end{Highlighting}
\end{Shaded}

Output:

\begin{verbatim}
Whiskers is eating.
Whiskers is sleeping.
Whiskers is taking a shit.
\end{verbatim}

I detta exempel skapar vi en instans av katten med namnet ``Whiskers''
och använder sedan dess metoder för att få katten att äta, sova och göra
sina behov. Resultatet skrivs ut i konsolen.

Detta är ett grundläggande exempel på hur gränssnitt kan användas för
att implementera polymorfism och separera implementationen av
komponenter från deras användning. Genom att använda gränssnitt kan vi
skapa flexibla och återanvändbara komponenter i våra program.

Gränssnitt är ett viktigt koncept inom objektorienterad programmering
och ger oss möjlighet att skapa en enhetlig och modulär kodstruktur.
Genom att definiera gränssnitt kan vi även underlätta samarbete mellan
olika utvecklare genom att specificera vilka metoder som förväntas
implementeras.

Användningen av gränssnitt är vanligt förekommande i Java och är en
viktig del av språkets design. Genom att använda gränssnitt kan vi skapa
mer flexibla och underhållbara program.

Kommentarer: - \texttt{Animal}: Gränssnittet \texttt{Animal} definierar
fyra metoder: \texttt{getName()}, \texttt{eat()}, \texttt{sleep()} och
\texttt{shit()}. Dessa metoder måste implementeras av alla klasser som
implementerar gränssnittet \texttt{Animal}. - \texttt{Cat}: Klassen
\texttt{Cat} implementerar gränssnittet \texttt{Animal} och
tillhandahåller implementationen av de fyra metoderna. Den har också en
privat egenskap \texttt{name} för att lagra namnet på katten. -
\texttt{Main}: Huvudklassen \texttt{Main} skapar en instans av
\texttt{Cat} med namnet ``Whiskers'' och använder sedan kattens metoder
för att simulera kattens beteende.

Genom att använda gränssnitt kan vi skapa en gemensam plattform för
olika klasser och möjliggöra utbytbarhet och polymorfism. Gränssnitt ger
oss flexibilitet och en tydlig struktur i vår kod.

Detta var en introduktion till gränssnitt i Java. Förhoppningsvis har du
fått en grundläggande förståelse för hur man skapar och använder
gränssnitt i dina Java-program. Fortsätt utforska och experimentera med
gränssnitt för att bygga mer flexibla och återanvändbara komponenter i
dina program!

\hypertarget{plugins}{%
\section{Plugins}\label{plugins}}

Plugins är en kraftfull mekanism för att lägga till ny funktionalitet i
en applikation. De ger en möjlighet att utöka och anpassa applikationen
utan att behöva modifiera dess kärnfunktionalitet. I denna artikel
kommer vi att utforska hur man implementerar ett pluginsystem med hjälp
av gränssnitt (interfaces) i Java.

\hypertarget{fuxf6rdelar-med-att-anvuxe4nda-plugins}{%
\subsection{Fördelar med att använda
Plugins}\label{fuxf6rdelar-med-att-anvuxe4nda-plugins}}

Att använda ett pluginsystem i din applikation kan ge flera fördelar:

\begin{itemize}
\item
  \textbf{Modularitet}: Ett pluginsystem främjar en modulär design där
  olika delar av funktionaliteten kan isoleras i separata plugins. Detta
  gör det enkelt att lägga till, ändra eller ta bort specifika
  funktioner utan att påverka resten av applikationen.
\item
  \textbf{Utbytbarhet}: Genom att använda ett pluginsystem kan användare
  eller utvecklare enkelt byta ut eller uppgradera plugins utan att
  påverka kärnfunktionaliteten i applikationen. Detta möjliggör snabb
  anpassning och utveckling av applikationen efter behov.
\item
  \textbf{Återanvändbarhet}: Plugins kan vara återanvändbara och
  användas i olika applikationer. Detta sparar tid och resurser genom
  att undvika att bygga samma funktionalitet från grunden i varje
  applikation.
\item
  \textbf{Flexibilitet}: Ett pluginsystem ger möjlighet att anpassa
  applikationen efter olika användares eller kunders krav och
  preferenser. Användare kan välja vilka plugins de vill använda och
  konfigurera applikationen enligt sina behov.
\end{itemize}

\hypertarget{hur-man-implementerar-ett-pluginsystem-med-gruxe4nssnitt-interfaces}{%
\subsection{Hur man implementerar ett pluginsystem med gränssnitt
(interfaces)}\label{hur-man-implementerar-ett-pluginsystem-med-gruxe4nssnitt-interfaces}}

För att implementera ett pluginsystem med hjälp av gränssnitt i Java
följer vi följande steg:

\hypertarget{steg-1-definiera-ett-gruxe4nssnitt-fuxf6r-plugins}{%
\subsubsection{Steg 1: Definiera ett gränssnitt för
Plugins}\label{steg-1-definiera-ett-gruxe4nssnitt-fuxf6r-plugins}}

Vi börjar med att definiera ett gränssnitt (interface) som beskriver den
gemensamma funktionaliteten som alla plugins måste implementera. Detta
gränssnitt kommer att fungera som en kontrakt för att säkerställa att
alla plugins följer samma struktur och regler.

\hypertarget{cb1}{}
\begin{Shaded}
\begin{Highlighting}[]
\KeywordTok{public} \KeywordTok{interface}\NormalTok{ Plugin }\OperatorTok{\{}
\DataTypeTok{void} \FunctionTok{run}\OperatorTok{();}
\OperatorTok{\}}
\end{Highlighting}
\end{Shaded}

I det här exemplet har vi definierat ett \texttt{Plugin}-gränssnitt med
en metod \texttt{run()}. Alla plugins i vårt system måste implementera
denna metod.

\hypertarget{steg-2-implementera-plugins}{%
\subsubsection{Steg 2: Implementera
Plugins}\label{steg-2-implementera-plugins}}

Nästa steg är att implementera olika plugins genom att skapa klasser som
implementerar det \texttt{Plugin}-gränssnitt som vi har definierat.
Varje pluginsklass kommer att innehålla den specifika funktionaliteten
för pluginet.

\hypertarget{cb2}{}
\begin{Shaded}
\begin{Highlighting}[]
\KeywordTok{public} \KeywordTok{class}\NormalTok{ MarvelPlugin }\KeywordTok{implements}\NormalTok{ Plugin }\OperatorTok{\{}
\AttributeTok{@Override}
\KeywordTok{public} \DataTypeTok{void} \FunctionTok{run}\OperatorTok{()} \OperatorTok{\{}
\BuiltInTok{System}\OperatorTok{.}\FunctionTok{out}\OperatorTok{.}\FunctionTok{println}\OperatorTok{(}\StringTok{"Kör Marvel{-}pluginet..."}\OperatorTok{);}
\CommentTok{// Implementera Marvel{-}specifik funktionalitet här}
\OperatorTok{\}}
\OperatorTok{\}}

\KeywordTok{public} \KeywordTok{class}\NormalTok{ DCPlugin }\KeywordTok{implements}\NormalTok{ Plugin }\OperatorTok{\{}
\AttributeTok{@Override}
\KeywordTok{public} \DataTypeTok{void} \FunctionTok{run}\OperatorTok{()} \OperatorTok{\{}
\BuiltInTok{System}\OperatorTok{.}\FunctionTok{out}\OperatorTok{.}\FunctionTok{println}\OperatorTok{(}\StringTok{"Kör DC{-}pluginet..."}\OperatorTok{);}
\CommentTok{// Implementera DC{-}specifik funktionalitet här}
\OperatorTok{\}}
\OperatorTok{\}}

\KeywordTok{public} \KeywordTok{class}\NormalTok{ StarWarsPlugin }\KeywordTok{implements}\NormalTok{ Plugin }\OperatorTok{\{}
\AttributeTok{@Override}
\KeywordTok{public} \DataTypeTok{void} \FunctionTok{run}\OperatorTok{()} \OperatorTok{\{}
\BuiltInTok{System}\OperatorTok{.}\FunctionTok{out}\OperatorTok{.}\FunctionTok{println}\OperatorTok{(}\StringTok{"Kör Star Wars{-}pluginet..."}\OperatorTok{);}
\CommentTok{// Implementera Star Wars{-}specifik funktionalitet här}
\OperatorTok{\}}
\OperatorTok{\}}
\end{Highlighting}
\end{Shaded}

I det här exemplet har vi implementerat tre olika plugins:
\texttt{MarvelPlugin}, \texttt{DCPlugin} och \texttt{StarWarsPlugin}.
Varje pluginsklass innehåller sin egen implementation av
\texttt{run()}-metoden.

\hypertarget{steg-3-ladda-och-kuxf6ra-plugins}{%
\subsubsection{Steg 3: Ladda och köra
Plugins}\label{steg-3-ladda-och-kuxf6ra-plugins}}

Det sista steget är att ladda och köra plugins i vår applikation. För
att göra detta kan vi använda Java Reflection API för att dynamiskt
ladda klasser och skapa instanser av dem.

\hypertarget{cb3}{}
\begin{Shaded}
\begin{Highlighting}[]
\KeywordTok{import} \ImportTok{java}\OperatorTok{.}\ImportTok{util}\OperatorTok{.}\ImportTok{ArrayList}\OperatorTok{;}
\KeywordTok{import} \ImportTok{java}\OperatorTok{.}\ImportTok{util}\OperatorTok{.}\ImportTok{List}\OperatorTok{;}

\KeywordTok{public} \KeywordTok{class}\NormalTok{ Application }\OperatorTok{\{}
\KeywordTok{private} \BuiltInTok{List}\OperatorTok{\textless{}}\NormalTok{Plugin}\OperatorTok{\textgreater{}}\NormalTok{ plugins}\OperatorTok{;}

\KeywordTok{public} \FunctionTok{Application}\OperatorTok{()} \OperatorTok{\{}
\NormalTok{plugins }\OperatorTok{=} \KeywordTok{new} \BuiltInTok{ArrayList}\OperatorTok{\textless{}\textgreater{}();}
\OperatorTok{\}}

\KeywordTok{public} \DataTypeTok{void} \FunctionTok{loadPlugins}\OperatorTok{()} \OperatorTok{\{}
\CommentTok{// Ladda alla plugins från en viss mapp eller konfiguration}
\CommentTok{// Till exempel kan vi lägga till MarvelPlugin, DCPlugin och StarWarsPlugin i plugins{-}listan}
\NormalTok{plugins}\OperatorTok{.}\FunctionTok{add}\OperatorTok{(}\KeywordTok{new} \FunctionTok{MarvelPlugin}\OperatorTok{());}
\NormalTok{plugins}\OperatorTok{.}\FunctionTok{add}\OperatorTok{(}\KeywordTok{new} \FunctionTok{DCPlugin}\OperatorTok{());}
\NormalTok{plugins}\OperatorTok{.}\FunctionTok{add}\OperatorTok{(}\KeywordTok{new} \FunctionTok{StarWarsPlugin}\OperatorTok{());}
\OperatorTok{\}}

\KeywordTok{public} \DataTypeTok{void} \FunctionTok{runPlugins}\OperatorTok{()} \OperatorTok{\{}
\ControlFlowTok{for} \OperatorTok{(}\NormalTok{Plugin plugin }\OperatorTok{:}\NormalTok{ plugins}\OperatorTok{)} \OperatorTok{\{}
\NormalTok{plugin}\OperatorTok{.}\FunctionTok{run}\OperatorTok{();}
\OperatorTok{\}}
\OperatorTok{\}}

\KeywordTok{public} \DataTypeTok{static} \DataTypeTok{void} \FunctionTok{main}\OperatorTok{(}\BuiltInTok{String}\OperatorTok{[]}\NormalTok{ args}\OperatorTok{)} \OperatorTok{\{}
\NormalTok{Application app }\OperatorTok{=} \KeywordTok{new} \FunctionTok{Application}\OperatorTok{();}
\NormalTok{app}\OperatorTok{.}\FunctionTok{loadPlugins}\OperatorTok{();}
\NormalTok{app}\OperatorTok{.}\FunctionTok{runPlugins}\OperatorTok{();}
\OperatorTok{\}}
\OperatorTok{\}}
\end{Highlighting}
\end{Shaded}

I det här exemplet har vi skapat en \texttt{Application}-klass som har
en lista av \texttt{Plugin}-instanser. Vi kan använda
\texttt{loadPlugins()}-metoden för att ladda alla plugins som vi vill
använda i vår applikation. Sedan kan vi använda
\texttt{runPlugins()}-metoden för att köra alla plugins i ordning.

\hypertarget{slutsats}{%
\subsection{Slutsats}\label{slutsats}}

Ett pluginsystem är en kraftfull mekanism som möjliggör enkel tillägg av
ny funktionalitet till en applikation. Genom att använda gränssnitt kan
vi definiera gemensam funktionalitet som alla plugins måste
implementera, vilket ger en flexibel och modulär arkitektur. Genom att
följa stegen som beskrivs ovan kan vi implementera ett pluginsystem i
Java och dra nytta av dess fördelar.

\hypertarget{anvuxe4ndningsomruxe5den-fuxf6r-pluginsystem}{%
\subsection{Användningsområden för
Pluginsystem}\label{anvuxe4ndningsomruxe5den-fuxf6r-pluginsystem}}

Pluginsystem kan vara användbara i en mängd olika scenarier och
applikationer. Här är några exempel på användningsområden:

\begin{itemize}
\tightlist
\item
  \textbf{Textredigerare}: Ett pluginsystem kan användas för att lägga
  till nya funktioner och redigeringsfunktion
\end{itemize}

er i en textredigerare, som stavningskontroll, syntaxhöjning eller
automatiserade uppgifter.

\begin{itemize}
\item
  \textbf{Webbläsare}: Webbläsare kan dra nytta av pluginsystem för att
  lägga till stöd för olika webbtekniker, tillägg för annonshantering
  eller användarskript.
\item
  \textbf{Grafikprogram}: Grafikprogram kan använda plugins för att
  lägga till nya filter, effekter eller verktyg för att utöka
  funktionaliteten.
\item
  \textbf{E-handelsplattformar}: Plugins kan användas för att lägga till
  betalningsportaler, fraktalternativ eller integration med externa
  tjänster i en e-handelsplattform.
\end{itemize}

Dessa är bara några exempel, och det finns många andra områden där
pluginsystem kan vara användbara.

\hypertarget{abstrakta-klasser}{%
\section{Abstrakta klasser}\label{abstrakta-klasser}}

Abstrakta klasser är klasser som innehåller både kod och abstrakta
metoder. De fungerar som en kombination av ett interface och en klass.

Innehållsförteckning

\{: .text-delta \} 1. TOC \{:toc\}

\hypertarget{beskrivning}{%
\subsection{Beskrivning}\label{beskrivning}}

En abstrakt klass används för att andra klasser ska kunna ärva från den.
Det går inte att skapa objekt av en abstrakt klass, utan den används
endast för att andra klasser ska kunna ärva från den.

\hypertarget{exempel}{%
\subsection{Exempel}\label{exempel}}

Vi skapar en klass kallad Shape, i den har vi metoden area(). Det ser
enkelt ut men det är effektfullt.

\hypertarget{cb1}{}
\begin{Shaded}
\begin{Highlighting}[]
\CommentTok{// Abstrakt klass för geometriska former}
\KeywordTok{public} \KeywordTok{abstract} \KeywordTok{class} \BuiltInTok{Shape} \OperatorTok{\{}
\CommentTok{// Abstrakt metod för att beräkna arean}
\KeywordTok{public} \KeywordTok{abstract} \DataTypeTok{double} \FunctionTok{area}\OperatorTok{();}
\OperatorTok{\}}

\CommentTok{// Underklass för cirkel som ärver från Shape}
\KeywordTok{public} \KeywordTok{class}\NormalTok{ Circle }\KeywordTok{extends} \BuiltInTok{Shape} \OperatorTok{\{}
\KeywordTok{private} \DataTypeTok{double}\NormalTok{ radius}\OperatorTok{;} \CommentTok{// Radien på cirkeln}

\CommentTok{// Konstruktor för att skapa en cirkel med given radie}
\KeywordTok{public} \FunctionTok{Circle}\OperatorTok{(}\DataTypeTok{double}\NormalTok{ radius}\OperatorTok{)} \OperatorTok{\{}
\KeywordTok{this}\OperatorTok{.}\FunctionTok{radius} \OperatorTok{=}\NormalTok{ radius}\OperatorTok{;}
\OperatorTok{\}}

\CommentTok{// Getter för att hämta radien på cirkeln}
\KeywordTok{public} \DataTypeTok{double} \FunctionTok{getRadius}\OperatorTok{()} \OperatorTok{\{}
\ControlFlowTok{return}\NormalTok{ radius}\OperatorTok{;}
\OperatorTok{\}}

\CommentTok{// Setter för att sätta radien på cirkeln}
\KeywordTok{public} \DataTypeTok{void} \FunctionTok{setRadius}\OperatorTok{(}\DataTypeTok{double}\NormalTok{ radius}\OperatorTok{)} \OperatorTok{\{}
\KeywordTok{this}\OperatorTok{.}\FunctionTok{radius} \OperatorTok{=}\NormalTok{ radius}\OperatorTok{;}
\OperatorTok{\}}

\CommentTok{// Överskuggning av area{-}metoden för att beräkna arean av cirkeln}
\AttributeTok{@Override}
\KeywordTok{public} \DataTypeTok{double} \FunctionTok{area}\OperatorTok{()} \OperatorTok{\{}
\ControlFlowTok{return} \BuiltInTok{Math}\OperatorTok{.}\FunctionTok{PI} \OperatorTok{*}\NormalTok{ radius }\OperatorTok{*}\NormalTok{ radius}\OperatorTok{;}
\OperatorTok{\}}
\OperatorTok{\}}

\CommentTok{// Underklass för rektangel som ärver från Shape}
\KeywordTok{public} \KeywordTok{class} \BuiltInTok{Rectangle} \KeywordTok{extends} \BuiltInTok{Shape} \OperatorTok{\{}
\KeywordTok{private} \DataTypeTok{double}\NormalTok{ width}\OperatorTok{;} \CommentTok{// Bredden på rektangeln}
\KeywordTok{private} \DataTypeTok{double}\NormalTok{ height}\OperatorTok{;} \CommentTok{// Höjden på rektangeln}

\CommentTok{// Konstruktor för att skapa en rektangel med given bredd och höjd}
\KeywordTok{public} \BuiltInTok{Rectangle}\OperatorTok{(}\DataTypeTok{double}\NormalTok{ width}\OperatorTok{,} \DataTypeTok{double}\NormalTok{ height}\OperatorTok{)} \OperatorTok{\{}
\KeywordTok{this}\OperatorTok{.}\FunctionTok{width} \OperatorTok{=}\NormalTok{ width}\OperatorTok{;}
\KeywordTok{this}\OperatorTok{.}\FunctionTok{height} \OperatorTok{=}\NormalTok{ height}\OperatorTok{;}
\OperatorTok{\}}

\CommentTok{// Getter för att hämta bredden på rektangeln}
\KeywordTok{public} \DataTypeTok{double} \FunctionTok{getWidth}\OperatorTok{()} \OperatorTok{\{}
\ControlFlowTok{return}\NormalTok{ width}\OperatorTok{;}
\OperatorTok{\}}

\CommentTok{// Setter för att sätta bredden på rektangeln}
\KeywordTok{public} \DataTypeTok{void} \FunctionTok{setWidth}\OperatorTok{(}\DataTypeTok{double}\NormalTok{ width}\OperatorTok{)} \OperatorTok{\{}
\KeywordTok{this}\OperatorTok{.}\FunctionTok{width} \OperatorTok{=}\NormalTok{ width}\OperatorTok{;}
\OperatorTok{\}}

\CommentTok{// Getter för att hämta höjden på rektangeln}
\KeywordTok{public} \DataTypeTok{double} \FunctionTok{getHeight}\OperatorTok{()} \OperatorTok{\{}
\ControlFlowTok{return}\NormalTok{ height}\OperatorTok{;}
\OperatorTok{\}}

\CommentTok{// Setter för att sätta höjden på rektangeln}
\KeywordTok{public} \DataTypeTok{void} \FunctionTok{setHeight}\OperatorTok{(}\DataTypeTok{double}\NormalTok{ height}\OperatorTok{)} \OperatorTok{\{}
\KeywordTok{this}\OperatorTok{.}\FunctionTok{height} \OperatorTok{=}\NormalTok{ height}\OperatorTok{;}
\OperatorTok{\}}

\CommentTok{// Överskuggning av area{-}metoden för att beräkna arean av rektangeln}
\AttributeTok{@Override}
\KeywordTok{public} \DataTypeTok{double} \FunctionTok{area}\OperatorTok{()} \OperatorTok{\{}
\ControlFlowTok{return}\NormalTok{ width }\OperatorTok{*}\NormalTok{ height}\OperatorTok{;}
\OperatorTok{\}}
\OperatorTok{\}}
\end{Highlighting}
\end{Shaded}

Vi kan nu skapa en array av typen Shape och lägga in objekt av typen
Circle och Rectangle i den. Vi kan sedan loopa igenom arrayen och anropa
area-metoden på varje objekt.

\hypertarget{cb2}{}
\begin{Shaded}
\begin{Highlighting}[]
\CommentTok{// Skapar en array av typen Shape}
\BuiltInTok{Shape}\OperatorTok{[]}\NormalTok{ shapes }\OperatorTok{=} \KeywordTok{new} \BuiltInTok{Shape}\OperatorTok{[}\DecValTok{2}\OperatorTok{];}

\CommentTok{// Lägger in ett objekt av typen Circle i arrayen}
\NormalTok{shapes}\OperatorTok{[}\DecValTok{0}\OperatorTok{]} \OperatorTok{=} \KeywordTok{new} \FunctionTok{Circle}\OperatorTok{(}\DecValTok{5}\OperatorTok{);}

\CommentTok{// Lägger in ett objekt av typen Rectangle i arrayen}
\NormalTok{shapes}\OperatorTok{[}\DecValTok{1}\OperatorTok{]} \OperatorTok{=} \KeywordTok{new} \BuiltInTok{Rectangle}\OperatorTok{(}\DecValTok{5}\OperatorTok{,} \DecValTok{10}\OperatorTok{);}

\CommentTok{// Loopar igenom arrayen och anropar area{-}metoden på varje objekt}
\ControlFlowTok{for} \OperatorTok{(}\BuiltInTok{Shape}\NormalTok{ shape }\OperatorTok{:}\NormalTok{ shapes}\OperatorTok{)} \OperatorTok{\{}
\BuiltInTok{System}\OperatorTok{.}\FunctionTok{out}\OperatorTok{.}\FunctionTok{println}\OperatorTok{(}\NormalTok{shape}\OperatorTok{.}\FunctionTok{area}\OperatorTok{());}
\OperatorTok{\}}
\end{Highlighting}
\end{Shaded}

Detta gör att vi kan använda oss av polymorfism för att beräkna arean av
olika geometriska former. Coolt va!

\hypertarget{referenser}{%
\subsection{Referenser}\label{referenser}}

\begin{itemize}
\tightlist
\item
  \href{https://docs.microsoft.com/en-us/dotnet/Java/programming-guide/classes-and-structs/abstract-classes}{Abstrakta
  klasser}
\item
  \href{https://www.w3schools.com/cs/cs_abstract.php}{Abstract class
  W3Schools}
\end{itemize}

\hypertarget{exempel}{%
\section{Exempel}\label{exempel}}

I den här artikeln ska vi titta på hur man skapar en abstrakt klass med
abstrakta och virtuella metoder i Java.

För detta behöver vi Maven libraryn JSOUP 1.14.3 eller nyare.

org.jsoup jsoup 1.14.3

\hypertarget{beskrivning}{%
\subsection{Beskrivning}\label{beskrivning}}

En abstrakt klass kan innehålla vanliga metoder, abstrakta och virtuella
metoder. Vi kommer att använda den här mallen för att skapa en abstrakt
klass som hanterar webbskrapning för olika sidor. Vi använder även
följande Java-paket:

\begin{itemize}
\tightlist
\item
  JSoup för att hantera HTML-dokument.
\end{itemize}

\hypertarget{fuxf6rdelar}{%
\subsection{Fördelar}\label{fuxf6rdelar}}

När du använder abstrakta klasser kan du dra nytta av följande fördelar:

\begin{enumerate}
\def\labelenumi{\arabic{enumi}.}
\tightlist
\item
  \textbf{Återanvändbarhet:} Abstrakta klasser kan fungera som grund för
  andra klasser och möjliggöra återanvändning av kod.
\item
  \textbf{Moduläritet:} Genom att använda abstrakta metoder kan du
  separera implementationen av en metod från själva klassen, vilket
  leder till en modulär design.
\item
  \textbf{Flexibilitet:} Abstrakta klasser kan användas som bas för
  olika implementationer och ge flexibilitet i utvecklingen.
\end{enumerate}

\hypertarget{begruxe4nsningar}{%
\subsection{Begränsningar}\label{begruxe4nsningar}}

Det finns några begränsningar att vara medveten om när du använder
abstrakta klasser:

\begin{enumerate}
\def\labelenumi{\arabic{enumi}.}
\tightlist
\item
  \textbf{Enkel arv:} En klass kan bara ärva från en enda abstrakt
  klass, vilket kan begränsa möjligheterna till hierarkiska strukturer.
\item
  \textbf{Ingen direkt instansiering:} Eftersom en abstrakt klass inte
  kan instansieras direkt måste den ärvas och implementeras i en konkret
  klass för att användas.
\end{enumerate}

\hypertarget{anvuxe4ndningsomruxe5den}{%
\subsection{Användningsområden}\label{anvuxe4ndningsomruxe5den}}

Abstrakta klasser är användbara i olika scenarier, inklusive:

\begin{enumerate}
\def\labelenumi{\arabic{enumi}.}
\tightlist
\item
  \textbf{Framtida utbyggbarhet:} Genom att definiera abstrakta metoder
  kan du planera för framtida utökning och implementation av specifika
  beteenden.
\item
  \textbf{Grundläggande mallar:} Abstrakta klasser kan fungera som grund
  för att skapa mallar eller ramverk med fördefinierad funktionalitet.
\item
  \textbf{Plugin-arkitektur:} Genom att använda abstrakta klasser kan du
  skapa en flexibel arkitektur för att lägga till och hantera plugins i
  din applikation.
\end{enumerate}

\hypertarget{kodexempel}{%
\subsection{Kodexempel}\label{kodexempel}}

Här är en implementation av en abstrakt klass för webbskrapning i Java:

\hypertarget{cb1}{}
\begin{Shaded}
\begin{Highlighting}[]
\KeywordTok{import} \ImportTok{org}\OperatorTok{.}\ImportTok{jsoup}\OperatorTok{.}\ImportTok{Jsoup}\OperatorTok{;}
\KeywordTok{import} \ImportTok{org}\OperatorTok{.}\ImportTok{jsoup}\OperatorTok{.}\ImportTok{nodes}\OperatorTok{.}\ImportTok{Document}\OperatorTok{;}
\KeywordTok{import} \ImportTok{org}\OperatorTok{.}\ImportTok{jsoup}\OperatorTok{.}\ImportTok{nodes}\OperatorTok{.}\ImportTok{Element}\OperatorTok{;}
\KeywordTok{import} \ImportTok{org}\OperatorTok{.}\ImportTok{jsoup}\OperatorTok{.}\ImportTok{select}\OperatorTok{.}\ImportTok{Elements}\OperatorTok{;}

\KeywordTok{import} \ImportTok{java}\OperatorTok{.}\ImportTok{io}\OperatorTok{.}\ImportTok{IOException}\OperatorTok{;}
\KeywordTok{import} \ImportTok{java}\OperatorTok{.}\ImportTok{util}\OperatorTok{.}\ImportTok{ArrayList}\OperatorTok{;}
\KeywordTok{import} \ImportTok{java}\OperatorTok{.}\ImportTok{util}\OperatorTok{.}\ImportTok{List}\OperatorTok{;}

\KeywordTok{public} \KeywordTok{abstract} \KeywordTok{class}\NormalTok{ WebScraper }\OperatorTok{\{}
\KeywordTok{private} \BuiltInTok{String}\NormalTok{ url }\OperatorTok{=} \StringTok{""}\OperatorTok{;}
\KeywordTok{private} \BuiltInTok{String}\NormalTok{ title }\OperatorTok{=} \StringTok{""}\OperatorTok{;}
\KeywordTok{private} \BuiltInTok{String}\NormalTok{ description }\OperatorTok{=} \StringTok{""}\OperatorTok{;}
\KeywordTok{private} \BuiltInTok{String}\NormalTok{ tags }\OperatorTok{=} \StringTok{""}\OperatorTok{;}
\KeywordTok{private} \BuiltInTok{Document}\NormalTok{ htmlDocument }\OperatorTok{=} \KeywordTok{null}\OperatorTok{;}

\KeywordTok{public} \BuiltInTok{String} \FunctionTok{getUrl}\OperatorTok{()} \OperatorTok{\{}
\ControlFlowTok{return}\NormalTok{ url}\OperatorTok{;}
\OperatorTok{\}}

\KeywordTok{public} \DataTypeTok{void} \FunctionTok{setUrl}\OperatorTok{(}\BuiltInTok{String}\NormalTok{ url}\OperatorTok{)} \OperatorTok{\{}
\KeywordTok{this}\OperatorTok{.}\FunctionTok{url} \OperatorTok{=}\NormalTok{ url}\OperatorTok{;}
\OperatorTok{\}}

\KeywordTok{public} \BuiltInTok{String} \FunctionTok{getTitle}\OperatorTok{()} \OperatorTok{\{}
\ControlFlowTok{return}\NormalTok{ title}\OperatorTok{;}
\OperatorTok{\}}

\KeywordTok{public} \DataTypeTok{void} \FunctionTok{setTitle}\OperatorTok{(}\BuiltInTok{String}\NormalTok{ title}\OperatorTok{)} \OperatorTok{\{}
\KeywordTok{this}\OperatorTok{.}\FunctionTok{title} \OperatorTok{=}\NormalTok{ title}\OperatorTok{;}
\OperatorTok{\}}

\KeywordTok{public} \BuiltInTok{String} \FunctionTok{getDescription}\OperatorTok{()} \OperatorTok{\{}
\ControlFlowTok{return}\NormalTok{ description}\OperatorTok{;}
\OperatorTok{\}}

\KeywordTok{public} \DataTypeTok{void} \FunctionTok{setDescription}\OperatorTok{(}\BuiltInTok{String}\NormalTok{ description}\OperatorTok{)} \OperatorTok{\{}
\KeywordTok{this}\OperatorTok{.}\FunctionTok{description} \OperatorTok{=}\NormalTok{ description}\OperatorTok{;}
\OperatorTok{\}}

\KeywordTok{public} \BuiltInTok{String} \FunctionTok{getTags}\OperatorTok{()} \OperatorTok{\{}
\ControlFlowTok{return}\NormalTok{ tags}\OperatorTok{;}
\OperatorTok{\}}

\KeywordTok{public} \DataTypeTok{void} \FunctionTok{setTags}\OperatorTok{(}\BuiltInTok{String}\NormalTok{ tags}\OperatorTok{)} \OperatorTok{\{}
\KeywordTok{this}\OperatorTok{.}\FunctionTok{tags} \OperatorTok{=}\NormalTok{ tags}\OperatorTok{;}
\OperatorTok{\}}

\KeywordTok{public} \BuiltInTok{Document} \FunctionTok{getHtmlDocument}\OperatorTok{()} \OperatorTok{\{}
\ControlFlowTok{return}\NormalTok{ htmlDocument}\OperatorTok{;}
\OperatorTok{\}}

\KeywordTok{public} \DataTypeTok{void} \FunctionTok{setHtmlDocument}\OperatorTok{(}\BuiltInTok{Document}\NormalTok{ htmlDocument}\OperatorTok{)} \OperatorTok{\{}
\KeywordTok{this}\OperatorTok{.}\FunctionTok{htmlDocument} \OperatorTok{=}\NormalTok{ htmlDocument}\OperatorTok{;}
\OperatorTok{\}}

\KeywordTok{public} \BuiltInTok{String} \FunctionTok{getHtml}\OperatorTok{()} \OperatorTok{\{}
\ControlFlowTok{return}\NormalTok{ htmlDocument }\OperatorTok{!=} \KeywordTok{null} \OperatorTok{?}\NormalTok{ htmlDocument}\OperatorTok{.}\FunctionTok{outerHtml}\OperatorTok{()} \OperatorTok{:;}
\OperatorTok{\}}

\KeywordTok{public} \BuiltInTok{String} \FunctionTok{getText}\OperatorTok{()} \OperatorTok{\{}
\ControlFlowTok{return}\NormalTok{ htmlDocument }\OperatorTok{!=} \KeywordTok{null} \OperatorTok{?}\NormalTok{ htmlDocument}\OperatorTok{.}\FunctionTok{text}\OperatorTok{()} \OperatorTok{:;}
\OperatorTok{\}}

\KeywordTok{public} \BuiltInTok{String} \FunctionTok{getHtmlDocument}\OperatorTok{(}\BuiltInTok{String}\NormalTok{ url}\OperatorTok{)} \OperatorTok{\{}
\FunctionTok{setUrl}\OperatorTok{(}\NormalTok{url}\OperatorTok{);}
\ControlFlowTok{try} \OperatorTok{\{}
\NormalTok{htmlDocument }\OperatorTok{=}\NormalTok{ Jsoup}\OperatorTok{.}\FunctionTok{connect}\OperatorTok{(}\NormalTok{url}\OperatorTok{).}\FunctionTok{get}\OperatorTok{();}
\NormalTok{title }\OperatorTok{=}\NormalTok{ htmlDocument}\OperatorTok{.}\FunctionTok{title}\OperatorTok{();}
\NormalTok{description }\OperatorTok{=}\NormalTok{ htmlDocument}\OperatorTok{.}\FunctionTok{select}\OperatorTok{(}\StringTok{"meta[name=description]"}\OperatorTok{).}\FunctionTok{attr}\OperatorTok{(}\StringTok{"content"}\OperatorTok{);}
\NormalTok{tags }\OperatorTok{=}\NormalTok{ htmlDocument}\OperatorTok{.}\FunctionTok{select}\OperatorTok{(}\StringTok{"meta[name=keywords]"}\OperatorTok{).}\FunctionTok{attr}\OperatorTok{(}\StringTok{"content"}\OperatorTok{);}
\OperatorTok{\}} \ControlFlowTok{catch} \OperatorTok{(}\BuiltInTok{IOException}\NormalTok{ e}\OperatorTok{)} \OperatorTok{\{}
\NormalTok{e}\OperatorTok{.}\FunctionTok{printStackTrace}\OperatorTok{();}
\OperatorTok{\}}
\ControlFlowTok{return} \FunctionTok{getHtml}\OperatorTok{();}
\OperatorTok{\}}

\KeywordTok{public} \BuiltInTok{String} \FunctionTok{getDivById}\OperatorTok{(}\BuiltInTok{String}\NormalTok{ id}\OperatorTok{)} \OperatorTok{\{}
\BuiltInTok{Element}\NormalTok{ div }\OperatorTok{=}\NormalTok{ htmlDocument }\OperatorTok{!=} \KeywordTok{null} \OperatorTok{?}\NormalTok{ htmlDocument}\OperatorTok{.}\FunctionTok{getElementById}\OperatorTok{(}\NormalTok{id}\OperatorTok{)} \OperatorTok{:} \KeywordTok{null}\OperatorTok{;}
\ControlFlowTok{return}\NormalTok{ div }\OperatorTok{!=} \KeywordTok{null} \OperatorTok{?}\NormalTok{ div}\OperatorTok{.}\FunctionTok{outerHtml}\OperatorTok{()} \OperatorTok{:;}
\OperatorTok{\}}

\KeywordTok{public} \BuiltInTok{List}\OperatorTok{\textless{}}\BuiltInTok{String}\OperatorTok{\textgreater{}} \FunctionTok{getDivByClass}\OperatorTok{(}\BuiltInTok{String}\NormalTok{ className}\OperatorTok{)} \OperatorTok{\{}
\BuiltInTok{Elements}\NormalTok{ divs }\OperatorTok{=}\NormalTok{ htmlDocument }\OperatorTok{!=} \KeywordTok{null} \OperatorTok{?}\NormalTok{ htmlDocument}\OperatorTok{.}\FunctionTok{getElementsByClass}\OperatorTok{(}\NormalTok{className}\OperatorTok{)} \OperatorTok{:} \KeywordTok{null}\OperatorTok{;}
\BuiltInTok{List}\OperatorTok{\textless{}}\BuiltInTok{String}\OperatorTok{\textgreater{}}\NormalTok{ divHtmlList }\OperatorTok{=} \KeywordTok{new} \BuiltInTok{ArrayList}\OperatorTok{\textless{}\textgreater{}();}
\ControlFlowTok{if} \OperatorTok{(}\NormalTok{divs }\OperatorTok{!=} \KeywordTok{null}\OperatorTok{)} \OperatorTok{\{}
\ControlFlowTok{for} \OperatorTok{(}\BuiltInTok{Element}\NormalTok{ div }\OperatorTok{:}\NormalTok{ divs}\OperatorTok{)} \OperatorTok{\{}
\NormalTok{divHtmlList}\OperatorTok{.}\FunctionTok{add}\OperatorTok{(}\NormalTok{div}\OperatorTok{.}\FunctionTok{outerHtml}\OperatorTok{());}
\OperatorTok{\}}
\OperatorTok{\}}
\ControlFlowTok{return}\NormalTok{ divHtmlList}\OperatorTok{;}
\OperatorTok{\}}

\KeywordTok{public} \BuiltInTok{List}\OperatorTok{\textless{}}\BuiltInTok{String}\OperatorTok{\textgreater{}} \FunctionTok{getElementByClass}\OperatorTok{(}\BuiltInTok{String}\NormalTok{ element}\OperatorTok{,} \BuiltInTok{String}\NormalTok{ className}\OperatorTok{)} \OperatorTok{\{}
\BuiltInTok{Elements}\NormalTok{ elements }\OperatorTok{=}\NormalTok{ htmlDocument }\OperatorTok{!=} \KeywordTok{null} \OperatorTok{?}\NormalTok{ htmlDocument}\OperatorTok{.}\FunctionTok{select}\OperatorTok{(}\NormalTok{element }\OperatorTok{+} \StringTok{"."} \OperatorTok{+}\NormalTok{ className}\OperatorTok{)} \OperatorTok{:} \KeywordTok{null}\OperatorTok{;}
\BuiltInTok{List}\OperatorTok{\textless{}}\BuiltInTok{String}\OperatorTok{\textgreater{}}\NormalTok{ elementHtmlList }\OperatorTok{=} \KeywordTok{new} \BuiltInTok{ArrayList}\OperatorTok{\textless{}\textgreater{}();}
\ControlFlowTok{if} \OperatorTok{(}\NormalTok{elements }\OperatorTok{!=} \KeywordTok{null}\OperatorTok{)} \OperatorTok{\{}
\ControlFlowTok{for} \OperatorTok{(}\BuiltInTok{Element}\NormalTok{ el }\OperatorTok{:}\NormalTok{ elements}\OperatorTok{)} \OperatorTok{\{}
\NormalTok{elementHtmlList}\OperatorTok{.}\FunctionTok{add}\OperatorTok{(}\NormalTok{el}\OperatorTok{.}\FunctionTok{outerHtml}\OperatorTok{());}
\OperatorTok{\}}
\OperatorTok{\}}
\ControlFlowTok{return}\NormalTok{ elementHtmlList}\OperatorTok{;}
\OperatorTok{\}}

\KeywordTok{public} \BuiltInTok{List}\OperatorTok{\textless{}}\BuiltInTok{String}\OperatorTok{\textgreater{}} \FunctionTok{getImages}\OperatorTok{()} \OperatorTok{\{}
\BuiltInTok{List}\OperatorTok{\textless{}}\BuiltInTok{String}\OperatorTok{\textgreater{}}\NormalTok{ imageUrls }\OperatorTok{=} \KeywordTok{new} \BuiltInTok{ArrayList}\OperatorTok{\textless{}\textgreater{}();}
\ControlFlowTok{if} \OperatorTok{(}\NormalTok{htmlDocument }\OperatorTok{!=} \KeywordTok{null}\OperatorTok{)} \OperatorTok{\{}
\BuiltInTok{Elements}\NormalTok{ images }\OperatorTok{=}\NormalTok{ htmlDocument}\OperatorTok{.}\FunctionTok{select}\OperatorTok{(}\StringTok{"img"}\OperatorTok{);}
\ControlFlowTok{for} \OperatorTok{(}\BuiltInTok{Element}\NormalTok{ image }\OperatorTok{:}\NormalTok{ images}\OperatorTok{)} \OperatorTok{\{}
\BuiltInTok{String}\NormalTok{ imageUrl }\OperatorTok{=}\NormalTok{ image}\OperatorTok{.}\FunctionTok{attr}\OperatorTok{(}\StringTok{"src"}\OperatorTok{);}
\ControlFlowTok{if} \OperatorTok{(}\NormalTok{imageUrl}\OperatorTok{.}\FunctionTok{endsWith}\OperatorTok{(}\StringTok{".jpg"}\OperatorTok{)} \OperatorTok{||}\NormalTok{ imageUrl}\OperatorTok{.}\FunctionTok{endsWith}\OperatorTok{(}\StringTok{".png"}\OperatorTok{))} \OperatorTok{\{}
\ControlFlowTok{if} \OperatorTok{(!}\NormalTok{imageUrl}\OperatorTok{.}\FunctionTok{startsWith}\OperatorTok{(}\StringTok{"http:"}\OperatorTok{)} \OperatorTok{\&\&} \OperatorTok{!}\NormalTok{imageUrl}\OperatorTok{.}\FunctionTok{startsWith}\OperatorTok{(}\StringTok{"https:"}\OperatorTok{))} \OperatorTok{\{}
\NormalTok{imageUrl }\OperatorTok{=}\NormalTok{ url}\OperatorTok{.}\FunctionTok{trim}\OperatorTok{()} \OperatorTok{+}\NormalTok{ imageUrl}\OperatorTok{;}
\OperatorTok{\}}
\NormalTok{imageUrls}\OperatorTok{.}\FunctionTok{add}\OperatorTok{(}\NormalTok{imageUrl}\OperatorTok{);}
\OperatorTok{\}}
\OperatorTok{\}}
\OperatorTok{\}}
\ControlFlowTok{return}\NormalTok{ imageUrls}\OperatorTok{;}
\OperatorTok{\}}

\KeywordTok{public} \KeywordTok{abstract} \DataTypeTok{void} \FunctionTok{scrape}\OperatorTok{(}\BuiltInTok{String}\NormalTok{ url}\OperatorTok{);}
\OperatorTok{\}}
\end{Highlighting}
\end{Shaded}

Nu ska vi göra en implementation av denna klass för att hämta innehållet
från en specifik sida.

\hypertarget{cb2}{}
\begin{Shaded}
\begin{Highlighting}[]
\KeywordTok{import} \ImportTok{org}\OperatorTok{.}\ImportTok{jsoup}\OperatorTok{.}\ImportTok{Jsoup}\OperatorTok{;}
\KeywordTok{import} \ImportTok{org}\OperatorTok{.}\ImportTok{jsoup}\OperatorTok{.}\ImportTok{nodes}\OperatorTok{.}\ImportTok{Document}\OperatorTok{;}
\KeywordTok{import} \ImportTok{org}\OperatorTok{.}\ImportTok{jsoup}\OperatorTok{.}\ImportTok{nodes}\OperatorTok{.}\ImportTok{Element}\OperatorTok{;}
\KeywordTok{import} \ImportTok{org}\OperatorTok{.}\ImportTok{jsoup}\OperatorTok{.}\ImportTok{select}\OperatorTok{.}\ImportTok{Elements}\OperatorTok{;}

\KeywordTok{import} \ImportTok{java}\OperatorTok{.}\ImportTok{io}\OperatorTok{.}\ImportTok{IOException}\OperatorTok{;}
\KeywordTok{import} \ImportTok{java}\OperatorTok{.}\ImportTok{util}\OperatorTok{.}\ImportTok{List}\OperatorTok{;}

\KeywordTok{public} \KeywordTok{class}\NormalTok{ GetKittens }\KeywordTok{extends}\NormalTok{ WebScraper }\OperatorTok{\{}
\KeywordTok{public} \BuiltInTok{String}

\FunctionTok{downloadKitten}\OperatorTok{()} \OperatorTok{\{}
\FunctionTok{scrape}\OperatorTok{(}\StringTok{"https://www.pinterest.com/katesaidy/cute{-}kitten{-}pics/"}\OperatorTok{);}
\BuiltInTok{String}\NormalTok{ imgFolder }\OperatorTok{=} \BuiltInTok{System}\OperatorTok{.}\FunctionTok{getProperty}\OperatorTok{(}\StringTok{"user.home"}\OperatorTok{)} \OperatorTok{+} \StringTok{"/Pictures"}\OperatorTok{;}
\BuiltInTok{String}\NormalTok{ filename }\OperatorTok{=}\NormalTok{ imgFolder }\OperatorTok{+} \StringTok{"/Daily kitten.jpg"}\OperatorTok{;}
\BuiltInTok{File}\NormalTok{ file }\OperatorTok{=} \KeywordTok{new} \BuiltInTok{File}\OperatorTok{(}\NormalTok{filename}\OperatorTok{);}
\ControlFlowTok{if} \OperatorTok{(}\NormalTok{file}\OperatorTok{.}\FunctionTok{exists}\OperatorTok{())} \OperatorTok{\{}
\ControlFlowTok{return}\NormalTok{ filename}\OperatorTok{;}
\OperatorTok{\}} \ControlFlowTok{else} \OperatorTok{\{}
\ControlFlowTok{return} \StringTok{""}\OperatorTok{;}
\OperatorTok{\}}
\OperatorTok{\}}

\AttributeTok{@Override}
\KeywordTok{public} \BuiltInTok{List}\OperatorTok{\textless{}}\BuiltInTok{String}\OperatorTok{\textgreater{}} \FunctionTok{getImages}\OperatorTok{()} \OperatorTok{\{}
\BuiltInTok{List}\OperatorTok{\textless{}}\BuiltInTok{String}\OperatorTok{\textgreater{}}\NormalTok{ imageUrls }\OperatorTok{=} \KeywordTok{super}\OperatorTok{.}\FunctionTok{getImages}\OperatorTok{();}

\CommentTok{// IMG SRC fungerade inte på Pinterest så vi söker med regex istället}
\CommentTok{// För att göra det mer genomskinnligt så överridar vi bara GetImages}
\ControlFlowTok{if} \OperatorTok{(}\FunctionTok{getHtmlDocument}\OperatorTok{()} \OperatorTok{!=} \KeywordTok{null}\OperatorTok{)} \OperatorTok{\{}
\BuiltInTok{String}\NormalTok{ html }\OperatorTok{=} \FunctionTok{getHtml}\OperatorTok{();}
\BuiltInTok{String}\NormalTok{ regex }\OperatorTok{=} \StringTok{"([\textquotesingle{}}\SpecialCharTok{\textbackslash{}"}\StringTok{])([\^{}\textquotesingle{}}\SpecialCharTok{\textbackslash{}"}\StringTok{]+}\SpecialCharTok{\textbackslash{}\textbackslash{}}\StringTok{.(jpg|png|bmp|gif))}\SpecialCharTok{\textbackslash{}\textbackslash{}}\StringTok{1"}\OperatorTok{;}
\BuiltInTok{Pattern}\NormalTok{ pattern }\OperatorTok{=} \BuiltInTok{Pattern}\OperatorTok{.}\FunctionTok{compile}\OperatorTok{(}\NormalTok{regex}\OperatorTok{);}
\BuiltInTok{Matcher}\NormalTok{ matcher }\OperatorTok{=}\NormalTok{ pattern}\OperatorTok{.}\FunctionTok{matcher}\OperatorTok{(}\NormalTok{html}\OperatorTok{);}
\ControlFlowTok{while} \OperatorTok{(}\NormalTok{matcher}\OperatorTok{.}\FunctionTok{find}\OperatorTok{())} \OperatorTok{\{}
\BuiltInTok{System}\OperatorTok{.}\FunctionTok{out}\OperatorTok{.}\FunctionTok{println}\OperatorTok{(}\NormalTok{matcher}\OperatorTok{.}\FunctionTok{group}\OperatorTok{(}\DecValTok{2}\OperatorTok{));}
\NormalTok{imageUrls}\OperatorTok{.}\FunctionTok{add}\OperatorTok{(}\NormalTok{matcher}\OperatorTok{.}\FunctionTok{group}\OperatorTok{(}\DecValTok{2}\OperatorTok{));}
\OperatorTok{\}}
\OperatorTok{\}}
\ControlFlowTok{return}\NormalTok{ imageUrls}\OperatorTok{;}
\OperatorTok{\}}

\CommentTok{// Här är scrape{-}metoden som vi måste implementera}
\AttributeTok{@Override}
\KeywordTok{public} \DataTypeTok{void} \FunctionTok{scrape}\OperatorTok{(}\BuiltInTok{String}\NormalTok{ url}\OperatorTok{)} \OperatorTok{\{}
\BuiltInTok{System}\OperatorTok{.}\FunctionTok{out}\OperatorTok{.}\FunctionTok{println}\OperatorTok{(}\StringTok{"Downloading page"}\OperatorTok{);}
\FunctionTok{getHtmlDocument}\OperatorTok{(}\NormalTok{url}\OperatorTok{);}
\BuiltInTok{System}\OperatorTok{.}\FunctionTok{out}\OperatorTok{.}\FunctionTok{println}\OperatorTok{(}\StringTok{"Getting image list"}\OperatorTok{);}
\BuiltInTok{List}\OperatorTok{\textless{}}\BuiltInTok{String}\OperatorTok{\textgreater{}}\NormalTok{ images }\OperatorTok{=} \FunctionTok{getImages}\OperatorTok{();}
\BuiltInTok{System}\OperatorTok{.}\FunctionTok{out}\OperatorTok{.}\FunctionTok{println}\OperatorTok{(}\StringTok{"Removing crap images"}\OperatorTok{);}
\NormalTok{images}\OperatorTok{.}\FunctionTok{removeIf}\OperatorTok{(}\NormalTok{pic }\OperatorTok{{-}\textgreater{}}\NormalTok{ pic}\OperatorTok{.}\FunctionTok{contains}\OperatorTok{(}\StringTok{"/images/user/"}\OperatorTok{));}
\ControlFlowTok{if} \OperatorTok{(}\NormalTok{images}\OperatorTok{.}\FunctionTok{isEmpty}\OperatorTok{())} \OperatorTok{\{}
\ControlFlowTok{return}\OperatorTok{;}
\OperatorTok{\}}

\CommentTok{// Hämtar en slumpmässig bild från listan}
\BuiltInTok{Random}\NormalTok{ random }\OperatorTok{=} \KeywordTok{new} \BuiltInTok{Random}\OperatorTok{();}
\BuiltInTok{System}\OperatorTok{.}\FunctionTok{out}\OperatorTok{.}\FunctionTok{println}\OperatorTok{(}\StringTok{"Selecting random image"}\OperatorTok{);}
\BuiltInTok{String}\NormalTok{ randomImage }\OperatorTok{=}\NormalTok{ images}\OperatorTok{.}\FunctionTok{get}\OperatorTok{(}\NormalTok{random}\OperatorTok{.}\FunctionTok{nextInt}\OperatorTok{(}\NormalTok{images}\OperatorTok{.}\FunctionTok{size}\OperatorTok{()));}
\BuiltInTok{System}\OperatorTok{.}\FunctionTok{out}\OperatorTok{.}\FunctionTok{println}\OperatorTok{(}\StringTok{"Downloading image"}\OperatorTok{);}
\BuiltInTok{System}\OperatorTok{.}\FunctionTok{out}\OperatorTok{.}\FunctionTok{println}\OperatorTok{(}\NormalTok{randomImage}\OperatorTok{);}
\ControlFlowTok{try} \OperatorTok{\{}
\BuiltInTok{URL}\NormalTok{ imageUrl }\OperatorTok{=} \KeywordTok{new} \BuiltInTok{URL}\OperatorTok{(}\NormalTok{randomImage}\OperatorTok{);}
\BuiltInTok{BufferedImage}\NormalTok{ image }\OperatorTok{=} \BuiltInTok{ImageIO}\OperatorTok{.}\FunctionTok{read}\OperatorTok{(}\NormalTok{imageUrl}\OperatorTok{);}
\BuiltInTok{System}\OperatorTok{.}\FunctionTok{out}\OperatorTok{.}\FunctionTok{println}\OperatorTok{(}\StringTok{"Saving image"}\OperatorTok{);}
\BuiltInTok{String}\NormalTok{ imgFolder }\OperatorTok{=} \BuiltInTok{System}\OperatorTok{.}\FunctionTok{getProperty}\OperatorTok{(}\StringTok{"user.home"}\OperatorTok{)} \OperatorTok{+} \StringTok{"/Pictures"}\OperatorTok{;}
\BuiltInTok{String}\NormalTok{ filename }\OperatorTok{=}\NormalTok{ imgFolder }\OperatorTok{+} \StringTok{"/Daily kitten.jpg"}\OperatorTok{;}
\BuiltInTok{File}\NormalTok{ file }\OperatorTok{=} \KeywordTok{new} \BuiltInTok{File}\OperatorTok{(}\NormalTok{filename}\OperatorTok{);}
\BuiltInTok{ImageIO}\OperatorTok{.}\FunctionTok{write}\OperatorTok{(}\NormalTok{image}\OperatorTok{,} \StringTok{"jpg"}\OperatorTok{,}\NormalTok{ file}\OperatorTok{);}
\OperatorTok{\}} \ControlFlowTok{catch} \OperatorTok{(}\BuiltInTok{IOException}\NormalTok{ e}\OperatorTok{)} \OperatorTok{\{}
\NormalTok{e}\OperatorTok{.}\FunctionTok{printStackTrace}\OperatorTok{();}
\OperatorTok{\}}
\OperatorTok{\}}
\OperatorTok{\}}
\end{Highlighting}
\end{Shaded}

Jag hoppas att du gillar exemplet och att det är till hjälp för dig.
Kattbilder är alltid en hit! Ha det roligt med webbscraping och
experimentera med olika sidor att hämta innehåll från. Lycka till! 😺

\hypertarget{termer}{%
\subsection{Termer}\label{termer}}

Här är en lista över några termer som används i koden:

\begin{itemize}
\tightlist
\item
  \texttt{WebScraper}: En abstrakt klass som hanterar webbskrapning och
  definierar olika metoder för att hämta och bearbeta webbsidor.
\item
  \texttt{url}: En variabel som representerar adressen till den webbsida
  som ska webscrapas.
\item
  \texttt{title}: En variabel som representerar titeln på webbsidan.
\item
  \texttt{description}: En variabel som representerar sidans egna
  beskrivning.
\item
  \texttt{tags}: En variabel som representerar meta-nyckelord från
  webbsidan.
\item
  \texttt{htmlDocument}: En variabel som representerar det HTML-dokument
  som hämtas från webbsidan.
\item
  \texttt{html}: En metod som returnerar HTML-dokumentet som en sträng.
\item
  \texttt{text}: En metod som returnerar all text på webbsidan utan
  HTML-taggar.
\item
  \texttt{getHtmlDocument}: En metod som öppnar en webbsida och hämtar
  dess HTML-dokument.
\item
  \texttt{getDivById}: En metod som hämtar innehållet i en div med en
  specifik id från webbsidan.
\item
  \texttt{getDivByClass}: En metod som hämtar innehållet i alla divar
  med en specifik klass från webbsidan.
\item
  \texttt{getElementByClass}: En metod som hämtar innehållet i alla
  element med en specifik klass från webbsidan.
\item
  \texttt{getImages}: En metod som hämtar URL:er till alla bilder (IMG
  SRC-länkar) på webbsidan.
\item
  \texttt{scrape}: En abstrakt metod som definierar hur webbsidan ska
  webscrapas.
\item
  \texttt{GetKittens}: En klass som ärver från \texttt{WebScraper} och
  implementerar \texttt{scrape}-metoden för att hämta kattbilder från en
  specifik sida.
\item
  \texttt{Abstrakt\ klass}: En klass som inte kan instansieras direkt
  utan måste ärvas och implementeras i en konkret klass.
\item
  \texttt{Virtuell\ metod}: En metod som kan överskridas i en subklass.
\item
  \texttt{Paket}: En samling av Java-klasser och resurser som kan
  användas för att tillhandahålla specifik funktionalitet.
\item
  \texttt{Grafik}: Bilder, ikoner och andra visuella element som används
  på en webbsida.
\item
  \texttt{HTML-dokument}: En textfil som innehåller HTML-kod för att
  skapa en webbsida.
\item
  \texttt{Överskrida}: Att skriva om en metod i en subklass.
\item
  \texttt{Utöka}: Att lägga till funktionalitet i en subklass.
\item
  \texttt{Klassimplementation}: En klass som ärver från en abstrakt
  klass och implementerar dess abstrakta metoder.
\item
  \texttt{Återanvändbar\ kod}: Kod som kan återanvändas i olika delar av
  ett projekt.
\item
  \texttt{Underhållbar\ kod}: Kod som är lätt att förstå och underhålla.
\end{itemize}

\hypertarget{slutsats}{%
\subsection{Slutsats}\label{slutsats}}

I den här artikeln har vi utforskat användningen av abstrakta klasser
och metoder i Java. Vi har sett hur man kan skapa en abstrakt klass med
abstrakta och virtuella metoder för att hantera webbskrapning av olika
sidor. Vi har också diskuterat användningen av JSoup för att underlätta
hanteringen av HTML-dokument.

En abstrakt klass ger oss en bra grundstruktur för att implementera
specifik funktionalitet och samtidigt möjliggöra flexibilitet genom att
tillåta subklasser att överskrida och utöka funktionaliteten genom att
implementera abstrakta metoder. Detta hjälper oss att skapa
återanvändbar och underhållbar kod.

Vi har också sett en konkret implementation av den abstrakta klassen för
att hämta kattbilder från en specifik sida. Genom att använda olika
metoder som \texttt{getImages} och \texttt{getHtmlDocument}, kan vi
hämta och bearbeta bilder från webbsidan.

Det är viktigt att nämna att det finns många andra sätt att använda
abstrakta klasser och metoder i Java. Det är en kraftfull mekanism som
kan användas för att skapa flexibla och modulära applikationer. Jag
hoppas att du har lärt dig något nytt och att du kommer att använda
denna kunskap i dina egna projekt. Lycka till! 😺

\hypertarget{referenser}{%
\subsection{Referenser}\label{referenser}}

\begin{itemize}
\tightlist
\item
  \href{https://jsoup.org/}{JSoup}
\item
  \href{https://www.java.com/}{Java}
\item
  \href{https://en.wikipedia.org/wiki/Abstract_type}{Wikipedia}
\end{itemize}

\hypertarget{delegater}{%
\section{delegater}\label{delegater}}

\hypertarget{exempel}{%
\section{Exempel}\label{exempel}}

klass som innehåller en metod för att utföra en beräkning med hjälp av
delegater. \textbar{} \textbar{} Addera \textbar{} En metod som tar två
heltal som parametrar och returnerar deras summa. \textbar{} \textbar{}
Subtrahera \textbar{} En metod som tar två heltal som parametrar och
returnerar deras differens. \textbar{} \textbar{} Multiplicera\textbar{}
En metod som tar två heltal som parametrar och returnerar deras produkt.
\textbar{} \textbar{} Dividera \textbar{} En metod som tar två heltal
som parametrar och returnerar deras kvot. \textbar{}

En delegat i C\# är en typ som representerar referenser till metoder med
en specifik signatur. Delegater används ofta för att implementera
händelsehantering och callback-funktioner. En delegat kan användas för
att skapa en referens till en metod och sedan användas för att anropa
den metoden. Här är ett exempel på hur man använder en delegat i C\#:

\hypertarget{cb1}{}
\begin{Shaded}
\begin{Highlighting}[]
\CommentTok{// Definiera en delegat med namnet CalcParams som tar två heltal som parametrar och returnerar inget.}
\NormalTok{delegate }\DataTypeTok{void} \FunctionTok{CalcParams}\OperatorTok{(}\DataTypeTok{int}\NormalTok{ a}\OperatorTok{,} \DataTypeTok{int}\NormalTok{ b}\OperatorTok{);}

\KeywordTok{public} \KeywordTok{class}\NormalTok{ Calculator }\OperatorTok{\{}
\CommentTok{// Metod som utför en beräkning med hjälp av delegaten CalcParams.}
\KeywordTok{public} \DataTypeTok{void} \FunctionTok{PerformCalculation}\OperatorTok{(}\DataTypeTok{int}\NormalTok{ a}\OperatorTok{,} \DataTypeTok{int}\NormalTok{ b}\OperatorTok{,}\NormalTok{ CalcParams calculation}\OperatorTok{)} \OperatorTok{\{}
\FunctionTok{calculation}\OperatorTok{(}\NormalTok{a}\OperatorTok{,}\NormalTok{ b}\OperatorTok{);}

\KeywordTok{public} \DataTypeTok{static} \DataTypeTok{void} \FunctionTok{main}\OperatorTok{(}\BuiltInTok{String}\OperatorTok{[]}\NormalTok{ args}\OperatorTok{)} \OperatorTok{\{}
\CommentTok{// Skapa en instans av Calculator{-}klassen.}
\NormalTok{Calculator calculator }\OperatorTok{=} \KeywordTok{new} \FunctionTok{Calculator}\OperatorTok{();}
\CommentTok{// Skapa en instans av delegaten CalcParams som refererar till metoden Addera.}
\NormalTok{CalcParams addDelegate }\OperatorTok{=} \KeywordTok{new} \FunctionTok{CalcParams}\OperatorTok{(}\NormalTok{calculator}\OperatorTok{.}\FunctionTok{Addera}\OperatorTok{);}
\CommentTok{// Anropa PerformCalculation{-}metoden och skicka med delegaten Addera.}
\NormalTok{calculator}\OperatorTok{.}\FunctionTok{PerformCalculation}\OperatorTok{(}\DecValTok{5}\OperatorTok{,} \DecValTok{3}\OperatorTok{,}\NormalTok{ addDelegate}\OperatorTok{);}
\OperatorTok{\}}
\OperatorTok{\}}
\end{Highlighting}
\end{Shaded}

I det här exemplet definierar vi en delegat med namnet CalcParams som
tar två heltal som parametrar och returnerar inget. Sedan skapar vi en
instans av Calculator-klassen och en instans av delegaten CalcParams som
refererar till metoden Addera. Vi anropar sedan
PerformCalculation-metoden på instansen av Calculator-klassen och
skickar med delegaten Addera. När PerformCalculation-metoden anropas,
kommer den att utföra beräkningen genom att anropa delegaten med de
angivna parametrarna.

Metoden Addera adderar två heltal och skriver ut resultatet till
konsolen. Detta är ett exempel på hur man använder en delegat i C\#. I
Java kan man använda liknande koncept med hjälp av gränssnitt och
lambda-uttryck. Nedan följer ett exempel på hur man skulle implementera
samma funktionalitet i Java:

\hypertarget{cb2}{}
\begin{Shaded}
\begin{Highlighting}[]
\CommentTok{// Definiera ett gränssnitt med namnet CalcParams som har en metod som tar två heltal som parametrar och returnerar inget.}
\KeywordTok{interface}\NormalTok{ CalcParams }\OperatorTok{\{}
\DataTypeTok{void} \FunctionTok{calculate}\OperatorTok{(}\DataTypeTok{int}\NormalTok{ a}\OperatorTok{,} \DataTypeTok{int}\NormalTok{ b}\OperatorTok{);}

\KeywordTok{public} \KeywordTok{class}\NormalTok{ Calculator }\OperatorTok{\{}
\CommentTok{// Metod som utför en beräkning med hjälp av gränssnittet CalcParams.}
\KeywordTok{public} \DataTypeTok{void} \FunctionTok{performCalculation}\OperatorTok{(}\DataTypeTok{int}\NormalTok{ a}\OperatorTok{,} \DataTypeTok{int}\NormalTok{ b}\OperatorTok{,}\NormalTok{ CalcParams calculation}\OperatorTok{)} \OperatorTok{\{}
\NormalTok{calculation}\OperatorTok{.}\FunctionTok{calculate}\OperatorTok{(}\NormalTok{a}\OperatorTok{,}\NormalTok{ b}\OperatorTok{);}

\KeywordTok{public} \KeywordTok{class}\NormalTok{ Main }\OperatorTok{\{}
\KeywordTok{public} \DataTypeTok{static} \DataTypeTok{void} \FunctionTok{main}\OperatorTok{(}\BuiltInTok{String}\OperatorTok{[]}\NormalTok{ args}\OperatorTok{)} \OperatorTok{\{}
\CommentTok{// Skapa en lambda{-}uttryck som implementerar gränssnittet CalcParams och adderar två heltal.}
\NormalTok{CalcParams addLambda }\OperatorTok{=} \OperatorTok{(}\DataTypeTok{int}\NormalTok{ a}\OperatorTok{,} \DataTypeTok{int}\NormalTok{ b}\OperatorTok{)} \OperatorTok{{-}\textgreater{}} \OperatorTok{\{}
\DataTypeTok{int}\NormalTok{ sum }\OperatorTok{=}\NormalTok{ a }\OperatorTok{+}\NormalTok{ b}\OperatorTok{;}
\BuiltInTok{System}\OperatorTok{.}\FunctionTok{out}\OperatorTok{.}\FunctionTok{println}\OperatorTok{(}\StringTok{"Summan av "} \OperatorTok{+}\NormalTok{ a }\OperatorTok{+} \StringTok{" och "} \OperatorTok{+}\NormalTok{ b }\OperatorTok{+} \StringTok{" är "} \OperatorTok{+}\NormalTok{ sum}\OperatorTok{);}
\OperatorTok{\};}
\CommentTok{// Anropa performCalculation{-}metoden och skicka med lambda{-}uttrycket.}
\NormalTok{calculator}\OperatorTok{.}\FunctionTok{performCalculation}\OperatorTok{(}\DecValTok{5}\OperatorTok{,} \DecValTok{3}\OperatorTok{,}\NormalTok{ addLambda}\OperatorTok{);}
\OperatorTok{\}}
\OperatorTok{\}}
\end{Highlighting}
\end{Shaded}

I det här exemplet definierar vi ett gränssnitt med namnet CalcParams
som har en metod calculate som tar två heltal som parametrar och
returnerar inget. Sedan skapar vi en instans av Calculator-klassen och
en lambda-uttryck som implementerar gränssnittet CalcParams och adderar
två heltal. Vi anropar sedan performCalculation-metoden på instansen av
Calculator-klassen och skickar med lambda-uttrycket. När
performCalculation-metoden anropas, kommer den att utföra beräkningen
genom att anropa calculate-metoden på lambda-uttrycket med de angivna
parametrarna.

Detta är ett exempel på hur man kan implementera samma funktionalitet
med hjälp av gränssnitt och lambda-uttryck i Java.

\hypertarget{events}{%
\section{Events}\label{events}}

I denna översikt kommer vi att utforska konceptet händelser i C\#.
Händelser är en funktion som låter oss reagera på händelser som
inträffar under körningen av vårt program. \#\# Vad är händelser?

Händelser kan ses som en mekanism för att anropa metoder när något
specifikt händer i programmet. Istället för att använda if-satser eller
liknande för att kontrollera och reagera på olika scenarier kan vi
definiera händelser som triggas när en händelse uppstår. En händelse kan
vara något som att en knapp klickas, en fil ändras, en timer slår till
eller ett objekt uppdateras. Genom att använda händelser kan vi separera
logiken för att hantera händelsen från den del av koden som utlöser
händelsen. \#\# Användningsområden för händelser

Händelser kan vara till nytta i en mängd olika scenarier, inklusive: 1.
Grafiska användargränssnitt (GUI): I GUI-applikationer kan händelser
användas för att hantera händelser som knappklickar, musinmatning eller
fönsterfokusändringar. 2. Spelprogrammering: I spel kan händelser
användas för att hantera händelser som spelarens interaktion,
karaktärsdöd eller nivåuppgradering. 3. Multitrådad programmering: I
flertrådade applikationer kan händelser användas för att synkronisera
och kommunicera mellan olika trådar. 4. Kommunikation med externa
enheter: Händelser kan användas för att hantera händelser som inträffar
vid kommunikation med externa enheter eller system, t.ex. sensorer eller
nätverksanslutningar. \#\# Kodexempel

För att ge dig en bättre förståelse för hur händelser fungerar i C\# låt
oss titta på ett exempel:

\hypertarget{cb1}{}
\begin{Shaded}
\begin{Highlighting}[]
\KeywordTok{import} \ImportTok{java}\OperatorTok{.}\ImportTok{util}\OperatorTok{.}\ImportTok{ArrayList}\OperatorTok{;}
\KeywordTok{import} \ImportTok{java}\OperatorTok{.}\ImportTok{util}\OperatorTok{.}\ImportTok{List}\OperatorTok{;}
\CommentTok{// Definierar en händelsehanterare}
\KeywordTok{interface} \BuiltInTok{EventHandler} \OperatorTok{\{}
\DataTypeTok{void} \FunctionTok{handleEvent}\OperatorTok{(}\BuiltInTok{String}\NormalTok{ event}\OperatorTok{);}
\OperatorTok{\}}
\CommentTok{// Definierar en händelsekälla}
\KeywordTok{class}\NormalTok{ EventSource }\OperatorTok{\{}
\KeywordTok{private} \BuiltInTok{List}\OperatorTok{\textless{}}\BuiltInTok{EventHandler}\OperatorTok{\textgreater{}}\NormalTok{ eventHandlers }\OperatorTok{=} \KeywordTok{new} \BuiltInTok{ArrayList}\OperatorTok{\textless{}\textgreater{}();}
\CommentTok{// Registrerar en händelsehanterare}
\KeywordTok{public} \DataTypeTok{void} \FunctionTok{registerEventHandler}\OperatorTok{(}\BuiltInTok{EventHandler}\NormalTok{ eventHandler}\OperatorTok{)} \OperatorTok{\{}
\NormalTok{eventHandlers}\OperatorTok{.}\FunctionTok{add}\OperatorTok{(}\NormalTok{eventHandler}\OperatorTok{);}
\OperatorTok{\}}
\CommentTok{// Utlöser händelsen}
\KeywordTok{public} \DataTypeTok{void} \FunctionTok{triggerEvent}\OperatorTok{(}\BuiltInTok{String}\NormalTok{ event}\OperatorTok{)} \OperatorTok{\{}
\ControlFlowTok{for} \OperatorTok{(}\BuiltInTok{EventHandler}\NormalTok{ eventHandler }\OperatorTok{:}\NormalTok{ eventHandlers}\OperatorTok{)} \OperatorTok{\{}
\NormalTok{eventHandler}\OperatorTok{.}\FunctionTok{handleEvent}\OperatorTok{(}\NormalTok{event}\OperatorTok{);}
\OperatorTok{\}}
\CommentTok{// En exempelklass som implementerar EventHandler}
\KeywordTok{class}\NormalTok{ ExampleEventHandler }\KeywordTok{implements} \BuiltInTok{EventHandler} \OperatorTok{\{}
\AttributeTok{@Override}
\KeywordTok{public} \DataTypeTok{void} \FunctionTok{handleEvent}\OperatorTok{(}\BuiltInTok{String}\NormalTok{ event}\OperatorTok{)} \OperatorTok{\{}
\BuiltInTok{System}\OperatorTok{.}\FunctionTok{out}\OperatorTok{.}\FunctionTok{println}\OperatorTok{(}\StringTok{"Händelse mottagen: "} \OperatorTok{+}\NormalTok{ event}\OperatorTok{);}
\KeywordTok{public} \KeywordTok{class}\NormalTok{ Main }\OperatorTok{\{}
\KeywordTok{public} \DataTypeTok{static} \DataTypeTok{void} \FunctionTok{main}\OperatorTok{(}\BuiltInTok{String}\OperatorTok{[]}\NormalTok{ args}\OperatorTok{)} \OperatorTok{\{}
\CommentTok{// Skapar en händelsekälla}
\NormalTok{EventSource eventSource }\OperatorTok{=} \KeywordTok{new} \FunctionTok{EventSource}\OperatorTok{();}
\CommentTok{// Skapar en händelsehanterare}
\BuiltInTok{EventHandler}\NormalTok{ eventHandler }\OperatorTok{=} \KeywordTok{new} \FunctionTok{ExampleEventHandler}\OperatorTok{();}
\CommentTok{// Registrerar händelsehanteraren hos händelsekällan}
\NormalTok{eventSource}\OperatorTok{.}\FunctionTok{registerEventHandler}\OperatorTok{(}\NormalTok{eventHandler}\OperatorTok{);}
\CommentTok{// Utlöser en händelse}
\NormalTok{eventSource}\OperatorTok{.}\FunctionTok{triggerEvent}\OperatorTok{(}\StringTok{"En händelse har inträffat"}\OperatorTok{);}
\end{Highlighting}
\end{Shaded}

I detta exempel definierar vi en händelsehanterare genom att skapa ett
gränssnitt \texttt{EventHandler} som har en metod \texttt{handleEvent}.
Vi definierar också en händelsekälla \texttt{EventSource} som kan
registrera händelsehanterare och utlösa händelser. Vi skapar en
exempelklass \texttt{ExampleEventHandler} som implementerar
\texttt{EventHandler} och skriver ut en meddelande när en händelse
mottas. I \texttt{Main}-klassen skapar vi en instans av
\texttt{EventSource}, en instans av \texttt{ExampleEventHandler} och
registrerar händelsehanteraren hos händelsekällan. Sedan utlöser vi en
händelse genom att anropa \texttt{triggerEvent}-metoden på
händelsekällan. När händelsen utlöses kommer
\texttt{ExampleEventHandler} att reagera genom att skriva ut ett
meddelande. Detta är bara ett grundläggande exempel på hur händelser kan
användas i Java. Det finns många andra möjligheter och
användningsområden för händelser i Java-programmering.

`````\texttt{java\ import\ java.util.*;\ public\ class\ Hero\ \{\ private\ List\textless{}LevelUpListener\textgreater{}\ levelUpListeners\ =\ new\ ArrayList\textless{}\textgreater{}();\ private\ List\textless{}DiedListener\textgreater{}\ diedListeners\ =\ new\ ArrayList\textless{}\textgreater{}();\ private\ int\ level\ =\ 1;\ private\ int\ XP\ =\ 0;\ public\ int\ getLevel()\ \{\ return\ level;\ public\ int\ getXP()\ \{\ return\ XP;\ public\ int\ getMaxXP()\ \{\ return\ level\ *\ 100;\ public\ void\ increaseXP(int\ amount)\ \{\ XP\ +=\ amount;\ if\ (XP\ \textgreater{}=\ getMaxXP())\ \{\ level++;\ onLevelUp();\ public\ void\ die()\ \{\ onDied();\ public\ void\ addLevelUpListener(LevelUpListener\ listener)\ \{\ levelUpListeners.add(listener);\ public\ void\ addDiedListener(DiedListener\ listener)\ \{\ diedListeners.add(listener);\ private\ void\ onLevelUp()\ \{\ for\ (LevelUpListener\ listener\ :\ levelUpListeners)\ \{\ listener.onLevelUp();\ private\ void\ onDied()\ \{\ for\ (DiedListener\ listener\ :\ diedListeners)\ \{\ listener.onDied();\ public\ interface\ LevelUpListener\ \{\ void\ onLevelUp();\ public\ interface\ DiedListener\ \{\ void\ onDied();\ public\ class\ Program\ \{\ Hero\ hero\ =\ new\ Hero();\ hero.addLevelUpListener(new\ LevelUpListener()\ \{\ @Override\ public\ void\ onLevelUp()\ \{\ System.out.println("Hjälten\ har\ nått\ en\ ny\ nivå!");\ \}\ \});\ hero.addDiedListener(new\ DiedListener()\ \{\ public\ void\ onDied()\ \{\ System.out.println("Hjälten\ har\ dött!");\ hero.increaseXP(100);\ hero.die();\ I\ detta\ exempel\ har\ vi\ en\ klass\ som\ heter}Hero\texttt{,\ som\ representerar\ en\ spelkaraktär.\ Klassen\ har\ två\ listor,}levelUpListeners\texttt{och}diedListeners\texttt{,\ som\ innehåller\ lyssnare\ för\ händelserna}LevelUp\texttt{och}Died\texttt{.\ När\ hjältens\ erfarenhetspoäng\ (}XP\texttt{)\ ökar,\ kontrolleras\ om\ XP\ överstiger\ det\ maximala\ värdet\ (}getMaxXP()\texttt{).\ Om\ så\ är\ fallet,\ ökar\ hjältens\ nivå\ och}onLevelUp()\texttt{-metoden\ anropas\ för\ varje\ lyssnare\ i}levelUpListeners\texttt{.\ På\ samma\ sätt,\ när\ hjältens\ hälsopoäng\ (}HP\texttt{)\ når\ 0\ eller\ mindre,\ anropas}onDied()\texttt{-metoden\ för\ varje\ lyssnare\ i}diedListeners\texttt{.\ I}Program\texttt{-klassen\ skapar\ vi\ en\ instans\ av}Hero\texttt{och\ lägger\ till\ anonyma\ klasser\ som\ implementerar}LevelUpListener\texttt{och}DiedListener`
för att definiera beteendet för varje händelse. När vi sedan ökar
hjältens XP eller simulerar dess död, kommer de tillhörande metoderna
att anropas och en meddelandetext skrivs ut i konsolen. Detta är bara
ett grundläggande exempel för att illustrera hur händelselyssnare kan
användas i Java. I praktiken kan de vara mycket kraftfulla och användas
för att skapa interaktion och dynamik i programmet. \#\# Summering

Händelselyssnare i Java möjliggör reaktioner på händelser under
programkörningen. Genom att separera händelsehanteringen från den del av
koden som utlöser händelsen, kan vi skapa mer modulära och underhållbara
program. Händelselyssnare kan användas i olika scenarier, inklusive
GUI-programmering, spelutveckling och multitrådad programmering. Genom
att följa principerna om ren kod (Clean Code) kan vi skriva läsbar och
effektiv kod som är lätt att förstå och underhålla.

\hypertarget{exempel}{%
\section{Exempel}\label{exempel}}

Vi ska skapa ett bankkonto där vi kan sätta in pengar och ta ut pengar.
Ett bankkonto är en vanlig komponent i finansiella system och används
för att hantera insättningar, uttag och saldo. \#\# Beskrivning

Vårt bankkonto kommer att ha flera events som informerar oss om olika
händelser som inträffar på kontot. Dessa events inkluderar: -
Insättning: När pengar sätts in på kontot. - Uttag: När pengar tas ut
från kontot. - Saldo mindre än 0: När saldot på kontot blir mindre än 0.
- Insättning med felaktig summa: När en insättning görs med en ogiltig
summa (t.ex. 0 eller negativt belopp). - Uttag med felaktig summa: När
ett uttag görs med en ogiltig summa. Vi kommer också att ha en särskild
händelse för insättningar över en viss gräns, till exempel 15000 kr.
Denna händelse kan användas för att varna för eventuell kriminell
aktivitet. För att implementera detta bankkonto använder vi Java och
skapar en klass som heter ``Account''. Klassen har en variabel
``balance'' för att hålla reda på kontots saldo, samt olika events för
att rapportera händelser till intresserade lyssnare. Här är Java-koden
som implementerar bankkontot:

\hypertarget{cb1}{}
\begin{Shaded}
\begin{Highlighting}[]
\KeywordTok{import} \ImportTok{java}\OperatorTok{.}\ImportTok{util}\OperatorTok{.}\ImportTok{ArrayList}\OperatorTok{;}
\KeywordTok{import} \ImportTok{java}\OperatorTok{.}\ImportTok{util}\OperatorTok{.}\ImportTok{List}\OperatorTok{;}
\KeywordTok{public} \KeywordTok{class}\NormalTok{ Account }\OperatorTok{\{}
\KeywordTok{private} \DataTypeTok{int}\NormalTok{ balance}\OperatorTok{;}
\KeywordTok{private} \BuiltInTok{List}\OperatorTok{\textless{}}\NormalTok{DepositListener}\OperatorTok{\textgreater{}}\NormalTok{ depositListeners}\OperatorTok{;}
\KeywordTok{private} \BuiltInTok{List}\OperatorTok{\textless{}}\NormalTok{WithdrawListener}\OperatorTok{\textgreater{}}\NormalTok{ withdrawListeners}\OperatorTok{;}
\KeywordTok{private} \BuiltInTok{List}\OperatorTok{\textless{}}\NormalTok{BalanceBelowZeroListener}\OperatorTok{\textgreater{}}\NormalTok{ balanceBelowZeroListeners}\OperatorTok{;}
\KeywordTok{private} \BuiltInTok{List}\OperatorTok{\textless{}}\NormalTok{DepositInvalidAmountListener}\OperatorTok{\textgreater{}}\NormalTok{ depositInvalidAmountListeners}\OperatorTok{;}
\KeywordTok{private} \BuiltInTok{List}\OperatorTok{\textless{}}\NormalTok{WithdrawInvalidAmountListener}\OperatorTok{\textgreater{}}\NormalTok{ withdrawInvalidAmountListeners}\OperatorTok{;}
\KeywordTok{private} \BuiltInTok{List}\OperatorTok{\textless{}}\NormalTok{DepositAboveLimitListener}\OperatorTok{\textgreater{}}\NormalTok{ depositAboveLimitListeners}\OperatorTok{;}
\KeywordTok{public} \FunctionTok{Account}\OperatorTok{()} \OperatorTok{\{}
\NormalTok{balance }\OperatorTok{=} \DecValTok{0}\OperatorTok{;}
\NormalTok{depositListeners }\OperatorTok{=} \KeywordTok{new} \BuiltInTok{ArrayList}\OperatorTok{\textless{}\textgreater{}();}
\NormalTok{withdrawListeners }\OperatorTok{=} \KeywordTok{new} \BuiltInTok{ArrayList}\OperatorTok{\textless{}\textgreater{}();}
\NormalTok{balanceBelowZeroListeners }\OperatorTok{=} \KeywordTok{new} \BuiltInTok{ArrayList}\OperatorTok{\textless{}\textgreater{}();}
\NormalTok{depositInvalidAmountListeners }\OperatorTok{=} \KeywordTok{new} \BuiltInTok{ArrayList}\OperatorTok{\textless{}\textgreater{}();}
\NormalTok{withdrawInvalidAmountListeners }\OperatorTok{=} \KeywordTok{new} \BuiltInTok{ArrayList}\OperatorTok{\textless{}\textgreater{}();}
\NormalTok{depositAboveLimitListeners }\OperatorTok{=} \KeywordTok{new} \BuiltInTok{ArrayList}\OperatorTok{\textless{}\textgreater{}();}
\OperatorTok{\}}
\KeywordTok{public} \DataTypeTok{void} \FunctionTok{addDepositListener}\OperatorTok{(}\NormalTok{DepositListener listener}\OperatorTok{)} \OperatorTok{\{}
\NormalTok{depositListeners}\OperatorTok{.}\FunctionTok{add}\OperatorTok{(}\NormalTok{listener}\OperatorTok{);}
\KeywordTok{public} \DataTypeTok{void} \FunctionTok{addWithdrawListener}\OperatorTok{(}\NormalTok{WithdrawListener listener}\OperatorTok{)} \OperatorTok{\{}
\NormalTok{withdrawListeners}\OperatorTok{.}\FunctionTok{add}\OperatorTok{(}\NormalTok{listener}\OperatorTok{);}
\KeywordTok{public} \DataTypeTok{void} \FunctionTok{addBalanceBelowZeroListener}\OperatorTok{(}\NormalTok{BalanceBelowZeroListener listener}\OperatorTok{)} \OperatorTok{\{}
\NormalTok{balanceBelowZeroListeners}\OperatorTok{.}\FunctionTok{add}\OperatorTok{(}\NormalTok{listener}\OperatorTok{);}
\KeywordTok{public} \DataTypeTok{void} \FunctionTok{addDepositInvalidAmountListener}\OperatorTok{(}\NormalTok{DepositInvalidAmountListener listener}\OperatorTok{)} \OperatorTok{\{}
\NormalTok{depositInvalidAmountListeners}\OperatorTok{.}\FunctionTok{add}\OperatorTok{(}\NormalTok{listener}\OperatorTok{);}
\KeywordTok{public} \DataTypeTok{void} \FunctionTok{addWithdrawInvalidAmountListener}\OperatorTok{(}\NormalTok{WithdrawInvalidAmountListener listener}\OperatorTok{)} \OperatorTok{\{}
\NormalTok{withdrawInvalidAmountListeners}\OperatorTok{.}\FunctionTok{add}\OperatorTok{(}\NormalTok{listener}\OperatorTok{);}
\KeywordTok{public} \DataTypeTok{void} \FunctionTok{addDepositAboveLimitListener}\OperatorTok{(}\NormalTok{DepositAboveLimitListener listener}\OperatorTok{)} \OperatorTok{\{}
\NormalTok{depositAboveLimitListeners}\OperatorTok{.}\FunctionTok{add}\OperatorTok{(}\NormalTok{listener}\OperatorTok{);}
\KeywordTok{public} \DataTypeTok{int} \FunctionTok{getBalance}\OperatorTok{()} \OperatorTok{\{}
\ControlFlowTok{return}\NormalTok{ balance}\OperatorTok{;}
\KeywordTok{public} \DataTypeTok{void} \FunctionTok{deposit}\OperatorTok{(}\DataTypeTok{int}\NormalTok{ amount}\OperatorTok{)} \OperatorTok{\{}
\ControlFlowTok{if} \OperatorTok{(}\NormalTok{amount }\OperatorTok{\textless{}=} \DecValTok{0}\OperatorTok{)} \OperatorTok{\{}
\ControlFlowTok{for} \OperatorTok{(}\NormalTok{DepositInvalidAmountListener listener }\OperatorTok{:}\NormalTok{ depositInvalidAmountListeners}\OperatorTok{)} \OperatorTok{\{}
\NormalTok{listener}\OperatorTok{.}\FunctionTok{onDepositInvalidAmount}\OperatorTok{(}\NormalTok{amount}\OperatorTok{);}
\OperatorTok{\}}
\ControlFlowTok{return}\OperatorTok{;}
\OperatorTok{\}}
\ControlFlowTok{if} \OperatorTok{(}\NormalTok{amount }\OperatorTok{\textgreater{}} \DecValTok{15000}\OperatorTok{)} \OperatorTok{\{}
\ControlFlowTok{for} \OperatorTok{(}\NormalTok{DepositAboveLimitListener listener }\OperatorTok{:}\NormalTok{ depositAboveLimitListeners}\OperatorTok{)} \OperatorTok{\{}
\NormalTok{listener}\OperatorTok{.}\FunctionTok{onDepositAboveLimit}\OperatorTok{(}\NormalTok{amount}\OperatorTok{);}
\NormalTok{balance }\OperatorTok{+=}\NormalTok{ amount}\OperatorTok{;}
\ControlFlowTok{for} \OperatorTok{(}\NormalTok{DepositListener listener }\OperatorTok{:}\NormalTok{ depositListeners}\OperatorTok{)} \OperatorTok{\{}
\NormalTok{listener}\OperatorTok{.}\FunctionTok{onDeposit}\OperatorTok{(}\NormalTok{amount}\OperatorTok{);}
\KeywordTok{public} \DataTypeTok{void} \FunctionTok{withdraw}\OperatorTok{(}\DataTypeTok{int}\NormalTok{ amount}\OperatorTok{)} \OperatorTok{\{}
\ControlFlowTok{for} \OperatorTok{(}\NormalTok{WithdrawInvalidAmountListener listener }\OperatorTok{:}\NormalTok{ withdrawInvalidAmountListeners}\OperatorTok{)} \OperatorTok{\{}
\NormalTok{listener}\OperatorTok{.}\FunctionTok{onWithdrawInvalidAmount}\OperatorTok{(}\NormalTok{amount}\OperatorTok{);}
\ControlFlowTok{if} \OperatorTok{(}\NormalTok{balance }\OperatorTok{{-}}\NormalTok{ amount }\OperatorTok{\textless{}} \DecValTok{0}\OperatorTok{)} \OperatorTok{\{}
\ControlFlowTok{for} \OperatorTok{(}\NormalTok{BalanceBelowZeroListener listener }\OperatorTok{:}\NormalTok{ balanceBelowZeroListeners}\OperatorTok{)} \OperatorTok{\{}
\NormalTok{listener}\OperatorTok{.}\FunctionTok{onBalanceBelowZero}\OperatorTok{(}\NormalTok{amount}\OperatorTok{);}
\NormalTok{balance }\OperatorTok{{-}=}\NormalTok{ amount}\OperatorTok{;}
\ControlFlowTok{for} \OperatorTok{(}\NormalTok{WithdrawListener listener }\OperatorTok{:}\NormalTok{ withdrawListeners}\OperatorTok{)} \OperatorTok{\{}
\NormalTok{listener}\OperatorTok{.}\FunctionTok{onWithdraw}\OperatorTok{(}\NormalTok{amount}\OperatorTok{);}
\OperatorTok{\}}
\KeywordTok{public} \KeywordTok{interface}\NormalTok{ DepositListener }\OperatorTok{\{}
\DataTypeTok{void} \FunctionTok{onDeposit}\OperatorTok{(}\DataTypeTok{int}\NormalTok{ amount}\OperatorTok{);}
\KeywordTok{public} \KeywordTok{interface}\NormalTok{ WithdrawListener }\OperatorTok{\{}
\DataTypeTok{void} \FunctionTok{onWithdraw}\OperatorTok{(}\DataTypeTok{int}\NormalTok{ amount}\OperatorTok{);}
\KeywordTok{public} \KeywordTok{interface}\NormalTok{ BalanceBelowZeroListener }\OperatorTok{\{}
\DataTypeTok{void} \FunctionTok{onBalanceBelowZero}\OperatorTok{(}\DataTypeTok{int}\NormalTok{ amount}\OperatorTok{);}
\KeywordTok{public} \KeywordTok{interface}\NormalTok{ DepositInvalidAmountListener }\OperatorTok{\{}
\DataTypeTok{void} \FunctionTok{onDepositInvalidAmount}\OperatorTok{(}\DataTypeTok{int}\NormalTok{ amount}\OperatorTok{);}
\KeywordTok{public} \KeywordTok{interface}\NormalTok{ WithdrawInvalidAmountListener }\OperatorTok{\{}
\DataTypeTok{void} \FunctionTok{onWithdrawInvalidAmount}\OperatorTok{(}\DataTypeTok{int}\NormalTok{ amount}\OperatorTok{);}
\KeywordTok{public} \KeywordTok{interface}\NormalTok{ DepositAboveLimitListener }\OperatorTok{\{}
\DataTypeTok{void} \FunctionTok{onDepositAboveLimit}\OperatorTok{(}\DataTypeTok{int}\NormalTok{ amount}\OperatorTok{);}
\end{Highlighting}
\end{Shaded}

public class DepositEventArgs \{ private int amount; public
DepositEventArgs(int amount) \{ this.amount = amount; public int
getAmount() \{ return amount; public class WithdrawEventArgs \{ public
WithdrawEventArgs(int amount) \{ public class BalanceBelowZeroEventArgs
\{ public BalanceBelowZeroEventArgs(int amount) \{ public class
DepositInvalidAmountEventArgs \{ public
DepositInvalidAmountEventArgs(int amount) \{ public class
WithdrawInvalidAmountEventArgs \{ public
WithdrawInvalidAmountEventArgs(int amount) \{ public class
DepositAboveLimitEventArgs \{ public DepositAboveLimitEventArgs(int
amount) \{ Nu kan vi använda vår ``Account''-klass för att skapa
bankkonton och hantera händelserna som inträffar på dem.\# Java import
java.util.EventListener; if (amount \textgreater{} 0) \{ balance +=
amount; \} else \{ DepositInvalidAmountEventArgs args = new
DepositInvalidAmountEventArgs(amount); onDepositInvalidAmount(args); if
(amount \textgreater{} 0 \&\& amount \textless= balance) \{ balance -=
amount; WithdrawInvalidAmountEventArgs args = new
WithdrawInvalidAmountEventArgs(amount); onWithdrawInvalidAmount(args);
if (balance \textless{} 0) \{ BalanceBelowZeroEventArgs args = new
BalanceBelowZeroEventArgs(balance); onBalanceBelowZero(args); public
event EventHandler Deposit; protected virtual void
OnDeposit(DepositEventArgs e) \{ Deposit?.Invoke(this, e); public event
EventHandler Withdraw; protected virtual void
OnWithdraw(WithdrawEventArgs e) Withdraw?.Invoke(this, e); public event
EventHandler BalanceBelowZero; protected virtual void
OnBalanceBelowZero(BalanceBelowZeroEventArgs e)
BalanceBelowZero?.Invoke(this, e); public event EventHandler
DepositInvalidAmount; protected virtual void
OnDepositInvalidAmount(DepositInvalidAmountEventArgs e)
DepositInvalidAmount?.Invoke(this, e); public event EventHandler
WithdrawInvalidAmount; protected virtual void
OnWithdrawInvalidAmount(WithdrawInvalidAmountEventArgs e)
WithdrawInvalidAmount?.Invoke(this, e); public event EventHandler
DepositAboveLimit; protected virtual void
OnDepositAboveLimit(DepositAboveLimitEventArgs e)
DepositAboveLimit?.Invoke(this, e); public class DepositEventArgs
extends EventArgs \{ public class WithdrawEventArgs extends EventArgs \{
public class BalanceBelowZeroEventArgs extends EventArgs \{ public class
DepositInvalidAmountEventArgs extends EventArgs \{ public class
WithdrawInvalidAmountEventArgs extends EventArgs \{ public class
DepositAboveLimitEventArgs extends EventArgs \{ public class Program \{
public static void main(String{[}{]} args) \{ Account account = new
Account(); account.addDepositListener(new DepositListener() \{
{@Override} public void onDeposit(DepositEventArgs e) \{
System.out.println(``Insättning:'' + e.getAmount()); \});
account.addWithdrawListener(new WithdrawListener() \{ public void
onWithdraw(WithdrawEventArgs e) \{ System.out.println(``Uttag:'' +
e.getAmount()); account.addBalanceBelowZeroListener(new
BalanceBelowZeroListener() \{ public void
onBalanceBelowZero(BalanceBelowZeroEventArgs e) \{
System.out.println(``Saldo under noll:'' + e.getAmount());
account.addDepositInvalidAmountListener(new
DepositInvalidAmountListener() \{ public void
onDepositInvalidAmount(DepositInvalidAmountEventArgs e) \{
System.out.println(``Ogiltig insättningsbelopp:'' + e.getAmount());
account.addWithdrawInvalidAmountListener(new
WithdrawInvalidAmountListener() \{ public void
onWithdrawInvalidAmount(WithdrawInvalidAmountEventArgs e) \{
System.out.println(``Ogiltig uttagsbelopp:'' + e.getAmount());
account.addDepositAboveLimitListener(new DepositAboveLimitListener() \{
public void onDepositAboveLimit(DepositAboveLimitEventArgs e) \{
System.out.println(``Insättning över gränsen:'' + e.getAmount());
account.deposit(1000); account.withdraw(500); account.withdraw(600);
account.deposit(-100); public interface DepositListener extends
EventListener \{ void onDeposit(DepositEventArgs e); public interface
WithdrawListener extends EventListener \{ void
onWithdraw(WithdrawEventArgs e); public interface
BalanceBelowZeroListener extends EventListener \{ void
onBalanceBelowZero(BalanceBelowZeroEventArgs e); public interface
DepositInvalidAmountListener extends EventListener \{ void
onDepositInvalidAmount(DepositInvalidAmountEventArgs e); public
interface WithdrawInvalidAmountListener extends EventListener \{ void
onWithdrawInvalidAmount(WithdrawInvalidAmountEventArgs e); public
interface DepositAboveLimitListener extends EventListener \{ void
onDepositAboveLimit(DepositAboveLimitEventArgs e);

```I det här exemplet skapas en instans av klassen ``Account'' och olika
händelselyssnare kopplas till dess events. Därefter utförs några
operationer på kontot, som insättningar och uttag, vilket resulterar i
att relevanta events triggas och meddelanden skrivs ut till konsolen.
Detta är bara en grundläggande implementation av ett bankkonto i Java.
Beroende på användningsområdet kan det finnas ytterligare funktioner och
logik som behöver läggas till. public class Main \{ public void
onDepositAboveLimit(DepositAboveLimitEvent event) \{
System.out.println(``Insättning på'' + event.getAmount() + '' kr är över
gränsen''); System.out.println(``Ring SÄPO!''); public void
onWithdrawInvalidAmount(WithdrawInvalidAmountEvent event) \{
System.out.println(``Uttag på'' + event.getAmount() + '' kr är inte
tillåtet''); public void
onDepositInvalidAmount(DepositInvalidAmountEvent event) \{
System.out.println(``Insättning på'' + event.getAmount() + '' kr är inte
tillåtet''); public void onBalanceBelowZero(BalanceBelowZeroEvent event)
\{ System.out.println(``Uttag på'' + event.getAmount() + '' kr är inte
tillåtet då det skulle sätta saldo under 0''); public void
onWithdraw(WithdrawEvent event) \{ System.out.println(``Uttag på'' +
event.getAmount() + '' kr''); public void onDeposit(DepositEvent event)
\{ System.out.println(``Insättning på'' + event.getAmount() + '' kr'');
account.withdraw(-100); account.deposit(20000); void
onDepositAboveLimit(DepositAboveLimitEvent event); public class
DepositAboveLimitEvent \{ private double amount; public
DepositAboveLimitEvent(double amount) \{ public double getAmount() \{
void onWithdrawInvalidAmount(WithdrawInvalidAmountEvent event); public
class WithdrawInvalidAmountEvent \{ public
WithdrawInvalidAmountEvent(double amount) \{ void
onDepositInvalidAmount(DepositInvalidAmountEvent event); public class
DepositInvalidAmountEvent \{ public DepositInvalidAmountEvent(double
amount) \{ void onBalanceBelowZero(BalanceBelowZeroEvent event); public
class BalanceBelowZeroEvent \{ public BalanceBelowZeroEvent(double
amount) \{ void onWithdraw(WithdrawEvent event); public class
WithdrawEvent \{ public WithdrawEvent(double amount) \{ void
onDeposit(DepositEvent event); public class DepositEvent \{ public
DepositEvent(double amount) \{ public void withdraw(double amount) \{ if
(amount \textless{} 0) \{ System.out.println(``Uttag på'' + amount + ''
kr är inte tillåtet''); WithdrawEvent event = new WithdrawEvent(amount);
listener.onWithdraw(event); if (amount \textgreater{} 10000) \{
listener.onDepositAboveLimit(new DepositAboveLimitEvent(amount)); public
void deposit(double amount) \{ System.out.println(``Insättning på'' +
amount + '' kr är inte tillåtet''); DepositEvent event = new
DepositEvent(amount); listener.onDeposit(event); public

\hypertarget{filhantering}{%
\section{Filhantering}\label{filhantering}}

Vi ska skapa en klass och spara den i hårddisken, sedan ska vi läsa in
den.

\hypertarget{skapa-en-klass}{%
\subsection{Skapa en klass}\label{skapa-en-klass}}

\hypertarget{cb1}{}
\begin{Shaded}
\begin{Highlighting}[]
\KeywordTok{class}\NormalTok{ Person}\OperatorTok{\{}
    \BuiltInTok{String}\NormalTok{ namn}\OperatorTok{;}
    \DataTypeTok{int}\NormalTok{ ålder}\OperatorTok{;}

    \KeywordTok{public} \FunctionTok{Person}\OperatorTok{(}\BuiltInTok{String}\NormalTok{ namn}\OperatorTok{,} \DataTypeTok{int}\NormalTok{ ålder}\OperatorTok{)\{}
        \KeywordTok{this}\OperatorTok{.}\FunctionTok{namn} \OperatorTok{=}\NormalTok{ namn}\OperatorTok{;}
        \KeywordTok{this}\OperatorTok{.}\NormalTok{ålder }\OperatorTok{=}\NormalTok{ ålder}\OperatorTok{;}
    \OperatorTok{\}}

    \KeywordTok{public} \BuiltInTok{String} \FunctionTok{toString}\OperatorTok{()\{}
        \ControlFlowTok{return} \StringTok{"Namn: "} \OperatorTok{+}\NormalTok{ namn }\OperatorTok{+} \StringTok{" Ålder:"} \OperatorTok{+}\NormalTok{ ålder}\OperatorTok{;}
    \OperatorTok{\}}

    \KeywordTok{public} \BuiltInTok{String} \FunctionTok{getNamn}\OperatorTok{()\{}
        \ControlFlowTok{return}\NormalTok{ namn}\OperatorTok{;}
    \OperatorTok{\}}

    \KeywordTok{public} \DataTypeTok{int} \FunctionTok{getÅlder}\OperatorTok{()\{}
        \ControlFlowTok{return}\NormalTok{ ålder}\OperatorTok{;}
    \OperatorTok{\}}

    \KeywordTok{public} \DataTypeTok{void} \FunctionTok{setNamn}\OperatorTok{(}\BuiltInTok{String}\NormalTok{ namn}\OperatorTok{)\{}
        \KeywordTok{this}\OperatorTok{.}\FunctionTok{namn} \OperatorTok{=}\NormalTok{ namn}\OperatorTok{;}
    \OperatorTok{\}}

    \KeywordTok{public} \DataTypeTok{void} \FunctionTok{setÅlder}\OperatorTok{(}\DataTypeTok{int}\NormalTok{ ålder}\OperatorTok{)\{}
        \KeywordTok{this}\OperatorTok{.}\NormalTok{ålder }\OperatorTok{=}\NormalTok{ ålder}\OperatorTok{;}
    \OperatorTok{\}}
\OperatorTok{\}}
\end{Highlighting}
\end{Shaded}

\hypertarget{luxe4sa-en-textfil}{%
\subsection{Läsa en textfil}\label{luxe4sa-en-textfil}}

För att läsa en textfil, använder vi \texttt{FileReader}-klassen. Vi
skapar ett nytt objekt av klassen och ger den filnamnet som argument.
Sedan kan vi använda \texttt{readLine()}-metoden för att läsa in rader
från filen.

\hypertarget{cb2}{}
\begin{Shaded}
\begin{Highlighting}[]
\BuiltInTok{FileReader}\NormalTok{ fr }\OperatorTok{=} \KeywordTok{new} \BuiltInTok{FileReader}\OperatorTok{(}\StringTok{"filnamn.txt"}\OperatorTok{);}
\BuiltInTok{String}\NormalTok{ rad }\OperatorTok{=}\NormalTok{ fr}\OperatorTok{.}\FunctionTok{readLine}\OperatorTok{();}
\ControlFlowTok{while}\OperatorTok{(}\NormalTok{rad }\OperatorTok{!=} \KeywordTok{null}\OperatorTok{)\{}
    \BuiltInTok{System}\OperatorTok{.}\FunctionTok{out}\OperatorTok{.}\FunctionTok{println}\OperatorTok{(}\NormalTok{rad}\OperatorTok{);}
\NormalTok{    rad }\OperatorTok{=}\NormalTok{ fr}\OperatorTok{.}\FunctionTok{readLine}\OperatorTok{();}
\OperatorTok{\}}
\end{Highlighting}
\end{Shaded}

\hypertarget{skriva-till-en-textfil}{%
\subsection{Skriva till en textfil}\label{skriva-till-en-textfil}}

För att skriva till en textfil, använder vi \texttt{FileWriter}-klassen.
Vi skapar ett nytt objekt av klassen och ger den filnamnet som argument.
Sedan kan vi använda \texttt{write()}-metoden för att skriva till filen.

\hypertarget{cb3}{}
\begin{Shaded}
\begin{Highlighting}[]
\CommentTok{// Skapa en person}
\NormalTok{Person p }\OperatorTok{=} \KeywordTok{new} \FunctionTok{Person}\OperatorTok{(}\StringTok{"Luke Skywalker"}\OperatorTok{,} \DecValTok{23}\OperatorTok{);}

\CommentTok{// Serialisera personen}
\BuiltInTok{String}\NormalTok{ person }\OperatorTok{=}\NormalTok{ p}\OperatorTok{.}\FunctionTok{toString}\OperatorTok{();}

\CommentTok{// Skapa en ny fil och skriv personen till den}
\BuiltInTok{FileWriter}\NormalTok{ fw }\OperatorTok{=} \KeywordTok{new} \BuiltInTok{FileWriter}\OperatorTok{(}\StringTok{"person.txt"}\OperatorTok{);}
\NormalTok{fw}\OperatorTok{.}\FunctionTok{write}\OperatorTok{(}\NormalTok{person}\OperatorTok{);}
\NormalTok{fw}\OperatorTok{.}\FunctionTok{close}\OperatorTok{();}

\CommentTok{// Läs in personen från filen}
\BuiltInTok{FileReader}\NormalTok{ fr }\OperatorTok{=} \KeywordTok{new} \BuiltInTok{FileReader}\OperatorTok{(}\StringTok{"person.txt"}\OperatorTok{);}
\BuiltInTok{String}\NormalTok{ rad }\OperatorTok{=}\NormalTok{ fr}\OperatorTok{.}\FunctionTok{readLine}\OperatorTok{();}
\NormalTok{ft}\OperatorTok{.}\FunctionTok{close}\OperatorTok{();}

\CommentTok{// Skapa en ny person från strängen}
\CommentTok{// genom att dela upp den på Namn: och Ålder:}
\BuiltInTok{String}\OperatorTok{[]}\NormalTok{ delar }\OperatorTok{=}\NormalTok{ rad}\OperatorTok{.}\FunctionTok{split}\OperatorTok{(}\StringTok{"Namn: | Ålder:"}\OperatorTok{);}
\NormalTok{Person p2 }\OperatorTok{=} \KeywordTok{new} \FunctionTok{Person}\OperatorTok{(}\NormalTok{delar}\OperatorTok{[}\DecValTok{1}\OperatorTok{],} \BuiltInTok{Integer}\OperatorTok{.}\FunctionTok{parseInt}\OperatorTok{(}\NormalTok{delar}\OperatorTok{[}\DecValTok{2}\OperatorTok{]));}
\BuiltInTok{System}\OperatorTok{.}\FunctionTok{out}\OperatorTok{.}\FunctionTok{println}\OperatorTok{(}\NormalTok{p2}\OperatorTok{);}
\end{Highlighting}
\end{Shaded}

\hypertarget{luxe4sa-och-skriva-till-en-json-fil}{%
\subsection{Läsa och skriva till en
Json-fil}\label{luxe4sa-och-skriva-till-en-json-fil}}

För att läsa och skriva till en Json-fil, använder vi
\texttt{Gson}-klassen. Vi skapar ett nytt objekt av klassen och ger den
filnamnet som argument. Sedan kan vi använda \texttt{toJson()}-metoden
för att skriva till filen och \texttt{fromJson()}-metoden för att läsa
från filen.

Lägg till Gson-biblioteket i Maven:

\hypertarget{cb4}{}
\begin{Shaded}
\begin{Highlighting}[]
\NormalTok{\textless{}}\KeywordTok{dependency}\NormalTok{\textgreater{}}
\NormalTok{    \textless{}}\KeywordTok{groupId}\NormalTok{\textgreater{}com.google.code.gson\textless{}/}\KeywordTok{groupId}\NormalTok{\textgreater{}}
\NormalTok{    \textless{}}\KeywordTok{artifactId}\NormalTok{\textgreater{}gson\textless{}/}\KeywordTok{artifactId}\NormalTok{\textgreater{}}
\NormalTok{    \textless{}}\KeywordTok{version}\NormalTok{\textgreater{}2.8.9\textless{}/}\KeywordTok{version}\NormalTok{\textgreater{}}
\NormalTok{\textless{}/}\KeywordTok{dependency}\NormalTok{\textgreater{}}
\end{Highlighting}
\end{Shaded}

\hypertarget{cb5}{}
\begin{Shaded}
\begin{Highlighting}[]
\CommentTok{// Skapa en person}
\NormalTok{Person p }\OperatorTok{=} \KeywordTok{new} \FunctionTok{Person}\OperatorTok{(}\StringTok{"Leia Skywalker"}\OperatorTok{,} \DecValTok{23}\OperatorTok{);}

\CommentTok{// Serialisera personen}
\NormalTok{Gson gson }\OperatorTok{=} \KeywordTok{new} \FunctionTok{Gson}\OperatorTok{();}
\BuiltInTok{String}\NormalTok{ person }\OperatorTok{=}\NormalTok{ gson}\OperatorTok{.}\FunctionTok{toJson}\OperatorTok{(}\NormalTok{p}\OperatorTok{);}

\CommentTok{// Skapa en ny fil och skriv personen till den}
\BuiltInTok{FileWriter}\NormalTok{ fw }\OperatorTok{=} \KeywordTok{new} \BuiltInTok{FileWriter}\OperatorTok{(}\StringTok{"person.json"}\OperatorTok{);}
\NormalTok{fw}\OperatorTok{.}\FunctionTok{write}\OperatorTok{(}\NormalTok{person}\OperatorTok{);}
\NormalTok{fw}\OperatorTok{.}\FunctionTok{close}\OperatorTok{();}

\CommentTok{// Läs in personen från filen}
\BuiltInTok{FileReader}\NormalTok{ fr }\OperatorTok{=} \KeywordTok{new} \BuiltInTok{FileReader}\OperatorTok{(}\StringTok{"person.json"}\OperatorTok{);}
\NormalTok{Person p2 }\OperatorTok{=}\NormalTok{ gson}\OperatorTok{.}\FunctionTok{fromJson}\OperatorTok{(}\NormalTok{fr}\OperatorTok{,}\NormalTok{ Person}\OperatorTok{.}\FunctionTok{class}\OperatorTok{);}
\BuiltInTok{System}\OperatorTok{.}\FunctionTok{out}\OperatorTok{.}\FunctionTok{println}\OperatorTok{(}\NormalTok{p2}\OperatorTok{);}
\end{Highlighting}
\end{Shaded}

\hypertarget{luxe4sa-och-skriva-binuxe4ra-filer}{%
\subsection{Läsa och skriva binära
filer}\label{luxe4sa-och-skriva-binuxe4ra-filer}}

För att läsa och skriva binära filer, använder vi
\texttt{FileInputStream} och \texttt{FileOutputStream}-klasserna. Vi
skapar ett nytt objekt av klassen och ger den filnamnet som argument.
Sedan kan vi använda \texttt{read()} och \texttt{write()}-metoderna för
att läsa och skriva till filen.

\hypertarget{cb6}{}
\begin{Shaded}
\begin{Highlighting}[]
\CommentTok{// Skapa en person}
\NormalTok{Person p }\OperatorTok{=} \KeywordTok{new} \FunctionTok{Person}\OperatorTok{(}\StringTok{"Han Solo"}\OperatorTok{,} \DecValTok{33}\OperatorTok{);}

\CommentTok{// Serialisera personen}
\BuiltInTok{ByteArrayOutputStream}\NormalTok{ baos }\OperatorTok{=} \KeywordTok{new} \BuiltInTok{ByteArrayOutputStream}\OperatorTok{();}
\BuiltInTok{ObjectOutputStream}\NormalTok{ oos }\OperatorTok{=} \KeywordTok{new} \BuiltInTok{ObjectOutputStream}\OperatorTok{(}\NormalTok{baos}\OperatorTok{);}
\NormalTok{oos}\OperatorTok{.}\FunctionTok{writeObject}\OperatorTok{(}\NormalTok{p}\OperatorTok{);}

\CommentTok{// Skapa en ny fil och skriv personen till den}
\BuiltInTok{FileOutputStream}\NormalTok{ fos }\OperatorTok{=} \KeywordTok{new} \BuiltInTok{FileOutputStream}\OperatorTok{(}\StringTok{"person.bin"}\OperatorTok{);}
\NormalTok{fos}\OperatorTok{.}\FunctionTok{write}\OperatorTok{(}\NormalTok{baos}\OperatorTok{.}\FunctionTok{toByteArray}\OperatorTok{());}
\NormalTok{fos}\OperatorTok{.}\FunctionTok{close}\OperatorTok{();}

\CommentTok{// Läs in personen från filen}
\BuiltInTok{FileInputStream}\NormalTok{ fis }\OperatorTok{=} \KeywordTok{new} \BuiltInTok{FileInputStream}\OperatorTok{(}\StringTok{"person.bin"}\OperatorTok{);}
\BuiltInTok{ObjectInputStream}\NormalTok{ ois }\OperatorTok{=} \KeywordTok{new} \BuiltInTok{ObjectInputStream}\OperatorTok{(}\NormalTok{fis}\OperatorTok{);}
\NormalTok{Person p2 }\OperatorTok{=} \OperatorTok{(}\NormalTok{Person}\OperatorTok{)}\NormalTok{ ois}\OperatorTok{.}\FunctionTok{readObject}\OperatorTok{();}
\BuiltInTok{System}\OperatorTok{.}\FunctionTok{out}\OperatorTok{.}\FunctionTok{println}\OperatorTok{(}\NormalTok{p2}\OperatorTok{);}
\end{Highlighting}
\end{Shaded}

\hypertarget{luxe4sa-och-skriva-binuxe4ra-filer-med-base64}{%
\subsubsection{Läsa och skriva binära filer med
Base64}\label{luxe4sa-och-skriva-binuxe4ra-filer-med-base64}}

För att läsa och skriva binära filer med Base64, använder vi
\texttt{Base64}-klassen. Vi skapar ett nytt objekt av klassen och ger
den filnamnet som argument. Sedan kan vi använda \texttt{encode()} och
\texttt{decode()}-metoderna för att läsa och skriva till filen.

\hypertarget{cb7}{}
\begin{Shaded}
\begin{Highlighting}[]
\CommentTok{// Skapa en person}
\NormalTok{Person p }\OperatorTok{=} \KeywordTok{new} \FunctionTok{Person}\OperatorTok{(}\StringTok{"Chewbacca"}\OperatorTok{,} \DecValTok{53}\OperatorTok{);}

\CommentTok{// Serialisera personen}
\BuiltInTok{ByteArrayOutputStream}\NormalTok{ baos }\OperatorTok{=} \KeywordTok{new} \BuiltInTok{ByteArrayOutputStream}\OperatorTok{();}
\BuiltInTok{ObjectOutputStream}\NormalTok{ oos }\OperatorTok{=} \KeywordTok{new} \BuiltInTok{ObjectOutputStream}\OperatorTok{(}\NormalTok{baos}\OperatorTok{);}
\NormalTok{oos}\OperatorTok{.}\FunctionTok{writeObject}\OperatorTok{(}\NormalTok{p}\OperatorTok{);}

\CommentTok{// Skapa en ny fil och skriv personen till den}
\BuiltInTok{FileOutputStream}\NormalTok{ fos }\OperatorTok{=} \KeywordTok{new} \BuiltInTok{FileOutputStream}\OperatorTok{(}\StringTok{"person.bin"}\OperatorTok{);}
\NormalTok{fos}\OperatorTok{.}\FunctionTok{write}\OperatorTok{(}\NormalTok{Base64}\OperatorTok{.}\FunctionTok{getEncoder}\OperatorTok{().}\FunctionTok{encode}\OperatorTok{(}\NormalTok{baos}\OperatorTok{.}\FunctionTok{toByteArray}\OperatorTok{()));}
\NormalTok{fos}\OperatorTok{.}\FunctionTok{close}\OperatorTok{();}

\CommentTok{// Läs in personen från filen}
\BuiltInTok{FileInputStream}\NormalTok{ fis }\OperatorTok{=} \KeywordTok{new} \BuiltInTok{FileInputStream}\OperatorTok{(}\StringTok{"person.bin"}\OperatorTok{);}
\BuiltInTok{ObjectInputStream}\NormalTok{ ois }\OperatorTok{=} \KeywordTok{new} \BuiltInTok{ObjectInputStream}\OperatorTok{(}\NormalTok{fis}\OperatorTok{);}
\NormalTok{Person p2 }\OperatorTok{=} \OperatorTok{(}\NormalTok{Person}\OperatorTok{)}\NormalTok{ ois}\OperatorTok{.}\FunctionTok{readObject}\OperatorTok{();}
\BuiltInTok{System}\OperatorTok{.}\FunctionTok{out}\OperatorTok{.}\FunctionTok{println}\OperatorTok{(}\NormalTok{p2}\OperatorTok{);}
\end{Highlighting}
\end{Shaded}

\hypertarget{luxe4sa-och-skriva-binuxe4ra-filer-med-base64-och-gzip}{%
\subsection{Läsa och skriva binära filer med Base64 och
GZIP}\label{luxe4sa-och-skriva-binuxe4ra-filer-med-base64-och-gzip}}

För att läsa och skriva binära filer med Base64 och GZIP, använder vi
\texttt{Base64}-klassen och \texttt{GZIPOutputStream}-klassen. Vi skapar
ett nytt objekt av klassen och ger den filnamnet som argument. Sedan kan
vi använda \texttt{encode()} och \texttt{decode()}-metoderna för att
läsa och skriva till filen.

\hypertarget{cb8}{}
\begin{Shaded}
\begin{Highlighting}[]
\CommentTok{// Skapa en person}
\NormalTok{Person p }\OperatorTok{=} \KeywordTok{new} \FunctionTok{Person}\OperatorTok{(}\StringTok{"Obi{-}Wan Kenobi"}\OperatorTok{,} \DecValTok{53}\OperatorTok{);}

\CommentTok{// Serialisera personen}
\BuiltInTok{ByteArrayOutputStream}\NormalTok{ baos }\OperatorTok{=} \KeywordTok{new} \BuiltInTok{ByteArrayOutputStream}\OperatorTok{();}
\BuiltInTok{ObjectOutputStream}\NormalTok{ oos }\OperatorTok{=} \KeywordTok{new} \BuiltInTok{ObjectOutputStream}\OperatorTok{(}\NormalTok{baos}\OperatorTok{);}
\NormalTok{oos}\OperatorTok{.}\FunctionTok{writeObject}\OperatorTok{(}\NormalTok{p}\OperatorTok{);}

\CommentTok{// Skapa en ny fil och skriv personen till den}
\BuiltInTok{FileOutputStream}\NormalTok{ fos }\OperatorTok{=} \KeywordTok{new} \BuiltInTok{FileOutputStream}\OperatorTok{(}\StringTok{"person.gzis"}\OperatorTok{);}
\BuiltInTok{GZIPOutputStream}\NormalTok{ gzos }\OperatorTok{=} \KeywordTok{new} \BuiltInTok{GZIPOutputStream}\OperatorTok{(}\NormalTok{fos}\OperatorTok{);}
\NormalTok{gos}\OperatorTok{.}\FunctionTok{write}\OperatorTok{(}\NormalTok{Base64}\OperatorTok{.}\FunctionTok{getEncoder}\OperatorTok{().}\FunctionTok{encode}\OperatorTok{(}\NormalTok{baos}\OperatorTok{.}\FunctionTok{toByteArray}\OperatorTok{()));}
\NormalTok{gos}\OperatorTok{.}\FunctionTok{close}\OperatorTok{();}

\CommentTok{// Läs in personen från filen}
\BuiltInTok{FileInputStream}\NormalTok{ fis }\OperatorTok{=} \KeywordTok{new} \BuiltInTok{FileInputStream}\OperatorTok{(}\StringTok{"person.gzis"}\OperatorTok{);}
\BuiltInTok{GZIPInputStream}\NormalTok{ gzis }\OperatorTok{=} \KeywordTok{new} \BuiltInTok{GZIPInputStream}\OperatorTok{(}\NormalTok{fis}\OperatorTok{);}
\BuiltInTok{ObjectInputStream}\NormalTok{ ois }\OperatorTok{=} \KeywordTok{new} \BuiltInTok{ObjectInputStream}\OperatorTok{(}\NormalTok{gzis}\OperatorTok{);}
\NormalTok{Person p2 }\OperatorTok{=} \OperatorTok{(}\NormalTok{Person}\OperatorTok{)}\NormalTok{ ois}\OperatorTok{.}\FunctionTok{readObject}\OperatorTok{();}
\BuiltInTok{System}\OperatorTok{.}\FunctionTok{out}\OperatorTok{.}\FunctionTok{println}\OperatorTok{(}\NormalTok{p2}\OperatorTok{);}
\end{Highlighting}
\end{Shaded}

\hypertarget{uxf6vrig-filhantering}{%
\subsection{Övrig filhantering}\label{uxf6vrig-filhantering}}

\hypertarget{kolla-om-en-fil-finns}{%
\subsubsection{Kolla om en fil finns}\label{kolla-om-en-fil-finns}}

\hypertarget{cb9}{}
\begin{Shaded}
\begin{Highlighting}[]
\BuiltInTok{File}\NormalTok{ f }\OperatorTok{=} \KeywordTok{new} \BuiltInTok{File}\OperatorTok{(}\StringTok{"person.txt"}\OperatorTok{);}
\ControlFlowTok{if}\OperatorTok{(}\NormalTok{f}\OperatorTok{.}\FunctionTok{exists}\OperatorTok{())\{}
    \BuiltInTok{System}\OperatorTok{.}\FunctionTok{out}\OperatorTok{.}\FunctionTok{println}\OperatorTok{(}\StringTok{"Filen finns"}\OperatorTok{);}
\OperatorTok{\}}
\end{Highlighting}
\end{Shaded}

\hypertarget{skapa-en-tom-fil-om-den-inte-finns}{%
\subsubsection{Skapa en tom fil om den inte
finns}\label{skapa-en-tom-fil-om-den-inte-finns}}

\hypertarget{cb10}{}
\begin{Shaded}
\begin{Highlighting}[]
\BuiltInTok{File}\NormalTok{ f }\OperatorTok{=} \KeywordTok{new} \BuiltInTok{File}\OperatorTok{(}\StringTok{"person.txt"}\OperatorTok{);}
\ControlFlowTok{if}\OperatorTok{(!}\NormalTok{f}\OperatorTok{.}\FunctionTok{exists}\OperatorTok{())\{}
\NormalTok{    f}\OperatorTok{.}\FunctionTok{createNewFile}\OperatorTok{();}
\OperatorTok{\}}
\end{Highlighting}
\end{Shaded}

\hypertarget{skapa-en-mapp-om-den-inte-finns}{%
\subsubsection{Skapa en mapp om den inte
finns}\label{skapa-en-mapp-om-den-inte-finns}}

\hypertarget{cb11}{}
\begin{Shaded}
\begin{Highlighting}[]
\BuiltInTok{File}\NormalTok{ f }\OperatorTok{=} \KeywordTok{new} \BuiltInTok{File}\OperatorTok{(}\StringTok{"mapp"}\OperatorTok{);}
\ControlFlowTok{if}\OperatorTok{(!}\NormalTok{f}\OperatorTok{.}\FunctionTok{exists}\OperatorTok{())\{}
\NormalTok{    f}\OperatorTok{.}\FunctionTok{mkdir}\OperatorTok{();}
\OperatorTok{\}}
\end{Highlighting}
\end{Shaded}

\hypertarget{radera-en-fil}{%
\subsection{Radera en fil}\label{radera-en-fil}}

\hypertarget{cb12}{}
\begin{Shaded}
\begin{Highlighting}[]
\BuiltInTok{File}\NormalTok{ f }\OperatorTok{=} \KeywordTok{new} \BuiltInTok{File}\OperatorTok{(}\StringTok{"person.txt"}\OperatorTok{);}
\ControlFlowTok{if}\OperatorTok{(}\NormalTok{f}\OperatorTok{.}\FunctionTok{exists}\OperatorTok{())\{}
\NormalTok{    f}\OperatorTok{.}\FunctionTok{delete}\OperatorTok{();}
\OperatorTok{\}}
\end{Highlighting}
\end{Shaded}

\hypertarget{radera-en-mapp-om-den-uxe4r-tom}{%
\subsection{Radera en mapp om den är
tom}\label{radera-en-mapp-om-den-uxe4r-tom}}

\hypertarget{cb13}{}
\begin{Shaded}
\begin{Highlighting}[]
\BuiltInTok{File}\NormalTok{ f }\OperatorTok{=} \KeywordTok{new} \BuiltInTok{File}\OperatorTok{(}\StringTok{"mapp"}\OperatorTok{);}
\ControlFlowTok{if}\OperatorTok{(}\NormalTok{f}\OperatorTok{.}\FunctionTok{exists}\OperatorTok{())\{}
    \ControlFlowTok{if} \OperatorTok{(}\NormalTok{f}\OperatorTok{.}\FunctionTok{isDirectory}\OperatorTok{()} \OperatorTok{\&\&}\NormalTok{ f}\OperatorTok{.}\FunctionTok{list}\OperatorTok{().}\FunctionTok{length} \OperatorTok{==} \DecValTok{0}\OperatorTok{)}
\NormalTok{        f}\OperatorTok{.}\FunctionTok{delete}\OperatorTok{();}
\OperatorTok{\}}
\end{Highlighting}
\end{Shaded}

\hypertarget{flytta-en-fil}{%
\subsection{Flytta en fil}\label{flytta-en-fil}}

\hypertarget{cb14}{}
\begin{Shaded}
\begin{Highlighting}[]
\BuiltInTok{File}\NormalTok{ f }\OperatorTok{=} \KeywordTok{new} \BuiltInTok{File}\OperatorTok{(}\StringTok{"person.txt"}\OperatorTok{);}
\ControlFlowTok{if}\OperatorTok{(}\NormalTok{f}\OperatorTok{.}\FunctionTok{exists}\OperatorTok{())\{}
\NormalTok{    f}\OperatorTok{.}\FunctionTok{renameTo}\OperatorTok{(}\KeywordTok{new} \BuiltInTok{File}\OperatorTok{(}\StringTok{"person2.txt"}\OperatorTok{));}
\OperatorTok{\}}
\end{Highlighting}
\end{Shaded}

\hypertarget{lista-alla-filer-i-en-mapp}{%
\subsection{Lista alla filer i en
mapp}\label{lista-alla-filer-i-en-mapp}}

\hypertarget{cb15}{}
\begin{Shaded}
\begin{Highlighting}[]
\BuiltInTok{File}\NormalTok{ f }\OperatorTok{=} \KeywordTok{new} \BuiltInTok{File}\OperatorTok{(}\StringTok{"mapp"}\OperatorTok{);}
\ControlFlowTok{if}\OperatorTok{(}\NormalTok{f}\OperatorTok{.}\FunctionTok{exists}\OperatorTok{())\{}
    \BuiltInTok{File}\OperatorTok{[]}\NormalTok{ files }\OperatorTok{=}\NormalTok{ f}\OperatorTok{.}\FunctionTok{listFiles}\OperatorTok{();}
    \ControlFlowTok{for}\OperatorTok{(}\BuiltInTok{File}\NormalTok{ file }\OperatorTok{:}\NormalTok{ files}\OperatorTok{)\{}
        \BuiltInTok{System}\OperatorTok{.}\FunctionTok{out}\OperatorTok{.}\FunctionTok{println}\OperatorTok{(}\NormalTok{file}\OperatorTok{.}\FunctionTok{getName}\OperatorTok{());}
    \OperatorTok{\}}
\OperatorTok{\}}
\end{Highlighting}
\end{Shaded}

\hypertarget{lista-alla-filer-i-en-mapp-och-dess-undermappar}{%
\subsection{Lista alla filer i en mapp och dess
undermappar}\label{lista-alla-filer-i-en-mapp-och-dess-undermappar}}

\hypertarget{cb16}{}
\begin{Shaded}
\begin{Highlighting}[]
\BuiltInTok{File}\NormalTok{ f }\OperatorTok{=} \KeywordTok{new} \BuiltInTok{File}\OperatorTok{(}\StringTok{"mapp"}\OperatorTok{);}

\ControlFlowTok{if}\OperatorTok{(}\NormalTok{f}\OperatorTok{.}\FunctionTok{exists}\OperatorTok{())\{}
    \BuiltInTok{File}\OperatorTok{[]}\NormalTok{ files }\OperatorTok{=}\NormalTok{ f}\OperatorTok{.}\FunctionTok{listFiles}\OperatorTok{();}
    \ControlFlowTok{for}\OperatorTok{(}\BuiltInTok{File}\NormalTok{ file }\OperatorTok{:}\NormalTok{ files}\OperatorTok{)\{}
        \ControlFlowTok{if}\OperatorTok{(}\NormalTok{file}\OperatorTok{.}\FunctionTok{isDirectory}\OperatorTok{())\{}
            \BuiltInTok{File}\OperatorTok{[]}\NormalTok{ subFiles }\OperatorTok{=}\NormalTok{ file}\OperatorTok{.}\FunctionTok{listFiles}\OperatorTok{();}
            \ControlFlowTok{for}\OperatorTok{(}\BuiltInTok{File}\NormalTok{ subFile }\OperatorTok{:}\NormalTok{ subFiles}\OperatorTok{)\{}
                \BuiltInTok{System}\OperatorTok{.}\FunctionTok{out}\OperatorTok{.}\FunctionTok{println}\OperatorTok{(}\NormalTok{subFile}\OperatorTok{.}\FunctionTok{getName}\OperatorTok{());}
            \OperatorTok{\}}
        \OperatorTok{\}}\ControlFlowTok{else}\OperatorTok{\{}
            \BuiltInTok{System}\OperatorTok{.}\FunctionTok{out}\OperatorTok{.}\FunctionTok{println}\OperatorTok{(}\NormalTok{file}\OperatorTok{.}\FunctionTok{getName}\OperatorTok{());}
        \OperatorTok{\}}
    \OperatorTok{\}}
\OperatorTok{\}}
\end{Highlighting}
\end{Shaded}

\hypertarget{hitta-en-fil-som-med-namn-som-innehuxe5ller-katt}{%
\subsection{Hitta en fil som med namn som innehåller
``Katt''}\label{hitta-en-fil-som-med-namn-som-innehuxe5ller-katt}}

\hypertarget{cb17}{}
\begin{Shaded}
\begin{Highlighting}[]
\BuiltInTok{File}\NormalTok{ f }\OperatorTok{=} \KeywordTok{new} \BuiltInTok{File}\OperatorTok{(}\StringTok{"mapp"}\OperatorTok{);}
\ControlFlowTok{if}\OperatorTok{(}\NormalTok{f}\OperatorTok{.}\FunctionTok{exists}\OperatorTok{())\{}
    \BuiltInTok{File}\OperatorTok{[]}\NormalTok{ files }\OperatorTok{=}\NormalTok{ f}\OperatorTok{.}\FunctionTok{listFiles}\OperatorTok{();}
    \ControlFlowTok{for}\OperatorTok{(}\BuiltInTok{File}\NormalTok{ file }\OperatorTok{:}\NormalTok{ files}\OperatorTok{)\{}
        \ControlFlowTok{if}\OperatorTok{(}\NormalTok{file}\OperatorTok{.}\FunctionTok{getName}\OperatorTok{().}\FunctionTok{contains}\OperatorTok{(}\StringTok{"Katt"}\OperatorTok{))\{}
            \BuiltInTok{System}\OperatorTok{.}\FunctionTok{out}\OperatorTok{.}\FunctionTok{println}\OperatorTok{(}\NormalTok{file}\OperatorTok{.}\FunctionTok{getName}\OperatorTok{());}
        \OperatorTok{\}}
    \OperatorTok{\}}
\OperatorTok{\}}
\end{Highlighting}
\end{Shaded}

\hypertarget{hitta-en-fil-som-med-namn-som-innehuxe5ller-hemlighet-och-radera-den}{%
\subsection{Hitta en fil som med namn som innehåller ``Hemlighet'' och
radera
den}\label{hitta-en-fil-som-med-namn-som-innehuxe5ller-hemlighet-och-radera-den}}

\hypertarget{cb18}{}
\begin{Shaded}
\begin{Highlighting}[]
\BuiltInTok{File}\NormalTok{ f }\OperatorTok{=} \KeywordTok{new} \BuiltInTok{File}\OperatorTok{(}\StringTok{"mapp"}\OperatorTok{);}
\ControlFlowTok{if}\OperatorTok{(}\NormalTok{f}\OperatorTok{.}\FunctionTok{exists}\OperatorTok{())\{}
    \BuiltInTok{File}\OperatorTok{[]}\NormalTok{ files }\OperatorTok{=}\NormalTok{ f}\OperatorTok{.}\FunctionTok{listFiles}\OperatorTok{();}
    \ControlFlowTok{for}\OperatorTok{(}\BuiltInTok{File}\NormalTok{ file }\OperatorTok{:}\NormalTok{ files}\OperatorTok{)\{}
        \ControlFlowTok{if}\OperatorTok{(}\NormalTok{file}\OperatorTok{.}\FunctionTok{getName}\OperatorTok{().}\FunctionTok{contains}\OperatorTok{(}\StringTok{"Hemlighet"}\OperatorTok{))\{}
\NormalTok{            file}\OperatorTok{.}\FunctionTok{delete}\OperatorTok{();}
        \OperatorTok{\}}
    \OperatorTok{\}}
\OperatorTok{\}}
\end{Highlighting}
\end{Shaded}

\hypertarget{hitta-en-fil-med-texten-luxf6senord-i-sig}{%
\subsection{Hitta en fil med texten ``Lösenord'' i
sig}\label{hitta-en-fil-med-texten-luxf6senord-i-sig}}

\hypertarget{cb19}{}
\begin{Shaded}
\begin{Highlighting}[]
\BuiltInTok{File}\NormalTok{ f }\OperatorTok{=} \KeywordTok{new} \BuiltInTok{File}\OperatorTok{(}\StringTok{"mapp"}\OperatorTok{);}

\ControlFlowTok{if}\OperatorTok{(}\NormalTok{f}\OperatorTok{.}\FunctionTok{exists}\OperatorTok{())\{}
    \BuiltInTok{File}\OperatorTok{[]}\NormalTok{ files }\OperatorTok{=}\NormalTok{ f}\OperatorTok{.}\FunctionTok{listFiles}\OperatorTok{();}
    \ControlFlowTok{for}\OperatorTok{(}\BuiltInTok{File}\NormalTok{ file }\OperatorTok{:}\NormalTok{ files}\OperatorTok{)\{}
        \ControlFlowTok{if}\OperatorTok{(}\NormalTok{file}\OperatorTok{.}\FunctionTok{isFile}\OperatorTok{())\{}
            \BuiltInTok{FileReader}\NormalTok{ fr }\OperatorTok{=} \KeywordTok{new} \BuiltInTok{FileReader}\OperatorTok{(}\NormalTok{file}\OperatorTok{);}
            \BuiltInTok{String}\NormalTok{ rad }\OperatorTok{=}\NormalTok{ fr}\OperatorTok{.}\FunctionTok{readLine}\OperatorTok{();}
            \ControlFlowTok{while}\OperatorTok{(}\NormalTok{rad }\OperatorTok{!=} \KeywordTok{null}\OperatorTok{)\{}
                \ControlFlowTok{if}\OperatorTok{(}\NormalTok{rad}\OperatorTok{.}\FunctionTok{contains}\OperatorTok{(}\StringTok{"Lösenord"}\OperatorTok{))\{}
                    \BuiltInTok{System}\OperatorTok{.}\FunctionTok{out}\OperatorTok{.}\FunctionTok{println}\OperatorTok{(}\NormalTok{file}\OperatorTok{.}\FunctionTok{getName}\OperatorTok{());}
                \OperatorTok{\}}
\NormalTok{                rad }\OperatorTok{=}\NormalTok{ fr}\OperatorTok{.}\FunctionTok{readLine}\OperatorTok{();}
            \OperatorTok{\}}
        \OperatorTok{\}}
    \OperatorTok{\}}
\OperatorTok{\}}
\end{Highlighting}
\end{Shaded}

\hypertarget{fuxe5-information-om-filen}{%
\subsection{Få information om filen}\label{fuxe5-information-om-filen}}

\hypertarget{cb20}{}
\begin{Shaded}
\begin{Highlighting}[]
\BuiltInTok{File}\NormalTok{ f }\OperatorTok{=} \KeywordTok{new} \BuiltInTok{File}\OperatorTok{(}\StringTok{"person.txt"}\OperatorTok{);}

\ControlFlowTok{if}\OperatorTok{(}\NormalTok{f}\OperatorTok{.}\FunctionTok{exists}\OperatorTok{())\{}
    \BuiltInTok{System}\OperatorTok{.}\FunctionTok{out}\OperatorTok{.}\FunctionTok{println}\OperatorTok{(}\StringTok{"Filnamn: "} \OperatorTok{+}\NormalTok{ f}\OperatorTok{.}\FunctionTok{getName}\OperatorTok{());}
    \BuiltInTok{System}\OperatorTok{.}\FunctionTok{out}\OperatorTok{.}\FunctionTok{println}\OperatorTok{(}\StringTok{"Filtyp: "} \OperatorTok{+}\NormalTok{ f}\OperatorTok{.}\FunctionTok{getType}\OperatorTok{());}
    \BuiltInTok{System}\OperatorTok{.}\FunctionTok{out}\OperatorTok{.}\FunctionTok{println}\OperatorTok{(}\StringTok{"Mapp: "} \OperatorTok{+}\NormalTok{ f}\OperatorTok{.}\FunctionTok{getParent}\OperatorTok{());}
    \BuiltInTok{System}\OperatorTok{.}\FunctionTok{out}\OperatorTok{.}\FunctionTok{println}\OperatorTok{(}\StringTok{"Sökväg: "} \OperatorTok{+}\NormalTok{ f}\OperatorTok{.}\FunctionTok{getAbsolutePath}\OperatorTok{());}
    \BuiltInTok{System}\OperatorTok{.}\FunctionTok{out}\OperatorTok{.}\FunctionTok{println}\OperatorTok{(}\StringTok{"Storlek: "} \OperatorTok{+}\NormalTok{ f}\OperatorTok{.}\FunctionTok{length}\OperatorTok{());}
    \BuiltInTok{System}\OperatorTok{.}\FunctionTok{out}\OperatorTok{.}\FunctionTok{println}\OperatorTok{(}\StringTok{"Senast ändrad: "} \OperatorTok{+}\NormalTok{ f}\OperatorTok{.}\FunctionTok{lastModified}\OperatorTok{());}
    \BuiltInTok{System}\OperatorTok{.}\FunctionTok{out}\OperatorTok{.}\FunctionTok{println}\OperatorTok{(}\StringTok{"Är fil: "} \OperatorTok{+}\NormalTok{ f}\OperatorTok{.}\FunctionTok{isFile}\OperatorTok{());}
    \BuiltInTok{System}\OperatorTok{.}\FunctionTok{out}\OperatorTok{.}\FunctionTok{println}\OperatorTok{(}\StringTok{"Är mapp: "} \OperatorTok{+}\NormalTok{ f}\OperatorTok{.}\FunctionTok{isDirectory}\OperatorTok{());}
    \BuiltInTok{System}\OperatorTok{.}\FunctionTok{out}\OperatorTok{.}\FunctionTok{println}\OperatorTok{(}\StringTok{"Är dold: "} \OperatorTok{+}\NormalTok{ f}\OperatorTok{.}\FunctionTok{isHidden}\OperatorTok{());}
    \BuiltInTok{System}\OperatorTok{.}\FunctionTok{out}\OperatorTok{.}\FunctionTok{println}\OperatorTok{(}\StringTok{"Är läsbar: "} \OperatorTok{+}\NormalTok{ f}\OperatorTok{.}\FunctionTok{canRead}\OperatorTok{());}
    \BuiltInTok{System}\OperatorTok{.}\FunctionTok{out}\OperatorTok{.}\FunctionTok{println}\OperatorTok{(}\StringTok{"Är skrivbar: "} \OperatorTok{+}\NormalTok{ f}\OperatorTok{.}\FunctionTok{canWrite}\OperatorTok{());}
    \BuiltInTok{System}\OperatorTok{.}\FunctionTok{out}\OperatorTok{.}\FunctionTok{println}\OperatorTok{(}\StringTok{"Är exekverbar: "} \OperatorTok{+}\NormalTok{ f}\OperatorTok{.}\FunctionTok{canExecute}\OperatorTok{());}
    \BuiltInTok{System}\OperatorTok{.}\FunctionTok{out}\OperatorTok{.}\FunctionTok{println}\OperatorTok{(}\StringTok{"Är samma fil: "} \OperatorTok{+}\NormalTok{ f}\OperatorTok{.}\FunctionTok{equals}\OperatorTok{(}\KeywordTok{new} \BuiltInTok{File}\OperatorTok{(}\StringTok{"person.txt"}\OperatorTok{)));}
\OperatorTok{\}}
\end{Highlighting}
\end{Shaded}

\hypertarget{termer-och-fuxf6rklaringar}{%
\subsection{Termer och förklaringar}\label{termer-och-fuxf6rklaringar}}

Här följer termer som använts i denna artikel och deras förklaringar.

\begin{longtable}[]{@{}
  >{\raggedright\arraybackslash}p{(\columnwidth - 2\tabcolsep) * \real{0.5000}}
  >{\raggedright\arraybackslash}p{(\columnwidth - 2\tabcolsep) * \real{0.5000}}@{}}
\toprule\noalign{}
\begin{minipage}[b]{\linewidth}\raggedright
Term
\end{minipage} & \begin{minipage}[b]{\linewidth}\raggedright
Förklaring
\end{minipage} \\
\midrule\noalign{}
\endhead
\bottomrule\noalign{}
\endlastfoot
\texttt{FileReader} & En klass som används för att läsa textfiler. \\
\texttt{FileWriter} & En klass som används för att skriva till
textfiler. \\
\texttt{Gson} & En klass som används för att läsa och skriva till
Json-filer. \\
\texttt{Base64} & En klass som används för att läsa och skriva binära
filer. \\
\texttt{GZIPOutputStream} & En klass som används för att läsa och skriva
binära filer med Base64 och GZIP. \\
\texttt{ObjectOutputStream} & En klass som används för att läsa och
skriva binära filer. \\
\texttt{ObjectInputStream} & En klass som används för att läsa och
skriva binära filer. \\
\texttt{ByteArrayOutputStream} & En klass som används för att läsa och
skriva binära filer. \\
\texttt{FileInputStream} & En klass som används för att läsa och skriva
binära filer. \\
\texttt{FileOutputStream} & En klass som används för att läsa och skriva
binära filer. \\
\texttt{GZIPInputStream} & En klass som används för att läsa och skriva
binära filer med Base64 och GZIP. \\
\texttt{readLine()} & En metod som används för att läsa en rad från en
textfil. \\
\texttt{write()} & En metod som används för att skriva till en
textfil. \\
\texttt{toJson()} & En metod som används för att skriva till en
Json-fil. \\
\texttt{fromJson()} & En metod som används för att läsa från en
Json-fil. \\
\texttt{read()} & En metod som används för att läsa från en binär
fil. \\
\texttt{write()} & En metod som används för att skriva till en binär
fil. \\
\texttt{encode()} & En metod som används för att läsa från en binär fil
med Base64. \\
\texttt{decode()} & En metod som används för att skriva till en binär
fil med Base64. \\
\texttt{close()} & En metod som används för att stänga en fil. \\
\texttt{split()} & En metod som används för att dela upp en sträng i en
array. \\
\texttt{toString()} & En metod som används för att konvertera ett objekt
till en sträng. \\
\texttt{getNamn()} & En metod som används för att hämta namnet på en
person. \\
\texttt{getÅlder()} & En metod som används för att hämta åldern på en
person. \\
\texttt{setNamn()} & En metod som används för att sätta namnet på en
person. \\
\texttt{setÅlder()} & En metod som används för att sätta åldern på en
person. \\
\texttt{Person} & En klass som används för att representera en
person. \\
\texttt{namn} & En variabel som används för att lagra namnet på en
person. \\
\texttt{ålder} & En variabel som används för att lagra åldern på en
person. \\
\texttt{p} & En variabel som används för att representera en person. \\
\texttt{p2} & En variabel som används för att representera en person. \\
\texttt{rad} & En variabel som används för att representera en rad i en
textfil. \\
\texttt{delar} & En variabel som används för att representera en array
av strängar. \\
\texttt{fr} & En variabel som används för att representera en
textfil. \\
\texttt{fw} & En variabel som används för att representera en
textfil. \\
\texttt{gson} & En variabel som används för att representera en
Json-fil. \\
\texttt{baos} & En variabel som används för att representera en binär
fil. \\
\texttt{oos} & En variabel som används för att representera en binär
fil. \\
\texttt{fos} & En variabel som används för att representera en binär
fil. \\
\texttt{fis} & En variabel som används för att representera en binär
fil. \\
\texttt{gzos} & En variabel som används för att representera en binär
fil med Base64 och GZIP. \\
\texttt{gzis} & En variabel som används för att representera en binär
fil med Base64 och GZIP. \\
\texttt{ois} & En variabel som används för att representera en binär
fil. \\
\texttt{fr} & En variabel som används för att representera en
textfil. \\
Serialize & Att konvertera ett objekt till en sträng. \\
Deserialize & Att konvertera en sträng till ett objekt. \\
\end{longtable}

\hypertarget{sammanfattning}{%
\subsection{Sammanfattning}\label{sammanfattning}}

Det finns många olika sätt att spara information på, välj den som känns
trevligast för dig helt enkelt.

\hypertarget{csv}{%
\section{CSV}\label{csv}}

Filformatet CSV (Comma Separated Values) är ett filformat som används
för att lagra data i en tabell. Det är ett vanligt filformat som används
för att lagra data i Excel. Det är ett textbaserat filformat som
använder kommatecken för att separera värdena i en rad. CSV-filer kan
enkelt läsas och skrivas med hjälp av Java-kod.

Innehållsförteckning

\{: .text-delta \} 1. TOC \{:toc\}

\hypertarget{fuxf6rdelar}{%
\subsection{Fördelar}\label{fuxf6rdelar}}

CSV-filformatet har flera fördelar när det gäller att lagra och hantera
data:

\begin{itemize}
\tightlist
\item
  \textbf{Enkelhet}: CSV-filer är lätta att skapa och förstå. Eftersom
  det är ett textbaserat format kan det enkelt redigeras med hjälp av en
  textredigerare eller ett kalkylbladsprogram.
\item
  \textbf{Kompatibilitet}: CSV-filer kan läsas och skrivas av de flesta
  programvaror och programmeringsspråk. Detta gör det enkelt att dela
  och utbyta data mellan olika system.
\item
  \textbf{Effektivitet}: CSV-filer är kompakta eftersom de använder en
  enkel textrepresentation för att lagra data. Detta minimerar
  filstorleken och gör det snabbt att läsa och skriva data.
\item
  \textbf{Tabellstruktur}: CSV-filer behåller tabellstrukturen, vilket
  gör det enkelt att representera relationella data. Varje rad i filen
  motsvarar en post i tabellen, och varje kolumn motsvarar en attribut
  eller en egenskap för posten.
\end{itemize}

\hypertarget{begruxe4nsningar}{%
\subsection{Begränsningar}\label{begruxe4nsningar}}

Trots sina fördelar har CSV-filer vissa begränsningar att vara medveten
om:

\begin{itemize}
\tightlist
\item
  \textbf{Ingen standard}: Det finns ingen officiell standard för
  CSV-filformatet, vilket kan leda till inkompatibiliteter mellan olika
  implementationer. Det är viktigt att vara medveten om vilka regler och
  konventioner som används för att skapa och tolka CSV-filer.
\item
  \textbf{Begränsad datatypsstöd}: CSV-filer hanterar bara textbaserade
  data. Om du behöver lagra mer komplexa datatyper som datum, tid eller
  binära data kan det kräva extra ansträngning för att konvertera och
  tolka dessa värden korrekt.
\item
  \textbf{Brister i struktur}: CSV-filer saknar en inbyggd mekanism för
  att beskriva datastrukturen. Det innebär att det kan vara svårt att
  veta vad varje kolumn representerar om det inte finns en tydlig
  dokumentation eller överenskommelse om filens innehåll.
\end{itemize}

\hypertarget{anvuxe4ndningsomruxe5den}{%
\subsection{Användningsområden}\label{anvuxe4ndningsomruxe5den}}

CSV-filer används inom olika områden och har många användningsområden:

\begin{itemize}
\tightlist
\item
  \textbf{Dataimport och -export}: CSV-filer används ofta för att
  importera och exportera data från olika program och system. Det kan
  vara användbart när du behöver överföra data mellan olika databaser,
  kalkylblad eller applikationer.
\item
  \textbf{Dataanalys}: CSV-filer är vanliga inom dataanalys och
  affärsintelligens. Genom att exportera data från olika källor till
  CSV-format kan du samla och analysera informationen med hjälp av
  specialiserade analysverktyg eller skript.
\item
  \textbf{Testdata}: CSV-filer kan användas för att skapa testdata för
  programvarutestning. Genom att generera CSV-filer med olika scenarier
  och värden kan du testa programmet med olika dataset och se till att
  det fungerar korrekt i olika situationer.
\item
  \textbf{Konfigurationsfiler}: CSV-filer kan användas som en enkel form
  av konfigurationsfiler. Du kan använda dem för att lagra och läsa in
  inställningar eller parametrar som behövs för att konfigurera en
  applikation eller ett system.
\end{itemize}

\hypertarget{kodexempel}{%
\subsection{Kodexempel}\label{kodexempel}}

Här är några kodexempel som visar hur du kan läsa och skriva CSV-filer i
Java:

\hypertarget{luxe4sning-av-csv-fil}{%
\subsubsection{Läsning av CSV-fil}\label{luxe4sning-av-csv-fil}}

\hypertarget{cb1}{}
\begin{Shaded}
\begin{Highlighting}[]
\KeywordTok{import} \ImportTok{java}\OperatorTok{.}\ImportTok{io}\OperatorTok{.}\ImportTok{BufferedReader}\OperatorTok{;}
\KeywordTok{import} \ImportTok{java}\OperatorTok{.}\ImportTok{io}\OperatorTok{.}\ImportTok{FileReader}\OperatorTok{;}
\KeywordTok{import} \ImportTok{java}\OperatorTok{.}\ImportTok{io}\OperatorTok{.}\ImportTok{IOException}\OperatorTok{;}

\KeywordTok{public} \KeywordTok{class}\NormalTok{ CSVReader }\OperatorTok{\{}
    \KeywordTok{public} \DataTypeTok{static} \DataTypeTok{void} \FunctionTok{main}\OperatorTok{(}\BuiltInTok{String}\OperatorTok{[]}\NormalTok{ args}\OperatorTok{)} \OperatorTok{\{}
        \BuiltInTok{String}\NormalTok{ file }\OperatorTok{=} \StringTok{"myfile.csv"}\OperatorTok{;}

        \ControlFlowTok{try} \OperatorTok{(}\BuiltInTok{BufferedReader}\NormalTok{ br }\OperatorTok{=} \KeywordTok{new} \BuiltInTok{BufferedReader}\OperatorTok{(}\KeywordTok{new} \BuiltInTok{FileReader}\OperatorTok{(}\NormalTok{file}\OperatorTok{)))} \OperatorTok{\{}
            \BuiltInTok{String}\NormalTok{ line}\OperatorTok{;}
            \ControlFlowTok{while} \OperatorTok{((}\NormalTok{line }\OperatorTok{=}\NormalTok{ br}\OperatorTok{.}\FunctionTok{readLine}\OperatorTok{())} \OperatorTok{!=} \KeywordTok{null}\OperatorTok{)} \OperatorTok{\{}
                \BuiltInTok{String}\OperatorTok{[]}\NormalTok{ values }\OperatorTok{=}\NormalTok{ line}\OperatorTok{.}\FunctionTok{split}\OperatorTok{(}\StringTok{","}\OperatorTok{);}
                \ControlFlowTok{for} \OperatorTok{(}\BuiltInTok{String}\NormalTok{ value }\OperatorTok{:}\NormalTok{ values}\OperatorTok{)} \OperatorTok{\{}
                    \BuiltInTok{System}\OperatorTok{.}\FunctionTok{out}\OperatorTok{.}\FunctionTok{print}\OperatorTok{(}\NormalTok{value }\OperatorTok{+} \StringTok{" "}\OperatorTok{);}
                \OperatorTok{\}}
                \BuiltInTok{System}\OperatorTok{.}\FunctionTok{out}\OperatorTok{.}\FunctionTok{println}\OperatorTok{();}
            \OperatorTok{\}}
        \OperatorTok{\}} \ControlFlowTok{catch} \OperatorTok{(}\BuiltInTok{IOException}\NormalTok{ e}\OperatorTok{)} \OperatorTok{\{}
\NormalTok{            e}\OperatorTok{.}\FunctionTok{printStackTrace}\OperatorTok{();}
        \OperatorTok{\}}
    \OperatorTok{\}}
\OperatorTok{\}}
\end{Highlighting}
\end{Shaded}

I det här exemplet används \texttt{BufferedReader} och
\texttt{FileReader} för att läsa in CSV-filen rad för rad. Sedan används
\texttt{String.split} för att dela upp varje rad i separata värden
baserat på kommatecken. Till sist skrivs varje värde ut på konsolen.

\hypertarget{skrivning-till-csv-fil}{%
\subsubsection{Skrivning till CSV-fil}\label{skrivning-till-csv-fil}}

\hypertarget{cb2}{}
\begin{Shaded}
\begin{Highlighting}[]
\KeywordTok{import} \ImportTok{java}\OperatorTok{.}\ImportTok{io}\OperatorTok{.}\ImportTok{BufferedWriter}\OperatorTok{;}
\KeywordTok{import} \ImportTok{java}\OperatorTok{.}\ImportTok{io}\OperatorTok{.}\ImportTok{FileWriter}\OperatorTok{;}
\KeywordTok{import} \ImportTok{java}\OperatorTok{.}\ImportTok{util}\OperatorTok{.}\ImportTok{ArrayList}\OperatorTok{;}
\KeywordTok{import} \ImportTok{java}\OperatorTok{.}\ImportTok{util}\OperatorTok{.}\ImportTok{List}\OperatorTok{;}

\KeywordTok{public} \KeywordTok{class}\NormalTok{ CSVWriter }\OperatorTok{\{}
    \KeywordTok{public} \DataTypeTok{static} \DataTypeTok{void} \FunctionTok{main}\OperatorTok{(}\BuiltInTok{String}\OperatorTok{[]}\NormalTok{ args}\OperatorTok{)} \OperatorTok{\{}
        \BuiltInTok{List}\OperatorTok{\textless{}}\NormalTok{Person}\OperatorTok{\textgreater{}}\NormalTok{ people }\OperatorTok{=} \KeywordTok{new} \BuiltInTok{ArrayList}\OperatorTok{\textless{}\textgreater{}();}
\NormalTok{        people}\OperatorTok{.}\FunctionTok{add}\OperatorTok{(}\KeywordTok{new} \FunctionTok{Person}\OperatorTok{(}\DecValTok{1}\OperatorTok{,} \StringTok{"Clark Kent"}\OperatorTok{,} \StringTok{"Superman"}\OperatorTok{,} \DecValTok{35}\OperatorTok{));}
\NormalTok{        people}\OperatorTok{.}\FunctionTok{add}\OperatorTok{(}\KeywordTok{new} \FunctionTok{Person}\OperatorTok{(}\DecValTok{2}\OperatorTok{,} \StringTok{"Bruce Wayne"}\OperatorTok{,} \StringTok{"Batman"}\OperatorTok{,} \DecValTok{40}\OperatorTok{));}
\NormalTok{        people}\OperatorTok{.}\FunctionTok{add}\OperatorTok{(}\KeywordTok{new} \FunctionTok{Person}\OperatorTok{(}\DecValTok{3}\OperatorTok{,} \StringTok{"Barry Allen"}\OperatorTok{,} \StringTok{"Flash"}\OperatorTok{,} \DecValTok{25}\OperatorTok{));}

        \ControlFlowTok{try} \OperatorTok{(}\BuiltInTok{BufferedWriter}\NormalTok{ bw }\OperatorTok{=} \KeywordTok{new} \BuiltInTok{BufferedWriter}\OperatorTok{(}\KeywordTok{new} \BuiltInTok{FileWriter}\OperatorTok{(}\NormalTok{file}\OperatorTok{)))} \OperatorTok{\{}
            \CommentTok{// Lägg till rubriker}
\NormalTok{            bw}\OperatorTok{.}\FunctionTok{write}\OperatorTok{(}\StringTok{"Id,Namn,Alias,Ålder"}\OperatorTok{);}
\NormalTok{            bw}\OperatorTok{.}\FunctionTok{newLine}\OperatorTok{();}
            \CommentTok{// Lägg till data}
            \ControlFlowTok{for} \OperatorTok{(}\NormalTok{Person person }\OperatorTok{:}\NormalTok{ people}\OperatorTok{)} \OperatorTok{\{}
\NormalTok{                bw}\OperatorTok{.}\FunctionTok{write}\OperatorTok{(}\NormalTok{person}\OperatorTok{.}\FunctionTok{getId}\OperatorTok{()} \OperatorTok{+} \StringTok{","} \OperatorTok{+}\NormalTok{ person}\OperatorTok{.}\FunctionTok{getName}\OperatorTok{()} \OperatorTok{+} \StringTok{","} \OperatorTok{+}\NormalTok{ person}\OperatorTok{.}\FunctionTok{getAlias}\OperatorTok{()} \OperatorTok{+} \StringTok{","} \OperatorTok{+}\NormalTok{ person}\OperatorTok{.}\FunctionTok{getAge}\OperatorTok{());}
\NormalTok{                bw}\OperatorTok{.}\FunctionTok{newLine}\OperatorTok{();}
            \OperatorTok{\}}
        \OperatorTok{\}} \ControlFlowTok{catch} \OperatorTok{(}\BuiltInTok{IOException}\NormalTok{ e}\OperatorTok{)} \OperatorTok{\{}
\NormalTok{            e}\OperatorTok{.}\FunctionTok{printStackTrace}\OperatorTok{();}
        \OperatorTok{\}}
    \OperatorTok{\}}
\OperatorTok{\}}

\KeywordTok{class}\NormalTok{ Person }\OperatorTok{\{}
    \KeywordTok{private} \DataTypeTok{int}\NormalTok{ id}\OperatorTok{;}
    \KeywordTok{private} \BuiltInTok{String}\NormalTok{ name}\OperatorTok{;}
    \KeywordTok{private} \BuiltInTok{String}\NormalTok{ alias}\OperatorTok{;}
    \KeywordTok{private} \DataTypeTok{int}\NormalTok{ age}\OperatorTok{;}

    \KeywordTok{public} \FunctionTok{Person}\OperatorTok{(}\DataTypeTok{int}\NormalTok{ id}\OperatorTok{,} \BuiltInTok{String}\NormalTok{ name}\OperatorTok{,} \BuiltInTok{String}\NormalTok{ alias}\OperatorTok{,} \DataTypeTok{int}\NormalTok{ age}\OperatorTok{)} \OperatorTok{\{}
        \KeywordTok{this}\OperatorTok{.}\FunctionTok{id} \OperatorTok{=}\NormalTok{ id}\OperatorTok{;}
        \KeywordTok{this}\OperatorTok{.}\FunctionTok{name} \OperatorTok{=}\NormalTok{ name}\OperatorTok{;}
        \KeywordTok{this}\OperatorTok{.}\FunctionTok{alias} \OperatorTok{=}\NormalTok{ alias}\OperatorTok{;}
        \KeywordTok{this}\OperatorTok{.}\FunctionTok{age} \OperatorTok{=}\NormalTok{ age}\OperatorTok{;}
    \OperatorTok{\}}

    \KeywordTok{public} \DataTypeTok{int} \FunctionTok{getId}\OperatorTok{()} \OperatorTok{\{}
        \ControlFlowTok{return}\NormalTok{ id}\OperatorTok{;}
    \OperatorTok{\}}

    \KeywordTok{public} \BuiltInTok{String} \FunctionTok{getName}\OperatorTok{()} \OperatorTok{\{}
        \ControlFlowTok{return}\NormalTok{ name}\OperatorTok{;}
    \OperatorTok{\}}

    \KeywordTok{public} \BuiltInTok{String} \FunctionTok{getAlias}\OperatorTok{()} \OperatorTok{\{}
        \ControlFlowTok{return}\NormalTok{ alias}\OperatorTok{;}
    \OperatorTok{\}}

    \KeywordTok{public} \DataTypeTok{int} \FunctionTok{getAge}\OperatorTok{()} \OperatorTok{\{}
        \ControlFlowTok{return}\NormalTok{ age}\OperatorTok{;}
    \OperatorTok{\}}
\OperatorTok{\}}
\end{Highlighting}
\end{Shaded}

I det här exemplet används \texttt{BufferedWriter} och
\texttt{FileWriter} för att skriva data till CSV-filen. Först läggs
rubrikerna till med \texttt{bw.write} och \texttt{bw.newLine}. Sedan
läggs varje person i listan till med \texttt{bw.write} och
\texttt{bw.newLine}.

I detta exempel används en lista av \texttt{Person}-objekt för att hålla
data. Sedan används \texttt{StringBuilder} för att bygga upp innehållet
i CSV-filen. Till sist skrivs strängen till filen med
\texttt{File.WriteAllText}.

\hypertarget{kodfuxf6rklaringar}{%
\subsection{Kodförklaringar}\label{kodfuxf6rklaringar}}

\hypertarget{files.readalllines}{%
\subsubsection{Files.readAllLines}\label{files.readalllines}}

\texttt{Files.readAllLines} läser in alla rader i en fil och lägger dem
i en lista. Varje rad blir ett element i listan.

\hypertarget{string.split}{%
\subsubsection{String.split}\label{string.split}}

\texttt{String.split} delar upp en sträng i en lista av strängar. Den
delar upp strängen vid varje komma.

\hypertarget{stringbuilder}{%
\subsubsection{StringBuilder}\label{stringbuilder}}

\texttt{StringBuilder} är en klass som används för att bygga upp en
sträng. Det är en effektivare metod än att använda strängkonkatenation.
Det är en klass som finns i \texttt{java.lang}.

\hypertarget{stringbuilder.append}{%
\subsubsection{StringBuilder.append}\label{stringbuilder.append}}

\texttt{StringBuilder.append} lägger till en rad till strängen. Den
lägger till en radbrytning efter raden.

\hypertarget{files.write}{%
\subsubsection{Files.write}\label{files.write}}

\texttt{Files.write} skriver en sträng till en fil. Den skriver över
allt som finns i filen.

\hypertarget{slutsats}{%
\subsection{Slutsats}\label{slutsats}}

CSV-filformatet är ett enkelt och flexibelt sätt att lagra och hantera
tabulära data. Det är vanligt förekommande inom olika områden, inklusive
datahantering, analys och testning. Med hjälp av Java kan du enkelt läsa
och skriva CSV-filer och manipulera data enligt dina behov. Var medveten
om dess begränsningar och regler för att få ut mesta möjliga av detta
filformat.

\hypertarget{termer}{%
\subsection{Termer}\label{termer}}

\begin{longtable}[]{@{}
  >{\raggedright\arraybackslash}p{(\columnwidth - 2\tabcolsep) * \real{0.2000}}
  >{\raggedright\arraybackslash}p{(\columnwidth - 2\tabcolsep) * \real{0.7900}}@{}}
\toprule\noalign{}
\begin{minipage}[b]{\linewidth}\raggedright
Term
\end{minipage} & \begin{minipage}[b]{\linewidth}\raggedright
Förklaring
\end{minipage} \\
\midrule\noalign{}
\endhead
\bottomrule\noalign{}
\endlastfoot
CSV & Comma Separated Values \\
Filformat & Ett sätt att lagra data i en fil \\
Textbaserat & Ett filformat som använder text för att representera
data \\
Kompatibilitet & Förmågan att fungera tillsammans med andra system \\
Prestanda & Hur snabbt eller effektivt ett system fungerar \\
Datatyp & En typ av data som kan lagras i en variabel \\
Binär & Ett filformat som använder binära tal för att representera
data \\
Datastruktur & Ett sätt att organisera data i en fil \\
Dataanalys & Att analysera data för att hitta mönster och trender \\
Affärsintelligens & Att använda data för att fatta beslut och förbättra
verksamheten \\
Testning & Att testa ett program för att hitta fel och problem \\
Konfigurationsfil & En fil som innehåller inställningar och parametrar
för ett program \\
Inställningar & Parametrar som används för att konfigurera ett
program \\
\end{longtable}

\hypertarget{luxe4nkar}{%
\subsection{Länkar}\label{luxe4nkar}}

\begin{itemize}
\tightlist
\item
  \href{https://en.wikipedia.org/wiki/Comma-separated_values}{CSV-filformatet}
\item
  \href{https://docs.oracle.com/en/java/javase/11/docs/api/java.base/java/nio/file/Files.html\#write(java.nio.file.Path,java.lang.Iterable,java.nio.file.OpenOption...)}{CSV-filformatet
  i Java}
\item
  \href{https://calendar.google.com/}{Google Kalender}
\item
  \href{https://outlook.live.com/calendar/0/view/month}{Outlook
  Kalender}
\item
  \href{https://support.google.com/calendar/answer/37118?hl=en&co=GENIE.Platform\%3DDesktop\#zippy=\%2Ccreate-or-edit-a-csv-file}{Google
  Kalender CSV}
\end{itemize}

\hypertarget{file-klassen}{%
\section{File-klassen}\label{file-klassen}}

File-klassen i Java, som tillhör paketet \texttt{java.io}, erbjuder en
mängd användbara metoder för att hantera filer. Du kan använda dessa
metoder för att skapa, skriva, läsa och ta bort filer. Exempel på några
av de vanligaste metoderna inkluderar \texttt{createNewFile},
\texttt{delete}, \texttt{exists} och \texttt{read}. Genom att använda
File-klassen kan du enkelt utföra filrelaterade operationer i dina
Java-program.

Innehållsförteckning

\{: .text-delta \}

\begin{enumerate}
\def\labelenumi{\arabic{enumi}.}
\tightlist
\item
  TOC \{:toc\}
\end{enumerate}

\hypertarget{tldr}{%
\subsection{TL;DR}\label{tldr}}

File-klassen i Java, som tillhör paketet \texttt{java.io}, erbjuder
användbara metoder för filhantering. Du kan använda dessa metoder för
att skapa, skriva, läsa och ta bort filer. Några av de vanligaste
metoderna inkluderar \texttt{createNewFile}, \texttt{delete},
\texttt{exists} och \texttt{renameTo}. Genom att använda File-klassen
kan du enkelt utföra filrelaterade operationer i dina Java-program.

\hypertarget{beskrivning}{%
\subsection{Beskrivning}\label{beskrivning}}

File-klassen innehåller ett brett utbud av användbara metoder för
filhantering. Här är några av de vanligaste metoderna och deras
funktioner:

\begin{itemize}
\tightlist
\item
  \texttt{createNewFile}: Skapar en ny fil på den angivna sökvägen.
\item
  \texttt{delete}: Tar bort en befintlig fil.
\item
  \texttt{exists}: Kontrollerar om en fil existerar.
\item
  \texttt{renameTo}: Byter namn på en fil.
\item
  \texttt{length}: Returnerar storleken på en fil i bytes.
\item
  \texttt{canRead}: Kontrollerar om en fil kan läsas.
\item
  \texttt{canWrite}: Kontrollerar om en fil kan skrivas till.
\item
  \texttt{canExecute}: Kontrollerar om en fil kan köras som ett program.
\item
  \texttt{isFile}: Kontrollerar om en fil är en vanlig fil.
\item
  \texttt{isDirectory}: Kontrollerar om en fil är en mapp.
\item
  \texttt{getParent}: Returnerar sökvägen till filens överordnade mapp.
\item
  \texttt{getName}: Returnerar filens namn.
\item
  \texttt{getPath}: Returnerar filens sökväg.
\item
  \texttt{lastModified}: Returnerar tidpunkten för senaste ändring av en
  fil.
\item
  \texttt{list}: Returnerar en array av filer och mappar i en mapp.
\item
  \texttt{mkdir}: Skapar en ny mapp på den angivna sökvägen.
\item
  \texttt{listFiles}: Returnerar en array av \texttt{File}-objekt som
  representerar filer och mappar i en mapp.
\item
  \texttt{FileReader}: Öppnar en befintlig textfil för läsning.
\item
  \texttt{FileWriter}: Öppnar en fil för skrivning.
\item
  \texttt{Files.readAllBytes}: Läser innehållet i en binär fil som en
  byte-array.
\item
  \texttt{Files.readAllLines}: Läser innehållet i en textfil och
  returnerar det som en array av strängar.
\item
  \texttt{Files.readAllText}: Läser innehållet i en textfil och
  returnerar det som en sträng.
\item
  \texttt{Files.lines}: Läser innehållet i en textfil och returnerar det
  som en uppräkningsbar sekvens av strängar.
\item
  \texttt{Files.writeAllBytes}: Skriver en byte-array till en fil.
\item
  \texttt{Files.writeAllLines}: Skriver en array av strängar till en
  textfil, en sträng per rad.
\item
  \texttt{Files.writeAllText}: Skriver en sträng till en textfil.
\end{itemize}

Dessa metoder ger oss flexibilitet att utföra olika åtgärder på filer,
inklusive skapande, läsning, skrivning, flyttning och radering.

\hypertarget{exempel}{%
\subsection{Exempel}\label{exempel}}

Här är några kodexempel som visar hur man använder några av metoderna i
File-klassen:

\hypertarget{cb1}{}
\begin{Shaded}
\begin{Highlighting}[]
\KeywordTok{import} \ImportTok{java}\OperatorTok{.}\ImportTok{io}\OperatorTok{.}\ImportTok{IOException}\OperatorTok{;}
\KeywordTok{import} \ImportTok{java}\OperatorTok{.}\ImportTok{nio}\OperatorTok{.}\ImportTok{file}\OperatorTok{.}\ImportTok{Files}\OperatorTok{;}
\KeywordTok{import} \ImportTok{java}\OperatorTok{.}\ImportTok{nio}\OperatorTok{.}\ImportTok{file}\OperatorTok{.}\ImportTok{Path}\OperatorTok{;}
\KeywordTok{import} \ImportTok{java}\OperatorTok{.}\ImportTok{nio}\OperatorTok{.}\ImportTok{file}\OperatorTok{.}\ImportTok{Paths}\OperatorTok{;}

\KeywordTok{public} \KeywordTok{class}\NormalTok{ Main }\OperatorTok{\{}
    \KeywordTok{public} \DataTypeTok{static} \DataTypeTok{void} \FunctionTok{main}\OperatorTok{(}\BuiltInTok{String}\OperatorTok{[]}\NormalTok{ args}\OperatorTok{)} \OperatorTok{\{}
        \BuiltInTok{String}\NormalTok{ file }\OperatorTok{=} \StringTok{"C:}\SpecialCharTok{\textbackslash{}\textbackslash{}}\StringTok{Temp}\SpecialCharTok{\textbackslash{}\textbackslash{}}\StringTok{test.txt"}\OperatorTok{;}

        \CommentTok{// Skapar en fil}
        \ControlFlowTok{try} \OperatorTok{\{}
            \ControlFlowTok{if} \OperatorTok{(!}\NormalTok{Files}\OperatorTok{.}\FunctionTok{exists}\OperatorTok{(}\NormalTok{Paths}\OperatorTok{.}\FunctionTok{get}\OperatorTok{(}\NormalTok{file}\OperatorTok{)))} \OperatorTok{\{}
\NormalTok{                Files}\OperatorTok{.}\FunctionTok{createFile}\OperatorTok{(}\NormalTok{Paths}\OperatorTok{.}\FunctionTok{get}\OperatorTok{(}\NormalTok{file}\OperatorTok{));}
            \OperatorTok{\}}
        \OperatorTok{\}} \ControlFlowTok{catch} \OperatorTok{(}\BuiltInTok{IOException}\NormalTok{ e}\OperatorTok{)} \OperatorTok{\{}
\NormalTok{            e}\OperatorTok{.}\FunctionTok{printStackTrace}\OperatorTok{();}
        \OperatorTok{\}}

        \CommentTok{// Skriver till en fil}
        \ControlFlowTok{try} \OperatorTok{\{}
\NormalTok{            Files}\OperatorTok{.}\FunctionTok{write}\OperatorTok{(}\NormalTok{Paths}\OperatorTok{.}\FunctionTok{get}\OperatorTok{(}\NormalTok{file}\OperatorTok{),} \StringTok{"Hello World!"}\OperatorTok{.}\FunctionTok{getBytes}\OperatorTok{());}
        \OperatorTok{\}} \ControlFlowTok{catch} \OperatorTok{(}\BuiltInTok{IOException}\NormalTok{ e}\OperatorTok{)} \OperatorTok{\{}
\NormalTok{            e}\OperatorTok{.}\FunctionTok{printStackTrace}\OperatorTok{();}
        \OperatorTok{\}}

        \CommentTok{// Läser innehållet i en fil}
        \ControlFlowTok{try} \OperatorTok{\{}
            \BuiltInTok{String}\NormalTok{ text }\OperatorTok{=}\NormalTok{ Files}\OperatorTok{.}\FunctionTok{readString}\OperatorTok{(}\NormalTok{Paths}\OperatorTok{.}\FunctionTok{get}\OperatorTok{(}\NormalTok{file}\OperatorTok{));}
            \BuiltInTok{System}\OperatorTok{.}\FunctionTok{out}\OperatorTok{.}\FunctionTok{println}\OperatorTok{(}\NormalTok{text}\OperatorTok{);}
        \OperatorTok{\}} \ControlFlowTok{catch} \OperatorTok{(}\BuiltInTok{IOException}\NormalTok{ e}\OperatorTok{)} \OperatorTok{\{}
\NormalTok{            e}\OperatorTok{.}\FunctionTok{printStackTrace}\OperatorTok{();}
        \OperatorTok{\}}

        \CommentTok{// Tar bort en fil}
        \ControlFlowTok{try} \OperatorTok{\{}
\NormalTok{            Files}\OperatorTok{.}\FunctionTok{delete}\OperatorTok{(}\NormalTok{Paths}\OperatorTok{.}\FunctionTok{get}\OperatorTok{(}\NormalTok{file}\OperatorTok{));}
        \OperatorTok{\}} \ControlFlowTok{catch} \OperatorTok{(}\BuiltInTok{IOException}\NormalTok{ e}\OperatorTok{)} \OperatorTok{\{}
\NormalTok{            e}\OperatorTok{.}\FunctionTok{printStackTrace}\OperatorTok{();}
        \OperatorTok{\}}
    \OperatorTok{\}}
\OperatorTok{\}}
\end{Highlighting}
\end{Shaded}

I detta exempel skapar vi en fil på sökvägen
``C:\textbackslash Temp\textbackslash test.txt'' om den inte redan
finns. Sedan använder vi metoden \texttt{Files.write} för att skriva
texten ``Hello World!'' till filen. Vi använder
\texttt{Files.readString} för att läsa in innehållet i filen och skriva
ut det till konsolen. Slutligen tar vi bort filen med
\texttt{Files.delete}.

Det finns många fler metoder i File-klassen som kan vara användbara
beroende på dina specifika behov. Genom att utforska dokumentationen för
File-klassen kan du lära dig mer om var och en av metoderna och hur de
kan användas.

\hypertarget{slutsats}{%
\subsection{Slutsats}\label{slutsats}}

File-klassen i Java erbjuder enkla och kraftfulla metoder för att
hantera filer. Oavsett om du behöver skapa, skriva, läsa eller ta bort
filer finns det en metod i File-klassen som kan hjälpa dig. Genom att
använda de olika metoderna kan du effektivt arbeta med filsystemet i
dina Java-applikationer.

\hypertarget{termer}{%
\subsection{Termer}\label{termer}}

\begin{itemize}
\tightlist
\item
  \textbf{File}: En klass i \texttt{java.io}-paketet som innehåller
  metoder för att hantera filer.
\item
  \textbf{appendText}: En metod i File-klassen som öppnar en befintlig
  fil för skrivning och lägger till text i slutet av filen.
\item
  \textbf{copy}: En metod i File-klassen som kopierar en fil till en ny
  plats.
\item
  \textbf{create}: En metod i File-klassen som skapar en ny fil.
\item
  \textbf{delete}: En metod i File-klassen som tar bort en befintlig
  fil.
\item
  \textbf{exists}: En metod i File-klassen som kontrollerar om en fil
  existerar.
\item
  \textbf{move}: En metod i File-klassen som flyttar en fil till en ny
  plats.
\item
  \textbf{open}: En metod i Java File-klassen som öppnar en fil i en
  \texttt{FileInputStream}.
\item
  \textbf{openRead}: Det finns ingen direkt motsvarighet till
  File-klassens \texttt{OpenRead}-metod i Java. Istället kan du använda
  \texttt{FileInputStream}-klassen för att öppna en fil för läsning.
\item
  \textbf{openText}: Det finns ingen direkt motsvarighet till
  File-klassens \texttt{OpenText}-metod i Java. Istället kan du använda
  \texttt{FileReader}-klassen för att öppna en befintlig textfil för
  läsning.
\item
  \textbf{openWrite}: En metod i Java File-klassen som öppnar en fil för
  skrivning med hjälp av \texttt{FileOutputStream}.
\item
  \textbf{readAllBytes}: En metod i Java File-klassen som läser
  innehållet i en binär fil som en byte-array med hjälp av
  \texttt{FileInputStream}.
\item
  \textbf{readAllLines}: En metod i Java File-klassen som läser
  innehållet i en textfil och returnerar det som en array av strängar.
  Du kan använda \texttt{BufferedReader}-klassen tillsammans med
  \texttt{FileReader}-klassen för att uppnå detta.
\item
  \textbf{readAllText}: En metod i Java File-klassen som läser
  innehållet i en textfil och returnerar det som en sträng med hjälp av
  \texttt{FileReader}.
\item
  \textbf{readLines}: En metod i Java File-klassen som läser innehållet
  i en textfil och returnerar det som en uppräkningsbar sekvens av
  strängar med hjälp av \texttt{BufferedReader} och \texttt{FileReader}.
\item
  \textbf{writeAllBytes}: En metod i Java File-klassen som skriver en
  byte-array till en fil med hjälp av \texttt{FileOutputStream}.
\item
  \textbf{writeAllLines}: En metod i Java File-klassen som skriver en
  lista av strängar till en textfil, en sträng per rad. Du kan använda
  \texttt{BufferedWriter} tillsammans med \texttt{FileWriter}-klassen
  för att uppnå detta.
\item
  \textbf{writeAllText}: En metod i Java File-klassen som skriver en
  sträng till en textfil.
\end{itemize}

\hypertarget{json}{%
\section{Json}\label{json}}

JSON är ett textbaserat filformat som används för att lagra och överföra
data på ett strukturerat sätt. Det är enklare att läsa och skriva än XML
och används ofta inom webbapplikationer. JSON används för att
representera data i form av objekt och arrayer, och det är ett populärt
filformat som har stöd i olika programmeringsspråk och plattformar.

\hypertarget{gson-biblioteket}{%
\subsection{Gson-biblioteket}\label{gson-biblioteket}}

I Java finns det inbyggda verktyg och bibliotek för att arbeta med
JSON-data. En av de mest populära biblioteken är Gson, som
tillhandahåller funktioner för att serialisera och deserialisera
JSON-data till Java-objekt och vice versa. Gson gör det enkelt att
konvertera en JSON-sträng till ett Java-objekt och att generera en
JSON-sträng från ett Java-objekt.

För att använda Gson-biblioteket behöver du lägga till beroendet i ditt
projekt. Du kan göra det genom att lägga till följande kod i din
pom.xml-fil om du använder Maven:

\hypertarget{cb1}{}
\begin{Shaded}
\begin{Highlighting}[]
\NormalTok{\textless{}}\KeywordTok{dependencies}\NormalTok{\textgreater{}}
\NormalTok{    \textless{}}\KeywordTok{dependency}\NormalTok{\textgreater{}}
\NormalTok{        \textless{}}\KeywordTok{groupId}\NormalTok{\textgreater{}com.google.code.gson\textless{}/}\KeywordTok{groupId}\NormalTok{\textgreater{}}
\NormalTok{        \textless{}}\KeywordTok{artifactId}\NormalTok{\textgreater{}gson\textless{}/}\KeywordTok{artifactId}\NormalTok{\textgreater{}}
\NormalTok{        \textless{}}\KeywordTok{version}\NormalTok{\textgreater{}2.8.7\textless{}/}\KeywordTok{version}\NormalTok{\textgreater{}}
\NormalTok{    \textless{}/}\KeywordTok{dependency}\NormalTok{\textgreater{}}
\NormalTok{\textless{}/}\KeywordTok{dependencies}\NormalTok{\textgreater{}}
\end{Highlighting}
\end{Shaded}

När du har lagt till beroendet kan du använda Gson för att serialisera
och deserialisera JSON-data. Här är några exempel på hur du kan använda
Gson:

\hypertarget{cb2}{}
\begin{Shaded}
\begin{Highlighting}[]
\KeywordTok{import} \ImportTok{com}\OperatorTok{.}\ImportTok{google}\OperatorTok{.}\ImportTok{gson}\OperatorTok{.}\ImportTok{Gson}\OperatorTok{;}

\KeywordTok{public} \KeywordTok{class}\NormalTok{ Main }\OperatorTok{\{}
    \KeywordTok{public} \DataTypeTok{static} \DataTypeTok{void} \FunctionTok{main}\OperatorTok{(}\BuiltInTok{String}\OperatorTok{[]}\NormalTok{ args}\OperatorTok{)} \OperatorTok{\{}
        \CommentTok{// Skapa en instans av Gson}
\NormalTok{        Gson gson }\OperatorTok{=} \KeywordTok{new} \FunctionTok{Gson}\OperatorTok{();}

        \CommentTok{// Serialisera ett Java{-}objekt till en JSON{-}sträng}
\NormalTok{        Person person }\OperatorTok{=} \KeywordTok{new} \FunctionTok{Person}\OperatorTok{(}\StringTok{"John Doe"}\OperatorTok{,} \DecValTok{30}\OperatorTok{);}
        \BuiltInTok{String}\NormalTok{ json }\OperatorTok{=}\NormalTok{ gson}\OperatorTok{.}\FunctionTok{toJson}\OperatorTok{(}\NormalTok{person}\OperatorTok{);}
        \BuiltInTok{System}\OperatorTok{.}\FunctionTok{out}\OperatorTok{.}\FunctionTok{println}\OperatorTok{(}\NormalTok{json}\OperatorTok{);} \CommentTok{// \{"name":"John Doe","age":30\}}

        \CommentTok{// Deserialisera en JSON{-}sträng till ett Java{-}objekt}
        \BuiltInTok{String}\NormalTok{ json2 }\OperatorTok{=} \StringTok{"\{}\SpecialCharTok{\textbackslash{}"}\StringTok{name}\SpecialCharTok{\textbackslash{}"}\StringTok{:}\SpecialCharTok{\textbackslash{}"}\StringTok{Jane Smith}\SpecialCharTok{\textbackslash{}"}\StringTok{,}\SpecialCharTok{\textbackslash{}"}\StringTok{age}\SpecialCharTok{\textbackslash{}"}\StringTok{:25\}"}\OperatorTok{;}
\NormalTok{        Person person2 }\OperatorTok{=}\NormalTok{ gson}\OperatorTok{.}\FunctionTok{fromJson}\OperatorTok{(}\NormalTok{json2}\OperatorTok{,}\NormalTok{ Person}\OperatorTok{.}\FunctionTok{class}\OperatorTok{);}
        \BuiltInTok{System}\OperatorTok{.}\FunctionTok{out}\OperatorTok{.}\FunctionTok{println}\OperatorTok{(}\NormalTok{person2}\OperatorTok{.}\FunctionTok{getName}\OperatorTok{());} \CommentTok{// Jane Smith}
        \BuiltInTok{System}\OperatorTok{.}\FunctionTok{out}\OperatorTok{.}\FunctionTok{println}\OperatorTok{(}\NormalTok{person2}\OperatorTok{.}\FunctionTok{getAge}\OperatorTok{());} \CommentTok{// 25}
    \OperatorTok{\}}
\OperatorTok{\}}

\KeywordTok{class}\NormalTok{ Person }\OperatorTok{\{}
    \KeywordTok{private} \BuiltInTok{String}\NormalTok{ name}\OperatorTok{;}
    \KeywordTok{private} \DataTypeTok{int}\NormalTok{ age}\OperatorTok{;}

    \KeywordTok{public} \FunctionTok{Person}\OperatorTok{(}\BuiltInTok{String}\NormalTok{ name}\OperatorTok{,} \DataTypeTok{int}\NormalTok{ age}\OperatorTok{)} \OperatorTok{\{}
        \KeywordTok{this}\OperatorTok{.}\FunctionTok{name} \OperatorTok{=}\NormalTok{ name}\OperatorTok{;}
        \KeywordTok{this}\OperatorTok{.}\FunctionTok{age} \OperatorTok{=}\NormalTok{ age}\OperatorTok{;}
    \OperatorTok{\}}

    \KeywordTok{public} \BuiltInTok{String} \FunctionTok{getName}\OperatorTok{()} \OperatorTok{\{}
        \ControlFlowTok{return}\NormalTok{ name}\OperatorTok{;}
    \OperatorTok{\}}

    \KeywordTok{public} \DataTypeTok{int} \FunctionTok{getAge}\OperatorTok{()} \OperatorTok{\{}
        \ControlFlowTok{return}\NormalTok{ age}\OperatorTok{;}
    \OperatorTok{\}}
\OperatorTok{\}}
\end{Highlighting}
\end{Shaded}

I exemplet ovan använder vi Gson för att serialisera en Person-objekt
till en JSON-sträng och deserialisera en JSON-sträng till ett
Person-objekt. Gson använder Java Reflection för att konvertera objekt
mellan Java och JSON och det gör det enkelt att hantera komplexa
datastrukturer.

\hypertarget{skapa-json-manuellt-med-stringbuilder}{%
\subsection{Skapa JSON-manuellt med
StringBuilder}\label{skapa-json-manuellt-med-stringbuilder}}

Förutom att använda bibliotek kan vi också skapa en JSON-sträng manuellt
med hjälp av StringBuilder-klassen. Det här kan vara användbart om du
behöver bygga en JSON-sträng dynamiskt eller om du inte vill använda ett
tredjepartsbibliotek för att hantera JSON.

Här är ett exempel på hur du kan använda StringBuilder för att skapa en
JSON-sträng som representerar en Star Wars-karaktär:

\hypertarget{cb3}{}
\begin{Shaded}
\begin{Highlighting}[]
\BuiltInTok{StringBuilder}\NormalTok{ sb }\OperatorTok{=} \KeywordTok{new} \BuiltInTok{StringBuilder}\OperatorTok{();}
\NormalTok{sb}\OperatorTok{.}\FunctionTok{append}\OperatorTok{(}\StringTok{"\{"}\OperatorTok{);}
\NormalTok{sb}\OperatorTok{.}\FunctionTok{append}\OperatorTok{(}\StringTok{"}\SpecialCharTok{\textbackslash{}"}\StringTok{name}\SpecialCharTok{\textbackslash{}"}\StringTok{: }\SpecialCharTok{\textbackslash{}"}\StringTok{Luke Skywalker}\SpecialCharTok{\textbackslash{}"}\StringTok{,"}\OperatorTok{);}
\NormalTok{sb}\OperatorTok{.}\FunctionTok{append}\OperatorTok{(}\StringTok{"}\SpecialCharTok{\textbackslash{}"}\StringTok{age}\SpecialCharTok{\textbackslash{}"}\StringTok{: 30"}\OperatorTok{);}
\NormalTok{sb}\OperatorTok{.}\FunctionTok{append}\OperatorTok{(}\StringTok{"\}"}\OperatorTok{);}
\BuiltInTok{String}\NormalTok{ json }\OperatorTok{=}\NormalTok{ sb}\OperatorTok{.}\FunctionTok{toString}\OperatorTok{();}
\end{Highlighting}
\end{Shaded}

I det här exemplet skapar vi en StringBuilder-instans och använder dess
\texttt{append()}-metod för att bygga upp JSON-strängen steg för steg.
Till slut använder vi \texttt{toString()}-metoden för att få den färdiga
JSON-strängen.

\hypertarget{luxe4sa-json-data}{%
\subsection{Läsa JSON-data}\label{luxe4sa-json-data}}

För att läsa JSON-data kan du använda Gson-biblioteket eller så kan du
använda Java's inbyggda verktyg för JSON-hantering. Här är ett exempel
på hur du kan läsa JSON-data utan att använda Gson:

\hypertarget{cb4}{}
\begin{Shaded}
\begin{Highlighting}[]
\KeywordTok{import} \ImportTok{javax}\OperatorTok{.}\ImportTok{json}\OperatorTok{.}\ImportTok{Json}\OperatorTok{;}
\KeywordTok{import} \ImportTok{javax}\OperatorTok{.}\ImportTok{json}\OperatorTok{.}\ImportTok{JsonObject}\OperatorTok{;}
\KeywordTok{import} \ImportTok{javax}\OperatorTok{.}\ImportTok{json}\OperatorTok{.}\ImportTok{JsonReader}\OperatorTok{;}
\KeywordTok{import} \ImportTok{java}\OperatorTok{.}\ImportTok{io}\OperatorTok{.}\ImportTok{StringReader}\OperatorTok{;}

\BuiltInTok{String}\NormalTok{ json }\OperatorTok{=} \StringTok{"\{}\SpecialCharTok{\textbackslash{}"}\StringTok{name}\SpecialCharTok{\textbackslash{}"}\StringTok{:}\SpecialCharTok{\textbackslash{}"}\StringTok{Luke Skywalker}\SpecialCharTok{\textbackslash{}"}\StringTok{,}\SpecialCharTok{\textbackslash{}"}\StringTok{age}\SpecialCharTok{\textbackslash{}"}\StringTok{:30\}"}\OperatorTok{;}

\NormalTok{JsonReader reader }\OperatorTok{=}\NormalTok{ Json}\OperatorTok{.}\FunctionTok{createReader}\OperatorTok{(}\KeywordTok{new} \BuiltInTok{StringReader}\OperatorTok{(}\NormalTok{json}\OperatorTok{));}
\NormalTok{JsonObject obj }\OperatorTok{=}\NormalTok{ reader}\OperatorTok{.}\FunctionTok{readObject}\OperatorTok{();}
\BuiltInTok{String}\NormalTok{ name }\OperatorTok{=}\NormalTok{ obj}\OperatorTok{.}\FunctionTok{getString}\OperatorTok{(}\StringTok{"name"}\OperatorTok{);}
\DataTypeTok{int}\NormalTok{ age }\OperatorTok{=}\NormalTok{ obj}\OperatorTok{.}\FunctionTok{getInt}\OperatorTok{(}\StringTok{"age"}\OperatorTok{);}

\BuiltInTok{System}\OperatorTok{.}\FunctionTok{out}\OperatorTok{.}\FunctionTok{println}\OperatorTok{(}\NormalTok{name}\OperatorTok{);} \CommentTok{// Luke Skywalker}
\BuiltInTok{System}\OperatorTok{.}\FunctionTok{out}\OperatorTok{.}\FunctionTok{println}\OperatorTok{(}\NormalTok{age}\OperatorTok{);} \CommentTok{// 30}
\end{Highlighting}
\end{Shaded}

I det här exemplet använder vi \texttt{javax.json}-paketet för att läsa
JSON-data. Vi skapar en JsonReader-instans och använder den för att läsa
in JSON-data från en sträng. Sedan använder vi JsonObject för att få
tillgång till de olika egenskaperna i JSON-objektet.

\hypertarget{fuxf6rdelar-och-nackdelar-med-json}{%
\subsection{Fördelar och nackdelar med
JSON}\label{fuxf6rdelar-och-nackdelar-med-json}}

Nu när du har lärt dig grunderna i JSON, låt oss titta på några fördelar
och nackdelar med att använda JSON:

\hypertarget{fuxf6rdelar}{%
\subsubsection{Fördelar:}\label{fuxf6rdelar}}

\begin{itemize}
\tightlist
\item
  JSON är lätt att läsa och skriva för både människor och maskiner.
\item
  Det är ett populärt filformat och har stort stöd i olika
  programmeringsspråk och plattformar.
\item
  JSON-data kan enkelt konverteras till objekt och vice versa, vilket
  gör det lätt att arbeta med.
\end{itemize}

\hypertarget{nackdelar}{%
\subsubsection{Nackdelar:}\label{nackdelar}}

\begin{itemize}
\tightlist
\item
  JSON saknar några av de mer avancerade funktionerna som XML erbjuder.
\item
  Det kan vara mindre lämpligt för strukturerade dokument med
  hierarkiska datastrukturer.
\item
  JSON kan bli svårt att hantera om filerna blir mycket stora.
\end{itemize}

\hypertarget{varfuxf6r-anvuxe4nda-gson-istuxe4llet-fuxf6r-org.json}{%
\subsection{Varför använda Gson istället för
org.json?}\label{varfuxf6r-anvuxe4nda-gson-istuxe4llet-fuxf6r-org.json}}

Du kanske undrar varför jag rekommenderar Gson istället för org.json.
Det finns flera skäl till detta:

\begin{itemize}
\tightlist
\item
  Gson har en mer intuitiv och enkel API-design.
\item
  Gson har bättre prestanda och lägre minnesanvändning än org.json.
\item
  Gson har stöd för avancerad typomvandling och anpassning av
  serialisering/deserialisering.
\item
  Gson har en aktiv och stor användarbas och har mer dokumentation och
  exempel att använda som referens.
\end{itemize}

\hypertarget{termer-och-fuxf6rklaringar}{%
\subsection{Termer och förklaringar}\label{termer-och-fuxf6rklaringar}}

Här är en tabell över termer som använts i denna artikel och deras
förklaringar:

\begin{longtable}[]{@{}
  >{\raggedright\arraybackslash}p{(\columnwidth - 2\tabcolsep) * \real{0.1200}}
  >{\raggedright\arraybackslash}p{(\columnwidth - 2\tabcolsep) * \real{0.8700}}@{}}
\toprule\noalign{}
\begin{minipage}[b]{\linewidth}\raggedright
Term
\end{minipage} & \begin{minipage}[b]{\linewidth}\raggedright
Förklaring
\end{minipage} \\
\midrule\noalign{}
\endhead
\bottomrule\noalign{}
\endlastfoot
JSON & JavaScript Object Notation, ett textbaserat filformat för att
lagra och transportera data \\
Serialisera & Konvertera Java-objekt till JSON-format \\
Deserialisera & Konvertera JSON till Java-objekt \\
Gson & Ett populärt bibliotek för JSON-hantering i Java \\
StringBuilder & En klass i Java för att bygga upp textsträngar
dynamiskt \\
javax.json & Ett paket i Java för JSON-hantering \\
org.json & Ett bibliotek för JSON-hantering i Java \\
\end{longtable}

För att fördjupa dina kunskaper och utforska mer om JSON och
Java-programmering rekommenderar jag följande resurser:

\begin{itemize}
\tightlist
\item
  \href{https://json.org/}{JSON.org}: Den officiella webbplatsen för
  JSON med massor av information och exempel.
\item
  \href{https://github.com/google/gson/blob/master/README.md}{Gson
  Dokumentation}: Officiell dokumentation för Gson-biblioteket, ett
  populärt bibliotek för att hantera JSON i Java.
\item
  \href{https://docs.oracle.com/javase/tutorial/}{Oracle Java
  Tutorials}: Officiella Java-tutorials från Oracle för att lära dig mer
  om Java-programmering.
\end{itemize}

\hypertarget{obligatorisk-dad-joke}{%
\subsubsection{Obligatorisk Dad Joke}\label{obligatorisk-dad-joke}}

Varför gick JSON till terapeuten? För att det hade problem med att parsa
sina känslor!

\hypertarget{mappar}{%
\section{Mappar}\label{mappar}}

I detta kodexempel använder vi klassen \texttt{File} från
\texttt{java.io}-paketet för att kontrollera, skapa och radera mappar i
Java-programmet.

\hypertarget{cb1}{}
\begin{Shaded}
\begin{Highlighting}[]
\KeywordTok{import} \ImportTok{java}\OperatorTok{.}\ImportTok{io}\OperatorTok{.}\ImportTok{File}\OperatorTok{;}

\KeywordTok{public} \KeywordTok{class}\NormalTok{ MapparExempel }\OperatorTok{\{}
    \KeywordTok{public} \DataTypeTok{static} \DataTypeTok{void} \FunctionTok{main}\OperatorTok{(}\BuiltInTok{String}\OperatorTok{[]}\NormalTok{ args}\OperatorTok{)} \OperatorTok{\{}
        \CommentTok{// Kontrollera om en mapp finns}
        \BuiltInTok{String}\NormalTok{ mappSokvag }\OperatorTok{=} \StringTok{"C:}\SpecialCharTok{\textbackslash{}\textbackslash{}}\StringTok{exempel}\SpecialCharTok{\textbackslash{}\textbackslash{}}\StringTok{mapp"}\OperatorTok{;}
        \BuiltInTok{File}\NormalTok{ mapp }\OperatorTok{=} \KeywordTok{new} \BuiltInTok{File}\OperatorTok{(}\NormalTok{mappSokvag}\OperatorTok{);}
        \ControlFlowTok{if} \OperatorTok{(}\NormalTok{mapp}\OperatorTok{.}\FunctionTok{exists}\OperatorTok{())} \OperatorTok{\{}
            \BuiltInTok{System}\OperatorTok{.}\FunctionTok{out}\OperatorTok{.}\FunctionTok{println}\OperatorTok{(}\StringTok{"Mappen finns."}\OperatorTok{);}
        \OperatorTok{\}} \ControlFlowTok{else} \OperatorTok{\{}
            \BuiltInTok{System}\OperatorTok{.}\FunctionTok{out}\OperatorTok{.}\FunctionTok{println}\OperatorTok{(}\StringTok{"Mappen finns inte."}\OperatorTok{);}
        \OperatorTok{\}}

        \CommentTok{// Skapa en mapp}
        \BuiltInTok{String}\NormalTok{ nyMappSokvag }\OperatorTok{=} \StringTok{"C:}\SpecialCharTok{\textbackslash{}\textbackslash{}}\StringTok{exempel}\SpecialCharTok{\textbackslash{}\textbackslash{}}\StringTok{ny\_mapp"}\OperatorTok{;}
        \BuiltInTok{File}\NormalTok{ nyMapp }\OperatorTok{=} \KeywordTok{new} \BuiltInTok{File}\OperatorTok{(}\NormalTok{nyMappSokvag}\OperatorTok{);}
        \ControlFlowTok{if} \OperatorTok{(!}\NormalTok{nyMapp}\OperatorTok{.}\FunctionTok{exists}\OperatorTok{())} \OperatorTok{\{}
            \ControlFlowTok{if} \OperatorTok{(}\NormalTok{nyMapp}\OperatorTok{.}\FunctionTok{mkdir}\OperatorTok{())} \OperatorTok{\{}
                \BuiltInTok{System}\OperatorTok{.}\FunctionTok{out}\OperatorTok{.}\FunctionTok{println}\OperatorTok{(}\StringTok{"Mappen har skapats."}\OperatorTok{);}
            \OperatorTok{\}} \ControlFlowTok{else} \OperatorTok{\{}
                \BuiltInTok{System}\OperatorTok{.}\FunctionTok{out}\OperatorTok{.}\FunctionTok{println}\OperatorTok{(}\StringTok{"Kunde inte skapa mappen."}\OperatorTok{);}
            \OperatorTok{\}}
        \OperatorTok{\}} \ControlFlowTok{else} \OperatorTok{\{}
            \BuiltInTok{System}\OperatorTok{.}\FunctionTok{out}\OperatorTok{.}\FunctionTok{println}\OperatorTok{(}\StringTok{"Mappen finns redan."}\OperatorTok{);}
        \OperatorTok{\}}

        \CommentTok{// Radera en mapp}
        \BuiltInTok{String}\NormalTok{ raderaMappSokvag }\OperatorTok{=} \StringTok{"C:}\SpecialCharTok{\textbackslash{}\textbackslash{}}\StringTok{exempel}\SpecialCharTok{\textbackslash{}\textbackslash{}}\StringTok{att\_radera"}\OperatorTok{;}
        \BuiltInTok{File}\NormalTok{ raderaMapp }\OperatorTok{=} \KeywordTok{new} \BuiltInTok{File}\OperatorTok{(}\NormalTok{raderaMappSokvag}\OperatorTok{);}
        \ControlFlowTok{if} \OperatorTok{(}\NormalTok{raderaMapp}\OperatorTok{.}\FunctionTok{exists}\OperatorTok{())} \OperatorTok{\{}
            \ControlFlowTok{if} \OperatorTok{(}\NormalTok{raderaMapp}\OperatorTok{.}\FunctionTok{delete}\OperatorTok{())} \OperatorTok{\{}
                \BuiltInTok{System}\OperatorTok{.}\FunctionTok{out}\OperatorTok{.}\FunctionTok{println}\OperatorTok{(}\StringTok{"Mappen har raderats."}\OperatorTok{);}
            \OperatorTok{\}} \ControlFlowTok{else} \OperatorTok{\{}
                \BuiltInTok{System}\OperatorTok{.}\FunctionTok{out}\OperatorTok{.}\FunctionTok{println}\OperatorTok{(}\StringTok{"Kunde inte radera mappen."}\OperatorTok{);}
            \OperatorTok{\}}
        \OperatorTok{\}}
    \OperatorTok{\}}
\OperatorTok{\}}
\end{Highlighting}
\end{Shaded}

I detta kodexempel utför vi följande operationer:

\begin{enumerate}
\def\labelenumi{\arabic{enumi}.}
\tightlist
\item
  Kontrollerar om en mapp existerar genom att skapa en
  \texttt{File}-instans med sökvägen till mappen och använda
  \texttt{exists()}-metoden.
\item
  Skapar en ny mapp genom att skapa en \texttt{File}-instans med
  sökvägen till den nya mappen och använda \texttt{mkdir()}-metoden. Vi
  kontrollerar först om mappen redan finns med hjälp av
  \texttt{exists()}-metoden.
\item
  Raderar en mapp genom att skapa en \texttt{File}-instans med sökvägen
  till mappen och använda \texttt{delete()}-metoden. Vi kontrollerar
  först om mappen existerar med \texttt{exists()}-metoden.
\end{enumerate}

Det är viktigt att vara försiktig när man arbetar med mappar och filer
och att hantera eventuella fel eller undantag som kan uppstå. Det kan
vara användbart att använda try-catch-block för att fånga och hantera
eventuella undantag som genereras vid arbete med mappar och filer.

\hypertarget{sammanfattning}{%
\subsubsection{Sammanfattning}\label{sammanfattning}}

I detta kodexempel har vi sett hur man kan arbeta med mappar i Java
genom att använda \texttt{File}-klassen från \texttt{java.io}-paketet.
Vi har utfört operationer som att kontrollera om en mapp existerar,
skapa en ny mapp och radera en mapp. Genom att använda dessa metoder kan
vi enkelt hantera mappar i våra Java-program.

Det är viktigt att vara försiktig när man arbetar med mappar och filer
och att hantera eventuella fel eller undantag som kan uppstå. Det kan
vara användbart att använda try-catch-block för att fånga och hantera
eventuella undantag som genereras vid arbete med mappar och filer.

Jag hoppas att denna artikel har varit användbar och gett dig en bättre
förståelse för hur man arbetar med mappar i Java. Om du har några frågor
eller vill lära dig mer, tveka inte att utforska dokumentationen och
resurserna från Oracle som jag har länkat till ovan. Ha kul med att
organisera dina filer och mappar i dina Java-projekt!

\textbf{Källor}:

\begin{itemize}
\tightlist
\item
  \href{https://docs.oracle.com/en/java/javase/17/docs/api/java.base/java/io/File.html}{Oracle
  Docs - File Class}
\end{itemize}

\hypertarget{obligatorisk-dad-joke}{%
\subsection{Obligatorisk Dad-joke}\label{obligatorisk-dad-joke}}

Varför älskar mappar att gå på festivaler?

För att de alltid får chansen att ``mappa'' upp sig!

\hypertarget{text}{%
\section{Text}\label{text}}

Textfiler är en vanlig filtyp som används för att lagra text. Det är ett
textbaserat filformat som används för att lagra data av olika slag.

Innehållsförteckning

\{: .text-delta \} 1. Innehållsförteckning \{:toc\}

\hypertarget{skapa-en-textfil}{%
\subsection{Skapa en textfil}\label{skapa-en-textfil}}

För att skapa en textfil i Java kan följande kod användas:

\hypertarget{cb1}{}
\begin{Shaded}
\begin{Highlighting}[]
\BuiltInTok{String}\NormalTok{ contents }\OperatorTok{=} \StringTok{"God morgon Mr Bond! Jag har ett meddelande till dig."}\OperatorTok{;}
\BuiltInTok{String}\NormalTok{ fileName }\OperatorTok{=} \StringTok{"Message.txt"}\OperatorTok{;}

\ControlFlowTok{try} \OperatorTok{\{}
    \BuiltInTok{FileWriter}\NormalTok{ fileWriter }\OperatorTok{=} \KeywordTok{new} \BuiltInTok{FileWriter}\OperatorTok{(}\NormalTok{fileName}\OperatorTok{);}
    \BuiltInTok{BufferedWriter}\NormalTok{ bufferedWriter }\OperatorTok{=} \KeywordTok{new} \BuiltInTok{BufferedWriter}\OperatorTok{(}\NormalTok{fileWriter}\OperatorTok{);}
\NormalTok{    bufferedWriter}\OperatorTok{.}\FunctionTok{write}\OperatorTok{(}\NormalTok{contents}\OperatorTok{);}
\NormalTok{    bufferedWriter}\OperatorTok{.}\FunctionTok{close}\OperatorTok{();}
    \BuiltInTok{System}\OperatorTok{.}\FunctionTok{out}\OperatorTok{.}\FunctionTok{println}\OperatorTok{(}\StringTok{"Textfilen har skapats."}\OperatorTok{);}
\OperatorTok{\}} \ControlFlowTok{catch} \OperatorTok{(}\BuiltInTok{IOException}\NormalTok{ e}\OperatorTok{)} \OperatorTok{\{}
    \BuiltInTok{System}\OperatorTok{.}\FunctionTok{out}\OperatorTok{.}\FunctionTok{println}\OperatorTok{(}\StringTok{"Ett fel inträffade vid skapandet av textfilen."}\OperatorTok{);}
\OperatorTok{\}}
\end{Highlighting}
\end{Shaded}

\hypertarget{luxe4s-in-en-textfil}{%
\subsection{Läs in en textfil}\label{luxe4s-in-en-textfil}}

För att läsa innehållet i en textfil kan följande kod användas:

\hypertarget{cb2}{}
\begin{Shaded}
\begin{Highlighting}[]
\ControlFlowTok{try} \OperatorTok{\{}
    \BuiltInTok{File}\NormalTok{ file }\OperatorTok{=} \KeywordTok{new} \BuiltInTok{File}\OperatorTok{(}\StringTok{"Message.txt"}\OperatorTok{);}
    \BuiltInTok{Scanner}\NormalTok{ scanner }\OperatorTok{=} \KeywordTok{new} \BuiltInTok{Scanner}\OperatorTok{(}\NormalTok{file}\OperatorTok{);}
    \ControlFlowTok{while} \OperatorTok{(}\NormalTok{scanner}\OperatorTok{.}\FunctionTok{hasNextLine}\OperatorTok{())} \OperatorTok{\{}
        \BuiltInTok{String}\NormalTok{ contents }\OperatorTok{=}\NormalTok{ scanner}\OperatorTok{.}\FunctionTok{nextLine}\OperatorTok{();}
        \BuiltInTok{System}\OperatorTok{.}\FunctionTok{out}\OperatorTok{.}\FunctionTok{println}\OperatorTok{(}\NormalTok{contents}\OperatorTok{);}
    \OperatorTok{\}}
\NormalTok{    scanner}\OperatorTok{.}\FunctionTok{close}\OperatorTok{();}
\OperatorTok{\}} \ControlFlowTok{catch} \OperatorTok{(}\BuiltInTok{FileNotFoundException}\NormalTok{ e}\OperatorTok{)} \OperatorTok{\{}
    \BuiltInTok{System}\OperatorTok{.}\FunctionTok{out}\OperatorTok{.}\FunctionTok{println}\OperatorTok{(}\StringTok{"Textfilen kunde inte hittas."}\OperatorTok{);}
\OperatorTok{\}}
\end{Highlighting}
\end{Shaded}

\hypertarget{spara-en-lista-i-en-textfil}{%
\subsection{Spara en lista i en
textfil}\label{spara-en-lista-i-en-textfil}}

För att spara en lista av strängar i en textfil kan följande kod
användas:

\hypertarget{cb3}{}
\begin{Shaded}
\begin{Highlighting}[]
\BuiltInTok{List}\OperatorTok{\textless{}}\BuiltInTok{String}\OperatorTok{\textgreater{}}\NormalTok{ names }\OperatorTok{=} \KeywordTok{new} \BuiltInTok{ArrayList}\OperatorTok{\textless{}\textgreater{}();}
\NormalTok{names}\OperatorTok{.}\FunctionTok{add}\OperatorTok{(}\StringTok{"Picard"}\OperatorTok{);}
\NormalTok{names}\OperatorTok{.}\FunctionTok{add}\OperatorTok{(}\StringTok{"Janeway"}\OperatorTok{);}
\NormalTok{names}\OperatorTok{.}\FunctionTok{add}\OperatorTok{(}\StringTok{"Kirk"}\OperatorTok{);}
\NormalTok{names}\OperatorTok{.}\FunctionTok{add}\OperatorTok{(}\StringTok{"Sisko"}\OperatorTok{);}
\NormalTok{names}\OperatorTok{.}\FunctionTok{add}\OperatorTok{(}\StringTok{"Archer"}\OperatorTok{);}

\ControlFlowTok{try} \OperatorTok{\{}
    \BuiltInTok{FileWriter}\NormalTok{ fileWriter }\OperatorTok{=} \KeywordTok{new} \BuiltInTok{FileWriter}\OperatorTok{(}\StringTok{"names.txt"}\OperatorTok{);}
    \BuiltInTok{BufferedWriter}\NormalTok{ bufferedWriter }\OperatorTok{=} \KeywordTok{new} \BuiltInTok{BufferedWriter}\OperatorTok{(}\NormalTok{fileWriter}\OperatorTok{);}
    \ControlFlowTok{for} \OperatorTok{(}\BuiltInTok{String}\NormalTok{ name }\OperatorTok{:}\NormalTok{ names}\OperatorTok{)} \OperatorTok{\{}
\NormalTok{        bufferedWriter}\OperatorTok{.}\FunctionTok{write}\OperatorTok{(}\NormalTok{name}\OperatorTok{);}
\NormalTok{        bufferedWriter}\OperatorTok{.}\FunctionTok{newLine}\OperatorTok{();}
    \OperatorTok{\}}
\NormalTok{    bufferedWriter}\OperatorTok{.}\FunctionTok{close}\OperatorTok{();}
    \BuiltInTok{System}\OperatorTok{.}\FunctionTok{out}\OperatorTok{.}\FunctionTok{println}\OperatorTok{(}\StringTok{"Listan har sparats i textfilen."}\OperatorTok{);}
\OperatorTok{\}} \ControlFlowTok{catch} \OperatorTok{(}\BuiltInTok{IOException}\NormalTok{ e}\OperatorTok{)} \OperatorTok{\{}
    \BuiltInTok{System}\OperatorTok{.}\FunctionTok{out}\OperatorTok{.}\FunctionTok{println}\OperatorTok{(}\StringTok{"Ett fel inträffade vid sparandet av listan i textfilen."}\OperatorTok{);}
\OperatorTok{\}}
\end{Highlighting}
\end{Shaded}

\hypertarget{luxe4s-in-en-lista-fruxe5n-en-textfil}{%
\subsection{Läs in en lista från en
textfil}\label{luxe4s-in-en-lista-fruxe5n-en-textfil}}

För att läsa innehållet från en textfil och spara det i en lista kan
följande kod användas:

\hypertarget{cb4}{}
\begin{Shaded}
\begin{Highlighting}[]
\BuiltInTok{List}\OperatorTok{\textless{}}\BuiltInTok{String}\OperatorTok{\textgreater{}}\NormalTok{ names }\OperatorTok{=} \KeywordTok{new} \BuiltInTok{ArrayList}\OperatorTok{\textless{}\textgreater{}();}

\ControlFlowTok{try} \OperatorTok{\{}
    \BuiltInTok{File}\NormalTok{ file }\OperatorTok{=} \KeywordTok{new} \BuiltInTok{File}\OperatorTok{(}\StringTok{"names.txt"}\OperatorTok{);}
    \BuiltInTok{Scanner}\NormalTok{ scanner }\OperatorTok{=} \KeywordTok{new} \BuiltInTok{Scanner}\OperatorTok{(}\NormalTok{file}\OperatorTok{);}
    \ControlFlowTok{while} \OperatorTok{(}\NormalTok{scanner}\OperatorTok{.}\FunctionTok{hasNextLine}\OperatorTok{())} \OperatorTok{\{}
        \BuiltInTok{String}\NormalTok{ name }\OperatorTok{=}\NormalTok{ scanner}\OperatorTok{.}\FunctionTok{nextLine}\OperatorTok{();}
\NormalTok{        names}\OperatorTok{.}\FunctionTok{add}\OperatorTok{(}\NormalTok{name}\OperatorTok{);}
    \OperatorTok{\}}
\NormalTok{    scanner}\OperatorTok{.}\FunctionTok{close}\OperatorTok{();}
\OperatorTok{\}} \ControlFlowTok{catch} \OperatorTok{(}\BuiltInTok{FileNotFoundException}\NormalTok{ e}\OperatorTok{)} \OperatorTok{\{}
    \BuiltInTok{System}\OperatorTok{.}\FunctionTok{out}\OperatorTok{.}\FunctionTok{println}\OperatorTok{(}\StringTok{"Textfilen kunde inte hittas."}\OperatorTok{);}
\OperatorTok{\}}
\end{Highlighting}
\end{Shaded}

Detta exempel visar hur man kan skapa, läsa, spara och läsa in textfiler
i Java. Genom att använda \texttt{FileWriter}, \texttt{BufferedWriter},
\texttt{FileReader}, \texttt{BufferedReader} och \texttt{Scanner} kan vi
utföra olika operationer på textfiler. Listor kan också enkelt sparas
och läsas från textfiler genom att använda dessa metoder.

Det är viktigt att hantera eventuella undantag (exceptions) som kan
uppstå vid hantering av filer, till exempel om filen inte hittas eller
om det uppstår problem med filåtkomst. Genom att använda try-catch-block
kan vi fånga och hantera dessa undantag på ett säkert sätt.

Det är också viktigt att stänga filanslutningar när de inte längre
behövs för att frigöra resurser och undvika minnesläckor. I kodexemplen
ovan används \texttt{close()}-metoden för att stänga anslutningarna
efter att de har använts.

\hypertarget{luxe4gg-till-text-i-en-textfil}{%
\subsection{Lägg till text i en
textfil}\label{luxe4gg-till-text-i-en-textfil}}

För att lägga till text i en textfil kan följande kod användas:

\hypertarget{cb5}{}
\begin{Shaded}
\begin{Highlighting}[]
\BuiltInTok{String}\NormalTok{ contents }\OperatorTok{=} \StringTok{"James Bond, din uppgift är enkel men avgörande: Sök upp den fiktiva elakingen och eliminera honom utan nåd. Låt inget stå i vägen för rättvisa och säkerhet. Var din vanliga självsäkra och eleganta själv. Tiden är knapp, agera snabbt och precist. Världen litar på dig, 007. Gör det som behöver göras."}\OperatorTok{;}
\BuiltInTok{String}\NormalTok{ fileName }\OperatorTok{=} \StringTok{"Message.txt"}\OperatorTok{;}

\ControlFlowTok{try} \OperatorTok{\{}
    \BuiltInTok{FileWriter}\NormalTok{ fileWriter }\OperatorTok{=} \KeywordTok{new} \BuiltInTok{FileWriter}\OperatorTok{(}\NormalTok{fileName}\OperatorTok{,} \KeywordTok{true}\OperatorTok{);}
    \BuiltInTok{BufferedWriter}\NormalTok{ bufferedWriter }\OperatorTok{=} \KeywordTok{new} \BuiltInTok{BufferedWriter}\OperatorTok{(}\NormalTok{fileWriter}\OperatorTok{);}
\NormalTok{    bufferedWriter}\OperatorTok{.}\FunctionTok{newLine}\OperatorTok{();}
\NormalTok{    bufferedWriter}\OperatorTok{.}\FunctionTok{write}\OperatorTok{(}\NormalTok{contents}\OperatorTok{);}
\NormalTok{    bufferedWriter}\OperatorTok{.}\FunctionTok{close}\OperatorTok{();}
    \BuiltInTok{System}\OperatorTok{.}\FunctionTok{out}\OperatorTok{.}\FunctionTok{println}\OperatorTok{(}\StringTok{"Texten har lagts till i textfilen."}\OperatorTok{);}
\OperatorTok{\}} \ControlFlowTok{catch} \OperatorTok{(}\BuiltInTok{IOException}\NormalTok{ e}\OperatorTok{)} \OperatorTok{\{}
    \BuiltInTok{System}\OperatorTok{.}\FunctionTok{out}\OperatorTok{.}\FunctionTok{println}\OperatorTok{(}\StringTok{"Ett fel inträffade vid tillägg av texten i textfilen."}\OperatorTok{);}
\OperatorTok{\}}
\end{Highlighting}
\end{Shaded}

I FileWriter-konstruktorn kan vi ange \texttt{true} som andra argument
för att lägga till text i en befintlig textfil. Detta gör att vi kan
lägga till text i en befintlig textfil utan att skriva över den
befintliga texten. Hade vi skrivit false istället för true hade den
befintliga texten skrivits över.

\hypertarget{skriv-uxf6ver-en-textfil}{%
\subsection{Skriv över en textfil}\label{skriv-uxf6ver-en-textfil}}

För att skriva över en textfil kan följande kod användas:

\hypertarget{cb6}{}
\begin{Shaded}
\begin{Highlighting}[]
\BuiltInTok{String}\NormalTok{ message }\OperatorTok{=} \StringTok{"Message deducted."}\OperatorTok{;}
\BuiltInTok{String}\NormalTok{ fileName }\OperatorTok{=} \StringTok{"Message.txt"}\OperatorTok{;}

\ControlFlowTok{try} \OperatorTok{\{}
    \BuiltInTok{FileWriter}\NormalTok{ fileWriter }\OperatorTok{=} \KeywordTok{new} \BuiltInTok{FileWriter}\OperatorTok{(}\NormalTok{fileName}\OperatorTok{);}
    \BuiltInTok{BufferedWriter}\NormalTok{ bufferedWriter }\OperatorTok{=} \KeywordTok{new} \BuiltInTok{BufferedWriter}\OperatorTok{(}\NormalTok{fileWriter}\OperatorTok{);}
\NormalTok{    bufferedWriter}\OperatorTok{.}\FunctionTok{write}\OperatorTok{(}\NormalTok{message}\OperatorTok{);}
\NormalTok{    bufferedWriter}\OperatorTok{.}\FunctionTok{close}\OperatorTok{();}
    \BuiltInTok{System}\OperatorTok{.}\FunctionTok{out}\OperatorTok{.}\FunctionTok{println}\OperatorTok{(}\StringTok{"Textfilen har skrivits över."}\OperatorTok{);}
\OperatorTok{\}} \ControlFlowTok{catch} \OperatorTok{(}\BuiltInTok{IOException}\NormalTok{ e}\OperatorTok{)} \OperatorTok{\{}
    \BuiltInTok{System}\OperatorTok{.}\FunctionTok{out}\OperatorTok{.}\FunctionTok{println}\OperatorTok{(}\StringTok{"Ett fel inträffade vid skrivning över textfilen."}\OperatorTok{);}
\OperatorTok{\}}
\end{Highlighting}
\end{Shaded}

\hypertarget{radera-en-textfil}{%
\subsection{Radera en textfil}\label{radera-en-textfil}}

För att radera en textfil kan följande kod användas:

\hypertarget{cb7}{}
\begin{Shaded}
\begin{Highlighting}[]
\BuiltInTok{String}\NormalTok{ fileName }\OperatorTok{=} \StringTok{"Message.txt"}\OperatorTok{;}

\ControlFlowTok{try} \OperatorTok{\{}
    \BuiltInTok{File}\NormalTok{ file }\OperatorTok{=} \KeywordTok{new} \BuiltInTok{File}\OperatorTok{(}\NormalTok{fileName}\OperatorTok{);}
    \ControlFlowTok{if} \OperatorTok{(}\NormalTok{file}\OperatorTok{.}\FunctionTok{delete}\OperatorTok{())} \OperatorTok{\{}
        \BuiltInTok{System}\OperatorTok{.}\FunctionTok{out}\OperatorTok{.}\FunctionTok{println}\OperatorTok{(}\StringTok{"Textfilen har raderats."}\OperatorTok{);}
    \OperatorTok{\}} \ControlFlowTok{else} \OperatorTok{\{}
        \BuiltInTok{System}\OperatorTok{.}\FunctionTok{out}\OperatorTok{.}\FunctionTok{println}

\OperatorTok{(}\StringTok{"Kunde inte radera textfilen."}\OperatorTok{);}
    \OperatorTok{\}}
\OperatorTok{\}} \ControlFlowTok{catch} \OperatorTok{(}\BuiltInTok{Exception}\NormalTok{ e}\OperatorTok{)} \OperatorTok{\{}
    \BuiltInTok{System}\OperatorTok{.}\FunctionTok{out}\OperatorTok{.}\FunctionTok{println}\OperatorTok{(}\StringTok{"Ett fel inträffade vid radering av textfilen."}\OperatorTok{);}
\OperatorTok{\}}
\end{Highlighting}
\end{Shaded}

\hypertarget{luxe4s-in-en-textfil-fruxe5n-en-url}{%
\subsection{Läs in en textfil från en
URL}\label{luxe4s-in-en-textfil-fruxe5n-en-url}}

För att läsa innehållet i en textfil från en URL kan följande kod
användas:

\hypertarget{cb8}{}
\begin{Shaded}
\begin{Highlighting}[]
\ControlFlowTok{try} \OperatorTok{\{}
    \BuiltInTok{URL}\NormalTok{ url }\OperatorTok{=} \KeywordTok{new} \BuiltInTok{URL}\OperatorTok{(}\StringTok{"https://minwebbsida/text.txt"}\OperatorTok{);}
    \BuiltInTok{Scanner}\NormalTok{ scanner }\OperatorTok{=} \KeywordTok{new} \BuiltInTok{Scanner}\OperatorTok{(}\NormalTok{url}\OperatorTok{.}\FunctionTok{openStream}\OperatorTok{());}
    \ControlFlowTok{while} \OperatorTok{(}\NormalTok{scanner}\OperatorTok{.}\FunctionTok{hasNextLine}\OperatorTok{())} \OperatorTok{\{}
        \BuiltInTok{String}\NormalTok{ contents }\OperatorTok{=}\NormalTok{ scanner}\OperatorTok{.}\FunctionTok{nextLine}\OperatorTok{();}
        \BuiltInTok{System}\OperatorTok{.}\FunctionTok{out}\OperatorTok{.}\FunctionTok{println}\OperatorTok{(}\NormalTok{contents}\OperatorTok{);}
    \OperatorTok{\}}
\NormalTok{    scanner}\OperatorTok{.}\FunctionTok{close}\OperatorTok{();}
\OperatorTok{\}} \ControlFlowTok{catch} \OperatorTok{(}\BuiltInTok{IOException}\NormalTok{ e}\OperatorTok{)} \OperatorTok{\{}
    \BuiltInTok{System}\OperatorTok{.}\FunctionTok{out}\OperatorTok{.}\FunctionTok{println}\OperatorTok{(}\StringTok{"Ett fel inträffade vid läsning av textfilen."}\OperatorTok{);}
\OperatorTok{\}}
\end{Highlighting}
\end{Shaded}

\hypertarget{spara-en-textfil-fruxe5n-en-url}{%
\subsection{Spara en textfil från en
URL}\label{spara-en-textfil-fruxe5n-en-url}}

För att spara en textfil från en URL kan följande kod användas:

\hypertarget{cb9}{}
\begin{Shaded}
\begin{Highlighting}[]
\ControlFlowTok{try} \OperatorTok{\{}
    \BuiltInTok{URL}\NormalTok{ url }\OperatorTok{=} \KeywordTok{new} \BuiltInTok{URL}\OperatorTok{(}\StringTok{"https://minwebbsida/text.txt"}\OperatorTok{);}
    \BuiltInTok{InputStream}\NormalTok{ inputStream }\OperatorTok{=}\NormalTok{ url}\OperatorTok{.}\FunctionTok{openStream}\OperatorTok{();}
\NormalTok{    Files}\OperatorTok{.}\FunctionTok{copy}\OperatorTok{(}\NormalTok{inputStream}\OperatorTok{,}\NormalTok{ Paths}\OperatorTok{.}\FunctionTok{get}\OperatorTok{(}\StringTok{"text.txt"}\OperatorTok{),}\NormalTok{ StandardCopyOption}\OperatorTok{.}\FunctionTok{REPLACE\_EXISTING}\OperatorTok{);}
    \BuiltInTok{System}\OperatorTok{.}\FunctionTok{out}\OperatorTok{.}\FunctionTok{println}\OperatorTok{(}\StringTok{"Textfilen har sparats."}\OperatorTok{);}
\OperatorTok{\}} \ControlFlowTok{catch} \OperatorTok{(}\BuiltInTok{IOException}\NormalTok{ e}\OperatorTok{)} \OperatorTok{\{}
    \BuiltInTok{System}\OperatorTok{.}\FunctionTok{out}\OperatorTok{.}\FunctionTok{println}\OperatorTok{(}\StringTok{"Ett fel inträffade vid sparande av textfilen."}\OperatorTok{);}
\OperatorTok{\}}
\end{Highlighting}
\end{Shaded}

\hypertarget{lista-puxe5-termer-i-artikeln}{%
\subsection{Lista på termer i
artikeln}\label{lista-puxe5-termer-i-artikeln}}

\begin{longtable}[]{@{}
  >{\raggedright\arraybackslash}p{(\columnwidth - 2\tabcolsep) * \real{0.2000}}
  >{\raggedright\arraybackslash}p{(\columnwidth - 2\tabcolsep) * \real{0.8000}}@{}}
\toprule\noalign{}
\begin{minipage}[b]{\linewidth}\raggedright
Ord
\end{minipage} & \begin{minipage}[b]{\linewidth}\raggedright
Förklaring
\end{minipage} \\
\midrule\noalign{}
\endhead
\bottomrule\noalign{}
\endlastfoot
Buffer & En buffert är en tillfällig lagringsplats för data. \\
Buffering & Att använda en buffert för att lagra data. \\
Byte & En byte är en enhet för digital information som består av åtta
bitar. \\
Byte array & En byte array är en samling av bytes. \\
Byte stream & En byte stream är en sekvens av bytes. \\
Character stream & En character stream är en sekvens av tecken. \\
Close & Att stänga en filanslutning. \\
Copy & Att kopiera en fil. \\
File & En fil är en samling av data som lagras på en enhet. \\
File path & En filväg är en sträng som identifierar en fil. \\
File writer & En filskrivare är en klass som används för att skriva till
filer. \\
File reader & En filskrivare är en klass som används för att läsa från
filer. \\
File input stream & En filinmatningsström är en ström som används för
att läsa från filer. \\
File output stream & En filutmatningsström är en ström som används för
att skriva till filer. \\
File writer & En filskrivare är en klass som används för att skriva till
filer. \\
File reader & En filskrivare är en klass som används för att läsa från
filer. \\
\end{longtable}

\hypertarget{obligatorisk-dad-joke}{%
\subsection{Obligatorisk Dad-joke}\label{obligatorisk-dad-joke}}

Varför gillar textfiler att gå på fester?

För att de älskar att dra skämt om ``line''-dans!

Path-klassen i Java erbjuder verktyg för att hantera sökvägar till filer
och mappar. Genom att använda Path-klassens olika metoder kan vi enkelt
utföra operationer som att slå ihop sökvägar, extrahera information om
filer och mappar, ändra filändelser och mycket mer. Detta gör det
lättare att arbeta med filsystemet inom våra program och skapa bättre
och mer robusta applikationer.

\hypertarget{cb1}{}
\begin{Shaded}
\begin{Highlighting}[]
\KeywordTok{import} \ImportTok{java}\OperatorTok{.}\ImportTok{nio}\OperatorTok{.}\ImportTok{file}\OperatorTok{.}\ImportTok{Path}\OperatorTok{;}
\KeywordTok{import} \ImportTok{java}\OperatorTok{.}\ImportTok{nio}\OperatorTok{.}\ImportTok{file}\OperatorTok{.}\ImportTok{Paths}\OperatorTok{;}

\KeywordTok{public} \KeywordTok{class}\NormalTok{ PathExample }\OperatorTok{\{}
    \KeywordTok{public} \DataTypeTok{static} \DataTypeTok{void} \FunctionTok{main}\OperatorTok{(}\BuiltInTok{String}\OperatorTok{[]}\NormalTok{ args}\OperatorTok{)} \OperatorTok{\{}
        \CommentTok{// Skapa en sökväg till en fil}
\NormalTok{        Path filePath }\OperatorTok{=}\NormalTok{ Paths}\OperatorTok{.}\FunctionTok{get}\OperatorTok{(}\StringTok{"C:}\SpecialCharTok{\textbackslash{}\textbackslash{}}\StringTok{Users}\SpecialCharTok{\textbackslash{}\textbackslash{}}\StringTok{User}\SpecialCharTok{\textbackslash{}\textbackslash{}}\StringTok{Documents}\SpecialCharTok{\textbackslash{}\textbackslash{}}\StringTok{file.txt"}\OperatorTok{);}

        \CommentTok{// Slå ihop sökvägar}
\NormalTok{        Path combinedPath }\OperatorTok{=}\NormalTok{ filePath}\OperatorTok{.}\FunctionTok{resolve}\OperatorTok{(}\StringTok{"subfolder}\SpecialCharTok{\textbackslash{}\textbackslash{}}\StringTok{file2.txt"}\OperatorTok{);}

        \CommentTok{// Extrahera filnamn}
        \BuiltInTok{String}\NormalTok{ fileName }\OperatorTok{=}\NormalTok{ filePath}\OperatorTok{.}\FunctionTok{getFileName}\OperatorTok{().}\FunctionTok{toString}\OperatorTok{();}

        \CommentTok{// Byta filändelse}
\NormalTok{        Path newFilePath }\OperatorTok{=}\NormalTok{ filePath}\OperatorTok{.}\FunctionTok{resolveSibling}\OperatorTok{(}\NormalTok{fileName}\OperatorTok{.}\FunctionTok{replace}\OperatorTok{(}\StringTok{".txt"}\OperatorTok{,} \StringTok{".csv"}\OperatorTok{));}

        \CommentTok{// Kontrollera om sökvägen har en filändelse}
        \DataTypeTok{boolean}\NormalTok{ hasExtension }\OperatorTok{=}\NormalTok{ filePath}\OperatorTok{.}\FunctionTok{toString}\OperatorTok{().}\FunctionTok{contains}\OperatorTok{(}\StringTok{"."}\OperatorTok{);}

        \CommentTok{// Skriv ut sökvägarna}
        \BuiltInTok{System}\OperatorTok{.}\FunctionTok{out}\OperatorTok{.}\FunctionTok{println}\OperatorTok{(}\StringTok{"File Path: "} \OperatorTok{+}\NormalTok{ filePath}\OperatorTok{);}
        \BuiltInTok{System}\OperatorTok{.}\FunctionTok{out}\OperatorTok{.}\FunctionTok{println}\OperatorTok{(}\StringTok{"Combined Path: "} \OperatorTok{+}\NormalTok{ combinedPath}\OperatorTok{);}
        \BuiltInTok{System}\OperatorTok{.}\FunctionTok{out}\OperatorTok{.}\FunctionTok{println}\OperatorTok{(}\StringTok{"File Name: "} \OperatorTok{+}\NormalTok{ fileName}\OperatorTok{);}
        \BuiltInTok{System}\OperatorTok{.}\FunctionTok{out}\OperatorTok{.}\FunctionTok{println}\OperatorTok{(}\StringTok{"New File Path: "} \OperatorTok{+}\NormalTok{ newFilePath}\OperatorTok{);}
        \BuiltInTok{System}\OperatorTok{.}\FunctionTok{out}\OperatorTok{.}\FunctionTok{println}\OperatorTok{(}\StringTok{"Has Extension: "} \OperatorTok{+}\NormalTok{ hasExtension}\OperatorTok{);}
    \OperatorTok{\}}
\OperatorTok{\}}
\end{Highlighting}
\end{Shaded}

I det här exemplet använder vi Path-klassen för att utföra olika
operationer på sökvägar. Vi skapar en sökväg till en fil och visar hur
man slår ihop sökvägar, extraherar filnamn, ändrar filändelser och
kontrollerar om en sökväg har en filändelse. Vi använder också metoden
\texttt{resolveSibling} för att byta filändelse och \texttt{resolve} för
att slå ihop sökvägar.

Resultatet av körningen av detta program är:

\begin{verbatim}
File Path: C:\Users\User\Documents\file.txt
Combined Path: C:\Users\User\Documents\subfolder\file2.txt
File Name: file.txt
New File Path: C:\Users\User\Documents\file.csv
Has Extension: true
\end{verbatim}

Som du kan se ger Path-klassen oss enkla och effektiva verktyg för att
hantera sökvägar till filer och mappar i Java-program.

Oavsett om du behöver slå ihop sökvägar, extrahera information om filer
och mappar eller ändra filändelser, kan Path-klassen vara till stor
hjälp. Genom att utforska dokumentationen för Path-klassen kan du lära
dig mer om de olika metoderna och hur de kan användas för att hantera
sökvägar i dina Java-program.

\hypertarget{path-klassen}{%
\section{Path-klassen}\label{path-klassen}}

\begin{itemize}
\tightlist
\item
  \href{https://docs.oracle.com/javase/8/docs/api/java/nio/file/Path.html}{Java
  Path Documentation}
\item
  \href{https://www.baeldung.com/java-nio-2-path}{Java Path Tutorial}
\item
\end{itemize}

\hypertarget{asynkron}{%
\section{Asynkron}\label{asynkron}}

Asynkrona metoder är metoder som inte blockerar tråden som kör metoden.
Detta gör det möjligt för oss att köra flera metoder samtidigt och
förbättra applikationens responsivitet.

\hypertarget{tldr}{%
\subsection{TL;DR}\label{tldr}}

Asynkrona metoder gör det möjligt att köra flera metoder samtidigt och
förbättra applikationens responsivitet genom att inte blockera tråden.

\hypertarget{nuxe4r-du-luxe4st-detta-ska-du-kunna}{%
\subsection{När du läst detta ska du
kunna}\label{nuxe4r-du-luxe4st-detta-ska-du-kunna}}

\begin{itemize}
\tightlist
\item
  Förstå vad asynkrona metoder är och hur de fungerar.
\item
  Veta när det är lämpligt att använda asynkrona metoder.
\item
  Kunna använda asynkrona metoder i Java.
\end{itemize}

\hypertarget{introduktion}{%
\subsection{Introduktion}\label{introduktion}}

I vissa situationer kan det vara nödvändigt att köra flera metoder
samtidigt i en applikation. Det kan handla om att kommunicera med andra
system och API:er eller att hantera tunga beräkningar. Traditionellt
sett skulle sådana metoder blockera huvudtråden i applikationen och göra
den otillgänglig för andra operationer tills den aktuella metoden är
klar. Detta kan resultera i en långsam och ineffektiv
användarupplevelse. För att undvika detta och göra applikationen mer
responsiv kan man använda sig av asynkrona metoder.

\hypertarget{beskrivning}{%
\subsection{Beskrivning}\label{beskrivning}}

Asynkrona metoder i Java gör det möjligt att utföra operationer utan att
blockera tråden som kör metoden. Istället för att vänta på att en
operation är klar innan man går vidare till nästa, kan man använda
nyckelordet \texttt{Runnable} tillsammans med \texttt{Thread} för att
ange att en metod ska köras asynkront. När en metod är markerad som
\texttt{Runnable} kommer den att returnera ett \texttt{Thread}-objekt
som representerar den asynkrona operationen. När man använder
\texttt{Thread} för att vänta på att en asynkron operation är klar,
kommer tråden som kör metoden att släppas och andra operationer kan
köras samtidigt.

\hypertarget{nuxe4r-ska-man-anvuxe4nda-asynkrona-metoder}{%
\subsection{När ska man använda asynkrona
metoder?}\label{nuxe4r-ska-man-anvuxe4nda-asynkrona-metoder}}

Asynkrona metoder är användbara när man vill förbättra applikationens
responsivitet och prestanda. De är också användbara när man vill köra
flera operationer samtidigt, som att kommunicera med andra system eller
hantera tunga beräkningar.

\hypertarget{sammanfattning}{%
\subsection{Sammanfattning}\label{sammanfattning}}

Asynkrona metoder är ett användbart verktyg för att förbättra
applikationens responsivitet och prestanda. De gör det möjligt att köra
flera metoder samtidigt och blockera inte tråden som kör metoden. I Java
kan man använda nyckelordet \texttt{Runnable} tillsammans med
\texttt{Thread} för att ange att en metod ska köras asynkront.3.
\textbf{Kodförenkling}: Asynkrona metoder kan göra koden mer läsbar och
förenkla kodstrukturen.

\hypertarget{java}{%
\subsection{Java}\label{java}}

Vanligtvis används asynkrona metoder när man inte vet hur lång tid en
operation kommer att ta, till exempel vid kommunikation över nätverk
eller när man arbetar med stora datamängder. Genom att köra sådana
operationer asynkront kan man göra applikationen mer responsiv och
tillåta att andra operationer utförs parallellt.

Det är viktigt att notera några viktiga punkter när man använder
asynkrona metoder:

\begin{itemize}
\tightlist
\item
  En asynkron metod måste alltid returnera antingen en \texttt{Future}
  eller en \texttt{Future\textless{}T\textgreater{}}. Detta är för att
  kunna köra flera metoder samtidigt.
\item
  Asynkrona metoder kan inte ha en \texttt{void}-returtyp. Istället bör
  de returnera en \texttt{Future} om inget värde behövs eller en
  \texttt{Future\textless{}T\textgreater{}} om det finns ett returvärde.
\item
  Asynkrona metoder kan inte ha \texttt{out}- eller
  \texttt{ref}-parametrar.
\item
  Asynkrona metoder kan inte ha parametrar som är av typen
  \texttt{varargs}.
\end{itemize}

\hypertarget{fuxf6rdelar}{%
\subsection{Fördelar}\label{fuxf6rdelar}}

Användningen av asynkrona metoder kan ge flera fördelar i en
applikation:

\begin{enumerate}
\def\labelenumi{\arabic{enumi}.}
\tightlist
\item
  \textbf{Responsivitet}: Genom att använda asynkrona metoder kan man
  undvika att blockera huvudtråden i applikationen och ge en mer
  responsiv användarupplevelse.
\item
  \textbf{Parallellism}: Genom att köra flera metoder samtidigt kan man
  utnyttja flera processorkärnor och därmed förbättra prestanda och
  genomströmning.
\item
  \textbf{Kodförenkling}: Asynkrona metoder kan göra koden mer läsbar
  och förenkla kodstrukturen.3. \textbf{Förbättrad skalbarhet}:
  Asynkrona metoder möjliggör parallell exekvering av operationer,
  vilket kan förbättra skalbarheten i en applikation och möjliggöra
  hantering av hög trafik eller arbetsbelastning.
\end{enumerate}

\hypertarget{begruxe4nsningar}{%
\subsection{Begränsningar}\label{begruxe4nsningar}}

Trots de fördelar som asynkrona metoder erbjuder finns det vissa
begränsningar och överväganden att tänka på:

\begin{enumerate}
\def\labelenumi{\arabic{enumi}.}
\tightlist
\item
  \textbf{Komplexitet}: Asynkron programmering kan vara mer komplext än
  synkron programmering, särskilt när det gäller att hantera fel och
  synchronisering mellan olika operationer.
\item
  \textbf{Ökad CPU-användning}: Asynkrona metoder kan resultera i en
  ökad CPU-användning på grund av överhead för att starta och hantera
  nya trådar eller arbetsenheter.
\item
  \textbf{Felhantering}: Det kan vara utmanande att hantera fel i
  asynkrona metoder, särskilt när flera operationer körs samtidigt.
  Korrekt felhantering och återställning är viktigt för att undvika
  buggar och felaktigt beteende.
\end{enumerate}

\hypertarget{anvuxe4ndningsomruxe5den}{%
\subsection{Användningsområden}\label{anvuxe4ndningsomruxe5den}}

Asynkrona metoder används i olika situationer och applikationsområden:

\begin{enumerate}
\def\labelenumi{\arabic{enumi}.}
\tightlist
\item
  \textbf{Webbapplikationer}: Vid kommunikation med webbtjänster och
  API:er kan asynkrona metoder användas för att göra anrop och hämta
  data på ett effektivt sätt utan att blockera huvudtråden.
\item
  \textbf{Databashantering}: Asynkrona metoder kan användas vid
  databashantering för att parallellt hämta och uppdatera data från
  flera källor samtidigt.Asynkrona metoder är ett användbart verktyg för
  att förbättra prestanda och responsivitet i applikationer som hanterar
  flera anrop samtidigt.
\end{enumerate}

\hypertarget{exempelkod---asynkrona-anrop-java}{%
\subsection{Exempelkod - Asynkrona anrop
(Java)}\label{exempelkod---asynkrona-anrop-java}}

Här är ett exempel på hur man kan använda asynkrona metoder i Java:

\hypertarget{cb1}{}
\begin{Shaded}
\begin{Highlighting}[]
\KeywordTok{import} \ImportTok{java}\OperatorTok{.}\ImportTok{io}\OperatorTok{.}\ImportTok{BufferedReader}\OperatorTok{;}
\KeywordTok{import} \ImportTok{java}\OperatorTok{.}\ImportTok{io}\OperatorTok{.}\ImportTok{IOException}\OperatorTok{;}
\KeywordTok{import} \ImportTok{java}\OperatorTok{.}\ImportTok{io}\OperatorTok{.}\ImportTok{InputStreamReader}\OperatorTok{;}
\KeywordTok{import} \ImportTok{java}\OperatorTok{.}\ImportTok{net}\OperatorTok{.}\ImportTok{URL}\OperatorTok{;}
\KeywordTok{import} \ImportTok{java}\OperatorTok{.}\ImportTok{net}\OperatorTok{.}\ImportTok{URLConnection}\OperatorTok{;}
\KeywordTok{import} \ImportTok{java}\OperatorTok{.}\ImportTok{util}\OperatorTok{.}\ImportTok{concurrent}\OperatorTok{.}\ImportTok{CompletableFuture}\OperatorTok{;}

\KeywordTok{public} \KeywordTok{class}\NormalTok{ Main }\OperatorTok{\{}
    \KeywordTok{public} \DataTypeTok{static}\NormalTok{ CompletableFuture}\OperatorTok{\textless{}}\BuiltInTok{String}\OperatorTok{\textgreater{}} \FunctionTok{getAsync}\OperatorTok{(}\BuiltInTok{String}\NormalTok{ url}\OperatorTok{)} \OperatorTok{\{}
        \CommentTok{// Starta en ny tråd och exekvera GET{-}begäran}
        \ControlFlowTok{return}\NormalTok{ CompletableFuture}\OperatorTok{.}\FunctionTok{supplyAsync}\OperatorTok{(()} \OperatorTok{{-}\textgreater{}} \OperatorTok{\{}
            \ControlFlowTok{try} \OperatorTok{\{}
                \CommentTok{// Skapa en URL{-}anslutning till webbplatsen}
                \BuiltInTok{URL}\NormalTok{ website }\OperatorTok{=} \KeywordTok{new} \BuiltInTok{URL}\OperatorTok{(}\NormalTok{url}\OperatorTok{);}
                \CommentTok{// Öppna anslutningen}
                \BuiltInTok{URLConnection}\NormalTok{ connection }\OperatorTok{=}\NormalTok{ website}\OperatorTok{.}\FunctionTok{openConnection}\OperatorTok{();}
                \CommentTok{// Skapa en BufferedReader för att läsa svaret}
                \BuiltInTok{BufferedReader}\NormalTok{ in }\OperatorTok{=} \KeywordTok{new} \BuiltInTok{BufferedReader}\OperatorTok{(}\KeywordTok{new} \BuiltInTok{InputStreamReader}\OperatorTok{(}\NormalTok{connection}\OperatorTok{.}\FunctionTok{getInputStream}\OperatorTok{()));}
                \CommentTok{// Läs svaret från anslutningen}
                \BuiltInTok{StringBuilder}\NormalTok{ response }\OperatorTok{=} \KeywordTok{new} \BuiltInTok{StringBuilder}\OperatorTok{();}

                \BuiltInTok{String}\NormalTok{ inputLine}\OperatorTok{;}
                \CommentTok{// Läs varje rad från svaret}
                \ControlFlowTok{while} \OperatorTok{((}\NormalTok{inputLine }\OperatorTok{=}\NormalTok{ in}\OperatorTok{.}\FunctionTok{readLine}\OperatorTok{())} \OperatorTok{!=} \KeywordTok{null}\OperatorTok{)}
\NormalTok{                    response}\OperatorTok{.}\FunctionTok{append}\OperatorTok{(}\NormalTok{inputLine}\OperatorTok{);}
\NormalTok{                in}\OperatorTok{.}\FunctionTok{close}\OperatorTok{();}
                \CommentTok{// Returnera svaret som en sträng}
                \ControlFlowTok{return}\NormalTok{ response}\OperatorTok{.}\FunctionTok{toString}\OperatorTok{();}
            \OperatorTok{\}} \ControlFlowTok{catch} \OperatorTok{(}\BuiltInTok{IOException}\NormalTok{ e}\OperatorTok{)} \OperatorTok{\{}
                \CommentTok{// Hantera eventuella IOExceptions}
\NormalTok{                e}\OperatorTok{.}\FunctionTok{printStackTrace}\OperatorTok{();}
                \ControlFlowTok{return} \KeywordTok{null}\OperatorTok{;}
            \OperatorTok{\}}
        \OperatorTok{\});}
    \OperatorTok{\}}

    \KeywordTok{public} \DataTypeTok{static} \DataTypeTok{void} \FunctionTok{main}\OperatorTok{(}\BuiltInTok{String}\OperatorTok{[]}\NormalTok{ args}\OperatorTok{)} \OperatorTok{\{}
        \CommentTok{// Utför en asynkron GET{-}begäran till "https://www.google.com"}
\NormalTok{        CompletableFuture}\OperatorTok{\textless{}}\BuiltInTok{String}\OperatorTok{\textgreater{}}\NormalTok{ result }\OperatorTok{=} \FunctionTok{getAsync}\OperatorTok{(}\StringTok{"https://www.google.com"}\OperatorTok{);}
        \ControlFlowTok{try} \OperatorTok{\{}
            \CommentTok{// Vänta på att resultatet ska bli tillgängligt}
            \BuiltInTok{System}\OperatorTok{.}\FunctionTok{out}\OperatorTok{.}\FunctionTok{println}\OperatorTok{(}\NormalTok{result}\OperatorTok{.}\FunctionTok{get}\OperatorTok{());}
        \OperatorTok{\}} \ControlFlowTok{catch} \OperatorTok{(}\BuiltInTok{Exception}\NormalTok{ e}\OperatorTok{)} \OperatorTok{\{}
            \CommentTok{// Hantera eventuella undantag}
\NormalTok{            e}\OperatorTok{.}\FunctionTok{printStackTrace}\OperatorTok{();}
        \OperatorTok{\}}
    \OperatorTok{\}}
\OperatorTok{\}}
\end{Highlighting}
\end{Shaded}

Här är förklaringar till vad varje del av koden gör:

\begin{enumerate}
\def\labelenumi{\arabic{enumi}.}
\tightlist
\item
  Importera nödvändiga klasser från Java API.
\item
  Skapa en klass \texttt{Main}.
\item
  Definiera \texttt{getAsync}-metoden som returnerar en
  \texttt{CompletableFuture} som innehåller resultatet av en asynkron
  HTTP GET-begäran.
\item
  Anropa \texttt{CompletableFuture.supplyAsync} för att starta en ny
  tråd och exekvera GET-begäran.
\item
  Inuti den nya tråden, försök att utföra GET-begäran.
\item
  Skapa en \texttt{URL}-instans baserat på den angivna URL-strängen.
\item
  Öppna en anslutning (\texttt{URLConnection}) till webbplatsen.
\item
  Skapa en \texttt{BufferedReader} för att läsa svaret från
  anslutningen.
\item
  Skapa en \texttt{StringBuilder} för att bygga upp hela svaret rad för
  rad.
\item
  Läs varje rad från svaret och lägg till den i
  \texttt{response}-objektet.
\item
  Stäng \texttt{BufferedReader}.
\item
  Returnera det kompletta svaret som en sträng.
\item
  Hantera eventuella \texttt{IOExceptions} genom att skriva ut
  felmeddelandet och returnera \texttt{null}.
\item
  Definiera \texttt{main}-metoden som startpunkt för programmet.
\item
  Anropa \texttt{getAsync}-metoden för att utföra en asynkron
  GET-begäran till ``https://www.google.com''.
\item
  Anropa \texttt{result.get()} för att vänta på att resultatet av
  GET-begäran ska vara tillgängligt.
\item
  Skriv ut resultatet till konsolen.
\item
  Hantera eventuella undantag genom att skriva ut felmeddelandet.```
\end{enumerate}

I exemplet ovan definieras en asynkron metod \texttt{getAsync} som
hämtar data från en given URL. Metoden använder
\texttt{CompletableFuture} för att göra ett asynkront HTTP-anrop och
returnerar innehållet som en sträng. I \texttt{main}-metoden anropas
\texttt{getAsync} och resultatet skrivs ut på konsolen.

\hypertarget{slutsats}{%
\subsection{Slutsats}\label{slutsats}}

Asynkrona metoder är ett användbart verktyg för att förbättra prestanda
och responsivitet i applikationer som hanterar flera anrop samtidigt.
Asynkrona metoder är ett kraftfullt verktyg för att göra applikationer
mer responsiva och effektiva. Genom att köra flera metoder samtidigt kan
man förbättra prestanda och användarupplevelse. Det är viktigt att
använda asynkrona metoder på rätt sätt och vara medveten om deras
begränsningar och överväganden. Med rätt användning kan asynkron
programmering vara ett värdefullt verktyg i utvecklingen av moderna
applikationer.

\hypertarget{termtabell}{%
\subsection{Termtabell}\label{termtabell}}

\begin{itemize}
\tightlist
\item
  \textbf{Asynkrona metoder}: Metoder som inte blockerar tråden och
  tillåter körsel av flera metoder samtidigt.
\item
  \textbf{Responsivitet}: Förmågan hos en applikation att snabbt svara
  på användarinteraktioner och andra händelser.
\item
  \textbf{Parallellism}: Exekvering av flera operationer samtidigt för
  att utnyttja flera processorkärnor och förbättra prestanda.
\item
  \textbf{Skalbarhet}: Förmågan hos en applikation att hantera en ökad
  arbetsbelastning och trafik utan att försämra prestanda och
  responsivitet.
\end{itemize}

\hypertarget{obligatorisk-dad-joke}{%
\subsection{Obligatorisk Dad-joke}\label{obligatorisk-dad-joke}}

Varför gillar programmerare att arbeta med asynkrona metoder?

För att de älskar att ``vänta'' på spänning!

\hypertarget{exempel}{%
\section{Exempel}\label{exempel}}

Yay! Nu ska vi titta på ett exempel där vi ska skapa en metod som söker
igenom alla filer i en mapp och returnerar en lista med filer som
innehåller en specifik text. Vi kommer att använda asynkrona metoder för
att kunna köra flera metoder samtidigt och därigenom effektivisera vår
sökning. Hur coolt låter inte det?

Innehållsförteckning

\{: .text-delta \} 1. Innehållsförteckning \{:toc\}

\hypertarget{kod}{%
\subsection{Kod}\label{kod}}

\hypertarget{cb1}{}
\begin{Shaded}
\begin{Highlighting}[]
\KeywordTok{import} \ImportTok{java}\OperatorTok{.}\ImportTok{io}\OperatorTok{.}\ImportTok{File}\OperatorTok{;}
\KeywordTok{import} \ImportTok{java}\OperatorTok{.}\ImportTok{io}\OperatorTok{.}\ImportTok{IOException}\OperatorTok{;}
\KeywordTok{import} \ImportTok{java}\OperatorTok{.}\ImportTok{nio}\OperatorTok{.}\ImportTok{file}\OperatorTok{.}\ImportTok{Files}\OperatorTok{;}
\KeywordTok{import} \ImportTok{java}\OperatorTok{.}\ImportTok{nio}\OperatorTok{.}\ImportTok{file}\OperatorTok{.}\ImportTok{Path}\OperatorTok{;}
\KeywordTok{import} \ImportTok{java}\OperatorTok{.}\ImportTok{nio}\OperatorTok{.}\ImportTok{file}\OperatorTok{.}\ImportTok{Paths}\OperatorTok{;}
\KeywordTok{import} \ImportTok{java}\OperatorTok{.}\ImportTok{util}\OperatorTok{.}\ImportTok{ArrayList}\OperatorTok{;}
\KeywordTok{import} \ImportTok{java}\OperatorTok{.}\ImportTok{util}\OperatorTok{.}\ImportTok{List}\OperatorTok{;}
\KeywordTok{import} \ImportTok{java}\OperatorTok{.}\ImportTok{util}\OperatorTok{.}\ImportTok{concurrent}\OperatorTok{.}\ImportTok{CompletableFuture}\OperatorTok{;}
\KeywordTok{import} \ImportTok{java}\OperatorTok{.}\ImportTok{util}\OperatorTok{.}\ImportTok{stream}\OperatorTok{.}\ImportTok{Collectors}\OperatorTok{;}
\KeywordTok{import} \ImportTok{java}\OperatorTok{.}\ImportTok{util}\OperatorTok{.}\ImportTok{stream}\OperatorTok{.}\ImportTok{Stream}\OperatorTok{;}

\KeywordTok{public} \KeywordTok{class}\NormalTok{ Main }\OperatorTok{\{}
    \KeywordTok{public} \DataTypeTok{static} \DataTypeTok{void} \FunctionTok{main}\OperatorTok{(}\BuiltInTok{String}\OperatorTok{[]}\NormalTok{ args}\OperatorTok{)} \KeywordTok{throws} \BuiltInTok{Exception} \OperatorTok{\{}
        \CommentTok{// Skapa en ny instans av DirReader{-}klassen}
\NormalTok{        DirReader dir }\OperatorTok{=} \KeywordTok{new} \FunctionTok{DirReader}\OperatorTok{();}

        \CommentTok{// Hämta användarens hemkatalog}
        \BuiltInTok{String}\NormalTok{ profile }\OperatorTok{=} \BuiltInTok{System}\OperatorTok{.}\FunctionTok{getProperty}\OperatorTok{(}\StringTok{"user.home"}\OperatorTok{);}

        \CommentTok{// Definiera sökvägen till källmappen}
        \BuiltInTok{String}\NormalTok{ source }\OperatorTok{=}\NormalTok{ Paths}\OperatorTok{.}\FunctionTok{get}\OperatorTok{(}\NormalTok{profile}\OperatorTok{,} \StringTok{"source"}\OperatorTok{).}\FunctionTok{toString}\OperatorTok{();}

        \CommentTok{// Sök efter alla .java{-}filer i källmappen asynkront}
        \BuiltInTok{List}\OperatorTok{\textless{}}\BuiltInTok{String}\OperatorTok{\textgreater{}}\NormalTok{ result }\OperatorTok{=}\NormalTok{ dir}\OperatorTok{.}\FunctionTok{searchFilesAsync}\OperatorTok{(}\NormalTok{source}\OperatorTok{,} \StringTok{"*.java"}\OperatorTok{,} \StringTok{"asynkron"}\OperatorTok{).}\FunctionTok{get}\OperatorTok{();}

        \CommentTok{// Skriv ut resultaten till konsolen}
        \BuiltInTok{System}\OperatorTok{.}\FunctionTok{out}\OperatorTok{.}\FunctionTok{println}\OperatorTok{(}\BuiltInTok{String}\OperatorTok{.}\FunctionTok{join}\OperatorTok{(}\BuiltInTok{System}\OperatorTok{.}\FunctionTok{lineSeparator}\OperatorTok{(),}\NormalTok{ result}\OperatorTok{));}
    \OperatorTok{\}}
\OperatorTok{\}}

\KeywordTok{class}\NormalTok{ DirReader }\OperatorTok{\{}
    \CommentTok{// Sök efter alla filer i en mapp som matchar ett angivet mönster och innehåller en viss text}
    \KeywordTok{public}\NormalTok{ CompletableFuture}\OperatorTok{\textless{}}\BuiltInTok{List}\OperatorTok{\textless{}}\BuiltInTok{String}\OperatorTok{\textgreater{}\textgreater{}} \FunctionTok{searchFilesAsync}\OperatorTok{(}\BuiltInTok{String}\NormalTok{ path}\OperatorTok{,} \BuiltInTok{String}\NormalTok{ searchPattern}\OperatorTok{,} \BuiltInTok{String}\NormalTok{ searchText}\OperatorTok{)} \OperatorTok{\{}
        \ControlFlowTok{try} \OperatorTok{(}\NormalTok{Stream}\OperatorTok{\textless{}}\NormalTok{Path}\OperatorTok{\textgreater{}}\NormalTok{ paths }\OperatorTok{=}\NormalTok{ Files}\OperatorTok{.}\FunctionTok{walk}\OperatorTok{(}\NormalTok{Paths}\OperatorTok{.}\FunctionTok{get}\OperatorTok{(}\NormalTok{path}\OperatorTok{)))} \OperatorTok{\{}
            \CommentTok{// Skapa en lista med uppgifter för att söka varje fil asynkront}
            \BuiltInTok{List}\OperatorTok{\textless{}}\NormalTok{CompletableFuture}\OperatorTok{\textless{}}\BuiltInTok{String}\OperatorTok{\textgreater{}\textgreater{}}\NormalTok{ tasks }\OperatorTok{=}\NormalTok{ paths}
                    \OperatorTok{.}\FunctionTok{filter}\OperatorTok{(}\NormalTok{Files}\OperatorTok{::}\NormalTok{isRegularFile}\OperatorTok{)} \CommentTok{// Endast vanliga filer beaktas}
                    \OperatorTok{.}\FunctionTok{filter}\OperatorTok{(}\NormalTok{p }\OperatorTok{{-}\textgreater{}}\NormalTok{ p}\OperatorTok{.}\FunctionTok{getFileName}\OperatorTok{().}\FunctionTok{toString}\OperatorTok{().}\FunctionTok{endsWith}\OperatorTok{(}\NormalTok{searchPattern}\OperatorTok{))} \CommentTok{// Endast filer som matchar sökmönstret beaktas}
                    \OperatorTok{.}\FunctionTok{map}\OperatorTok{(}\NormalTok{p }\OperatorTok{{-}\textgreater{}} \FunctionTok{searchFileAsync}\OperatorTok{(}\NormalTok{p}\OperatorTok{.}\FunctionTok{toString}\OperatorTok{(),}\NormalTok{ searchText}\OperatorTok{))} \CommentTok{// Sök i varje fil asynkront}
                    \OperatorTok{.}\FunctionTok{collect}\OperatorTok{(}\NormalTok{Collectors}\OperatorTok{.}\FunctionTok{toList}\OperatorTok{());}

            \CommentTok{// Vänta på att alla uppgifter ska slutföras och samla resultaten}
\NormalTok{            CompletableFuture}\OperatorTok{\textless{}}\BuiltInTok{Void}\OperatorTok{\textgreater{}}\NormalTok{ allTasks }\OperatorTok{=}\NormalTok{ CompletableFuture}\OperatorTok{.}\FunctionTok{allOf}\OperatorTok{(}\NormalTok{tasks}\OperatorTok{.}\FunctionTok{toArray}\OperatorTok{(}\KeywordTok{new}\NormalTok{ CompletableFuture}\OperatorTok{[}\NormalTok{tasks}\OperatorTok{.}\FunctionTok{size}\OperatorTok{()]));}
            \ControlFlowTok{return}\NormalTok{ allTasks}\OperatorTok{.}\FunctionTok{thenApply}\OperatorTok{(}\NormalTok{v }\OperatorTok{{-}\textgreater{}}\NormalTok{ tasks}\OperatorTok{.}\FunctionTok{stream}\OperatorTok{()}
                    \OperatorTok{.}\FunctionTok{map}\OperatorTok{(}\NormalTok{CompletableFuture}\OperatorTok{::}\NormalTok{join}\OperatorTok{)}
                    \OperatorTok{.}\FunctionTok{filter}\OperatorTok{(}\NormalTok{x }\OperatorTok{{-}\textgreater{}}\NormalTok{ x }\OperatorTok{!=} \KeywordTok{null}\OperatorTok{)}
                    \OperatorTok{.}\FunctionTok{collect}\OperatorTok{(}\NormalTok{Collectors}\OperatorTok{.}\FunctionTok{toList}\OperatorTok{()));}
        \OperatorTok{\}} \ControlFlowTok{catch} \OperatorTok{(}\BuiltInTok{IOException}\NormalTok{ e}\OperatorTok{)} \OperatorTok{\{}
            \CommentTok{// Om ett fel inträffar, returnera ett misslyckat CompletableFuture med undantaget}
\NormalTok{            e}\OperatorTok{.}\FunctionTok{printStackTrace}\OperatorTok{();}
            \ControlFlowTok{return}\NormalTok{ CompletableFuture}\OperatorTok{.}\FunctionTok{failedFuture}\OperatorTok{(}\NormalTok{e}\OperatorTok{);}
        \OperatorTok{\}}
    \OperatorTok{\}}

    \CommentTok{// Sök efter en given text i en enskild fil och returnera resultatet}
    \KeywordTok{public}\NormalTok{ CompletableFuture}\OperatorTok{\textless{}}\BuiltInTok{String}\OperatorTok{\textgreater{}} \FunctionTok{searchFileAsync}\OperatorTok{(}\BuiltInTok{String}\NormalTok{ path}\OperatorTok{,} \BuiltInTok{String}\NormalTok{ searchText}\OperatorTok{)} \OperatorTok{\{}
        \ControlFlowTok{return}\NormalTok{ CompletableFuture}\OperatorTok{.}\FunctionTok{supplyAsync}\OperatorTok{(()} \OperatorTok{{-}\textgreater{}} \OperatorTok{\{}
            \ControlFlowTok{try} \OperatorTok{\{}
                \CommentTok{// Läs in alla rader i filen}
                \BuiltInTok{List}\OperatorTok{\textless{}}\BuiltInTok{String}\OperatorTok{\textgreater{}}\NormalTok{ lines }\OperatorTok{=}\NormalTok{ Files}\OperatorTok{.}\FunctionTok{readAllLines}\OperatorTok{(}\NormalTok{Paths}\OperatorTok{.}\FunctionTok{get}\OperatorTok{(}\NormalTok{path}\OperatorTok{));}

                \CommentTok{// Sök efter texten i filen}
                \DataTypeTok{int}\NormalTok{ index }\OperatorTok{=}\NormalTok{ lines}\OperatorTok{.}\FunctionTok{indexOf}\OperatorTok{(}\NormalTok{searchText}\OperatorTok{);}

                \CommentTok{// Om texten hittas, returnera filens sökväg och radnummer}
                \ControlFlowTok{return}\NormalTok{ index }\OperatorTok{\textgreater{}=} \DecValTok{0} \OperatorTok{?}\NormalTok{ path }\OperatorTok{+} \StringTok{" {-} "} \OperatorTok{+}\NormalTok{ index }\OperatorTok{:} \KeywordTok{null}\OperatorTok{;}
            \OperatorTok{\}} \ControlFlowTok{catch} \OperatorTok{(}\BuiltInTok{IOException}\NormalTok{ e}\OperatorTok{)} \OperatorTok{\{}
                \CommentTok{// Om ett fel inträffar, returnera null}
\NormalTok{                e}\OperatorTok{.}\FunctionTok{printStackTrace}\OperatorTok{();}
                \ControlFlowTok{return} \KeywordTok{null}\OperatorTok{;}
            \OperatorTok{\}}
        \OperatorTok{\});}
    \OperatorTok{\}}
\OperatorTok{\}}
\end{Highlighting}
\end{Shaded}

\begin{enumerate}
\def\labelenumi{\arabic{enumi}.}
\tightlist
\item
  Vi börjar med att skapa en metod som heter \texttt{SearchFilesAsync}.
\item
  Metoden tar emot en sökväg, sökmönster och söktext.
\item
  Vi använder \texttt{Files.walk} för att hämta en ström med alla filer
  i den angivna sökvägen.
\item
  Därefter skapar vi en lista med tasks som kommer att köra metoden
  \texttt{SearchFileAsync}.
\item
  Vi loopar igenom varje fil och lägger till en task för varje fil i vår
  lista.
\item
  Genom att använda \texttt{CompletableFuture.allOf} väntar vi på att
  alla tasks ska bli klara.
\item
  Vi returnerar en lista med resultatet från varje task som inte är
  null.
\item
  Därefter skapar vi en metod som heter \texttt{SearchFileAsync}.
\item
  Metoden tar emot en sökväg och en söktext.
\item
  Vi använder \texttt{Files.readAllLines} för att läsa innehållet i
  filen och sparar det i en lista.
\item
  Vi skapar en variabel som håller koll på vilken rad vi är på i filen.
\item
  Sedan loopar vi igenom varje rad i filen.
\item
  Om raden innehåller den sökta texten så returnerar vi sökvägen till
  filen och vilken rad det var.
\item
  Vi läser in nästa rad från filen.
\item
  Vi ökar radnumret med 1.
\item
  Om vi inte hittar någon rad som innehåller den sökta texten så
  returnerar vi null.
\item
  Slutligen har vi vår \texttt{Main}-metod där vi skapar en instans av
  \texttt{DirReader} och definierar sökvägen till vår mapp.
\item
  Vi kör metoden \texttt{SearchFilesAsync} och sparar resultatet i en
  variabel.
\item
  Till sist skriver vi ut alla resultat i konsolen.
\end{enumerate}

Wow! Nu kan vi söka igenom filer asynkront och få en lista med de filer
som innehåller vår sökta text. Det är riktigt coolt att kunna köra flera
metoder samtidigt och på så sätt förbättra vår effektivitet. Fortsätt
det fantastiska arbetet!

\hypertarget{slutsats}{%
\subsection{Slutsats}\label{slutsats}}

Att kunna utföra asynkrona operationer är en kraftfull teknik som gör
att våra applikationer kan vara mer responsiva och effektiva. Genom att
använda asynkrona metoder kan vi undvika att blockera huvudtråden och
istället köra flera operationer samtidigt. Detta är särskilt användbart
vid uppgifter som tar tid, som att söka igenom filer eller kommunicera
med externa system. Genom att utnyttja asynkron programmering kan vi
skapa mer responsiva och skalbara applikationer. Fortsätt kodning med
glädje!

\hypertarget{termtabell}{%
\subsection{Termtabell}\label{termtabell}}

\begin{itemize}
\tightlist
\item
  \textbf{Asynkrona metoder}: Metoder som inte blockerar tråden och
  tillåter parallell exekvering av flera metoder samtidigt.
\item
  \textbf{Responsivitet}: Förmågan hos en applikation att snabbt svara
  på användarinteraktioner och andra händelser.
\item
  \textbf{Parallellism}: Exekvering av flera operationer samtidigt för
  att utnyttja flera processorkärnor och förbättra prestanda.
\item
  \textbf{Skalbarhet}: Förmågan hos en applikation att hantera en ökad
  arbetsbelastning och trafik utan att försämra prestanda och
  responsivitet.
\end{itemize}

\hypertarget{obligatorisk-dad-joke}{%
\subsection{Obligatorisk Dad-joke}\label{obligatorisk-dad-joke}}

Varför kallas asynkrona metoder för ``coola''?

För att de jobbar i sin egen ``takt''!

\hypertarget{api}{%
\section{API}\label{api}}

Nu ska vi prata om API, eller Application Programming Interface, som är
ett viktigt gränssnitt för att underlätta kommunikationen med
webbapplikationer. API:er fungerar som en slags mellanhand mellan två
applikationer, vilket gör det möjligt för oss att ställa frågor och få
svar från webbapplikationer samt skicka och ta emot data. Genom API:er
blir klienten och servern oberoende av varandra och kan enkelt utbyta
information.

Innehållsförteckning

\{: .text-delta \}

\begin{enumerate}
\def\labelenumi{\arabic{enumi}.}
\tightlist
\item
  TOC \{:toc\}
\end{enumerate}

\hypertarget{tl-dr}{%
\subsection{TL; DR}\label{tl-dr}}

API:er är ett viktigt gränssnitt för att underlätta kommunikationen med
webbapplikationer. De gör det möjligt för oss att ställa frågor och få
svar från webbapplikationer samt skicka och ta emot data. Genom API:er
blir klienten och servern oberoende av varandra och kan enkelt utbyta
information.

\hypertarget{fuxf6rdjupning}{%
\subsection{Fördjupning}\label{fuxf6rdjupning}}

API:er gör det möjligt för oss att använda webbapplikationer utan att
behöva interagera med deras grafiska gränssnitt (GUI). Istället kan vi
kommunicera med applikationen genom att ställa frågor och begära data
via en URL. Svaren från API:et kommer vanligtvis i formatet XML eller
JSON, vilket gör det lättare att läsa och tolka informationen.

Genom att använda API:er kan vi skapa mer flexibla och dynamiska
applikationer. Vi kan anpassa datautbytet mellan klienten och servern
och se till att endast de delar av servern som API:et tillåter blir
synliga för klienten. Detta underlättar kommunikationen och ger oss
möjlighet att skapa mer effektiva och modulära system.

\hypertarget{exempel}{%
\subsection{Exempel}\label{exempel}}

Här är ett exempel på hur vi kan använda ett API i Java för att hämta
data från The Movie Database (TMDb):

\hypertarget{cb1}{}
\begin{Shaded}
\begin{Highlighting}[]
\KeywordTok{import} \ImportTok{java}\OperatorTok{.}\ImportTok{io}\OperatorTok{.}\ImportTok{BufferedReader}\OperatorTok{;}
\KeywordTok{import} \ImportTok{java}\OperatorTok{.}\ImportTok{io}\OperatorTok{.}\ImportTok{IOException}\OperatorTok{;}
\KeywordTok{import} \ImportTok{java}\OperatorTok{.}\ImportTok{io}\OperatorTok{.}\ImportTok{InputStreamReader}\OperatorTok{;}
\KeywordTok{import} \ImportTok{java}\OperatorTok{.}\ImportTok{net}\OperatorTok{.}\ImportTok{HttpURLConnection}\OperatorTok{;}
\KeywordTok{import} \ImportTok{java}\OperatorTok{.}\ImportTok{net}\OperatorTok{.}\ImportTok{URL}\OperatorTok{;}

\KeywordTok{public} \KeywordTok{class}\NormalTok{ TmdbApiExample }\OperatorTok{\{}
    \KeywordTok{public} \DataTypeTok{static} \DataTypeTok{void} \FunctionTok{main}\OperatorTok{(}\BuiltInTok{String}\OperatorTok{[]}\NormalTok{ args}\OperatorTok{)} \OperatorTok{\{}
        \ControlFlowTok{try} \OperatorTok{\{}
            \CommentTok{// API{-}nyckel från The Movie Database (TMDb)}
            \BuiltInTok{String}\NormalTok{ apiKey }\OperatorTok{=} \StringTok{"DIN\_API\_NYCKEL\_HÄR"}\OperatorTok{;}

            \CommentTok{// URL för att fråga efter filmen "Happy Death Day"}
            \BuiltInTok{String}\NormalTok{ url }\OperatorTok{=} \StringTok{"https://api.themoviedb.org/3/search/movie?api\_key="} \OperatorTok{+}\NormalTok{ apiKey }\OperatorTok{+} \StringTok{"\&query=Happy+death+day"}\OperatorTok{;}

            \CommentTok{// Skapa en URL{-}objekt från strängen}
            \BuiltInTok{URL}\NormalTok{ apiURL }\OperatorTok{=} \KeywordTok{new} \BuiltInTok{URL}\OperatorTok{(}\NormalTok{url}\OperatorTok{);}

            \CommentTok{// Öppna en anslutning till URL:en}
            \BuiltInTok{HttpURLConnection}\NormalTok{ connection }\OperatorTok{=} \OperatorTok{(}\BuiltInTok{HttpURLConnection}\OperatorTok{)}\NormalTok{ apiURL}\OperatorTok{.}\FunctionTok{openConnection}\OperatorTok{();}

            \CommentTok{// Ange att vi vill göra en GET{-}begäran}
\NormalTok{            connection}\OperatorTok{.}\FunctionTok{setRequestMethod}\OperatorTok{(}\StringTok{"GET"}\OperatorTok{);}

            \CommentTok{// Hämta svarskod från anslutningen}
            \DataTypeTok{int}\NormalTok{ responseCode }\OperatorTok{=}\NormalTok{ connection}\OperatorTok{.}\FunctionTok{getResponseCode}\OperatorTok{();}
            \BuiltInTok{System}\OperatorTok{.}\FunctionTok{out}\OperatorTok{.}\FunctionTok{println}\OperatorTok{(}\StringTok{"Svarskod: "} \OperatorTok{+}\NormalTok{ responseCode}\OperatorTok{);}

            \CommentTok{// Läs svaret från anslutningen}
            \BuiltInTok{BufferedReader}\NormalTok{ reader }\OperatorTok{=} \KeywordTok{new} \BuiltInTok{BufferedReader}\OperatorTok{(}\KeywordTok{new} \BuiltInTok{InputStreamReader}\OperatorTok{(}\NormalTok{connection}\OperatorTok{.}\FunctionTok{getInputStream}\OperatorTok{()));}
            \BuiltInTok{String}\NormalTok{ inputLine}\OperatorTok{;}
            \BuiltInTok{StringBuilder}\NormalTok{ response }\OperatorTok{=} \KeywordTok{new} \BuiltInTok{StringBuilder}\OperatorTok{();}
            \ControlFlowTok{while} \OperatorTok{((}\NormalTok{inputLine }\OperatorTok{=}\NormalTok{ reader}\OperatorTok{.}\FunctionTok{readLine}\OperatorTok{())} \OperatorTok{!=} \KeywordTok{null}\OperatorTok{)} \OperatorTok{\{}
\NormalTok{                response}\OperatorTok{.}\FunctionTok{append}\OperatorTok{(}\NormalTok{inputLine}\OperatorTok{);}
            \OperatorTok{\}}
\NormalTok{            reader}\OperatorTok{.}\FunctionTok{close}\OperatorTok{();}

            \CommentTok{// Visa svaret}
            \BuiltInTok{System}\OperatorTok{.}\FunctionTok{out}\OperatorTok{.}\FunctionTok{println}\OperatorTok{(}\NormalTok{response}\OperatorTok{.}\FunctionTok{toString}\OperatorTok{());}

            \CommentTok{// Stäng anslutningen}
\NormalTok{            connection}\OperatorTok{.}\FunctionTok{disconnect}\OperatorTok{();}
        \OperatorTok{\}} \ControlFlowTok{catch} \OperatorTok{(}\BuiltInTok{IOException}\NormalTok{ e}\OperatorTok{)} \OperatorTok{\{}
\NormalTok{            e}\OperatorTok{.}\FunctionTok{printStackTrace}\OperatorTok{();}
        \OperatorTok{\}}
    \OperatorTok{\}}
\OperatorTok{\}}
\end{Highlighting}
\end{Shaded}

Koden ovan gör följande:

\begin{enumerate}
\def\labelenumi{\arabic{enumi}.}
\tightlist
\item
  Den skapar en URL för att fråga The Movie Database API efter filmen
  ``Happy Death Day'' med hjälp av en API-nyckel.
\item
  Den skapar en HTTP-anslutning (HttpURLConnection) till URL:en och
  specifierar att vi vill göra en GET-begäran.
\item
  Den hämtar svarskoden från anslutningen och visar den.
\item
  Den läser svaret från anslutningen och sparar det i en sträng.
\item
  Den visar svaret på konsolen.
\item
  Den stänger anslutningen.
\end{enumerate}

\hypertarget{referenser}{%
\subsection{Referenser}\label{referenser}}

\begin{itemize}
\tightlist
\item
  \href{https://www.themoviedb.org/documentation/api}{The Movie Database
  API}
\item
  \href{https://docs.oracle.com/en/java/javase/11/docs/api/java.net/java/net/HttpURLConnection.html}{Oracle
  - HttpURLConnection}
\end{itemize}

Jag hoppas att du har lärt er något nytt om API:er och hur de kan
underlätta kommunikationen mellan webbapplikationer. För mer information
och djupgående dokumentation rekommenderar jag er att besöka
referenserna ovan. Ha en fortsatt bra dag och lycka till med era
programmeringsprojekt!

\hypertarget{obligatorisk-dad-joke}{%
\subsection{Obligatorisk Dad-joke}\label{obligatorisk-dad-joke}}

Varför älskar utvecklare att arbeta med APIer?

För att de ger dem en chans att ``PUT'' sina skämt ``GET'' ersättningar!
😄

\hypertarget{uxf6vrigt}{%
\section{Övrigt}\label{uxf6vrigt}}

Här finns några andra ämnen som inte riktigt passar någon annanstans.
Detta kommer att fyllas på med tiden\ldots{}

\hypertarget{sql}{%
\section{SQL}\label{sql}}

SQL står för Structured Query Language och är ett språk som används för
att kommunicera med databaser. Det är ett av de mest använda språken i
världen.

\emph{Filen editerades senast 2022-12-09} \#\# Referenser

\begin{itemize}
\tightlist
\item
  \href{https://www.w3schools.com/sql/default.asp}{SQL}
\item
  \href{https://www.w3schools.com/sql/sql_ref_mysql.asp}{SQL Cheat
  Sheet}
\item
  \href{https://sqlzoo.net/}{SQL Zoo}
\item
  \href{https://sqlfiddle.com/}{SQL Fiddle}
\end{itemize}

\hypertarget{tabeller}{%
\section{Tabeller}\label{tabeller}}

Tabeller är som ett rutnät av data som lagras i en databas. Varje tabell
har en eller flera kolumner, och varje rad i tabellen är en post. Varje
post har en värde för varje kolumn. En tabell kan ha en primärnyckel,
som är en kolumn som är unik för varje post. En tabell kan också ha en
eller flera främmande nycklar, som är en kolumn som är kopplad till en
annan tabell. Exemplen här nedanför är för MySQL. Det kan skilja på
andra servrar \{: warning\}

Innehållsförteckning

\{: .text-delta \}

\begin{enumerate}
\def\labelenumi{\arabic{enumi}.}
\tightlist
\item
  TOC \{:toc\}
\end{enumerate}

\emph{Filen editerades senast 2022-12-09} \#\#\# Skapa en tabell om den
inte finns

Här skapar vi en tabell med namnet \texttt{users} och tre kolumner:
\texttt{id}, \texttt{name} och \texttt{email}.

\hypertarget{cb1}{}
\begin{Shaded}
\begin{Highlighting}[]
\KeywordTok{CREATE} \KeywordTok{TABLE} \ControlFlowTok{IF} \KeywordTok{NOT} \KeywordTok{EXISTS}\NormalTok{ users (}
\KeywordTok{id} \DataTypeTok{INT}\NormalTok{(}\DecValTok{6}\NormalTok{) UNSIGNED AUTO\_INCREMENT }\KeywordTok{PRIMARY} \KeywordTok{KEY}\NormalTok{,}
\NormalTok{name }\DataTypeTok{VARCHAR}\NormalTok{(}\DecValTok{30}\NormalTok{) }\KeywordTok{NOT} \KeywordTok{NULL}\NormalTok{,}
\NormalTok{email }\DataTypeTok{VARCHAR}\NormalTok{(}\DecValTok{50}\NormalTok{)}
\NormalTok{);}
\end{Highlighting}
\end{Shaded}

\hypertarget{luxe4gg-till-en-kolumn}{%
\subsubsection{Lägg till en kolumn}\label{luxe4gg-till-en-kolumn}}

\hypertarget{cb2}{}
\begin{Shaded}
\begin{Highlighting}[]
\KeywordTok{ALTER} \KeywordTok{TABLE}\NormalTok{ users }\KeywordTok{ADD} \KeywordTok{COLUMN}\NormalTok{ age }\DataTypeTok{INT}\NormalTok{(}\DecValTok{3}\NormalTok{);}
\end{Highlighting}
\end{Shaded}

\hypertarget{ta-bort-en-kolumn}{%
\subsubsection{Ta bort en kolumn}\label{ta-bort-en-kolumn}}

\hypertarget{cb3}{}
\begin{Shaded}
\begin{Highlighting}[]
\KeywordTok{ALTER} \KeywordTok{TABLE}\NormalTok{ users }\KeywordTok{DROP} \KeywordTok{COLUMN}\NormalTok{ age;}
\end{Highlighting}
\end{Shaded}

\hypertarget{referenser}{%
\subsection{Referenser}\label{referenser}}

\hypertarget{luxe4gg-till-en-kolumn-och-guxf6r-den-till-en-foreignkey}{%
\subsubsection{Lägg till en kolumn och gör den till en
ForeignKey}\label{luxe4gg-till-en-kolumn-och-guxf6r-den-till-en-foreignkey}}

Vi skapar en länk till en annan tabell. I detta fall till tabellen
\texttt{cars} och kolumnen \texttt{CarId}.

\hypertarget{cb4}{}
\begin{Shaded}
\begin{Highlighting}[]
\KeywordTok{ALTER} \KeywordTok{TABLE}\NormalTok{ users }\KeywordTok{ADD} \KeywordTok{COLUMN}\NormalTok{ CarId }\DataTypeTok{INT}\NormalTok{(}\DecValTok{6}\NormalTok{) UNSIGNED;}
\KeywordTok{ALTER} \KeywordTok{TABLE}\NormalTok{ users }\KeywordTok{ADD} \KeywordTok{FOREIGN} \KeywordTok{KEY}\NormalTok{ (CarId) }\KeywordTok{REFERENCES}\NormalTok{ cars(CarId);}
\end{Highlighting}
\end{Shaded}

\hypertarget{ta-bort-en-foreignkey}{%
\subsection{Ta bort en ForeignKey}\label{ta-bort-en-foreignkey}}

Vi tar bort referensen till Cars.CarId

\hypertarget{cb5}{}
\begin{Shaded}
\begin{Highlighting}[]
\KeywordTok{ALTER} \KeywordTok{TABLE}\NormalTok{ users }\KeywordTok{DROP} \KeywordTok{FOREIGN} \KeywordTok{KEY}\NormalTok{ users\_ibfk\_1;}
\end{Highlighting}
\end{Shaded}

users\_ibfk\_1 är den första ForeignKey som finns i tabellen users.

\hypertarget{skapa-en-tabell-med-en-foreignkey}{%
\subsection{Skapa en tabell med en
ForeignKey}\label{skapa-en-tabell-med-en-foreignkey}}

\hypertarget{cb6}{}
\begin{Shaded}
\begin{Highlighting}[]
\KeywordTok{CREATE} \KeywordTok{TABLE} \ControlFlowTok{IF} \KeywordTok{NOT} \KeywordTok{EXISTS}\NormalTok{ cars (}
\NormalTok{CarId }\DataTypeTok{INT}\NormalTok{(}\DecValTok{6}\NormalTok{) UNSIGNED AUTO\_INCREMENT }\KeywordTok{PRIMARY} \KeywordTok{KEY}\NormalTok{,}
\NormalTok{Model }\DataTypeTok{VARCHAR}\NormalTok{(}\DecValTok{30}\NormalTok{) }\KeywordTok{NOT} \KeywordTok{NULL}\NormalTok{,}
\NormalTok{Brand }\DataTypeTok{VARCHAR}\NormalTok{(}\DecValTok{50}\NormalTok{),}
\NormalTok{UserId }\DataTypeTok{INT}\NormalTok{(}\DecValTok{6}\NormalTok{) UNSIGNED,}
\KeywordTok{FOREIGN} \KeywordTok{KEY}\NormalTok{ (UserId) }\KeywordTok{REFERENCES}\NormalTok{ users(}\KeywordTok{id}\NormalTok{)}
\NormalTok{);}
\end{Highlighting}
\end{Shaded}

\hypertarget{skapa-en-tabell-med-en-foreignkey-och-en-primuxe4rnyckel}{%
\subsection{Skapa en tabell med en ForeignKey och en
primärnyckel}\label{skapa-en-tabell-med-en-foreignkey-och-en-primuxe4rnyckel}}

\hypertarget{cb7}{}
\begin{Shaded}
\begin{Highlighting}[]
\KeywordTok{CREATE} \KeywordTok{TABLE} \ControlFlowTok{IF} \KeywordTok{NOT} \KeywordTok{EXISTS}\NormalTok{ cars (}
\NormalTok{CarId }\DataTypeTok{INT}\NormalTok{(}\DecValTok{6}\NormalTok{) UNSIGNED AUTO\_INCREMENT }\KeywordTok{PRIMARY} \KeywordTok{KEY}\NormalTok{,}
\NormalTok{Model }\DataTypeTok{VARCHAR}\NormalTok{(}\DecValTok{30}\NormalTok{) }\KeywordTok{NOT} \KeywordTok{NULL}\NormalTok{,}
\NormalTok{Brand }\DataTypeTok{VARCHAR}\NormalTok{(}\DecValTok{50}\NormalTok{),}
\NormalTok{UserId }\DataTypeTok{INT}\NormalTok{(}\DecValTok{6}\NormalTok{) UNSIGNED,}
\KeywordTok{PRIMARY} \KeywordTok{KEY}\NormalTok{ (CarId, UserId),}
\KeywordTok{FOREIGN} \KeywordTok{KEY}\NormalTok{ (UserId) }\KeywordTok{REFERENCES}\NormalTok{ users(}\KeywordTok{id}\NormalTok{)}
\NormalTok{);}
\end{Highlighting}
\end{Shaded}

\hypertarget{referenser-1}{%
\subsection{Referenser}\label{referenser-1}}

\begin{itemize}
\tightlist
\item
  \href{https://www.w3schools.com/sql/sql_foreignkey.asp}{W3Schools}
\item
  \href{https://dev.mysql.com/doc/refman/8.0/en/create-table-foreign-keys.html}{MySQL}
\end{itemize}

\hypertarget{databaser}{%
\section{Databaser}\label{databaser}}

Innehållsförteckning

\{: .text-delta \}

\begin{enumerate}
\def\labelenumi{\arabic{enumi}.}
\tightlist
\item
  TOC \{:toc\}
\end{enumerate}

\emph{Filen editerades senast 2022-12-09} \#\#\# Skapa en databas

\hypertarget{cb1}{}
\begin{Shaded}
\begin{Highlighting}[]
\KeywordTok{CREATE} \KeywordTok{DATABASE}\NormalTok{ databasnamn;}
\end{Highlighting}
\end{Shaded}

\hypertarget{skapa-en-databas-om-den-inte-finns}{%
\subsubsection{Skapa en databas om den inte
finns}\label{skapa-en-databas-om-den-inte-finns}}

\hypertarget{cb2}{}
\begin{Shaded}
\begin{Highlighting}[]
\KeywordTok{CREATE} \KeywordTok{DATABASE} \ControlFlowTok{IF} \KeywordTok{NOT} \KeywordTok{EXISTS}\NormalTok{ databasnamn;}
\end{Highlighting}
\end{Shaded}

\hypertarget{ta-bort-en-databas}{%
\subsubsection{Ta bort en databas}\label{ta-bort-en-databas}}

Tänk på att backuppa först, annars är allt borta.\{: warning\}

\hypertarget{cb3}{}
\begin{Shaded}
\begin{Highlighting}[]
\KeywordTok{DROP} \KeywordTok{DATABASE}\NormalTok{ databasnamn;}
\end{Highlighting}
\end{Shaded}

\hypertarget{uxe4ndra-namn-puxe5-databasen}{%
\subsubsection{Ändra namn på
databasen}\label{uxe4ndra-namn-puxe5-databasen}}

Farligt! Detta kan skapa problem om du har kopplat dig till databasen
med ett program.\{: warning\}

\hypertarget{cb4}{}
\begin{Shaded}
\begin{Highlighting}[]
\KeywordTok{ALTER} \KeywordTok{DATABASE}\NormalTok{ databasnamn }\KeywordTok{RENAME} \KeywordTok{TO}\NormalTok{ nytt\_databasnamn;}
\end{Highlighting}
\end{Shaded}

\hypertarget{constraints}{%
\section{Constraints}\label{constraints}}

SQL-Constraints är ett sätt att begränsa datan i en tabell. Det finns
två typer av constraints, primary key och foreign key. Primary key är en
constraint som gör att en kolumn i en tabell kan ha unika värden.
Foreign key är en constraint som gör att en kolumn i en tabell kan ha
värden som finns i en annan tabell.

Innehållsförteckning

\{: .text-delta \}

\begin{enumerate}
\def\labelenumi{\arabic{enumi}.}
\tightlist
\item
  TOC \{:toc\}
\end{enumerate}

\emph{Filen editerades senast 2022-12-09} \#\# Primary Key

Primary Key är det värde som är unikt för varje rad i en tabell. Det är
det värde som används för att identifiera en rad i en tabell. Det är
också det värde som används för att koppla ihop tabeller med varandra.
Det är inte möjligt att ha fler än en primary key i en tabell. Det är
också inte möjligt att ha en primary key som är null. Det är också
möjligt att ha en primary key som är en kombination av flera kolumner.

\hypertarget{foreign-key}{%
\subsection{Foreign Key}\label{foreign-key}}

Foreign key är ett värde som finns i en tabell som är kopplat till ett
värde i en annan tabell. Det är det värde som används för att koppla
ihop tabeller med varandra. Det är inte möjligt att ha fler än en
foreign key i en tabell. Det är också inte möjligt att ha en foreign key
som är null. Det är också möjligt att ha en foreign key som är en
kombination av flera kolumner.

\hypertarget{exempel}{%
\subsection{Exempel}\label{exempel}}

\hypertarget{cb1}{}
\begin{Shaded}
\begin{Highlighting}[]
\KeywordTok{CREATE} \KeywordTok{TABLE}\NormalTok{ Person (}
\NormalTok{PersonID }\DataTypeTok{int} \KeywordTok{NOT} \KeywordTok{NULL} \KeywordTok{PRIMARY} \KeywordTok{KEY}\NormalTok{,}
\NormalTok{LastName }\DataTypeTok{varchar}\NormalTok{(}\DecValTok{255}\NormalTok{) }\KeywordTok{NOT} \KeywordTok{NULL}\NormalTok{,}
\NormalTok{FirstName }\DataTypeTok{varchar}\NormalTok{(}\DecValTok{255}\NormalTok{),}
\NormalTok{Address }\DataTypeTok{varchar}\NormalTok{(}\DecValTok{255}\NormalTok{),}
\NormalTok{City }\DataTypeTok{varchar}\NormalTok{(}\DecValTok{255}\NormalTok{),}
\KeywordTok{PRIMARY} \KeywordTok{KEY}\NormalTok{ (PersonID)}
\NormalTok{);}
\KeywordTok{CREATE} \KeywordTok{TABLE}\NormalTok{ Orders (}
\NormalTok{OrderID }\DataTypeTok{int} \KeywordTok{NOT} \KeywordTok{NULL} \KeywordTok{PRIMARY} \KeywordTok{KEY}\NormalTok{,}
\NormalTok{OrderDate }\DataTypeTok{date} \KeywordTok{NOT} \KeywordTok{NULL}\NormalTok{,}
\NormalTok{PersonID }\DataTypeTok{int} \KeywordTok{NOT} \KeywordTok{NULL}\NormalTok{,}
\KeywordTok{PRIMARY} \KeywordTok{KEY}\NormalTok{ (OrderID),}
\KeywordTok{FOREIGN} \KeywordTok{KEY}\NormalTok{ (PersonID) }\KeywordTok{REFERENCES}\NormalTok{ Person(PersonID)}
\NormalTok{);}
\end{Highlighting}
\end{Shaded}

Man använder Primary key för att identifiera en rad i en tabell, och
Foreign key för att koppla ihop tabeller med varandra.

\hypertarget{exempel-puxe5-hur-man-kopplar-ihop-tabellerna}{%
\subsection{Exempel på hur man kopplar ihop
tabellerna}\label{exempel-puxe5-hur-man-kopplar-ihop-tabellerna}}

\hypertarget{cb2}{}
\begin{Shaded}
\begin{Highlighting}[]
\KeywordTok{INSERT} \KeywordTok{INTO}\NormalTok{ Person (PersonID, LastName, FirstName, Address, City)}
\KeywordTok{VALUES}\NormalTok{ (}\DecValTok{1}\NormalTok{, }\StringTok{\textquotesingle{}Peter\textquotesingle{}}\NormalTok{, }\StringTok{\textquotesingle{}Petersson\textquotesingle{}}\NormalTok{, }\StringTok{\textquotesingle{}St Petersgatan 1\textquotesingle{}}\NormalTok{, }\StringTok{\textquotesingle{}Peterstad\textquotesingle{}}\NormalTok{);}
\KeywordTok{INSERT} \KeywordTok{INTO}\NormalTok{ Orders (OrderID, OrderDate, PersonID)}
\KeywordTok{VALUES}\NormalTok{ (}\DecValTok{1}\NormalTok{, }\StringTok{\textquotesingle{}2021{-}01{-}01\textquotesingle{}}\NormalTok{, }\DecValTok{1}\NormalTok{);}
\end{Highlighting}
\end{Shaded}

och vi kan sedan hämta ut informationen från båda tabellerna med
följande kod:

\hypertarget{cb3}{}
\begin{Shaded}
\begin{Highlighting}[]
\KeywordTok{SELECT} \OperatorTok{*} \KeywordTok{FROM}\NormalTok{ Person}
\KeywordTok{INNER} \KeywordTok{JOIN}\NormalTok{ Orders }\KeywordTok{ON}\NormalTok{ Person.PersonID }\OperatorTok{=}\NormalTok{ Orders.PersonID;}
\end{Highlighting}
\end{Shaded}

\hypertarget{referenser}{%
\subsection{Referenser}\label{referenser}}

\href{https://www.w3schools.com/sql/sql_constraints.asp}{SQL
Constraints}

\hypertarget{ordlista}{%
\section{Ordlista}\label{ordlista}}

Här finns några ordlistor med vanliga ord och uttryck som används inom
IT världen, med korta förklaringar.

\hypertarget{git}{%
\section{Git}\label{git}}

\begin{longtable}[]{@{}
  >{\raggedright\arraybackslash}p{(\columnwidth - 2\tabcolsep) * \real{0.5000}}
  >{\raggedright\arraybackslash}p{(\columnwidth - 2\tabcolsep) * \real{0.5000}}@{}}
\toprule\noalign{}
\begin{minipage}[b]{\linewidth}\raggedright
Ord
\end{minipage} & \begin{minipage}[b]{\linewidth}\raggedright
Förklaring
\end{minipage} \\
\midrule\noalign{}
\endhead
\bottomrule\noalign{}
\endlastfoot
Repository & En mapp som innehåller alla filer och historiken för en
projekt. \\
Branch & En gren i repot. Det är en ögonblickskopia av repot. \\
Commit & En commit är en punkt i historiken. Det är en punkt där du
sparar dina ändringar. \\
Push & Att skicka upp ändringar till servern. \\
Pull & Att hämta ner ändringar från servern. \\
Merge & Att slå ihop två grenar. \\
Rebase & Att flytta en gren till en annan gren. \\
Stash & Att spara undan ändringar som inte är klara. \\
Checkout & Att byta gren. \\
Clone & Att kopiera ett repo till din dator. \\
Fork & Att kopiera ett repo till din github. \\
Pull request & Att be servern merga två brancher. \\
Merge conflict & När två grenar har ändrat samma fil på samma ställe. \\
Remote & En koppling till ett annat repo. \\
Origin & Det vanligaste namnet på en remote. \\
Upstream & Det vanligaste namnet på en remote som inte är din. \\
Fork & Att kopiera ett repo till din github. \\
Pull request & Att be servern merga två brancher. \\
\end{longtable}

\hypertarget{programmering}{%
\section{Programmering}\label{programmering}}

\begin{longtable}[]{@{}
  >{\raggedright\arraybackslash}p{(\columnwidth - 2\tabcolsep) * \real{0.5000}}
  >{\raggedright\arraybackslash}p{(\columnwidth - 2\tabcolsep) * \real{0.5000}}@{}}
\toprule\noalign{}
\begin{minipage}[b]{\linewidth}\raggedright
Ord
\end{minipage} & \begin{minipage}[b]{\linewidth}\raggedright
Förklaring
\end{minipage} \\
\midrule\noalign{}
\endhead
\bottomrule\noalign{}
\endlastfoot
Algorithm & En algoritm är lista av kommandon i steg för steg-
beskrivning, för att lösa ett problem. \\
Attribut & En variabel som tillhör ett objekt. \\
Deserialisera & Att göra om en sträng till ett objekt. \\
Get & En metod som returnerar en variabel. \\
Instansvariabel & En variabel som tillhör ett objekt. \\
Interface & En mall för en klass. \\
Klass & En mall för ett objekt. \\
Klassvariabel & En variabel som tillhör en klass. \\
Konstruktor & En metod som körs när ett objekt skapas. \\
Metod & En funktion som tillhör ett objekt. \\
Modifierare & En modifierare är en modifierare som ändrar hur en
variabel eller metod fungerar. \\
Namespace & En grupp av klasser. \\
Namespace alias & Att ge en namespace ett alias. \\
Namespace import & Att importera en namespace. \\
Namespace using & Att använda en namespace. \\
Objekt & En instans av en klass. \\
Property & En metod som sätter och returnerar en variabel. \\
Serialisera & Att göra om ett objekt till en sträng. \\
Set & En metod som sätter en variabel. \\
Static & En modifierare som gör att en variabel tillhör klassen och inte
objektet. \\
using static & Att använda en klass som namespace och dess statiska
metoder. \\
\end{longtable}

\hypertarget{programmeringsspruxe5k}{%
\section{Programmeringsspråk}\label{programmeringsspruxe5k}}

--- \textbar{} --- \textbar{}\\
Angular \textbar{} Ett JavaScript-bibliotek. \textbar{}\\
Composer \textbar{} Ett pakethanteringsverktyg. \textbar{}\\
CSS \textbar{} Cascading Style Sheets. \textbar{}\\
Docker \textbar{} Ett program som gör det enkelt att köra program i en
container.\textbar{}\\
Docker Compose \textbar{} Ett verktyg för att skapa och köra
fleraDocker-containers. \textbar{}\\
Docker Desktop \textbar{} Ett program som gör det enkelt att köra
Docker. \textbar{}\\
Docker Hub \textbar{} Ett webbaserat verktyg för att lagra
Docker-containers.\textbar{}\\
Docker Machine \textbar{} Ett verktyg för att skapa och köra
fleraDocker-containers. \textbar{}\\
Docker Registry \textbar{} Ett webbaserat verktyg för att
lagraDocker-containers. \textbar{}\\
Docker Swarm \textbar{} Ett verktyg för att skapa och köra
fleraDocker-containers. \textbar{}\\
Docker Toolbox \textbar{} Ett program som gör det enkelt att köra
Docker. \textbar{}\\
Gatsby \textbar{} Ett statiskt webbplatsbyggverktyg. \textbar{}\\
Git \textbar{} Ett versionshanteringssystem. \textbar{}\\
GitHub \textbar{} Ett webbaserat versionshanteringssystem. \textbar{}\\
GitLab \textbar{} Ett webbaserat versionshanteringssystem. \textbar{}\\
HTML \textbar{} Hypertext Markup Language. \textbar{}\\
Hugo \textbar{} Ett statiskt webbplatsbyggverktyg. \textbar{}\\
JavaScript \textbar{} Ett programmeringsspråk. \textbar{}\\
Jekyll \textbar{} Ett statiskt webbplatsbyggverktyg. \textbar{}\\
JSON \textbar{} JavaScript Object Notation. \textbar{}\\
Markdown \textbar{} Ett textformat. \textbar{}\\
Node \textbar{} Ett JavaScript-bibliotek. \textbar{}\\
NPM \textbar{} Ett pakethanteringsverktyg. \textbar{}\\
Ord \textbar{} Förklaring \textbar{}\\
PHP \textbar{} Ett programmeringsspråk. \textbar{}\\
React \textbar{} Ett JavaScript-bibliotek. \textbar{}\\
SQL \textbar{} Structured Query Language. \textbar{}\\
Vue \textbar{} Ett JavaScript-bibliotek. \textbar{}\\
XML \textbar{} Extensible Markup Language. \textbar{}\\
YAML \textbar{} Yet Another Markup Language. \textbar{}\\
Yarn \textbar{} Ett pakethanteringsverktyg. \textbar{}

\hypertarget{webb}{%
\section{Webb}\label{webb}}

\begin{longtable}[]{@{}ll@{}}
\toprule\noalign{}
Ord & Förklaring \\
\midrule\noalign{}
\endhead
\bottomrule\noalign{}
\endlastfoot
CDN & Content Delivery Network. \\
DNS & Domain Name System. \\
DNS & Domain Name System. \\
FTP & File Transfer Protocol. \\
FTPS & File Transfer Protocol Secure. \\
HTTP & Hypertext Transfer Protocol. \\
HTTPS & Hypertext Transfer Protocol Secure. \\
IMAP & Internet Message Access Protocol. \\
IP & Internet Protocol. \\
POP & Post Office Protocol. \\
SFTP & Secure File Transfer Protocol. \\
SMTP & Simple Mail Transfer Protocol. \\
SSH & Secure Shell. \\
SSL & Secure Sockets Layer. \\
TLS & Transport Layer Security. \\
URL & Uniform Resource Locator. \\
VPN & Virtual Private Network. \\
VPS & Virtual Private Server. \\
\end{longtable}

\hypertarget{yrken}{%
\section{Yrken}\label{yrken}}

\begin{longtable}[]{@{}
  >{\raggedright\arraybackslash}p{(\columnwidth - 2\tabcolsep) * \real{0.5000}}
  >{\raggedright\arraybackslash}p{(\columnwidth - 2\tabcolsep) * \real{0.5000}}@{}}
\toprule\noalign{}
\begin{minipage}[b]{\linewidth}\raggedright
Ord
\end{minipage} & \begin{minipage}[b]{\linewidth}\raggedright
Förklaring
\end{minipage} \\
\midrule\noalign{}
\endhead
\bottomrule\noalign{}
\endlastfoot
Backend & Logiken bakom ett program. \\
Backend utvecklare & En person som hjälper företag med backend. \\
Cloud & En molntjänst. \\
Cloud Architect & En person som planerar och utvecklar molntjänster. \\
Cloud Consultant & En person som hjälper företag med molntjänster. \\
Cloud Engineer & En person som hjälper företag med molntjänster. \\
Cloud Security & En person som hjälper företag med molntjänster. \\
Cloud utvecklare & En person som hjälper företag med molntjänster. \\
DevOps & En metod för att utveckla och driftsätta program. \\
DevOps Architect & En person som planerar och utvecklar DevOps. \\
DevOps Consultant & En person som hjälper företag med DevOps. \\
DevOps Engineer & En person som hjälper företag med DevOps. \\
Frontend & Det som användaren ser. Användargränssnittet. \\
Frontend utvecklare & En person som hjälper företag med frontend. \\
Fullstack & En person som kan utveckla både frontend och backend. \\
IT & Informationsteknik. \\
IT-arkitekt & En person som planerar och utvecklar IT-system. \\
IT-konsult & En person som hjälper företag med IT. \\
IT-säkerhet & En person som hjälper företag med IT-säkerhet. \\
IT-support & En person som hjälper användare med IT. \\
IT-tekniker & En person som hjälper användare med IT. \\
Systemarkitekt & En person som planerar och utvecklar system. \\
Systemutvecklare & En person som hjälper företag med
systemutveckling. \\
\end{longtable}
